\documentclass[12pt]{article}
\usepackage{float}
\usepackage{xcolor}
\usepackage{newpxtext,euler}
\usepackage[T1]{fontenc}
\usepackage[spanish]{babel}
\usepackage{amsfonts,amsmath,amssymb,amsthm}
\usepackage{hyperref}
\usepackage{geometry}

\newcommand\N{\ensuremath{\mathbb{N}}}
\newcommand\R{\ensuremath{\mathbb{R}}}
\newcommand\Z{\ensuremath{\mathbb{Z}}}
\renewcommand\O{\ensuremath{\emptyset}}
\newcommand\Q{\ensuremath{\mathbb{Q}}}
\newcommand\C{\ensuremath{\mathbb{C}}}
\newcommand\T{\mathbb{T}}
\renewcommand{\epsilon}{\varepsilon}
\renewcommand{\hat}{\widehat}
\newcommand\jk{\langle k\rangle}


\setlength{\parindent}{0pt}


\input{boxes}
\usepackage{lipsum}
 \geometry{
 a4paper,
 total={170mm,260mm},
 left=20mm,
 top=15mm,
 }
 
\title{\vspace{-2cm}\par\noindent\rule{16cm}{1pt}\large
\\\bfseries Teoría Hiperbólica de Nudos
\vspace{-0.34cm}\par\noindent\hspace{0.15cm}\rule{16cm}{1pt}
\vspace{-0.6cm}
}
\author{\small \bfseries \quad \small Edgar Santiago Ochoa Quiroga\\ \small \quad \quad \texttt{eochoa@unal.edu.co}\quad\quad \quad\\}

\usepackage{titling}
\predate{\hspace{6.8cm}\small}
\postdate{}

\begin{document}
\maketitle
\begin{abstract}
    Abordaremos un primer acercamiento a la teoría hiperbólica de nudos, las bases de esta teoría y de la geometría hiperbólica enfocándonos en particular en el la estructura hiperbólica del complemento del 8-Nudo respecto a $S^3$ y de su construcción, incluyendo los ejemplos del nudo $6_1$ y del enlace de Whitehead que no son estándar encontrarlos hechos y un breve comentario acerca de la conjetura del volumen.
\end{abstract}
\section{Introducción}
    La primera vez que uno cruza caminos con el concepto de nudo, muy ocasionalmente no es desde el ámbito formal matemático al que la mayoría de los posibles lectores de este proyecto estén acostumbrados, sino es desde el ámbito mas geométrico y tangible, curiosamente al pasar al formalismo pareciera que toda la intuición geométrica usual para trabajar con ellos resulta inefectiva y las herramientas que se desarrollan se alejan de la misma y pasan a ser topológicas y algebraicas.\\

    La idea de este proyecto sera reencontrarnos con la geometría, en particular con la geometría hiperbólica, por lo que en primera instancia exploraremos los conceptos básicos necesarios de teoría de nudos y de la geometría hiperbólica, para posteriormente estudiar la construcción de Thurston del complemento del 8-Nudo sobre $S^3$ y algunos ejemplos adicionales mas comentarios acerca de un proceso general.\\

    Por el trabajo de Thurston un complemento de nudo tiene una de tres formas: o bien es un nudo toroidal, que puede dibujarse como una curva incrustada en un toro de Heegaard en $S^3$, o bien es un nudo satélite, que se encuentra en una vecindad tubular de otro nudo, o bien es hiperbólico, curiosamente esta ultima clase de nudos es la que contiene la mayor cantidad de ejemplos y ademas es de la cual se entiende menos su estructura. Estudiaremos como se dota al complemento de un nudo de estructura hiperbólica con ayuda de la descompocision y veremos las condiciones para que esta estructura este bien definida y sea completa, sin olvidarnos de hacer ejemplos al respecto y finalizaremos con un comentario acerca del camino que queda para entender los nudos hiperbólicos por medio de invariantes geométricos y la conjetura del volumen.
    \section{Preliminares}
    \subsection{Teoría de Nudos}
    %!TEX root = ./main.tex

Dentro del contexto de los posibles lectores de este documento, todo habremos estado en sintonía con el curso por lo que unos preliminares de teoría de nudos excesivamente extensos resultarían innecesarios, pero por formalidad y el propósito mas geométrico de este trabajo, resulta imperante dejar en claro cual va a ser el conjunto en el que vamos a embeber los nudos para dotarlos de esta geometría.

\begin{definition}
    Un nudo $K\subset S^3$ es un subconjunto de puntos homeomorfico a $S^1$ bajo un homeomorfismo PL (Lineal a trozos), también podemos pensar en un nudo como un embedimiento PL $K:S^1\to S^3$, y usaremos $K$ para referirnos tanto a la función como a $K(S^1)$.
\end{definition}

Al igual que siempre este homeomorfismo PL puede ser cambiado por un homeomorfismo suave en el sentido usual y en caso de enlace se define una función de la unión de copias disyuntas de $S^1$, la principal diferencia radica en cambiar $\R^3$ por $S^3$, pensando netamente en los espacios, la única ventaja aparente es el hecho de que el espacio escogido en este documento es compacto, pero mas adelante se vera el por que de esta elección desde un punto de vista geométrico, por el momento continuemos con las nociones usuales.

\begin{definition}
  Decimos que dos nudos (o enlaces) $K_1$ y $K_2$ son equivalentes si existe una isotopia ambiente, es decir una función $h:S^3\times [0,1]\to S^3$  continua  tal que $h_t:=h(-,t):S^3\to S^3$ es un homeomorfismo para cada $t$ fijo. Ademas $h_0(K_1)=K_1$ y $h_1(K_1)=K_2.$ 
\end{definition}

Si bien esta definición la volvemos a dar sobre $S^3$, es útil seguir pensando en este espacio como $\R^3\cup\{\infty\}$ al momento de visualizar la isotopia ambiente.

\begin{definition}
  Para un nudo $K$, denotamos por $N(K)$ una vecindad regular abierta de $K$ en $S^3$. El exterior del nudo es la variedad $S^3\setminus N(K)$ y el complemento del nudo es $S^3\setminus K.$
\end{definition}
A partir de estas tres definiciones incluidas en una sola debemos hacer varios comentarios. El primero de ellos es que  que tanto la vecindad como el exterior son 3-variedades, donde la primera es abierta por definición, mientras la segunda al ser el complemento de un abierto, es cerrada, y al ser $S^3$ compacto, es compacta (uno de los beneficios de estar en $S^3$).\\

Si bien no es inmediato, la frontera del exterior del nudo, es homeomorfa a un toro, esto es mas claro si pensamos en un tubo donde el nudo fluye por dentro. Normalmente en vez de vecindad regular se habla directamente de vecindad tubular.\\

Por ultimo para el caso de enlaces se toma unión disyunta de vecindades tubulares adecuadas para cada componente y definimos las mismas nociones.

Como siempre el propósito principal es clasificar y poder distinguir nudos (enlaces) entre si, para esto Thurston \cite{thurston_3d_geometry_topology} desarrollo invariantes con herramientas de la geometría hiperbólica, pero para poder si quiera entender un poco de ello tenemos que ver un poco de la misma.
    \subsection{Geometría Hiperbólica}
    %!TEX root = ./main.tex

Al ser preliminares no se entrara tanto en detalle en algunos resultados, mas se presentaran cálculos básicos necesarios para posteriormente poder calcular directamente en los complementos de los Nudos, por lo que al trabajar con las proyecciones de nudos en $S^2$ (Plano) es natural empezar entendiendo nociones básicas de la geometría hiperbólica en dimensión 2.

\begin{definition}
    El espacio hiperbólico 2-dimensional $\mathbb{H}^2$ es el conjunto de puntos
    $$\mathbb{H}^2:=\{z\in\C:\Im(z)>0\},$$
    Donde si $z=x+iy$, dotamos al espacio de una métrica dada por
    $$ds^2=\frac{dx^2+dy^2}{y^2}$$
\end{definition}
Esta métrica se puede interpretar como la métrica usual de el plano $dx^2+dy^2$ re escalada por el factor. Note en particular que al ser puntos del plano complejo tenemos dos interpretaciones, como un punto $z\in\C$ o como $(x,y)\in\R^2$, ambas serán útiles dependiendo del contexto por lo que cambiaremos entre ellas según convenga.\\

Esta idea de métrica y del espacio hiperbólico se pueden ver mas a fondo desde la perspectiva de la geometría Riemanniana, mas debido a que esto es un curso completo nos quedaremos con las formulas fundamentales, primero la que nos permite calcular longitudes.
\begin{definition}
    Sea $\gamma:[a.b]\to \mathbb{H}^2$ una curva diferenciable, definimos la longitud de la curva como 
    $$|\gamma|=\int_a^b\frac{1}{\gamma_y(s)}\sqrt{(\gamma^{\prime}_x(s))^2+(\gamma^{\prime}_y(s))^2}\,ds$$
    donde $\gamma(s)=(\gamma_x(s),\gamma_y(s)).$
\end{definition}


\begin{eg}
    Las longitudes mas básicas de calcular en la geometría euclidiana usual son las rectas horizontales y verticales, así que hagamos lo mismo. Considere el segmento que va del punto $(0,h)$ a el punto $(1,h)$, este segmento esta parametrizado por $\gamma(t)=(t,h)$ con $t\in[0,1]$, así obtenemos
    $$|\gamma|=\int_0^1\frac{1}{h}\sqrt{1^2+0^2}\,ds=\frac{1}{h}$$
    Que es lo esperado la distancia usual del segmento euclidiano (1) re escalada por la altura. Observe que esto nos dice que cuando $h\to\infty$ los puntos a pesar de que en el semiplano superior se ven a una distancia fija en realidad se están acercando, de manera similar si $h\to 0^+$ los puntos se alejan.\\

    Ahora consideremos una recta vertical, es decir una recta parametrizada por $\gamma(t)=(x,t)$, donde $x$ esta fijo y $t\in[a,b]$, por la formula tenemos
    $$|\gamma|=\int_a^b\frac{1}{s}\sqrt{0^2+1^2}\,ds=\log\left(\frac{b}{a}\right).$$
    De aquí podemos observar que si fijamos $a$ y $b\to\infty$ tenemos que estos se alejan, y si $b$ es fijo y $a\to 0^+$ la distancia también tiende a infinito.
    \end{eg}

    Este comportamiento que se sale de la intuición del dibujo euclideo también se traslada a la siguiente formula para calcular áreas y el calculo particular.

    \begin{definition}
        Dada una región $R\subset \bb{H}^2$, definimos el área como
        $$\text{area}(R)=\int_R\frac{1}{y^2}dR.$$
    \end{definition}
    \begin{eg}
        Como ejemplo de calculo considere la siguiente región
         $$R=\{(x,y)\in \bb{R}^2: 0\leq x\leq 1,y\geq 1\},$$
         Esto es una franja que se extiende al infinito, usando la formula
         $$\text{area}(R)=\int_0^1\int_1^\infty\frac{1}{y^2}\,dydx=\int_0^11\,dx=1.$$
         Esto nos muestra como regiones euclideas con área infinita pueden tener área finita, esto sera excesivamente relevante cuando hablemos del volumen hiperbólico de un nudo.
    \end{eg}

    Observe que a pesar de no pertenecer al espacio hiperbólico, la recta real y acercarse al infinito nos da información geométrica del mismo, por esto es común otorgarles nombre propio para referirse a ambos en conjunto.

    \begin{definition}
        Llamamos a $\bb{R}\cup\{\infty\}$ la frontera en el infinito de $\bb{H}^2$, y la denotamos por $\partial_\infty\bb{H}^2$
    \end{definition}

    Una observación particular es que al igual que antes por proyección estereográfica tenemos que $\partial_\infty\bb{H}^2\cong S^1$.\\

    Entre estos puntos de la frontera y todos los puntos que se encuentran en $\bb{H}^2$ nos gustaría encontrar aquellas curvas que minimizan la distancia entre puntos, en el caso euclideo estas siempre son rectas, pero en general no tendrían por que serlo en otras métricas, así definimos la siguiente nocion
    \begin{definition}
        Dados puntos  $p,q\in\bb{H}^2$ decimos que una geodésica entre estos puntos es una curva $\gamma$ que justamente minimiza $|\gamma|$. Una geodésica infinita es una curva $\gamma:\bb{R}\to \bb{H}^2$ tal que dado un intervalo $[s,t]$, la curva $\gamma([s,t])$ es una geodésica entre $\gamma(s)$ y $\gamma(t).$
    \end{definition}

    Al igual que en el caso euclideo las geodésicas tienen una forma particular en el caso hiperbólico, mas son diferentes

    \begin{theorem}
        Las geodésicas infinitas en $\bb{H}^2$ consiste de lineas verticales cuando dos puntos comparten parte real y de semicírculos con centro en la recta real para el otro caso.
    \end{theorem}

    Algo particular a notar es que estas geodésicas se intersectan con $\partial_\infty\bb{H}^2$ en ángulos rectos.\\

    \begin{definition}
        Una isometria $\phi:\bb{H}^2\to\bb{H}^2$ es un difeomorfismo tal que preserva la métrica de $\bb{H}^2$
    \end{definition}

    Claramente estas dos ultimas definiciones se pueden dar en contextos mas generales pero nos quedaremos en nuestro caso hiperbólico, dicho esto en particular nos interesan la isometrias que preservan la orientación, si juntamos todas estas isometrias estas forman un grupo con una descripción particular.
    \begin{theorem}
        El grupo de isometrias de $\bb{H}^2$ es generado por reflexiones a través geodésicas, en particular si miramos aquellas que preservan la orientación es el grupo dado por las siguientes transformaciones
        $$z\mapsto\frac{az+b}{cz+d},$$
        donde $a,b,c,d\in\bb{R}$ y $ad-bc>0.$
    \end{theorem}
    Estas transformaciones se conocen como transformaciones lineales fraccionarias. Note en particular si cada coeficiente lo normalizamos por $\sqrt{ad-bc}$ (que no afecta la transformación) podemos entender estas transformaciones como la acción del grupo $PSL(2,\bb{R})$ sobre $\bb{H}^2$ donde la acción esta dada por
    $$\begin{pmatrix}
        a&b\\
        c&d
    \end{pmatrix}\cdot z=\frac{az+b}{cz+d}.$$
    Este tipo de transformaciones son particularmente útiles ya que permiten escoger a donde llevamos un punto particular, esto se ve en el siguiente lema que es bastante útil.
    \begin{lemma}
        Dados tres puntos diferentes $z_1,z_2z_3\in\partial_\infty\bb{H}^2$, existe una isometria que preserva la orientación tal que tal que 
        $$z_i\mapsto\begin{cases}
            1, & i=1\\
            0, & i=2\\
            \infty & i=3
        \end{cases}$$
    \end{lemma}
    \begin{proof}
       El detalle mas importante que hay que cuidar es el tema de la orientación (mantener el signo del determinante positivo), si ningún punto es $\infty$ observe que son tres puntos en la recta real, así considere
       $$z\mapsto \frac{z_1-z_3}{z_1-z_2}\frac{z-z_2}{z-z_3}$$ 
       Note que 
       $$z_1\mapsto \frac{z_1-z_3}{z_1-z_2}\frac{z_1-z_2}{z_1-z_3}=1,\quad z_2\mapsto \frac{z_1-z_3}{z_1-z_2}\frac{z_2-z_2}{z_2-z_3}=0,\quad z_3\mapsto \frac{z_1-z_3}{z_1-z_2}\frac{z_3-z_2}{z_3-z_3}=\infty.$$
       Ahora note que el producto de coeficientes es $$(z_1-z_3)(-z_3)(z_1-z_2)-(-z_2)(z_1-z_3)(z_1-z_2)=(z_1-z_3)(z_1-z_2)(z_2-z_3).$$
       Observe que el signo depende completamente de el orden en el que uno etiquete los puntos, en particular si uno los etiqueta en un sentido antihorario partiendo desde $z_1$ claramente es positivo.\\

       En el caso que $z_i=\infty$ para alguno de los $i$, la isometria es la siguiente para cada caso
       $$z\mapsto\begin{cases}
           \dfrac{z-z_2}{z-z_3},&i=1\\
           \\
           \dfrac{z_1-z_3}{z-z_3},&i=2\\
           \\
           \dfrac{z-z_2}{z_1-z_2},&i=3
       \end{cases}$$
       Este ultimo es fácil verificar nuevamente que enviá los puntos a lo enunciado y el determinante resulta siempre positivo por la manera en que organizamos los puntos.

    \end{proof}
    Estas isometrias dan mucha facilidad no solo para cálculos sino también para algunos hechos de la geometría del espacio.
    \begin{corollary}
        Siempre existe una isometria tal que dados tres puntos diferentes $z_1,z_2,z_3$ en la frontera al infinito pueden ser enviados a otros tres diferentes $z_1^\prime,z_2^\prime,z_3^\prime$ también en la frontera al infinito.
    \end{corollary}
    \begin{corollary}
        Dadas dos geodésicas $l_1,l_2$ en $\bb{H}^2$ solo hay tres posibilidades
        \begin{itemize}
            \item $l_1\cap l_2=\{x\}\subset \bb{H}^2.$
            \item $l_1\cap l_2=\{x\}\subset \partial_\infty\bb{H}^2.$
            \item $ l_1\cap l_2=\varnothing.$ 
        \end{itemize}
    \end{corollary}
    \begin{eg}
       Previamente habíamos hallado la distancia entre dos puntos con la misma parte real (esto ya que la recta vertical si es una geodésica), pero por ejemplo para por ejemplo el punto $(0,h)$ y $(1,h)$ la longitud de ese segmento es claro que no es la mas corta, la mas corta es la que esta dada por el segmento del semicírculo que contiene a estos dos puntos, este circulo es 
       $\left(x-\dfrac{1}{2}\right)^2+y^2=h^2+\dfrac{1}{4}.$
       Este circulo interseca a la frontera en 
       $x=\dfrac{1}{2}\pm\sqrt{h^2+\dfrac{1}{4}},$ 
       esto es haciendo $y=0$, luego hay dos estrategias para calcular la distancia, una seria hallar la integral paramétrica, la otra es usar la isometria hallada previamente para convertir esta geodésica en la recta vertical de $0$ a $\infty$, en el lenguaje de la prueba si tomamos
       $z_1=\infty,\quad z_2=\dfrac{1}{2}+\sqrt{h^2+\dfrac{1}{4}},\quad z_3=\dfrac{1}{2}-\sqrt{h^2+\dfrac{1}{4}},$
       la transformación 
       $z\mapsto \dfrac{z-z_2}{z-z_3},$ hace justamente lo que queremos, Luego la imagen de $ih$ y $1+ih$ es respectivamente
       $a=\dfrac{ih-z_2}{ih-z_3},\quad b=\dfrac{1+ih-z_2}{1+ih-z_3},$
      Así por el calculo hecho previamente sabemos que la distancia entre estos dos puntos es $\pm\log\left(\dfrac{\Im(b)}{\Im(a)}\right)$, el $\pm$ esta ya que depende de que parte imaginaria sea mayor, y como las isometrias preservan distancias tenemos la distancia explicita entre los puntos originales. 
    \end{eg}

    \begin{figure}[H]
    \centering
    
%!TEX root = ./main.tex



\tikzset{every picture/.style={line width=0.75pt}} %set default line width to 0.75pt        

\begin{tikzpicture}[x=0.75pt,y=0.75pt,yscale=-1,xscale=1]
%uncomment if require: \path (0,877); %set diagram left start at 0, and has height of 877

%Shape: Arc [id:dp3323452367299544] 
\draw  [draw opacity=0] (247.5,190) .. controls (247.5,158.24) and (273.24,132.5) .. (305,132.5) .. controls (336.76,132.5) and (362.5,158.24) .. (362.5,190) -- (305,190) -- cycle ; \draw   (247.5,190) .. controls (247.5,158.24) and (273.24,132.5) .. (305,132.5) .. controls (336.76,132.5) and (362.5,158.24) .. (362.5,190) ;  
%Straight Lines [id:da5286576515069806] 
\draw    (133,190) -- (477,190) ;
\draw [shift={(480,190)}, rotate = 180] [fill={rgb, 255:red, 0; green, 0; blue, 0 }  ][line width=0.08]  [draw opacity=0] (8.93,-4.29) -- (0,0) -- (8.93,4.29) -- cycle    ;
\draw [shift={(130,190)}, rotate = 0] [fill={rgb, 255:red, 0; green, 0; blue, 0 }  ][line width=0.08]  [draw opacity=0] (8.93,-4.29) -- (0,0) -- (8.93,4.29) -- cycle    ;
%Straight Lines [id:da6944501524708641] 
\draw    (280,53) -- (280,190) ;
\draw [shift={(280,50)}, rotate = 90] [fill={rgb, 255:red, 0; green, 0; blue, 0 }  ][line width=0.08]  [draw opacity=0] (8.93,-4.29) -- (0,0) -- (8.93,4.29) -- cycle    ;
%Straight Lines [id:da6550119981386697] 
\draw    (130,140) -- (460,140) ;
%Shape: Arc [id:dp7496478031646566] 
\draw  [draw opacity=0][line width=2.25]  (279.69,138.36) .. controls (287.33,134.61) and (295.92,132.5) .. (305,132.5) .. controls (314.52,132.5) and (323.49,134.81) .. (331.4,138.9) -- (305,190) -- cycle ; \draw  [color={rgb, 255:red, 74; green, 144; blue, 226 }  ,draw opacity=1 ][line width=2.25]  (279.69,138.36) .. controls (287.33,134.61) and (295.92,132.5) .. (305,132.5) .. controls (314.52,132.5) and (323.49,134.81) .. (331.4,138.9) ;  

% Text Node
\draw (352,190.4) node [anchor=north west][inner sep=0.75pt]    {$z_{2}$};
% Text Node
\draw (242,190.4) node [anchor=north west][inner sep=0.75pt]    {$z_{3}$};
% Text Node
\draw (261,122.4) node [anchor=north west][inner sep=0.75pt]    {$ih$};
% Text Node
\draw (337,122.4) node [anchor=north west][inner sep=0.75pt]    {$1+ih$};


\end{tikzpicture}
        \caption{Geodésica que contiene a los puntos}
    \end{figure}

    Con estas herramientas se ve un camino mejor hacia los cálculos, pero recordemos hacia donde vamos, la idea es poder calcular de alguna forma sobre las descomposiciones de los complementos en $S^3$ de nudos/enlaces, estos complementos generaran regiones y estructuras particulares en $S^3$ por lo que necesitamos subir de dimensión e intentar ver áreas y volúmenes, las siguientes definiciones empiezan a ir en esa dirección.

    \begin{definition}
        Un triangulo ideal en $\bb{H}^2$ es un triangulo donde sus tres aristas son geodésicas y su vértices se encuentran en $\partial_\infty\bb{H}^2.$
    \end{definition}

        \begin{figure}[H]
        \centering
        
%!TEX root = ./main.tex



\tikzset{every picture/.style={line width=0.75pt}} %set default line width to 0.75pt        

\begin{tikzpicture}[x=0.75pt,y=0.75pt,yscale=-1,xscale=1]
%uncomment if require: \path (0,877); %set diagram left start at 0, and has height of 877

%Shape: Arc [id:dp3323452367299544] 
\draw  [draw opacity=0][line width=1.5]  (140,200) .. controls (140,200) and (140,200) .. (140,200) .. controls (140,168.24) and (165.74,142.5) .. (197.5,142.5) .. controls (229.26,142.5) and (255,168.24) .. (255,200) -- (197.5,200) -- cycle ; \draw  [color={rgb, 255:red, 74; green, 144; blue, 226 }  ,draw opacity=1 ][line width=1.5]  (140,200) .. controls (140,200) and (140,200) .. (140,200) .. controls (140,168.24) and (165.74,142.5) .. (197.5,142.5) .. controls (229.26,142.5) and (255,168.24) .. (255,200) ;  
%Straight Lines [id:da5286576515069806] 
\draw    (83,200) -- (297,200) ;
\draw [shift={(300,200)}, rotate = 180] [fill={rgb, 255:red, 0; green, 0; blue, 0 }  ][line width=0.08]  [draw opacity=0] (8.93,-4.29) -- (0,0) -- (8.93,4.29) -- cycle    ;
\draw [shift={(80,200)}, rotate = 0] [fill={rgb, 255:red, 0; green, 0; blue, 0 }  ][line width=0.08]  [draw opacity=0] (8.93,-4.29) -- (0,0) -- (8.93,4.29) -- cycle    ;
%Straight Lines [id:da6944501524708641] 
\draw [color={rgb, 255:red, 74; green, 144; blue, 226 }  ,draw opacity=1 ][line width=1.5]    (140,30) -- (140,200) ;
%Straight Lines [id:da7691230862880515] 
\draw [color={rgb, 255:red, 74; green, 144; blue, 226 }  ,draw opacity=1 ][line width=1.5]    (255,30) -- (255,200) ;
%Shape: Arc [id:dp4258690964395112] 
\draw  [draw opacity=0][line width=1.5]  (400,202.5) .. controls (400,202.5) and (400,202.5) .. (400,202.5) .. controls (400,188.69) and (413.43,177.5) .. (430,177.5) .. controls (446.57,177.5) and (460,188.69) .. (460,202.5) -- (430,202.5) -- cycle ; \draw  [color={rgb, 255:red, 74; green, 144; blue, 226 }  ,draw opacity=1 ][line width=1.5]  (400,202.5) .. controls (400,202.5) and (400,202.5) .. (400,202.5) .. controls (400,188.69) and (413.43,177.5) .. (430,177.5) .. controls (446.57,177.5) and (460,188.69) .. (460,202.5) ;  
%Straight Lines [id:da8977320996545466] 
\draw    (353,202.5) -- (567,202.5) ;
\draw [shift={(570,202.5)}, rotate = 180] [fill={rgb, 255:red, 0; green, 0; blue, 0 }  ][line width=0.08]  [draw opacity=0] (8.93,-4.29) -- (0,0) -- (8.93,4.29) -- cycle    ;
\draw [shift={(350,202.5)}, rotate = 0] [fill={rgb, 255:red, 0; green, 0; blue, 0 }  ][line width=0.08]  [draw opacity=0] (8.93,-4.29) -- (0,0) -- (8.93,4.29) -- cycle    ;
%Shape: Arc [id:dp38871258976776124] 
\draw  [draw opacity=0][line width=1.5]  (460,202.5) .. controls (460,202.5) and (460,202.5) .. (460,202.5) .. controls (460,188.69) and (473.43,177.5) .. (490,177.5) .. controls (506.57,177.5) and (520,188.69) .. (520,202.5) -- (490,202.5) -- cycle ; \draw  [color={rgb, 255:red, 74; green, 144; blue, 226 }  ,draw opacity=1 ][line width=1.5]  (460,202.5) .. controls (460,202.5) and (460,202.5) .. (460,202.5) .. controls (460,188.69) and (473.43,177.5) .. (490,177.5) .. controls (506.57,177.5) and (520,188.69) .. (520,202.5) ;  
%Shape: Arc [id:dp5339401693658827] 
\draw  [draw opacity=0][line width=1.5]  (400,200) .. controls (400,166.86) and (426.86,140) .. (460,140) .. controls (493.14,140) and (520,166.86) .. (520,200) -- (460,200) -- cycle ; \draw  [color={rgb, 255:red, 74; green, 144; blue, 226 }  ,draw opacity=1 ][line width=1.5]  (400,200) .. controls (400,166.86) and (426.86,140) .. (460,140) .. controls (493.14,140) and (520,166.86) .. (520,200) ;  
%Straight Lines [id:da11869494313797047] 
\draw [color={rgb, 255:red, 0; green, 0; blue, 0 }  ,draw opacity=1 ][line width=0.75]    (380.67,34.33) -- (380.67,201.33) ;
\draw [shift={(380.67,31.33)}, rotate = 90] [fill={rgb, 255:red, 0; green, 0; blue, 0 }  ,fill opacity=1 ][line width=0.08]  [draw opacity=0] (8.93,-4.29) -- (0,0) -- (8.93,4.29) -- cycle    ;

% Text Node
\draw (134,202.4) node [anchor=north west][inner sep=0.75pt]    {$0$};
% Text Node
\draw (251,202.4) node [anchor=north west][inner sep=0.75pt]    {$1$};


\end{tikzpicture}
            \caption{Algunos triángulos ideales}
        \end{figure}

    \begin{note}
        Una observación importante es que todos los triángulos ideales pueden ser llevados a el primer triangulo ideal del ejemplo anterior por una isometria y por tanto todos los triángulos ideales son isometricos entre si en $\bb{H}^2$.
    \end{note}

    \begin{definition}
        Un horociclo centrado en un punto $p\in\partial_\infty\bb{H}^2$ es definido como una curva que es perpendicular a todas las geodesicas que pasan a traves de $p$, en particular cuando $p$ esta sobre la recta real,los horociclos son círculos tangentes a $p$, en cambio cuando $p=\infty$ los horociclos son rectas paralelas al eje real. La región interior al horociclo se conoce como horobola, en el primer caso es la región al interior del circulo
    \end{definition}

    Este lenguaje nos permite hablar con precisión de algunas regiones particulares, ademas de que nos permite probar el hecho de que los triángulos ideales tienen área finita, aunque aquí seremos un poco mas precisos y usaremos calculo para mostrar esa área.

    \begin{lemma}
        El área de un triangulo ideal es $\pi$
    \end{lemma}

    \begin{proof}
        Dado un triangulo ideal en $\bb{H}^2$ ya sabemos que por medio de una isometria se puede llevar al triangulo ideal con vértices $0,1,\infty.$ Llamemoslos $T$, esta región se puede parametrizar de la siguiente forma 
        $$T=\left\{(x,y)\in\bb{H}^2: y\geq \sqrt{\frac{1}{4}-\left(x-\frac{1}{2}\right)^2},\,0\leq x\leq 1\right\}.$$
        Así por un ejercicio de calculo elemental tenemos
        \begin{align*}
            \text{area}(T)&=\int_0^1\int_{\sqrt{\frac{1}{4}-\left(x-\frac{1}{2}\right)^2}}^\infty\frac{1}{y^2}\,dydx\\
            &=\int_0^1\frac{1}{\sqrt{\frac{1}{4}-\left(x-\frac{1}{2}\right)^2}}\,dx\\
            &=\int_{-1}^1\frac{1}{\sqrt{1-u^2}}\,du\\
            &=\int_0^\pi\,d\theta=\pi
        \end{align*}
        Esto con la sustituciones $u=2x-1$ y $u=\cos \theta$. Como las isometrias mantienen el área llegamos a la conclusión de que todos tienen la misma área y es $\pi$.
    
    \end{proof}

    Ahora naturalmente queremos subir una dimensión ya que $S^3$ de base es una 3-variedad, por lo que para terminar esta sección definiremos el espacio hiperbólico 3-dimensional y se enunciaran algunos hechos relevantes análogos a algunos previos de el caso 2 dimensional.

    \begin{definition}
        El espacio hiperbolico 3-dimensional es
        $$\bb{H}^3:=\{(x+iy,t)\in\bb{C}\times \bb{R}:t>0\}.$$
        donde la metrica esta dada por
        $$ds^2=\frac{dx^2+dy^2+dt^2}{t^2}$$
    \end{definition}

    \begin{theorem}
    Las geodesicas en $\bb{H}^3$ son lineas verticales y semicirculos ortogonales a $\partial\bb{H}^3:=\bb{C}\cup\{\infty\}$. Los planos totalmente geodesicos son de manera similar planos verticales y las partes superiores de esferas centradas en algún punto de frontera.
    \end{theorem}
    \begin{theorem}
        El grupo de isometrias de $\bb{H}^3$ es generado por reflexiones a través de planos geodésicos, mientras que aquellas que preservan orientación es $PSL(2,\bb{C})$. Particularmente este grupo actuá sobre la frontera $\partial\bb{H}^3$ a través de la acción usual de transformaciones lineales fraccionarias, esta acción se extiende al interior.
    \end{theorem}
    \begin{definition}
        Un tetraedro ideal, es un tetraedro en $\bb{H}^3$ donde sus cuatro vértices se encuentran en $\partial\bb{H}^3$ y sus caras y aristas son geodésicas
    \end{definition}
  \begin{figure}[H]
        \centering
        
%!TEX root = ./main.tex

\tikzset{every picture/.style={line width=0.75pt}} %set default line width to 0.75pt        

\begin{tikzpicture}[x=0.75pt,y=0.75pt,yscale=-1,xscale=1]
%uncomment if require: \path (0,877); %set diagram left start at 0, and has height of 877

%Straight Lines [id:da5480235268340935] 
\draw    (200,200) -- (200,40) ;
%Shape: Arc [id:dp5722566538107826] 
\draw  [draw opacity=0] (200,200) .. controls (200,183.43) and (213.43,170) .. (230,170) .. controls (246.57,170) and (260,183.43) .. (260,200) -- (230,200) -- cycle ; \draw   (200,200) .. controls (200,183.43) and (213.43,170) .. (230,170) .. controls (246.57,170) and (260,183.43) .. (260,200) ;  
%Straight Lines [id:da05435319973401176] 
\draw    (260,200) -- (260,40) ;
%Shape: Arc [id:dp7827074833336561] 
\draw  [draw opacity=0] (260,200) .. controls (260,200) and (260,200) .. (260,200) .. controls (260,175.15) and (273.43,155) .. (290,155) .. controls (297.78,155) and (304.86,159.44) .. (310.19,166.72) -- (290,200) -- cycle ; \draw   (260,200) .. controls (260,200) and (260,200) .. (260,200) .. controls (260,175.15) and (273.43,155) .. (290,155) .. controls (297.78,155) and (304.86,159.44) .. (310.19,166.72) ;  
%Straight Lines [id:da6791216861853826] 
\draw    (310.19,166.72) -- (310,40) ;
%Shape: Arc [id:dp006836970713729906] 
\draw  [draw opacity=0][dash pattern={on 4.5pt off 4.5pt}] (200,187.5) .. controls (205.52,155.74) and (234.62,130) .. (265,130) .. controls (287.67,130) and (304.65,144.34) .. (309.5,164.82) -- (255,187.5) -- cycle ; \draw  [dash pattern={on 4.5pt off 4.5pt}] (200,187.5) .. controls (205.52,155.74) and (234.62,130) .. (265,130) .. controls (287.67,130) and (304.65,144.34) .. (309.5,164.82) ;  
%Straight Lines [id:da2490317153466941] 
\draw    (200,100) -- (260,100) ;
%Straight Lines [id:da6214491647211758] 
\draw    (260,100) -- (310,70) ;
%Straight Lines [id:da2566992937994781] 
\draw  [dash pattern={on 4.5pt off 4.5pt}]  (200,100) -- (310,70) ;
%Straight Lines [id:da19794337864554612] 
\draw    (120,200) -- (360,200) ;
%Straight Lines [id:da16808622174577048] 
\draw    (160,150) -- (120,200) ;
%Straight Lines [id:da6537649860646593] 
\draw    (160,150) -- (200,150) ;
%Straight Lines [id:da9885141829311728] 
\draw    (310,150) -- (400,150) ;
%Straight Lines [id:da010456063455708042] 
\draw    (400,150) -- (360,200) ;

% Text Node
\draw (197,202.4) node [anchor=north west][inner sep=0.75pt]    {$0$};
% Text Node
\draw (257,202.4) node [anchor=north west][inner sep=0.75pt]    {$1$};
% Text Node
\draw (311,162.4) node [anchor=north west][inner sep=0.75pt]    {$z$};


\end{tikzpicture}
            \caption{Un tetraedro ideal}
        \end{figure}
    \begin{note}
Es un poco difícil dibujar un tetraedro ideal donde ningun punto es el infinito, pero observe que por transformaciones al igual que en el caso de $\bb{H}^2$ podemos llevar cualesquiera tres puntos a $0,1,\infty$, por lo que un tetraedro ideal sera la figura del ejemplo anterior donde el cuarto punto es algún $z\in\bb{C}\setminus\{0,1\}$. 
\end{note}
\begin{definition}
    Una horoesfera alrededor de un punto ideal $p\in\partial\bb{H}^3$ es una superficie perpendicular a todas las geodésicas que pasan por ese punto, similar a antes serán esferas tangentes a el plano $\bb{C}$ o planos paralelos a este mismo.
\end{definition}







    \section{Descomposición del complemento de un nudo/enlace}
    En el mismo espíritu de la bibliografía presentaremos la construcción principal por medio del ejemplo estrella, para posteriormente verla en otro casos. El propósito es presentar una construcción geométrica que nos permita de manera local, entender como se descompone la 3-variedad $S^3\setminus K$ en bloques 3-dimensionales mucho mas sencillos, para así poder pasar entre diagramas de nudos/enlaces a complementos de nudos/enlaces.
    \subsection{Descomposición del nudo de 8}
    %!TEX root = ./main.tex



\begin{definition}
    Un poliedro es una 3-bola cerrada, tal que en su frontera se encuentre marcada con un grafo finito, tal que la regiones complementarias llamadas caras son simplemente conexas.
\end{definition}
Prácticamente cualquier dibujo clásico que se puede hacer (de manera natural sin estar buscando no ejemplos) resulta ser un ejemplo de poliedro
\begin{figure}[H]
    \centering
    
%!TEX root = ./main.tex





\tikzset{every picture/.style={line width=0.75pt}} %set default line width to 0.75pt        

\begin{tikzpicture}[x=0.75pt,y=0.75pt,yscale=-0.8,xscale=0.8]
%uncomment if require: \path (0,877); %set diagram left start at 0, and has height of 877

%Shape: Ellipse [id:dp0928071191771559] 
\draw  [fill={rgb, 255:red, 155; green, 155; blue, 155 }  ,fill opacity=0.27 ] (220,122.5) .. controls (220,65.89) and (262.53,20) .. (315,20) .. controls (367.47,20) and (410,65.89) .. (410,122.5) .. controls (410,179.11) and (367.47,225) .. (315,225) .. controls (262.53,225) and (220,179.11) .. (220,122.5) -- cycle ;
%Curve Lines [id:da13418485162745786] 
\draw    (220,122.5) .. controls (285.87,142.09) and (348.36,143.91) .. (410,115.67) ;
%Curve Lines [id:da5341258541218854] 
\draw  [dash pattern={on 4.5pt off 4.5pt}]  (220,122.5) .. controls (284.18,100.18) and (350.04,103.82) .. (410,115.67) ;
%Straight Lines [id:da3427735018941048] 
\draw    (280,90) -- (310,70) ;
%Straight Lines [id:da6109826278772268] 
\draw    (310,70) -- (340,90) ;
%Straight Lines [id:da31209815320417755] 
\draw    (310,30) -- (310,70) ;
%Curve Lines [id:da4696149880079775] 
\draw    (280,90) .. controls (279.67,65) and (284.33,45) .. (310,30) ;
%Curve Lines [id:da0013412730978942244] 
\draw    (310,30) .. controls (335,41) and (340.33,60.33) .. (340,90) ;
%Curve Lines [id:da5680668023106804] 
\draw    (280,90) .. controls (305.67,96.33) and (302.33,99.67) .. (340,90) ;




\end{tikzpicture}
    \caption{Ejemplo de poliedro}
\end{figure}
La condición final esta puesta para asegurar que las caras si se comporten como discos poligonales para tener un análogo de la imagen usual que uno suele tener de poliedro, mas si no la tuviéramos podemos hacer uno de esos dibujos que mencionábamos antes.
\begin{figure}[H]
    \centering
    
%!TEX root = ./main.tex



\tikzset{every picture/.style={line width=0.75pt}} %set default line width to 0.75pt        

\begin{tikzpicture}[x=0.75pt,y=0.75pt,yscale=-0.8,xscale=0.8]
%uncomment if require: \path (0,877); %set diagram left start at 0, and has height of 877

%Shape: Ellipse [id:dp0928071191771559] 
\draw  [fill={rgb, 255:red, 155; green, 155; blue, 155 }  ,fill opacity=0.27 ] (230,122.5) .. controls (230,65.89) and (272.53,20) .. (325,20) .. controls (377.47,20) and (420,65.89) .. (420,122.5) .. controls (420,179.11) and (377.47,225) .. (325,225) .. controls (272.53,225) and (230,179.11) .. (230,122.5) -- cycle ;
%Curve Lines [id:da13418485162745786] 
\draw    (230,122.5) .. controls (295.87,142.09) and (358.36,143.91) .. (420,115.67) ;
%Curve Lines [id:da5341258541218854] 
\draw  [dash pattern={on 4.5pt off 4.5pt}]  (230,122.5) .. controls (294.18,100.18) and (360.04,103.82) .. (420,115.67) ;
%Shape: Ellipse [id:dp9147901422013078] 
\draw   (285,45) .. controls (285,36.72) and (300.67,30) .. (320,30) .. controls (339.33,30) and (355,36.72) .. (355,45) .. controls (355,53.28) and (339.33,60) .. (320,60) .. controls (300.67,60) and (285,53.28) .. (285,45) -- cycle ;
%Shape: Ellipse [id:dp5524598870994889] 
\draw   (285,200) .. controls (285,191.72) and (300.67,185) .. (320,185) .. controls (339.33,185) and (355,191.72) .. (355,200) .. controls (355,208.28) and (339.33,215) .. (320,215) .. controls (300.67,215) and (285,208.28) .. (285,200) -- cycle ;




\end{tikzpicture}
    \caption{Grafo que genera una región que no es simplemente conexa}
\end{figure}
Como nuestro propósito es estudiar los complementos de los nudos en $S^3$ si lo pensamos bien la información que nos da el nudo es aquello que falta, es decir son huecos en la descomposición, esto motiva la siguiente definición. 
\begin{definition}
    Un poliedro ideal es un poliedro en el cual removemos todos los puntos que antes habíamos marcado como vértices.
\end{definition}

Nuestra idea sera cortar $S^3\setminus K$ en dos poliedros ideales, así tendríamos  una descripción del complemento como el pegado de dos estructuras mucho mas sencillas de comprender (en particular de visualizar).\\

Una primera intuición geométrica que nos da la construcción es la siguiente, Consideremos el siguiente diagrama del nudo de 8
\begin{figure}[H]
    \centering
    
%!TEX root = ./main.tex



\tikzset{every picture/.style={line width=0.75pt}} %set default line width to 0.75pt        

\begin{tikzpicture}[x=0.75pt,y=0.75pt,yscale=-0.85,xscale=0.85]
%uncomment if require: \path (0,877); %set diagram left start at 0, and has height of 877

%Shape: Free Drawing [id:dp0005319926756554016] 
\draw  [line width=2.25] [line join = round][line cap = round] (298.71,113.07) .. controls (290.26,111.76) and (290.37,96.06) .. (287.28,87.1) .. controls (278.97,62.96) and (283.59,34.1) .. (300.2,17.86) .. controls (303.92,14.22) and (310.73,9.55) .. (316.2,10.03) .. controls (409.49,18.29) and (466.49,178.81) .. (402.12,247.44) .. controls (394.25,255.83) and (381.75,257.34) .. (371.76,249.89) .. controls (363.9,244.03) and (356.41,229.52) .. (347.94,232.75) ;
%Shape: Free Drawing [id:dp4987515250733162] 
\draw  [line width=2.25] [line join = round][line cap = round] (332.24,229.4) .. controls (321.28,221.4) and (294.93,233.34) .. (285.13,238.66) .. controls (282.8,239.92) and (281.19,242.57) .. (278.71,243.48) .. controls (249.19,254.28) and (225.41,237.38) .. (209.61,206.31) .. controls (203.85,194.99) and (199.42,182.36) .. (198.44,169.63) .. controls (195.36,129.43) and (201.51,68.2) .. (227.52,37.92) .. controls (237.95,25.79) and (256.3,15.25) .. (272.26,14.4) .. controls (278.18,14.08) and (290.4,21.7) .. (287.75,21.28) ;
%Shape: Free Drawing [id:dp6231462806336996] 
\draw  [line width=2.25] [line join = round][line cap = round] (300.29,29.52) .. controls (313.33,39.03) and (314.36,55.83) .. (316.87,71.08) .. controls (318.17,78.94) and (320.99,87.94) .. (319.38,96.11) .. controls (317.59,105.21) and (311.25,114.25) .. (307.55,122.05) .. controls (298.18,141.77) and (288.6,158.03) .. (282.65,180.47) .. controls (280.81,187.4) and (278.84,202.37) .. (279.89,211.4) .. controls (280.38,215.6) and (281.27,219.79) .. (282.38,223.89) .. controls (283.01,226.24) and (287.39,230.94) .. (285.22,230.6) ;
%Shape: Free Drawing [id:dp34435703317452426] 
\draw  [line width=2.25] [line join = round][line cap = round] (293.48,242.2) .. controls (297.43,245.08) and (296.48,249.26) .. (302.16,253.4) .. controls (305.44,255.79) and (309.19,257.21) .. (312.85,258.65) .. controls (338.56,268.73) and (345.04,219.7) .. (343.78,201.62) .. controls (342.05,176.65) and (340.06,161.14) .. (328.08,139.58) .. controls (326.46,136.67) and (320.37,128.91) .. (319.8,128.44) .. controls (316.42,125.63) and (309.39,124.58) .. (312.66,125.09) ;
%Shape: Free Drawing [id:dp9595901252421886] 
\draw  [color={rgb, 255:red, 255; green, 255; blue, 255 }  ,draw opacity=1 ][line width=2.25] [line join = round][line cap = round] (295.15,117.89) .. controls (294.9,117.98) and (294.33,118.06) .. (294.34,117.76) ;
%Shape: Free Drawing [id:dp1969537050250879] 
\draw  [draw opacity=0][line width=2.25] [line join = round][line cap = round] (426.53,66.18) .. controls (426.65,66.26) and (427.07,66.26) .. (426.94,66.24) ;
%Shape: Free Drawing [id:dp026195381751755398] 
\draw  [draw opacity=0][line width=2.25] [line join = round][line cap = round] (300.82,112.05) .. controls (298.16,113.06) and (297.03,115.75) .. (296.34,118.52) ;
%Shape: Free Drawing [id:dp45027773036340213] 
\draw  [draw opacity=0][line width=2.25] [line join = round][line cap = round] (299.5,113.63) .. controls (292.93,121.91) and (291.72,122.28) .. (294.04,122.64) ;
%Shape: Free Drawing [id:dp42616248151326175] 
\draw  [draw opacity=0][line width=2.25] [line join = round][line cap = round] (299.96,112.81) .. controls (297.93,115.24) and (295.36,117.35) .. (294.18,120.42) ;
%Shape: Free Drawing [id:dp4649779193663184] 
\draw  [draw opacity=0][line width=2.25] [line join = round][line cap = round] (288.84,231.61) .. controls (287.29,232.56) and (285.87,233.82) .. (284.18,234.47) .. controls (283.73,234.64) and (284.58,233.19) .. (285.07,233.26) ;
%Shape: Free Drawing [id:dp48154694518391594] 
\draw  [draw opacity=0][line width=2.25] [line join = round][line cap = round] (281.37,233.59) .. controls (284.36,231.84) and (286.8,227.71) .. (290.28,228.25) ;
%Shape: Free Drawing [id:dp8989735988593472] 
\draw  [draw opacity=0][line width=2.25] [line join = round][line cap = round] (285.58,231.55) .. controls (287.76,231.89) and (289.57,229.92) .. (291.46,228.88) ;
%Shape: Free Drawing [id:dp8922120977301046] 
\draw  [draw opacity=0][line width=2.25] [line join = round][line cap = round] (284.02,230.41) .. controls (285.23,230.16) and (290.6,229.69) .. (290.74,227.42) ;
%Shape: Free Drawing [id:dp5528703258533438] 
\draw  [draw opacity=0][line width=2.25] [line join = round][line cap = round] (284.37,231.36) .. controls (286.58,230.33) and (292.29,228.72) .. (292.52,225.01) ;
%Shape: Free Drawing [id:dp9013988558965794] 
\draw  [draw opacity=0][line width=2.25] [line join = round][line cap = round] (284.31,232.25) .. controls (287.24,229.93) and (290.03,227.39) .. (292.55,224.57) ;
%Shape: Free Drawing [id:dp8372444458576085] 
\draw  [draw opacity=0][line width=2.25] [line join = round][line cap = round] (283.51,232.13) .. controls (288.08,228.65) and (290.79,224.57) .. (293.69,219.37) ;
%Shape: Free Drawing [id:dp8237859994697239] 
\draw  [draw opacity=0][line width=2.25] [line join = round][line cap = round] (191.28,50.96) .. controls (190.75,50.66) and (190.28,50.16) .. (190,49.59) ;




\end{tikzpicture}
    \caption{Proyección del nudo de 8}
\end{figure}
La idea subyacente sera considerar el diagrama en un plano proyección con cierto volumen, para que de esta forma si bien sea prácticamente una proyección en términos de visualización en realidad sea el mismo nudo original embebido en $S^2\times[-\varepsilon,\varepsilon]$, para $\varepsilon>0$, tal que todo eso este contenido en $S^3$, una ilustración que describe eso es la siguiente
\begin{figure}[H]
    \centering
    
%!TEX root = ./main.tex



\tikzset{every picture/.style={line width=0.75pt}} %set default line width to 0.75pt        

\begin{tikzpicture}[x=0.75pt,y=0.75pt,yscale=-0.8,xscale=0.8]
%uncomment if require: \path (0,877); %set diagram left start at 0, and has height of 877

%Shape: Free Drawing [id:dp7347232652227033] 
\draw  [line width=2.25] [line join = round][line cap = round] (340.64,155.74) .. controls (336.75,153.74) and (342.11,145.88) .. (343.56,140.89) .. controls (347.45,127.45) and (359.56,113.7) .. (373.56,108.19) .. controls (376.69,106.96) and (381.76,105.7) .. (384.4,106.81) .. controls (429.43,125.8) and (404.44,215.43) .. (348.27,239.63) .. controls (341.4,242.59) and (334.48,241.36) .. (331.88,236.03) .. controls (329.83,231.84) and (330.89,223.37) .. (325.46,223.64) ;
%Shape: Free Drawing [id:dp9263733489395964] 
\draw  [line width=2.25] [line join = round][line cap = round] (318.53,219.46) .. controls (315.61,213.7) and (298.07,215.5) .. (291.25,216.61) .. controls (289.63,216.88) and (287.91,217.95) .. (286.33,218.01) .. controls (267.56,218.73) and (261.07,206.47) .. (263.47,188.36) .. controls (264.34,181.76) and (266.33,174.72) .. (270.13,168.17) .. controls (282.13,147.51) and (305.96,117.75) .. (329.52,106.7) .. controls (338.96,102.27) and (351.92,99.9) .. (360.39,102.01) .. controls (363.53,102.79) and (367.22,108.56) .. (366,107.93) ;
%Shape: Free Drawing [id:dp06457306751650549] 
\draw  [line width=2.25] [line join = round][line cap = round] (369.66,114.06) .. controls (373.13,120.9) and (367.98,129.5) .. (364.12,137.55) .. controls (362.13,141.7) and (360.53,146.67) .. (356.95,150.51) .. controls (352.96,154.8) and (346.66,158.33) .. (342.12,161.65) .. controls (330.67,170.06) and (320.27,176.7) .. (309.64,187.02) .. controls (306.36,190.2) and (300.29,197.41) .. (297.78,202.11) .. controls (296.61,204.29) and (295.66,206.53) .. (294.84,208.77) .. controls (294.37,210.05) and (295.02,213.1) .. (294.03,212.59) ;
%Shape: Free Drawing [id:dp21552600534807442] 
\draw  [line width=2.25] [line join = round][line cap = round] (294.34,219.72) .. controls (295.39,221.79) and (293.49,223.74) .. (295.01,226.72) .. controls (295.88,228.44) and (297.32,229.75) .. (298.71,231.06) .. controls (308.48,240.2) and (328.36,216.63) .. (333.82,207.35) .. controls (341.37,194.54) and (345.59,186.45) .. (346.73,173.72) .. controls (346.88,172) and (346.38,167.14) .. (346.25,166.81) .. controls (345.46,164.86) and (342.22,163.22) .. (343.72,163.99) ;
%Shape: Free Drawing [id:dp7149381563661258] 
\draw  [color={rgb, 255:red, 255; green, 255; blue, 255 }  ,draw opacity=1 ][line width=2.25] [line join = round][line cap = round] (337.18,157.59) .. controls (337.02,157.6) and (336.7,157.55) .. (336.81,157.4) ;
%Shape: Free Drawing [id:dp22964569689630154] 
\draw  [draw opacity=0][line width=2.25] [line join = round][line cap = round] (422.02,152.55) .. controls (422.05,152.61) and (422.27,152.68) .. (422.2,152.64) ;
%Shape: Free Drawing [id:dp10065941420052993] 
\draw  [draw opacity=0][line width=2.25] [line join = round][line cap = round] (342.06,155.56) .. controls (340.35,155.65) and (338.87,156.82) .. (337.58,158.1) ;
%Shape: Free Drawing [id:dp3497860993106495] 
\draw  [draw opacity=0][line width=2.25] [line join = round][line cap = round] (340.85,156.15) .. controls (334.68,159.26) and (333.94,159.25) .. (335,159.8) ;
%Shape: Free Drawing [id:dp579495330613402] 
\draw  [draw opacity=0][line width=2.25] [line join = round][line cap = round] (341.36,155.81) .. controls (339.5,156.71) and (337.47,157.35) .. (335.82,158.71) ;
%Shape: Free Drawing [id:dp7142046075462641] 
\draw  [draw opacity=0][line width=2.25] [line join = round][line cap = round] (295.53,213.66) .. controls (294.41,213.89) and (293.26,214.3) .. (292.17,214.36) .. controls (291.88,214.37) and (292.81,213.78) .. (293.03,213.89) ;
%Shape: Free Drawing [id:dp36059894867560727] 
\draw  [draw opacity=0][line width=2.25] [line join = round][line cap = round] (291.03,213.47) .. controls (293.15,213.07) and (295.8,211.38) .. (297.4,212.2) ;
%Shape: Free Drawing [id:dp6910674550747714] 
\draw  [draw opacity=0][line width=2.25] [line join = round][line cap = round] (293.88,213.12) .. controls (294.87,213.63) and (296.47,212.93) .. (297.79,212.71) ;
%Shape: Free Drawing [id:dp5927274501309536] 
\draw  [draw opacity=0][line width=2.25] [line join = round][line cap = round] (293.46,212.3) .. controls (294.16,212.36) and (297.07,212.98) .. (297.91,211.86) ;
%Shape: Free Drawing [id:dp5667950092933809] 
\draw  [draw opacity=0][line width=2.25] [line join = round][line cap = round] (293.32,212.83) .. controls (294.8,212.66) and (298.27,212.76) .. (299.64,210.94) ;
%Shape: Free Drawing [id:dp23387142797305527] 
\draw  [draw opacity=0][line width=2.25] [line join = round][line cap = round] (292.99,213.26) .. controls (295.27,212.57) and (297.56,211.74) .. (299.81,210.72) ;
%Shape: Free Drawing [id:dp009814050173147182] 
\draw  [draw opacity=0][line width=2.25] [line join = round][line cap = round] (292.62,213.07) .. controls (296.13,212.06) and (298.9,210.44) .. (302.15,208.29) ;
%Shape: Free Drawing [id:dp1531990322186717] 
\draw  [draw opacity=0][line width=2.25] [line join = round][line cap = round] (306.53,107.47) .. controls (306.36,107.24) and (306.29,106.91) .. (306.34,106.58) ;

%Shape: Parallelogram [id:dp602722874823763] 
\draw   (275.1,90) -- (490,90) -- (397.9,255.35) -- (183,255.35) -- cycle ;
%Straight Lines [id:da33965271518677576] 
\draw    (183,255.35) -- (183,270) ;
%Straight Lines [id:da5238934226743897] 
\draw    (397.9,255.35) -- (397.9,270) ;
%Straight Lines [id:da46592613323086984] 
\draw    (490,90) -- (490,104.65) ;
%Straight Lines [id:da03324567448990312] 
\draw    (183,270) -- (397.97,270) ;
%Straight Lines [id:da05494519914066254] 
\draw    (397.9,270) -- (490,104.65) ;
%Shape: Circle [id:dp6719648380552748] 
\draw  [fill={rgb, 255:red, 155; green, 155; blue, 155 }  ,fill opacity=0.27 ] (355,42.5) .. controls (355,21.79) and (371.79,5) .. (392.5,5) .. controls (413.21,5) and (430,21.79) .. (430,42.5) .. controls (430,63.21) and (413.21,80) .. (392.5,80) .. controls (371.79,80) and (355,63.21) .. (355,42.5) -- cycle ;
%Curve Lines [id:da7420886829816177] 
\draw    (355,42.5) .. controls (381,49.67) and (405.67,50.33) .. (430,40) ;
%Curve Lines [id:da9534326817781857] 
\draw  [dash pattern={on 4.5pt off 4.5pt}]  (355,42.5) .. controls (380.33,34.33) and (406.33,35.67) .. (430,40) ;
%Shape: Circle [id:dp6798222831973407] 
\draw  [fill={rgb, 255:red, 155; green, 155; blue, 155 }  ,fill opacity=0.27 ] (265,317.5) .. controls (265,296.79) and (281.79,280) .. (302.5,280) .. controls (323.21,280) and (340,296.79) .. (340,317.5) .. controls (340,338.21) and (323.21,355) .. (302.5,355) .. controls (281.79,355) and (265,338.21) .. (265,317.5) -- cycle ;
%Curve Lines [id:da8454704520111102] 
\draw    (265,317.5) .. controls (291,324.67) and (315.67,325.33) .. (340,315) ;
%Curve Lines [id:da3154394219916593] 
\draw  [dash pattern={on 4.5pt off 4.5pt}]  (265,317.5) .. controls (290.33,309.33) and (316.33,310.67) .. (340,315) ;

% Text Node
\draw (305.43,143.97) node [anchor=north west][inner sep=0.75pt]  [rotate=-17.24]  {$A$};
% Text Node
\draw (371.51,168.5) node [anchor=north west][inner sep=0.75pt]  [rotate=-17.24]  {$B$};
% Text Node
\draw (320.48,185.73) node [anchor=north west][inner sep=0.75pt]  [rotate=-17.24]  {$C$};
% Text Node
\draw (246.92,194.77) node [anchor=north west][inner sep=0.75pt]  [rotate=-17.24]  {$D$};
% Text Node
\draw (353.07,129.16) node [anchor=north west][inner sep=0.75pt]  [rotate=-17.24]  {$E$};
% Text Node
\draw (304.31,214.98) node [anchor=north west][inner sep=0.75pt]  [rotate=-17.24]  {$F$};


\end{tikzpicture}
    \caption{Dos 3 bolas inflándose hasta chocar con el nudo}
\end{figure}
Los globos que se encuentran sobre y debajo del plano representaran las 3-bolas que identificaremos para que por medio de la función de pegado codifiquen $S^3\setminus K$. La idea sera inflar estas hasta que se encuentren las fronteras de ambas, la pregunta que surge es estas donde se encontraran y de que manera se producirá la identificación. Observemos que en una región local del diagrama donde no halla un cruce al expandirse los globos estos se verán de la siguiente manera
\begin{figure}[H]
    \centering
    
%!TEX root = ./main.tex



\tikzset{every picture/.style={line width=0.75pt}} %set default line width to 0.75pt        

\begin{tikzpicture}[x=0.75pt,y=0.75pt,yscale=-1,xscale=1]
%uncomment if require: \path (0,877); %set diagram left start at 0, and has height of 877

%Straight Lines [id:da9114133242483868] 
\draw    (130,190) -- (230,190) ;
%Shape: Arc [id:dp76715373468173] 
\draw  [draw opacity=0] (230,190) .. controls (230,173.43) and (243.43,160) .. (260,160) .. controls (276.57,160) and (290,173.43) .. (290,190) -- (260,190) -- cycle ; \draw   (230,190) .. controls (230,173.43) and (243.43,160) .. (260,160) .. controls (276.57,160) and (290,173.43) .. (290,190) ;  
%Straight Lines [id:da5480526040795924] 
\draw    (290,190) -- (360,190) ;
%Straight Lines [id:da8548422928953872] 
\draw    (240,70) -- (130,190) ;
%Straight Lines [id:da9009737644518513] 
\draw    (470,70) -- (360,190) ;
%Shape: Arc [id:dp7906930930707468] 
\draw  [draw opacity=0] (340,70) .. controls (340,53.43) and (353.43,40) .. (370,40) .. controls (386.57,40) and (400,53.43) .. (400,70) -- (370,70) -- cycle ; \draw   (340,70) .. controls (340,53.43) and (353.43,40) .. (370,40) .. controls (386.57,40) and (400,53.43) .. (400,70) ;  
%Straight Lines [id:da7284237450228865] 
\draw    (240,70) -- (340,70) ;
%Straight Lines [id:da9557077871704964] 
\draw    (400,70) -- (470,70) ;
%Shape: Arc [id:dp1517473778529409] 
\draw  [draw opacity=0] (310,100) .. controls (310,83.43) and (323.43,70) .. (340,70) .. controls (356.57,70) and (370,83.43) .. (370,100) -- (340,100) -- cycle ; \draw   (310,100) .. controls (310,83.43) and (323.43,70) .. (340,70) .. controls (356.57,70) and (370,83.43) .. (370,100) ;  
%Shape: Arc [id:dp9007707592571377] 
\draw  [draw opacity=0] (285,130) .. controls (285,113.43) and (298.43,100) .. (315,100) .. controls (331.57,100) and (345,113.43) .. (345,130) -- (315,130) -- cycle ; \draw   (285,130) .. controls (285,113.43) and (298.43,100) .. (315,100) .. controls (331.57,100) and (345,113.43) .. (345,130) ;  
%Shape: Arc [id:dp2562761265277008] 
\draw  [draw opacity=0] (260,160) .. controls (260,143.43) and (273.43,130) .. (290,130) .. controls (306.57,130) and (320,143.43) .. (320,160) -- (290,160) -- cycle ; \draw   (260,160) .. controls (260,143.43) and (273.43,130) .. (290,130) .. controls (306.57,130) and (320,143.43) .. (320,160) ;  
%Straight Lines [id:da12803577030308289] 
\draw    (390,30) -- (240,190) ;
%Shape: Circle [id:dp9817645653301902] 
\draw   (239,196.5) .. controls (239,189.04) and (245.04,183) .. (252.5,183) .. controls (259.96,183) and (266,189.04) .. (266,196.5) .. controls (266,203.96) and (259.96,210) .. (252.5,210) .. controls (245.04,210) and (239,203.96) .. (239,196.5) -- cycle ;
%Straight Lines [id:da8931818284701049] 
\draw    (410,40) -- (263,205) ;
%Curve Lines [id:da20985507457895636] 
\draw    (390,30) .. controls (401,23) and (411.8,29.4) .. (410,40) ;
%Straight Lines [id:da6630076546842583] 
\draw    (130,200) -- (230,200) ;
%Straight Lines [id:da7747329083523316] 
\draw    (140,190) -- (130,200) ;
%Straight Lines [id:da7546070667858217] 
\draw    (290,200) -- (360,200) ;
%Straight Lines [id:da46420675335921213] 
\draw    (470,80) -- (360,200) ;
%Straight Lines [id:da5784934546012858] 
\draw    (460,80) -- (470,80) ;
%Shape: Arc [id:dp14611856141981894] 
\draw  [draw opacity=0] (290,200) .. controls (290,216.57) and (276.57,230) .. (260,230) .. controls (243.43,230) and (230,216.57) .. (230,200) -- (260,200) -- cycle ; \draw   (290,200) .. controls (290,216.57) and (276.57,230) .. (260,230) .. controls (243.43,230) and (230,216.57) .. (230,200) ;  




\end{tikzpicture}
    \caption{Vista local alrededor de una vecindad tubular}
\end{figure}
Al ver esto nos podemos dar una idea de cuales serán las caras de nuestro poliedro para la identificación, precisamente serán las regiones que delimita el diagrama del nudo, así ya tenemos la primera parte identificada, las caras de nuestro poliedro ideal serán las regiones que delimita del diagrama del nudo.\\

La segunda parte que ya tenemos determinada son los vértices ideales del poliedro ya que como mencionamos antes la única información que no esta en nuestra 3-variedad es el nudo en si, entonces los vértices ideales serán las cuerdas del diagrama, mas cuando entremos en detalle del algoritmo veremos como funciona eso.\\

Lo ultimo que nos falta es determinar cuales serán las aristas que rodearan las caras ya determinadas, esto lo podemos ver con la única parte del diagrama del nudo a la cual no hemos prestado tanta atención, la zona de los cruces, para esto presentamos la siguiente ilustracion para dar una intuición de lo que ocurre en esas regiones al inflar los globos.
\begin{figure}[H]
    \centering
    
%!TEX root = ./main.tex



\tikzset{every picture/.style={line width=0.75pt}} %set default line width to 0.75pt        

\begin{tikzpicture}[x=0.75pt,y=0.75pt,yscale=-1,xscale=1]
%uncomment if require: \path (0,877); %set diagram left start at 0, and has height of 877

%Straight Lines [id:da6932355685480536] 
\draw [draw opacity=0][fill={rgb, 255:red, 155; green, 155; blue, 155 }  ,fill opacity=1 ]   (412.94,200) -- (440,160) -- (325,100) -- (325,145.87) -- cycle ;
%Straight Lines [id:da1384787451044316] 
\draw [draw opacity=0][fill={rgb, 255:red, 128; green, 128; blue, 128 }  ,fill opacity=1 ]   (210,160) -- (237.06,200) -- (315.46,151.74) -- (325,145.87) -- (325,100) -- cycle ;
%Straight Lines [id:da0741744233702265] 
\draw [line width=1.5]    (210,160) -- (440,40) ;
%Straight Lines [id:da5955578937587962] 
\draw  [dash pattern={on 4.5pt off 4.5pt}]  (210,40) -- (440,160) ;
%Straight Lines [id:da85493541645859] 
\draw    (210,160) -- (237.06,200) ;
%Straight Lines [id:da7842887989072671] 
\draw    (440,160) -- (412.94,200) ;
%Straight Lines [id:da2669884694205853] 
\draw    (210,40) -- (210,160) ;
%Straight Lines [id:da637854901562851] 
\draw    (440,40) -- (440,160) ;
%Straight Lines [id:da7206676036750514] 
\draw [color={rgb, 255:red, 74; green, 144; blue, 226 }  ,draw opacity=1 ]   (325,100) -- (325,142.87) ;
\draw [shift={(325,145.87)}, rotate = 270] [fill={rgb, 255:red, 74; green, 144; blue, 226 }  ,fill opacity=1 ][line width=0.08]  [draw opacity=0] (8.93,-4.29) -- (0,0) -- (8.93,4.29) -- cycle    ;
%Straight Lines [id:da12493325383084064] 
\draw [line width=1.5]    (325,145.87) -- (412.94,200) ;
%Straight Lines [id:da0009226395893477957] 
\draw  [dash pattern={on 4.5pt off 4.5pt}]  (237.06,200) -- (325,145.87) ;
%Straight Lines [id:da06537491608635204] 
\draw    (237.06,200) -- (412.94,200) ;
%Straight Lines [id:da9541213893361323] 
\draw    (237.06,53.33) -- (412.94,53.33) ;

% Text Node
\draw (316.29,59.07) node [anchor=north west][inner sep=0.75pt]    {$S$};
% Text Node
\draw (232.76,99.07) node [anchor=north west][inner sep=0.75pt]    {$V$};
% Text Node
\draw (402.88,85.73) node [anchor=north west][inner sep=0.75pt]    {$T$};
% Text Node
\draw (312.76,165.73) node [anchor=north west][inner sep=0.75pt]    {$U$};


\end{tikzpicture}
    \caption{Dibujo local como segmentos}
\end{figure}
Claramente estas intersecciones están des-suavizadas por propósitos prácticos, pero observe como en la región que se genera entre las cuerdas es una región donde se diferencian tanto las caras del globo de arriba como las de abajo por lo que a cada cruce del diagrama se dotara de una arista con una orientación. Con toda esta información ya podemos plantear un paso a paso completamente basado en el diagrama del nudo (combinatorio) para hacer la descomposición posterior en los dos poliedros.

\begin{block}
Dado un diagrama de un nudo $K$. Este divide el plano en múltiples regiones, designe cada una de estas por algún identificador. Posterior a esto en cada cruce que encuentre en el diagrama dibujaremos aristas (segmentos de recta con orientación) como se muestra a continuación
\begin{figure}[H]
    \centering
    
%!TEX root = ./main.tex



\tikzset{every picture/.style={line width=0.75pt}} %set default line width to 0.75pt        

\begin{tikzpicture}[x=0.75pt,y=0.75pt,yscale=-1,xscale=1]
%uncomment if require: \path (0,877); %set diagram left start at 0, and has height of 877

%Straight Lines [id:da46690641073913] 
\draw [line width=2.25]    (210,40) -- (350,160) ;
%Straight Lines [id:da10554492911680236] 
\draw [line width=2.25]    (300,90) -- (360,40) ;
%Straight Lines [id:da8098815627289185] 
\draw [line width=2.25]    (210,160) -- (270,110) ;
%Straight Lines [id:da05962605728576298] 
\draw [color={rgb, 255:red, 74; green, 144; blue, 226 }  ,draw opacity=1 ][line width=1.5]    (329,142) -- (251.62,129.63) ;
\draw [shift={(247.67,129)}, rotate = 9.08] [fill={rgb, 255:red, 74; green, 144; blue, 226 }  ,fill opacity=1 ][line width=0.08]  [draw opacity=0] (13.4,-6.43) -- (0,0) -- (13.4,6.44) -- (8.9,0) -- cycle    ;
%Straight Lines [id:da4600200086141276] 
\draw [color={rgb, 255:red, 74; green, 144; blue, 226 }  ,draw opacity=1 ][line width=1.5]    (240,64) -- (247.2,125.03) ;
\draw [shift={(247.67,129)}, rotate = 263.27] [fill={rgb, 255:red, 74; green, 144; blue, 226 }  ,fill opacity=1 ][line width=0.08]  [draw opacity=0] (13.4,-6.43) -- (0,0) -- (13.4,6.44) -- (8.9,0) -- cycle    ;
%Straight Lines [id:da4794416836174995] 
\draw [color={rgb, 255:red, 74; green, 144; blue, 226 }  ,draw opacity=1 ][line width=1.5]    (240,64) -- (314.39,76.35) ;
\draw [shift={(318.33,77)}, rotate = 189.42] [fill={rgb, 255:red, 74; green, 144; blue, 226 }  ,fill opacity=1 ][line width=0.08]  [draw opacity=0] (13.4,-6.43) -- (0,0) -- (13.4,6.44) -- (8.9,0) -- cycle    ;
%Straight Lines [id:da9958661012600074] 
\draw [color={rgb, 255:red, 74; green, 144; blue, 226 }  ,draw opacity=1 ][line width=1.5]    (329,142) -- (318.98,80.95) ;
\draw [shift={(318.33,77)}, rotate = 80.68] [fill={rgb, 255:red, 74; green, 144; blue, 226 }  ,fill opacity=1 ][line width=0.08]  [draw opacity=0] (13.4,-6.43) -- (0,0) -- (13.4,6.44) -- (8.9,0) -- cycle    ;




\end{tikzpicture}
    \caption{Arista en un cruce}
\end{figure}
\end{block}
Algunas observaciones antes de ver como se visualiza esto en el nudo de $8$, recordemos que antes había dicho que por cada cruce en el nudo surgía una arista mas en la ilustración anterior se dibujaron 4, en primer lugar observe que las cuatro tienen toda la misma orientación (ya sea desde la cuerda que pasa por abajo hacia la de arriba o viceversa) esto se debe a que estas aristas son isotopicas en $S^3$, esto se consigue deslizando tanto la cola o la cabeza de cada flecha a través de la cuerda, luego el motivo de dibujarlas las 4 es para determinar desde un principio cuales son las aristas que rodean cada región y convierten en una cara. Con esto en mente veamos en la proyección del nudo de $8$ como se ve este primer paso
\begin{figure}[H]


%!TEX root = ./main.tex

\tikzset{every picture/.style={line width=0.75pt}} %set default line width to 0.75pt        

\begin{tikzpicture}[x=0.85pt,y=0.75pt,yscale=-0.9,xscale=0.9]
%uncomment if require: \path (0,877); %set diagram left start at 0, and has height of 877

%Shape: Free Drawing [id:dp052719054339937954] 
\draw  [line width=2.25] [line join = round][line cap = round] (287.39,145.57) .. controls (277.01,144.32) and (278.03,126.78) .. (274.73,116.85) .. controls (265.82,90.11) and (273.17,57.75) .. (294.63,39.17) .. controls (299.44,35.01) and (308.14,29.62) .. (314.87,30.02) .. controls (429.8,36.81) and (491.19,214.64) .. (407.67,292.97) .. controls (397.46,302.54) and (381.91,304.56) .. (369.98,296.5) .. controls (360.6,290.15) and (352.15,274.14) .. (341.49,277.97) ;
%Shape: Free Drawing [id:dp25610451731952133] 
\draw  [line width=2.25] [line join = round][line cap = round] (322.25,274.65) .. controls (309.14,266) and (275.87,280.02) .. (263.45,286.22) .. controls (260.49,287.69) and (258.35,290.7) .. (255.22,291.78) .. controls (218.11,304.6) and (189.64,286.34) .. (171.87,252.05) .. controls (165.39,239.55) and (160.62,225.56) .. (160.14,211.36) .. controls (158.61,166.54) and (169.69,97.98) .. (203.58,63.49) .. controls (217.17,49.66) and (240.46,37.41) .. (260.26,36.05) .. controls (267.59,35.55) and (282.28,43.73) .. (279.03,43.34) ;
%Shape: Free Drawing [id:dp5679641581314395] 
\draw  [line width=2.25] [line join = round][line cap = round] (294.08,52.2) .. controls (309.66,62.48) and (309.98,81.23) .. (312.22,98.19) .. controls (313.38,106.94) and (316.36,116.92) .. (313.9,126.09) .. controls (311.17,136.3) and (302.82,146.56) .. (297.8,155.38) .. controls (285.1,177.65) and (272.32,196.06) .. (263.69,221.28) .. controls (261.03,229.08) and (257.74,245.85) .. (258.53,255.91) .. controls (258.89,260.59) and (259.75,265.24) .. (260.89,269.79) .. controls (261.54,272.4) and (266.69,277.54) .. (264.02,277.22) ;
%Shape: Free Drawing [id:dp6189964124099764] 
\draw  [line width=2.25] [line join = round][line cap = round] (273.57,289.95) .. controls (278.3,293.07) and (276.89,297.76) .. (283.68,302.24) .. controls (287.6,304.83) and (292.16,306.32) .. (296.61,307.83) .. controls (327.83,318.42) and (338.62,263.49) .. (338.1,243.31) .. controls (337.37,215.47) and (335.79,198.2) .. (322.2,174.43) .. controls (320.36,171.21) and (313.27,162.71) .. (312.59,162.19) .. controls (308.57,159.15) and (299.94,158.15) .. (303.95,158.64) ;
%Shape: Free Drawing [id:dp23463200731808187] 
\draw  [color={rgb, 255:red, 255; green, 255; blue, 255 }  ,draw opacity=1 ][line width=2.25] [line join = round][line cap = round] (282.71,151.05) .. controls (282.39,151.16) and (281.67,151.26) .. (281.71,150.93) ;
%Shape: Free Drawing [id:dp6791599529562201] 
\draw  [draw opacity=0][line width=2.25] [line join = round][line cap = round] (448.19,89.85) .. controls (448.33,89.94) and (448.85,89.93) .. (448.68,89.91) ;
%Shape: Free Drawing [id:dp4292482997182422] 
\draw  [draw opacity=0][line width=2.25] [line join = round][line cap = round] (290.05,144.38) .. controls (286.7,145.58) and (285.16,148.61) .. (284.14,151.72) ;
%Shape: Free Drawing [id:dp6085326410058335] 
\draw  [draw opacity=0][line width=2.25] [line join = round][line cap = round] (288.33,146.18) .. controls (279.73,155.6) and (278.21,156.05) .. (281.06,156.39) ;
%Shape: Free Drawing [id:dp9043644710987669] 
\draw  [draw opacity=0][line width=2.25] [line join = round][line cap = round] (288.94,145.25) .. controls (286.29,148.02) and (282.99,150.44) .. (281.35,153.91) ;
%Shape: Free Drawing [id:dp21505097221297298] 
\draw  [draw opacity=0][line width=2.25] [line join = round][line cap = round] (268.43,278.25) .. controls (266.45,279.36) and (264.63,280.8) .. (262.49,281.57) .. controls (261.93,281.77) and (263.07,280.13) .. (263.67,280.2) ;
%Shape: Free Drawing [id:dp28830093222917674] 
\draw  [draw opacity=0][line width=2.25] [line join = round][line cap = round] (259.08,280.66) .. controls (262.87,278.63) and (266.13,273.95) .. (270.4,274.46) ;
%Shape: Free Drawing [id:dp9547251910894086] 
\draw  [draw opacity=0][line width=2.25] [line join = round][line cap = round] (264.4,278.27) .. controls (267.07,278.59) and (269.43,276.35) .. (271.83,275.14) ;
%Shape: Free Drawing [id:dp719981311992137] 
\draw  [draw opacity=0][line width=2.25] [line join = round][line cap = round] (262.53,277.04) .. controls (264.04,276.73) and (270.71,276.07) .. (271.01,273.53) ;
%Shape: Free Drawing [id:dp29834365303629773] 
\draw  [draw opacity=0][line width=2.25] [line join = round][line cap = round] (262.91,278.09) .. controls (265.7,276.89) and (272.86,274.94) .. (273.36,270.79) ;
%Shape: Free Drawing [id:dp42617767268148665] 
\draw  [draw opacity=0][line width=2.25] [line join = round][line cap = round] (262.79,279.09) .. controls (266.54,276.42) and (270.14,273.51) .. (273.42,270.29) ;
%Shape: Free Drawing [id:dp5503601990364082] 
\draw  [draw opacity=0][line width=2.25] [line join = round][line cap = round] (261.8,278.97) .. controls (267.65,274.97) and (271.24,270.34) .. (275.13,264.45) ;
%Shape: Free Drawing [id:dp12708063524714575] 
\draw  [draw opacity=0][line width=2.25] [line join = round][line cap = round] (29,80) .. controls (28.37,79.68) and (27.82,79.13) .. (27.5,78.5) ;
%Straight Lines [id:da7851146054141257] 
\draw [color={rgb, 255:red, 74; green, 144; blue, 226 }  ,draw opacity=1 ][line width=1.5]    (274.5,69.5) -- (301.08,63.19) ;
\draw [shift={(304,62.5)}, rotate = 166.65] [color={rgb, 255:red, 74; green, 144; blue, 226 }  ,draw opacity=1 ][line width=1.5]    (14.21,-4.28) .. controls (9.04,-1.82) and (4.3,-0.39) .. (0,0) .. controls (4.3,0.39) and (9.04,1.82) .. (14.21,4.28)   ;
%Straight Lines [id:da4904183098904815] 
\draw [color={rgb, 255:red, 74; green, 144; blue, 226 }  ,draw opacity=1 ][line width=1.5]    (274.5,69.5) -- (263.05,38.81) ;
\draw [shift={(262,36)}, rotate = 69.54] [color={rgb, 255:red, 74; green, 144; blue, 226 }  ,draw opacity=1 ][line width=1.5]    (14.21,-4.28) .. controls (9.04,-1.82) and (4.3,-0.39) .. (0,0) .. controls (4.3,0.39) and (9.04,1.82) .. (14.21,4.28)   ;
%Straight Lines [id:da21498342913887913] 
\draw [color={rgb, 255:red, 74; green, 144; blue, 226 }  ,draw opacity=1 ][line width=1.5]    (301,34.5) -- (303.68,59.52) ;
\draw [shift={(304,62.5)}, rotate = 263.88] [color={rgb, 255:red, 74; green, 144; blue, 226 }  ,draw opacity=1 ][line width=1.5]    (14.21,-4.28) .. controls (9.04,-1.82) and (4.3,-0.39) .. (0,0) .. controls (4.3,0.39) and (9.04,1.82) .. (14.21,4.28)   ;
%Straight Lines [id:da5005258690622153] 
\draw [color={rgb, 255:red, 74; green, 144; blue, 226 }  ,draw opacity=1 ][line width=1.5]    (301,34.5) -- (265,35.88) ;
\draw [shift={(262,36)}, rotate = 357.8] [color={rgb, 255:red, 74; green, 144; blue, 226 }  ,draw opacity=1 ][line width=1.5]    (14.21,-4.28) .. controls (9.04,-1.82) and (4.3,-0.39) .. (0,0) .. controls (4.3,0.39) and (9.04,1.82) .. (14.21,4.28)   ;
%Straight Lines [id:da6750405427654139] 
\draw [color={rgb, 255:red, 74; green, 144; blue, 226 }  ,draw opacity=1 ][line width=1.5]    (278,132.5) -- (286.34,169.57) ;
\draw [shift={(287,172.5)}, rotate = 257.32] [color={rgb, 255:red, 74; green, 144; blue, 226 }  ,draw opacity=1 ][line width=1.5]    (14.21,-4.28) .. controls (9.04,-1.82) and (4.3,-0.39) .. (0,0) .. controls (4.3,0.39) and (9.04,1.82) .. (14.21,4.28)   ;
%Straight Lines [id:da3934982113036074] 
\draw [color={rgb, 255:red, 74; green, 144; blue, 226 }  ,draw opacity=1 ][line width=1.5]    (278,132.5) -- (306.01,135.21) ;
\draw [shift={(309,135.5)}, rotate = 185.53] [color={rgb, 255:red, 74; green, 144; blue, 226 }  ,draw opacity=1 ][line width=1.5]    (14.21,-4.28) .. controls (9.04,-1.82) and (4.3,-0.39) .. (0,0) .. controls (4.3,0.39) and (9.04,1.82) .. (14.21,4.28)   ;
%Straight Lines [id:da8163214263178915] 
\draw [color={rgb, 255:red, 74; green, 144; blue, 226 }  ,draw opacity=1 ][line width=1.5]    (320,172) -- (290,172.45) ;
\draw [shift={(287,172.5)}, rotate = 359.13] [color={rgb, 255:red, 74; green, 144; blue, 226 }  ,draw opacity=1 ][line width=1.5]    (14.21,-4.28) .. controls (9.04,-1.82) and (4.3,-0.39) .. (0,0) .. controls (4.3,0.39) and (9.04,1.82) .. (14.21,4.28)   ;
%Straight Lines [id:da9808570839490202] 
\draw [color={rgb, 255:red, 74; green, 144; blue, 226 }  ,draw opacity=1 ][line width=1.5]    (320,172) -- (309.87,138.37) ;
\draw [shift={(309,135.5)}, rotate = 73.23] [color={rgb, 255:red, 74; green, 144; blue, 226 }  ,draw opacity=1 ][line width=1.5]    (14.21,-4.28) .. controls (9.04,-1.82) and (4.3,-0.39) .. (0,0) .. controls (4.3,0.39) and (9.04,1.82) .. (14.21,4.28)   ;
%Straight Lines [id:da9411841582233399] 
\draw [color={rgb, 255:red, 74; green, 144; blue, 226 }  ,draw opacity=1 ][line width=1.5]    (258,261) -- (244.2,292.75) ;
\draw [shift={(243,295.5)}, rotate = 293.5] [color={rgb, 255:red, 74; green, 144; blue, 226 }  ,draw opacity=1 ][line width=1.5]    (14.21,-4.28) .. controls (9.04,-1.82) and (4.3,-0.39) .. (0,0) .. controls (4.3,0.39) and (9.04,1.82) .. (14.21,4.28)   ;
%Straight Lines [id:da003586680143401688] 
\draw [color={rgb, 255:red, 74; green, 144; blue, 226 }  ,draw opacity=1 ][line width=1.5]    (285,303) -- (245.95,296.03) ;
\draw [shift={(243,295.5)}, rotate = 10.12] [color={rgb, 255:red, 74; green, 144; blue, 226 }  ,draw opacity=1 ][line width=1.5]    (14.21,-4.28) .. controls (9.04,-1.82) and (4.3,-0.39) .. (0,0) .. controls (4.3,0.39) and (9.04,1.82) .. (14.21,4.28)   ;
%Straight Lines [id:da8295997211316378] 
\draw [color={rgb, 255:red, 74; green, 144; blue, 226 }  ,draw opacity=1 ][line width=1.5]    (258,261) -- (282.42,275.47) ;
\draw [shift={(285,277)}, rotate = 210.65] [color={rgb, 255:red, 74; green, 144; blue, 226 }  ,draw opacity=1 ][line width=1.5]    (14.21,-4.28) .. controls (9.04,-1.82) and (4.3,-0.39) .. (0,0) .. controls (4.3,0.39) and (9.04,1.82) .. (14.21,4.28)   ;
%Straight Lines [id:da4207009837534177] 
\draw [color={rgb, 255:red, 74; green, 144; blue, 226 }  ,draw opacity=1 ][line width=1.5]    (285,303) -- (285,280) ;
\draw [shift={(285,277)}, rotate = 90] [color={rgb, 255:red, 74; green, 144; blue, 226 }  ,draw opacity=1 ][line width=1.5]    (14.21,-4.28) .. controls (9.04,-1.82) and (4.3,-0.39) .. (0,0) .. controls (4.3,0.39) and (9.04,1.82) .. (14.21,4.28)   ;
%Straight Lines [id:da7105295943160274] 
\draw [color={rgb, 255:red, 74; green, 144; blue, 226 }  ,draw opacity=1 ][line width=1.5]    (336.5,256.5) -- (313.61,269.52) ;
\draw [shift={(311,271)}, rotate = 330.38] [color={rgb, 255:red, 74; green, 144; blue, 226 }  ,draw opacity=1 ][line width=1.5]    (14.21,-4.28) .. controls (9.04,-1.82) and (4.3,-0.39) .. (0,0) .. controls (4.3,0.39) and (9.04,1.82) .. (14.21,4.28)   ;
%Straight Lines [id:da6405444868178161] 
\draw [color={rgb, 255:red, 74; green, 144; blue, 226 }  ,draw opacity=1 ][line width=1.5]    (324,293.5) -- (312.5,273.6) ;
\draw [shift={(311,271)}, rotate = 59.98] [color={rgb, 255:red, 74; green, 144; blue, 226 }  ,draw opacity=1 ][line width=1.5]    (14.21,-4.28) .. controls (9.04,-1.82) and (4.3,-0.39) .. (0,0) .. controls (4.3,0.39) and (9.04,1.82) .. (14.21,4.28)   ;
%Straight Lines [id:da5651565239928864] 
\draw [color={rgb, 255:red, 74; green, 144; blue, 226 }  ,draw opacity=1 ][line width=1.5]    (324,293.5) -- (354.09,285.75) ;
\draw [shift={(357,285)}, rotate = 165.56] [color={rgb, 255:red, 74; green, 144; blue, 226 }  ,draw opacity=1 ][line width=1.5]    (14.21,-4.28) .. controls (9.04,-1.82) and (4.3,-0.39) .. (0,0) .. controls (4.3,0.39) and (9.04,1.82) .. (14.21,4.28)   ;
%Straight Lines [id:da48389743721592937] 
\draw [color={rgb, 255:red, 74; green, 144; blue, 226 }  ,draw opacity=1 ][line width=1.5]    (336.5,256.5) -- (355.25,282.56) ;
\draw [shift={(357,285)}, rotate = 234.27] [color={rgb, 255:red, 74; green, 144; blue, 226 }  ,draw opacity=1 ][line width=1.5]    (14.21,-4.28) .. controls (9.04,-1.82) and (4.3,-0.39) .. (0,0) .. controls (4.3,0.39) and (9.04,1.82) .. (14.21,4.28)   ;

% Text Node
\draw (210.5,152.4) node [anchor=north west][inner sep=0.75pt]    {$A$};
% Text Node
\draw (373,156.4) node [anchor=north west][inner sep=0.75pt]    {$B$};
% Text Node
\draw (292.5,218.9) node [anchor=north west][inner sep=0.75pt]    {$C$};
% Text Node
\draw (155.5,280.4) node [anchor=north west][inner sep=0.75pt]    {$D$};
% Text Node
\draw (286.5,96.4) node [anchor=north west][inner sep=0.75pt]    {$E$};
% Text Node
\draw (295,281.9) node [anchor=north west][inner sep=0.75pt]    {$F$};
% Text Node
\draw (258.5,50.9) node [anchor=north west][inner sep=0.75pt]  [font=\small]  {$0$};
% Text Node
\draw (268.5,141.9) node [anchor=north west][inner sep=0.75pt]  [font=\small]  {$1$};
% Text Node
\draw (260.5,300.4) node [anchor=north west][inner sep=0.75pt]  [font=\small]  {$2$};
% Text Node
\draw (332.5,293.4) node [anchor=north west][inner sep=0.75pt]  [font=\small]  {$3$};


\end{tikzpicture}

\caption{Paso 1 del algoritmo para el nudo de 8}

\end{figure}
Observe que por ejemplo la cara identificada por $A$ las aristas que la rodean son la $0,1$ y $2$. lo cual en la ilustración con solo $1$ arista por cruce es bastante difícil de ver
\begin{block}
    Contraiga cada cuerda que sea un cruce superior a un punto, este punto sera el vértice ideal de nuestro poliedro superior, dado por la esfera de arriba.
\end{block}

Una observación relevante es que contraer estas cuerdas a un punto es posible ya que el complemento de un punto en $S^2$ es homeomorfo al complemento de un segmento de recta en $S^2$, esto se puede mostrar fácilmente con ayuda de la proyección estereográfica, la otra observación que es netamente de visualización es que estamos viendo el grafo desde dentro de la bola superior, esto lo haremos con un paso intermedio que se puede ver en el siguiente diagrama
\begin{figure}[H]
    \centering
    
%!TEX root = ./main.tex


\tikzset{every picture/.style={line width=0.75pt}} %set default line width to 0.75pt        

\begin{tikzpicture}[x=0.75pt,y=0.75pt,yscale=-0.85,xscale=0.85]
%uncomment if require: \path (0,877); %set diagram left start at 0, and has height of 877

%Shape: Free Drawing [id:dp052719054339937954] 
\draw  [line width=2.25] [line join = round][line cap = round] (144.39,147.9) .. controls (134.01,146.65) and (135.03,129.11) .. (131.73,119.19) .. controls (122.82,92.44) and (130.17,60.08) .. (151.63,41.51) .. controls (156.44,37.34) and (165.14,31.96) .. (171.87,32.35) .. controls (286.8,39.15) and (348.19,216.97) .. (264.67,295.3) .. controls (254.46,304.88) and (238.91,306.9) .. (226.98,298.83) .. controls (217.6,292.48) and (209.15,276.48) .. (198.49,280.3) ;
%Shape: Free Drawing [id:dp25610451731952133] 
\draw  [line width=2.25] [line join = round][line cap = round] (179.25,276.98) .. controls (166.14,268.33) and (132.87,282.35) .. (120.45,288.55) .. controls (117.49,290.03) and (115.35,293.03) .. (112.22,294.11) .. controls (75.11,306.93) and (46.64,288.67) .. (28.87,254.38) .. controls (22.39,241.88) and (17.62,227.89) .. (17.14,213.69) .. controls (15.61,168.88) and (26.69,100.32) .. (60.58,65.82) .. controls (74.17,52) and (97.46,39.75) .. (117.26,38.38) .. controls (124.59,37.88) and (139.28,46.06) .. (136.03,45.67) ;
%Shape: Free Drawing [id:dp5679641581314395] 
\draw  [line width=2.25] [line join = round][line cap = round] (151.08,54.53) .. controls (166.66,64.82) and (166.98,83.56) .. (169.22,100.52) .. controls (170.38,109.27) and (173.36,119.25) .. (170.9,128.42) .. controls (168.17,138.64) and (159.82,148.9) .. (154.8,157.71) .. controls (142.1,179.99) and (129.32,198.39) .. (120.69,223.62) .. controls (118.03,231.41) and (114.74,248.18) .. (115.53,258.25) .. controls (115.89,262.93) and (116.75,267.57) .. (117.89,272.13) .. controls (118.54,274.73) and (123.69,279.87) .. (121.02,279.55) ;
%Shape: Free Drawing [id:dp6189964124099764] 
\draw  [line width=2.25] [line join = round][line cap = round] (130.57,292.28) .. controls (135.3,295.4) and (133.89,300.1) .. (140.68,304.58) .. controls (144.6,307.16) and (149.16,308.65) .. (153.61,310.16) .. controls (184.83,320.75) and (195.62,265.82) .. (195.1,245.65) .. controls (194.37,217.8) and (192.79,200.53) .. (179.2,176.76) .. controls (177.36,173.55) and (170.27,165.04) .. (169.59,164.53) .. controls (165.57,161.48) and (156.94,160.49) .. (160.95,160.97) ;
%Shape: Free Drawing [id:dp23463200731808187] 
\draw  [color={rgb, 255:red, 255; green, 255; blue, 255 }  ,draw opacity=1 ][line width=2.25] [line join = round][line cap = round] (139.71,153.38) .. controls (139.39,153.49) and (138.67,153.59) .. (138.71,153.26) ;
%Shape: Free Drawing [id:dp6791599529562201] 
\draw  [draw opacity=0][line width=2.25] [line join = round][line cap = round] (305.19,92.18) .. controls (305.33,92.27) and (305.85,92.26) .. (305.68,92.24) ;
%Shape: Free Drawing [id:dp4292482997182422] 
\draw  [draw opacity=0][line width=2.25] [line join = round][line cap = round] (147.05,146.71) .. controls (143.7,147.91) and (142.16,150.95) .. (141.14,154.06) ;
%Shape: Free Drawing [id:dp6085326410058335] 
\draw  [draw opacity=0][line width=2.25] [line join = round][line cap = round] (145.33,148.52) .. controls (136.73,157.93) and (135.21,158.38) .. (138.06,158.72) ;
%Shape: Free Drawing [id:dp9043644710987669] 
\draw  [draw opacity=0][line width=2.25] [line join = round][line cap = round] (145.94,147.58) .. controls (143.29,150.36) and (139.99,152.77) .. (138.35,156.24) ;
%Shape: Free Drawing [id:dp21505097221297298] 
\draw  [draw opacity=0][line width=2.25] [line join = round][line cap = round] (125.43,280.59) .. controls (123.45,281.69) and (121.63,283.14) .. (119.49,283.9) .. controls (118.93,284.11) and (120.07,282.46) .. (120.67,282.53) ;
%Shape: Free Drawing [id:dp28830093222917674] 
\draw  [draw opacity=0][line width=2.25] [line join = round][line cap = round] (116.08,282.99) .. controls (119.87,280.96) and (123.13,276.28) .. (127.4,276.79) ;
%Shape: Free Drawing [id:dp9547251910894086] 
\draw  [draw opacity=0][line width=2.25] [line join = round][line cap = round] (121.4,280.61) .. controls (124.07,280.93) and (126.43,278.68) .. (128.83,277.47) ;
%Shape: Free Drawing [id:dp719981311992137] 
\draw  [draw opacity=0][line width=2.25] [line join = round][line cap = round] (119.53,279.38) .. controls (121.04,279.07) and (127.71,278.4) .. (128.01,275.86) ;
%Shape: Free Drawing [id:dp29834365303629773] 
\draw  [draw opacity=0][line width=2.25] [line join = round][line cap = round] (119.91,280.43) .. controls (122.7,279.22) and (129.86,277.27) .. (130.36,273.12) ;
%Shape: Free Drawing [id:dp42617767268148665] 
\draw  [draw opacity=0][line width=2.25] [line join = round][line cap = round] (119.79,281.42) .. controls (123.54,278.76) and (127.14,275.85) .. (130.42,272.63) ;
%Shape: Free Drawing [id:dp5503601990364082] 
\draw  [draw opacity=0][line width=2.25] [line join = round][line cap = round] (118.8,281.3) .. controls (124.65,277.3) and (128.24,272.67) .. (132.13,266.79) ;
%Shape: Free Drawing [id:dp12708063524714575] 
\draw  [draw opacity=0][line width=2.25] [line join = round][line cap = round] (15,81.33) .. controls (14.37,81.02) and (13.82,80.47) .. (13.5,79.83) ;
%Straight Lines [id:da4904183098904815] 
\draw [color={rgb, 255:red, 74; green, 144; blue, 226 }  ,draw opacity=1 ][line width=1.5]    (142.5,49.83) -- (121.69,39.65) ;
\draw [shift={(119,38.33)}, rotate = 26.08] [color={rgb, 255:red, 74; green, 144; blue, 226 }  ,draw opacity=1 ][line width=1.5]    (14.21,-4.28) .. controls (9.04,-1.82) and (4.3,-0.39) .. (0,0) .. controls (4.3,0.39) and (9.04,1.82) .. (14.21,4.28)   ;
%Straight Lines [id:da21498342913887913] 
\draw [color={rgb, 255:red, 74; green, 144; blue, 226 }  ,draw opacity=1 ][line width=1.5]    (142.5,49.83) -- (158.67,62.94) ;
\draw [shift={(161,64.83)}, rotate = 219.04] [color={rgb, 255:red, 74; green, 144; blue, 226 }  ,draw opacity=1 ][line width=1.5]    (14.21,-4.28) .. controls (9.04,-1.82) and (4.3,-0.39) .. (0,0) .. controls (4.3,0.39) and (9.04,1.82) .. (14.21,4.28)   ;
%Straight Lines [id:da6750405427654139] 
\draw [color={rgb, 255:red, 74; green, 144; blue, 226 }  ,draw opacity=1 ][line width=1.5]    (135,134.83) -- (150.58,153.53) ;
\draw [shift={(152.5,155.83)}, rotate = 230.19] [color={rgb, 255:red, 74; green, 144; blue, 226 }  ,draw opacity=1 ][line width=1.5]    (14.21,-4.28) .. controls (9.04,-1.82) and (4.3,-0.39) .. (0,0) .. controls (4.3,0.39) and (9.04,1.82) .. (14.21,4.28)   ;
%Straight Lines [id:da9808570839490202] 
\draw [color={rgb, 255:red, 74; green, 144; blue, 226 }  ,draw opacity=1 ][line width=1.5]    (177,174.33) -- (160.68,158.89) ;
\draw [shift={(158.5,156.83)}, rotate = 43.41] [color={rgb, 255:red, 74; green, 144; blue, 226 }  ,draw opacity=1 ][line width=1.5]    (14.21,-4.28) .. controls (9.04,-1.82) and (4.3,-0.39) .. (0,0) .. controls (4.3,0.39) and (9.04,1.82) .. (14.21,4.28)   ;
%Straight Lines [id:da8295997211316378] 
\draw [color={rgb, 255:red, 74; green, 144; blue, 226 }  ,draw opacity=1 ][line width=1.5]    (115.5,256.33) -- (124.55,283.49) ;
\draw [shift={(125.5,286.33)}, rotate = 251.57] [color={rgb, 255:red, 74; green, 144; blue, 226 }  ,draw opacity=1 ][line width=1.5]    (14.21,-4.28) .. controls (9.04,-1.82) and (4.3,-0.39) .. (0,0) .. controls (4.3,0.39) and (9.04,1.82) .. (14.21,4.28)   ;
%Straight Lines [id:da4207009837534177] 
\draw [color={rgb, 255:red, 74; green, 144; blue, 226 }  ,draw opacity=1 ][line width=1.5]    (147.5,308.83) -- (127.6,288.48) ;
\draw [shift={(125.5,286.33)}, rotate = 45.64] [color={rgb, 255:red, 74; green, 144; blue, 226 }  ,draw opacity=1 ][line width=1.5]    (14.21,-4.28) .. controls (9.04,-1.82) and (4.3,-0.39) .. (0,0) .. controls (4.3,0.39) and (9.04,1.82) .. (14.21,4.28)   ;
%Straight Lines [id:da7105295943160274] 
\draw [color={rgb, 255:red, 74; green, 144; blue, 226 }  ,draw opacity=1 ][line width=1.5]    (189.5,279.33) -- (170.89,274.14) ;
\draw [shift={(168,273.33)}, rotate = 15.59] [color={rgb, 255:red, 74; green, 144; blue, 226 }  ,draw opacity=1 ][line width=1.5]    (14.21,-4.28) .. controls (9.04,-1.82) and (4.3,-0.39) .. (0,0) .. controls (4.3,0.39) and (9.04,1.82) .. (14.21,4.28)   ;
%Straight Lines [id:da5651565239928864] 
\draw [color={rgb, 255:red, 74; green, 144; blue, 226 }  ,draw opacity=1 ][line width=1.5]    (189.5,279.33) -- (211.15,286.4) ;
\draw [shift={(214,287.33)}, rotate = 198.08] [color={rgb, 255:red, 74; green, 144; blue, 226 }  ,draw opacity=1 ][line width=1.5]    (14.21,-4.28) .. controls (9.04,-1.82) and (4.3,-0.39) .. (0,0) .. controls (4.3,0.39) and (9.04,1.82) .. (14.21,4.28)   ;
%Shape: Circle [id:dp7563944317650207] 
\draw  [color={rgb, 255:red, 74; green, 144; blue, 226 }  ,draw opacity=1 ][fill={rgb, 255:red, 74; green, 144; blue, 226 }  ,fill opacity=1 ] (136.5,49.83) .. controls (136.5,46.52) and (139.19,43.83) .. (142.5,43.83) .. controls (145.81,43.83) and (148.5,46.52) .. (148.5,49.83) .. controls (148.5,53.15) and (145.81,55.83) .. (142.5,55.83) .. controls (139.19,55.83) and (136.5,53.15) .. (136.5,49.83) -- cycle ;
%Shape: Circle [id:dp8058641002050105] 
\draw  [color={rgb, 255:red, 74; green, 144; blue, 226 }  ,draw opacity=1 ][fill={rgb, 255:red, 74; green, 144; blue, 226 }  ,fill opacity=1 ] (150,157.33) .. controls (150,154.02) and (152.69,151.33) .. (156,151.33) .. controls (159.31,151.33) and (162,154.02) .. (162,157.33) .. controls (162,160.65) and (159.31,163.33) .. (156,163.33) .. controls (152.69,163.33) and (150,160.65) .. (150,157.33) -- cycle ;
%Shape: Circle [id:dp3830896703653248] 
\draw  [color={rgb, 255:red, 74; green, 144; blue, 226 }  ,draw opacity=1 ][fill={rgb, 255:red, 74; green, 144; blue, 226 }  ,fill opacity=1 ] (119.5,286.33) .. controls (119.5,283.02) and (122.19,280.33) .. (125.5,280.33) .. controls (128.81,280.33) and (131.5,283.02) .. (131.5,286.33) .. controls (131.5,289.65) and (128.81,292.33) .. (125.5,292.33) .. controls (122.19,292.33) and (119.5,289.65) .. (119.5,286.33) -- cycle ;
%Shape: Circle [id:dp568354357328672] 
\draw  [color={rgb, 255:red, 74; green, 144; blue, 226 }  ,draw opacity=1 ][fill={rgb, 255:red, 74; green, 144; blue, 226 }  ,fill opacity=1 ] (183.5,279.33) .. controls (183.5,276.02) and (186.19,273.33) .. (189.5,273.33) .. controls (192.81,273.33) and (195.5,276.02) .. (195.5,279.33) .. controls (195.5,282.65) and (192.81,285.33) .. (189.5,285.33) .. controls (186.19,285.33) and (183.5,282.65) .. (183.5,279.33) -- cycle ;
%Shape: Free Drawing [id:dp4081098483122335] 
\draw  [color={rgb, 255:red, 74; green, 144; blue, 226 }  ,draw opacity=1 ][line width=2.25] [line join = round][line cap = round] (492.73,148.24) .. controls (482.34,146.99) and (483.37,129.44) .. (480.06,119.52) .. controls (471.15,92.78) and (478.5,60.42) .. (499.97,41.84) .. controls (504.78,37.68) and (513.47,32.29) .. (520.2,32.69) .. controls (635.13,39.48) and (696.52,217.31) .. (613,295.63) .. controls (602.79,305.21) and (587.24,307.23) .. (575.31,299.16) .. controls (565.93,292.82) and (557.49,276.81) .. (546.82,280.63) ;
%Shape: Free Drawing [id:dp3946430254773102] 
\draw  [color={rgb, 255:red, 74; green, 144; blue, 226 }  ,draw opacity=1 ][line width=2.25] [line join = round][line cap = round] (527.58,277.31) .. controls (514.47,268.66) and (481.2,282.69) .. (468.78,288.88) .. controls (465.82,290.36) and (463.69,293.36) .. (460.56,294.44) .. controls (423.44,307.26) and (394.97,289.01) .. (377.2,254.72) .. controls (370.72,242.22) and (365.96,228.23) .. (365.47,214.03) .. controls (363.94,169.21) and (375.02,100.65) .. (408.92,66.15) .. controls (422.5,52.33) and (445.8,40.08) .. (465.59,38.72) .. controls (472.93,38.21) and (487.61,46.4) .. (484.36,46.01) ;
%Shape: Free Drawing [id:dp031095390456670646] 
\draw  [color={rgb, 255:red, 74; green, 144; blue, 226 }  ,draw opacity=1 ][line width=2.25] [line join = round][line cap = round] (499.41,54.87) .. controls (515,65.15) and (515.31,83.9) .. (517.55,100.86) .. controls (518.71,109.6) and (521.69,119.59) .. (519.24,128.76) .. controls (516.51,138.97) and (508.15,149.23) .. (503.13,158.05) .. controls (490.43,180.32) and (477.65,198.73) .. (469.03,223.95) .. controls (466.36,231.74) and (463.07,248.52) .. (463.86,258.58) .. controls (464.22,263.26) and (465.08,267.91) .. (466.22,272.46) .. controls (466.87,275.07) and (472.03,280.21) .. (469.36,279.89) ;
%Shape: Free Drawing [id:dp4420640700184618] 
\draw  [color={rgb, 255:red, 74; green, 144; blue, 226 }  ,draw opacity=1 ][line width=2.25] [line join = round][line cap = round] (478.91,292.62) .. controls (483.63,295.74) and (482.22,300.43) .. (489.01,304.91) .. controls (492.93,307.5) and (497.49,308.98) .. (501.94,310.49) .. controls (533.16,321.08) and (543.96,266.15) .. (543.43,245.98) .. controls (542.7,218.14) and (541.13,200.87) .. (527.53,177.09) .. controls (525.7,173.88) and (518.6,165.37) .. (517.92,164.86) .. controls (513.9,161.81) and (505.27,160.82) .. (509.29,161.3) ;
%Shape: Free Drawing [id:dp7555303909696401] 
\draw  [color={rgb, 255:red, 255; green, 255; blue, 255 }  ,draw opacity=1 ][line width=2.25] [line join = round][line cap = round] (488.04,153.72) .. controls (487.73,153.83) and (487.01,153.93) .. (487.05,153.6) ;
%Shape: Free Drawing [id:dp37169312634770046] 
\draw  [draw opacity=0][line width=2.25] [line join = round][line cap = round] (653.52,92.52) .. controls (653.66,92.61) and (654.18,92.59) .. (654.02,92.57) ;
%Shape: Free Drawing [id:dp9746607410675733] 
\draw  [draw opacity=0][line width=2.25] [line join = round][line cap = round] (495.39,147.04) .. controls (492.03,148.25) and (490.49,151.28) .. (489.47,154.39) ;
%Shape: Free Drawing [id:dp1538069765158795] 
\draw  [draw opacity=0][line width=2.25] [line join = round][line cap = round] (493.66,148.85) .. controls (485.06,158.27) and (483.54,158.71) .. (486.39,159.06) ;
%Shape: Free Drawing [id:dp7648430884866061] 
\draw  [draw opacity=0][line width=2.25] [line join = round][line cap = round] (494.27,147.92) .. controls (491.62,150.69) and (488.33,153.11) .. (486.69,156.57) ;
%Shape: Free Drawing [id:dp06814306363407108] 
\draw  [draw opacity=0][line width=2.25] [line join = round][line cap = round] (473.77,280.92) .. controls (471.79,282.03) and (469.96,283.47) .. (467.83,284.24) .. controls (467.26,284.44) and (468.4,282.79) .. (469,282.87) ;
%Shape: Free Drawing [id:dp9432314319008595] 
\draw  [draw opacity=0][line width=2.25] [line join = round][line cap = round] (464.41,283.32) .. controls (468.2,281.29) and (471.46,276.62) .. (475.73,277.13) ;
%Shape: Free Drawing [id:dp868875743192955] 
\draw  [draw opacity=0][line width=2.25] [line join = round][line cap = round] (469.73,280.94) .. controls (472.4,281.26) and (474.76,279.02) .. (477.16,277.8) ;
%Shape: Free Drawing [id:dp38260252168946307] 
\draw  [draw opacity=0][line width=2.25] [line join = round][line cap = round] (467.87,279.71) .. controls (469.37,279.4) and (476.04,278.73) .. (476.35,276.19) ;
%Shape: Free Drawing [id:dp16315732791425064] 
\draw  [draw opacity=0][line width=2.25] [line join = round][line cap = round] (468.24,280.76) .. controls (471.04,279.55) and (478.19,277.6) .. (478.69,273.45) ;
%Shape: Free Drawing [id:dp9749908385719288] 
\draw  [draw opacity=0][line width=2.25] [line join = round][line cap = round] (468.13,281.75) .. controls (471.87,279.09) and (475.47,276.18) .. (478.75,272.96) ;
%Shape: Free Drawing [id:dp34001164497982994] 
\draw  [draw opacity=0][line width=2.25] [line join = round][line cap = round] (467.13,281.63) .. controls (472.99,277.64) and (476.58,273.01) .. (480.46,267.12) ;
%Shape: Circle [id:dp954115660091245] 
\draw  [color={rgb, 255:red, 74; green, 144; blue, 226 }  ,draw opacity=1 ][fill={rgb, 255:red, 255; green, 255; blue, 255 }  ,fill opacity=1 ][line width=2.25]  (484.83,50.17) .. controls (484.83,46.85) and (487.52,44.17) .. (490.83,44.17) .. controls (494.15,44.17) and (496.83,46.85) .. (496.83,50.17) .. controls (496.83,53.48) and (494.15,56.17) .. (490.83,56.17) .. controls (487.52,56.17) and (484.83,53.48) .. (484.83,50.17) -- cycle ;
%Shape: Circle [id:dp12690402095779774] 
\draw  [color={rgb, 255:red, 74; green, 144; blue, 226 }  ,draw opacity=1 ][fill={rgb, 255:red, 74; green, 144; blue, 226 }  ,fill opacity=1 ] (498.33,157.67) .. controls (498.33,154.35) and (501.02,151.67) .. (504.33,151.67) .. controls (507.65,151.67) and (510.33,154.35) .. (510.33,157.67) .. controls (510.33,160.98) and (507.65,163.67) .. (504.33,163.67) .. controls (501.02,163.67) and (498.33,160.98) .. (498.33,157.67) -- cycle ;
%Shape: Circle [id:dp06506143308802792] 
\draw  [color={rgb, 255:red, 74; green, 144; blue, 226 }  ,draw opacity=1 ][fill={rgb, 255:red, 74; green, 144; blue, 226 }  ,fill opacity=1 ] (467.83,286.67) .. controls (467.83,283.35) and (470.52,280.67) .. (473.83,280.67) .. controls (477.15,280.67) and (479.83,283.35) .. (479.83,286.67) .. controls (479.83,289.98) and (477.15,292.67) .. (473.83,292.67) .. controls (470.52,292.67) and (467.83,289.98) .. (467.83,286.67) -- cycle ;
%Shape: Circle [id:dp9939141058309052] 
\draw  [color={rgb, 255:red, 74; green, 144; blue, 226 }  ,draw opacity=1 ][fill={rgb, 255:red, 74; green, 144; blue, 226 }  ,fill opacity=1 ] (531.83,279.67) .. controls (531.83,276.35) and (534.52,273.67) .. (537.83,273.67) .. controls (541.15,273.67) and (543.83,276.35) .. (543.83,279.67) .. controls (543.83,282.98) and (541.15,285.67) .. (537.83,285.67) .. controls (534.52,285.67) and (531.83,282.98) .. (531.83,279.67) -- cycle ;
%Shape: Circle [id:dp006791002980818472] 
\draw  [color={rgb, 255:red, 74; green, 144; blue, 226 }  ,draw opacity=1 ][fill={rgb, 255:red, 255; green, 255; blue, 255 }  ,fill opacity=1 ][line width=2.25]  (498.33,157.67) .. controls (498.33,154.35) and (501.02,151.67) .. (504.33,151.67) .. controls (507.65,151.67) and (510.33,154.35) .. (510.33,157.67) .. controls (510.33,160.98) and (507.65,163.67) .. (504.33,163.67) .. controls (501.02,163.67) and (498.33,160.98) .. (498.33,157.67) -- cycle ;
%Shape: Circle [id:dp919820840892326] 
\draw  [color={rgb, 255:red, 74; green, 144; blue, 226 }  ,draw opacity=1 ][fill={rgb, 255:red, 255; green, 255; blue, 255 }  ,fill opacity=1 ][line width=2.25]  (467.83,286.67) .. controls (467.83,283.35) and (470.52,280.67) .. (473.83,280.67) .. controls (477.15,280.67) and (479.83,283.35) .. (479.83,286.67) .. controls (479.83,289.98) and (477.15,292.67) .. (473.83,292.67) .. controls (470.52,292.67) and (467.83,289.98) .. (467.83,286.67) -- cycle ;
%Shape: Circle [id:dp5154014024748038] 
\draw  [color={rgb, 255:red, 74; green, 144; blue, 226 }  ,draw opacity=1 ][fill={rgb, 255:red, 255; green, 255; blue, 255 }  ,fill opacity=1 ][line width=2.25]  (531.83,279.67) .. controls (531.83,276.35) and (534.52,273.67) .. (537.83,273.67) .. controls (541.15,273.67) and (543.83,276.35) .. (543.83,279.67) .. controls (543.83,282.98) and (541.15,285.67) .. (537.83,285.67) .. controls (534.52,285.67) and (531.83,282.98) .. (531.83,279.67) -- cycle ;
%Shape: Free Drawing [id:dp7512942112724584] 
\draw  [color={rgb, 255:red, 74; green, 144; blue, 226 }  ,draw opacity=1 ][line width=3] [line join = round][line cap = round] (469,279.67) .. controls (464.11,279.67) and (459.21,279.44) .. (454.33,279) ;
%Shape: Free Drawing [id:dp34845684588155923] 
\draw  [color={rgb, 255:red, 74; green, 144; blue, 226 }  ,draw opacity=1 ][line width=3] [line join = round][line cap = round] (469.67,277.67) .. controls (469.67,272.56) and (472.73,266.46) .. (474.33,261.67) ;
%Shape: Free Drawing [id:dp6085374109120528] 
\draw  [color={rgb, 255:red, 74; green, 144; blue, 226 }  ,draw opacity=1 ][line width=3] [line join = round][line cap = round] (484.33,274.33) .. controls (483.44,277) and (482.92,279.82) .. (481.67,282.33) .. controls (481.17,283.33) and (478.63,284.61) .. (479.67,285) .. controls (480.71,285.39) and (490.91,285.67) .. (492.33,285.67) ;
%Shape: Free Drawing [id:dp7749044227211802] 
\draw  [color={rgb, 255:red, 74; green, 144; blue, 226 }  ,draw opacity=1 ][line width=3] [line join = round][line cap = round] (455,287) .. controls (459,287) and (463.06,286.3) .. (467,287) .. controls (467.79,287.14) and (465.76,288.2) .. (465.67,289) .. controls (464.82,296.64) and (464.33,299.97) .. (464.33,306.33) ;
%Shape: Free Drawing [id:dp10346459869848679] 
\draw  [color={rgb, 255:red, 74; green, 144; blue, 226 }  ,draw opacity=1 ][line width=3] [line join = round][line cap = round] (477,303.67) .. controls (477.46,301.39) and (478.44,299.25) .. (479,297) .. controls (479.43,295.26) and (478.11,292.55) .. (479.67,291.67) .. controls (486.4,287.88) and (487.85,291.67) .. (492.33,291.67) ;
%Shape: Free Drawing [id:dp9094306621747547] 
\draw  [color={rgb, 255:red, 74; green, 144; blue, 226 }  ,draw opacity=1 ][line width=3] [line join = round][line cap = round] (498.33,34.33) .. controls (498.33,38.78) and (496.58,43.58) .. (498.33,47.67) .. controls (498.76,48.66) and (509.78,47.67) .. (510.33,47.67) ;
%Shape: Free Drawing [id:dp9371550972142336] 
\draw  [color={rgb, 255:red, 74; green, 144; blue, 226 }  ,draw opacity=1 ][line width=3] [line join = round][line cap = round] (505.67,143) .. controls (505.67,146.11) and (503.98,149.72) .. (505.67,152.33) .. controls (506.77,154.06) and (518.03,150.33) .. (518.33,150.33) ;
%Shape: Free Drawing [id:dp16012343838757115] 
\draw  [color={rgb, 255:red, 74; green, 144; blue, 226 }  ,draw opacity=1 ][line width=3] [line join = round][line cap = round] (488.33,154.33) .. controls (500.55,154.33) and (499,151.21) .. (499,141) ;
%Shape: Free Drawing [id:dp9681856106265141] 
\draw  [color={rgb, 255:red, 74; green, 144; blue, 226 }  ,draw opacity=1 ][line width=3] [line join = round][line cap = round] (525.67,157.67) .. controls (522.88,157.67) and (513.67,157.22) .. (511,159) .. controls (509.65,159.9) and (512.23,162.05) .. (512.33,163.67) .. controls (512.57,167.22) and (513,170.77) .. (513,174.33) ;

% Text Node
\draw (67.5,154.73) node [anchor=north west][inner sep=0.75pt]    {$A$};
% Text Node
\draw (230,158.73) node [anchor=north west][inner sep=0.75pt]    {$B$};
% Text Node
\draw (149.5,221.23) node [anchor=north west][inner sep=0.75pt]    {$C$};
% Text Node
\draw (12.5,282.73) node [anchor=north west][inner sep=0.75pt]    {$D$};
% Text Node
\draw (143.5,98.73) node [anchor=north west][inner sep=0.75pt]    {$E$};
% Text Node
\draw (152,284.23) node [anchor=north west][inner sep=0.75pt]    {$F$};
% Text Node
\draw (115.5,53.23) node [anchor=north west][inner sep=0.75pt]  [font=\small]  {$0$};
% Text Node
\draw (125.5,144.23) node [anchor=north west][inner sep=0.75pt]  [font=\small]  {$1$};
% Text Node
\draw (103,270.23) node [anchor=north west][inner sep=0.75pt]  [font=\small]  {$2$};
% Text Node
\draw (189.5,295.73) node [anchor=north west][inner sep=0.75pt]  [font=\small]  {$3$};
% Text Node
\draw (415.83,155.07) node [anchor=north west][inner sep=0.75pt]    {$A$};
% Text Node
\draw (578.33,159.07) node [anchor=north west][inner sep=0.75pt]    {$B$};
% Text Node
\draw (497.83,221.57) node [anchor=north west][inner sep=0.75pt]    {$C$};
% Text Node
\draw (360.83,283.07) node [anchor=north west][inner sep=0.75pt]    {$D$};
% Text Node
\draw (491.83,99.07) node [anchor=north west][inner sep=0.75pt]    {$E$};
% Text Node
\draw (500.33,284.57) node [anchor=north west][inner sep=0.75pt]    {$F$};
% Text Node
\draw (383.83,62.9) node [anchor=north west][inner sep=0.75pt]  [font=\small]  {$0$};
% Text Node
\draw (459.83,99.9) node [anchor=north west][inner sep=0.75pt]  [font=\small]  {$1$};
% Text Node
\draw (448.67,219.23) node [anchor=north west][inner sep=0.75pt]  [font=\small]  {$2$};
% Text Node
\draw (499.83,257.4) node [anchor=north west][inner sep=0.75pt]  [font=\small]  {$3$};
% Text Node
\draw (549.83,211.23) node [anchor=north west][inner sep=0.75pt]  [font=\small]  {$1$};
% Text Node
\draw (523.83,91.57) node [anchor=north west][inner sep=0.75pt]  [font=\small]  {$0$};
% Text Node
\draw (505.33,315.9) node [anchor=north west][inner sep=0.75pt]  [font=\small]  {$2$};
% Text Node
\draw (613.17,67.4) node [anchor=north west][inner sep=0.75pt]  [font=\small]  {$3$};


\end{tikzpicture}

    \caption{Paso intermedio y paso 2 para el nudo de 8}
\end{figure}
\begin{block}
    De manera similar a antes contraiga ahora cada cruce por debajo como si fuera una cuerda continua a un punto, estos serán los vértices ideales del poliedro inferior, dado por la esfera de abajo.
\end{block}
Observe que para el poliedro de abajo lo que para nosotros desde el diagrama es una cuerda que pasa por debajo, para este sera una cuerda que pasa por encima, por eso ahora contraemos los cruces así
\begin{figure}[H]
    \centering
    
%!TEX root = ./main.tex



\tikzset{every picture/.style={line width=0.75pt}} %set default line width to 0.75pt        

\begin{tikzpicture}[x=0.75pt,y=0.75pt,yscale=-0.85,xscale=0.85]
%uncomment if require: \path (0,877); %set diagram left start at 0, and has height of 877

%Shape: Free Drawing [id:dp052719054339937954] 
\draw  [line width=2.25] [line join = round][line cap = round] (144.39,147.9) .. controls (134.01,146.65) and (135.03,129.11) .. (131.73,119.19) .. controls (122.82,92.44) and (130.17,60.08) .. (151.63,41.51) .. controls (156.44,37.34) and (165.14,31.96) .. (171.87,32.35) .. controls (286.8,39.15) and (348.19,216.97) .. (264.67,295.3) .. controls (254.46,304.88) and (238.91,306.9) .. (226.98,298.83) .. controls (217.6,292.48) and (209.15,276.48) .. (198.49,280.3) ;
%Shape: Free Drawing [id:dp25610451731952133] 
\draw  [line width=2.25] [line join = round][line cap = round] (179.25,276.98) .. controls (166.14,268.33) and (132.87,282.35) .. (120.45,288.55) .. controls (117.49,290.03) and (115.35,293.03) .. (112.22,294.11) .. controls (75.11,306.93) and (46.64,288.67) .. (28.87,254.38) .. controls (22.39,241.88) and (17.62,227.89) .. (17.14,213.69) .. controls (15.61,168.88) and (26.69,100.32) .. (60.58,65.82) .. controls (74.17,52) and (97.46,39.75) .. (117.26,38.38) .. controls (124.59,37.88) and (139.28,46.06) .. (136.03,45.67) ;
%Shape: Free Drawing [id:dp5679641581314395] 
\draw  [line width=2.25] [line join = round][line cap = round] (151.08,54.53) .. controls (166.66,64.82) and (166.98,83.56) .. (169.22,100.52) .. controls (170.38,109.27) and (173.36,119.25) .. (170.9,128.42) .. controls (168.17,138.64) and (159.82,148.9) .. (154.8,157.71) .. controls (142.1,179.99) and (129.32,198.39) .. (120.69,223.62) .. controls (118.03,231.41) and (114.74,248.18) .. (115.53,258.25) .. controls (115.89,262.93) and (116.75,267.57) .. (117.89,272.13) .. controls (118.54,274.73) and (123.69,279.87) .. (121.02,279.55) ;
%Shape: Free Drawing [id:dp6189964124099764] 
\draw  [line width=2.25] [line join = round][line cap = round] (130.57,292.28) .. controls (135.3,295.4) and (133.89,300.1) .. (140.68,304.58) .. controls (144.6,307.16) and (149.16,308.65) .. (153.61,310.16) .. controls (184.83,320.75) and (195.62,265.82) .. (195.1,245.65) .. controls (194.37,217.8) and (192.79,200.53) .. (179.2,176.76) .. controls (177.36,173.55) and (170.27,165.04) .. (169.59,164.53) .. controls (165.57,161.48) and (156.94,160.49) .. (160.95,160.97) ;
%Shape: Free Drawing [id:dp23463200731808187] 
\draw  [color={rgb, 255:red, 255; green, 255; blue, 255 }  ,draw opacity=1 ][line width=2.25] [line join = round][line cap = round] (139.71,153.38) .. controls (139.39,153.49) and (138.67,153.59) .. (138.71,153.26) ;
%Shape: Free Drawing [id:dp6791599529562201] 
\draw  [draw opacity=0][line width=2.25] [line join = round][line cap = round] (305.19,92.18) .. controls (305.33,92.27) and (305.85,92.26) .. (305.68,92.24) ;
%Shape: Free Drawing [id:dp4292482997182422] 
\draw  [draw opacity=0][line width=2.25] [line join = round][line cap = round] (147.05,146.71) .. controls (143.7,147.91) and (142.16,150.95) .. (141.14,154.06) ;
%Shape: Free Drawing [id:dp6085326410058335] 
\draw  [draw opacity=0][line width=2.25] [line join = round][line cap = round] (145.33,148.52) .. controls (136.73,157.93) and (135.21,158.38) .. (138.06,158.72) ;
%Shape: Free Drawing [id:dp9043644710987669] 
\draw  [draw opacity=0][line width=2.25] [line join = round][line cap = round] (145.94,147.58) .. controls (143.29,150.36) and (139.99,152.77) .. (138.35,156.24) ;
%Shape: Free Drawing [id:dp21505097221297298] 
\draw  [draw opacity=0][line width=2.25] [line join = round][line cap = round] (125.43,280.59) .. controls (123.45,281.69) and (121.63,283.14) .. (119.49,283.9) .. controls (118.93,284.11) and (120.07,282.46) .. (120.67,282.53) ;
%Shape: Free Drawing [id:dp28830093222917674] 
\draw  [draw opacity=0][line width=2.25] [line join = round][line cap = round] (116.08,282.99) .. controls (119.87,280.96) and (123.13,276.28) .. (127.4,276.79) ;
%Shape: Free Drawing [id:dp9547251910894086] 
\draw  [draw opacity=0][line width=2.25] [line join = round][line cap = round] (121.4,280.61) .. controls (124.07,280.93) and (126.43,278.68) .. (128.83,277.47) ;
%Shape: Free Drawing [id:dp719981311992137] 
\draw  [draw opacity=0][line width=2.25] [line join = round][line cap = round] (119.53,279.38) .. controls (121.04,279.07) and (127.71,278.4) .. (128.01,275.86) ;
%Shape: Free Drawing [id:dp29834365303629773] 
\draw  [draw opacity=0][line width=2.25] [line join = round][line cap = round] (119.91,280.43) .. controls (122.7,279.22) and (129.86,277.27) .. (130.36,273.12) ;
%Shape: Free Drawing [id:dp42617767268148665] 
\draw  [draw opacity=0][line width=2.25] [line join = round][line cap = round] (119.79,281.42) .. controls (123.54,278.76) and (127.14,275.85) .. (130.42,272.63) ;
%Shape: Free Drawing [id:dp5503601990364082] 
\draw  [draw opacity=0][line width=2.25] [line join = round][line cap = round] (118.8,281.3) .. controls (124.65,277.3) and (128.24,272.67) .. (132.13,266.79) ;
%Shape: Free Drawing [id:dp12708063524714575] 
\draw  [draw opacity=0][line width=2.25] [line join = round][line cap = round] (15,81.33) .. controls (14.37,81.02) and (13.82,80.47) .. (13.5,79.83) ;
%Straight Lines [id:da4904183098904815] 
\draw [color={rgb, 255:red, 74; green, 144; blue, 226 }  ,draw opacity=1 ][line width=1.5]    (165,31.67) -- (144.83,47.95) ;
\draw [shift={(142.5,49.83)}, rotate = 321.08] [color={rgb, 255:red, 74; green, 144; blue, 226 }  ,draw opacity=1 ][line width=1.5]    (14.21,-4.28) .. controls (9.04,-1.82) and (4.3,-0.39) .. (0,0) .. controls (4.3,0.39) and (9.04,1.82) .. (14.21,4.28)   ;
%Straight Lines [id:da21498342913887913] 
\draw [color={rgb, 255:red, 74; green, 144; blue, 226 }  ,draw opacity=1 ][line width=1.5]    (129,79.67) -- (141.26,52.57) ;
\draw [shift={(142.5,49.83)}, rotate = 114.35] [color={rgb, 255:red, 74; green, 144; blue, 226 }  ,draw opacity=1 ][line width=1.5]    (14.21,-4.28) .. controls (9.04,-1.82) and (4.3,-0.39) .. (0,0) .. controls (4.3,0.39) and (9.04,1.82) .. (14.21,4.28)   ;
%Straight Lines [id:da6750405427654139] 
\draw [color={rgb, 255:red, 74; green, 144; blue, 226 }  ,draw opacity=1 ][line width=1.5]    (152.5,155.83) -- (170.11,126.9) ;
\draw [shift={(171.67,124.33)}, rotate = 121.32] [color={rgb, 255:red, 74; green, 144; blue, 226 }  ,draw opacity=1 ][line width=1.5]    (14.21,-4.28) .. controls (9.04,-1.82) and (4.3,-0.39) .. (0,0) .. controls (4.3,0.39) and (9.04,1.82) .. (14.21,4.28)   ;
%Straight Lines [id:da9808570839490202] 
\draw [color={rgb, 255:red, 74; green, 144; blue, 226 }  ,draw opacity=1 ][line width=1.5]    (152.5,155.83) -- (139.05,181.67) ;
\draw [shift={(137.67,184.33)}, rotate = 297.5] [color={rgb, 255:red, 74; green, 144; blue, 226 }  ,draw opacity=1 ][line width=1.5]    (14.21,-4.28) .. controls (9.04,-1.82) and (4.3,-0.39) .. (0,0) .. controls (4.3,0.39) and (9.04,1.82) .. (14.21,4.28)   ;
%Straight Lines [id:da8295997211316378] 
\draw [color={rgb, 255:red, 74; green, 144; blue, 226 }  ,draw opacity=1 ][line width=1.5]    (122.83,283.67) -- (94.36,297.68) ;
\draw [shift={(91.67,299)}, rotate = 333.8] [color={rgb, 255:red, 74; green, 144; blue, 226 }  ,draw opacity=1 ][line width=1.5]    (14.21,-4.28) .. controls (9.04,-1.82) and (4.3,-0.39) .. (0,0) .. controls (4.3,0.39) and (9.04,1.82) .. (14.21,4.28)   ;
%Straight Lines [id:da4207009837534177] 
\draw [color={rgb, 255:red, 74; green, 144; blue, 226 }  ,draw opacity=1 ][line width=1.5]    (125.5,286.33) -- (150.18,277.36) ;
\draw [shift={(153,276.33)}, rotate = 160.02] [color={rgb, 255:red, 74; green, 144; blue, 226 }  ,draw opacity=1 ][line width=1.5]    (14.21,-4.28) .. controls (9.04,-1.82) and (4.3,-0.39) .. (0,0) .. controls (4.3,0.39) and (9.04,1.82) .. (14.21,4.28)   ;
%Straight Lines [id:da7105295943160274] 
\draw [color={rgb, 255:red, 74; green, 144; blue, 226 }  ,draw opacity=1 ][line width=1.5]    (197,243.67) -- (191.48,277.37) ;
\draw [shift={(191,280.33)}, rotate = 279.29] [color={rgb, 255:red, 74; green, 144; blue, 226 }  ,draw opacity=1 ][line width=1.5]    (14.21,-4.28) .. controls (9.04,-1.82) and (4.3,-0.39) .. (0,0) .. controls (4.3,0.39) and (9.04,1.82) .. (14.21,4.28)   ;
%Straight Lines [id:da5651565239928864] 
\draw [color={rgb, 255:red, 74; green, 144; blue, 226 }  ,draw opacity=1 ][line width=1.5]    (173.67,310.33) -- (188.14,282.01) ;
\draw [shift={(189.5,279.33)}, rotate = 117.06] [color={rgb, 255:red, 74; green, 144; blue, 226 }  ,draw opacity=1 ][line width=1.5]    (14.21,-4.28) .. controls (9.04,-1.82) and (4.3,-0.39) .. (0,0) .. controls (4.3,0.39) and (9.04,1.82) .. (14.21,4.28)   ;
%Shape: Circle [id:dp7563944317650207] 
\draw  [color={rgb, 255:red, 74; green, 144; blue, 226 }  ,draw opacity=1 ][fill={rgb, 255:red, 74; green, 144; blue, 226 }  ,fill opacity=1 ] (136.5,49.83) .. controls (136.5,46.52) and (139.19,43.83) .. (142.5,43.83) .. controls (145.81,43.83) and (148.5,46.52) .. (148.5,49.83) .. controls (148.5,53.15) and (145.81,55.83) .. (142.5,55.83) .. controls (139.19,55.83) and (136.5,53.15) .. (136.5,49.83) -- cycle ;
%Shape: Circle [id:dp8058641002050105] 
\draw  [color={rgb, 255:red, 74; green, 144; blue, 226 }  ,draw opacity=1 ][fill={rgb, 255:red, 74; green, 144; blue, 226 }  ,fill opacity=1 ] (146.5,155.83) .. controls (146.5,152.52) and (149.19,149.83) .. (152.5,149.83) .. controls (155.81,149.83) and (158.5,152.52) .. (158.5,155.83) .. controls (158.5,159.15) and (155.81,161.83) .. (152.5,161.83) .. controls (149.19,161.83) and (146.5,159.15) .. (146.5,155.83) -- cycle ;
%Shape: Circle [id:dp3830896703653248] 
\draw  [color={rgb, 255:red, 74; green, 144; blue, 226 }  ,draw opacity=1 ][fill={rgb, 255:red, 74; green, 144; blue, 226 }  ,fill opacity=1 ] (119.5,286.33) .. controls (119.5,283.02) and (122.19,280.33) .. (125.5,280.33) .. controls (128.81,280.33) and (131.5,283.02) .. (131.5,286.33) .. controls (131.5,289.65) and (128.81,292.33) .. (125.5,292.33) .. controls (122.19,292.33) and (119.5,289.65) .. (119.5,286.33) -- cycle ;
%Shape: Circle [id:dp568354357328672] 
\draw  [color={rgb, 255:red, 74; green, 144; blue, 226 }  ,draw opacity=1 ][fill={rgb, 255:red, 74; green, 144; blue, 226 }  ,fill opacity=1 ] (183.5,279.33) .. controls (183.5,276.02) and (186.19,273.33) .. (189.5,273.33) .. controls (192.81,273.33) and (195.5,276.02) .. (195.5,279.33) .. controls (195.5,282.65) and (192.81,285.33) .. (189.5,285.33) .. controls (186.19,285.33) and (183.5,282.65) .. (183.5,279.33) -- cycle ;
%Shape: Free Drawing [id:dp4081098483122335] 
\draw  [color={rgb, 255:red, 74; green, 144; blue, 226 }  ,draw opacity=1 ][line width=2.25] [line join = round][line cap = round] (492.73,148.24) .. controls (482.34,146.99) and (483.37,129.44) .. (480.06,119.52) .. controls (471.15,92.78) and (478.5,60.42) .. (499.97,41.84) .. controls (504.78,37.68) and (513.47,32.29) .. (520.2,32.69) .. controls (635.13,39.48) and (696.52,217.31) .. (613,295.63) .. controls (602.79,305.21) and (587.24,307.23) .. (575.31,299.16) .. controls (565.93,292.82) and (557.49,276.81) .. (546.82,280.63) ;
%Shape: Free Drawing [id:dp3946430254773102] 
\draw  [color={rgb, 255:red, 74; green, 144; blue, 226 }  ,draw opacity=1 ][line width=2.25] [line join = round][line cap = round] (527.58,277.31) .. controls (514.47,268.66) and (481.2,282.69) .. (468.78,288.88) .. controls (465.82,290.36) and (463.69,293.36) .. (460.56,294.44) .. controls (423.44,307.26) and (394.97,289.01) .. (377.2,254.72) .. controls (370.72,242.22) and (365.96,228.23) .. (365.47,214.03) .. controls (363.94,169.21) and (375.02,100.65) .. (408.92,66.15) .. controls (422.5,52.33) and (445.8,40.08) .. (465.59,38.72) .. controls (472.93,38.21) and (487.61,46.4) .. (484.36,46.01) ;
%Shape: Free Drawing [id:dp031095390456670646] 
\draw  [color={rgb, 255:red, 74; green, 144; blue, 226 }  ,draw opacity=1 ][line width=2.25] [line join = round][line cap = round] (499.41,54.87) .. controls (515,65.15) and (515.31,83.9) .. (517.55,100.86) .. controls (518.71,109.6) and (521.69,119.59) .. (519.24,128.76) .. controls (516.51,138.97) and (508.15,149.23) .. (503.13,158.05) .. controls (490.43,180.32) and (477.65,198.73) .. (469.03,223.95) .. controls (466.36,231.74) and (463.07,248.52) .. (463.86,258.58) .. controls (464.22,263.26) and (465.08,267.91) .. (466.22,272.46) .. controls (466.87,275.07) and (472.03,280.21) .. (469.36,279.89) ;
%Shape: Free Drawing [id:dp4420640700184618] 
\draw  [color={rgb, 255:red, 74; green, 144; blue, 226 }  ,draw opacity=1 ][line width=2.25] [line join = round][line cap = round] (478.91,292.62) .. controls (483.63,295.74) and (482.22,300.43) .. (489.01,304.91) .. controls (492.93,307.5) and (497.49,308.98) .. (501.94,310.49) .. controls (533.16,321.08) and (543.96,266.15) .. (543.43,245.98) .. controls (542.7,218.14) and (541.13,200.87) .. (527.53,177.09) .. controls (525.7,173.88) and (518.6,165.37) .. (517.92,164.86) .. controls (513.9,161.81) and (505.27,160.82) .. (509.29,161.3) ;
%Shape: Free Drawing [id:dp7555303909696401] 
\draw  [color={rgb, 255:red, 255; green, 255; blue, 255 }  ,draw opacity=1 ][line width=2.25] [line join = round][line cap = round] (488.04,153.72) .. controls (487.73,153.83) and (487.01,153.93) .. (487.05,153.6) ;
%Shape: Free Drawing [id:dp37169312634770046] 
\draw  [draw opacity=0][line width=2.25] [line join = round][line cap = round] (653.52,92.52) .. controls (653.66,92.61) and (654.18,92.59) .. (654.02,92.57) ;
%Shape: Free Drawing [id:dp9746607410675733] 
\draw  [draw opacity=0][line width=2.25] [line join = round][line cap = round] (495.39,147.04) .. controls (492.03,148.25) and (490.49,151.28) .. (489.47,154.39) ;
%Shape: Free Drawing [id:dp1538069765158795] 
\draw  [draw opacity=0][line width=2.25] [line join = round][line cap = round] (493.66,148.85) .. controls (485.06,158.27) and (483.54,158.71) .. (486.39,159.06) ;
%Shape: Free Drawing [id:dp7648430884866061] 
\draw  [draw opacity=0][line width=2.25] [line join = round][line cap = round] (494.27,147.92) .. controls (491.62,150.69) and (488.33,153.11) .. (486.69,156.57) ;
%Shape: Free Drawing [id:dp06814306363407108] 
\draw  [draw opacity=0][line width=2.25] [line join = round][line cap = round] (473.77,280.92) .. controls (471.79,282.03) and (469.96,283.47) .. (467.83,284.24) .. controls (467.26,284.44) and (468.4,282.79) .. (469,282.87) ;
%Shape: Free Drawing [id:dp9432314319008595] 
\draw  [draw opacity=0][line width=2.25] [line join = round][line cap = round] (464.41,283.32) .. controls (468.2,281.29) and (471.46,276.62) .. (475.73,277.13) ;
%Shape: Free Drawing [id:dp868875743192955] 
\draw  [draw opacity=0][line width=2.25] [line join = round][line cap = round] (469.73,280.94) .. controls (472.4,281.26) and (474.76,279.02) .. (477.16,277.8) ;
%Shape: Free Drawing [id:dp38260252168946307] 
\draw  [draw opacity=0][line width=2.25] [line join = round][line cap = round] (467.87,279.71) .. controls (469.37,279.4) and (476.04,278.73) .. (476.35,276.19) ;
%Shape: Free Drawing [id:dp16315732791425064] 
\draw  [draw opacity=0][line width=2.25] [line join = round][line cap = round] (468.24,280.76) .. controls (471.04,279.55) and (478.19,277.6) .. (478.69,273.45) ;
%Shape: Free Drawing [id:dp9749908385719288] 
\draw  [draw opacity=0][line width=2.25] [line join = round][line cap = round] (468.13,281.75) .. controls (471.87,279.09) and (475.47,276.18) .. (478.75,272.96) ;
%Shape: Free Drawing [id:dp34001164497982994] 
\draw  [draw opacity=0][line width=2.25] [line join = round][line cap = round] (467.13,281.63) .. controls (472.99,277.64) and (476.58,273.01) .. (480.46,267.12) ;
%Shape: Circle [id:dp954115660091245] 
\draw  [color={rgb, 255:red, 74; green, 144; blue, 226 }  ,draw opacity=1 ][fill={rgb, 255:red, 255; green, 255; blue, 255 }  ,fill opacity=1 ][line width=2.25]  (484.83,50.17) .. controls (484.83,46.85) and (487.52,44.17) .. (490.83,44.17) .. controls (494.15,44.17) and (496.83,46.85) .. (496.83,50.17) .. controls (496.83,53.48) and (494.15,56.17) .. (490.83,56.17) .. controls (487.52,56.17) and (484.83,53.48) .. (484.83,50.17) -- cycle ;
%Shape: Circle [id:dp12690402095779774] 
\draw  [color={rgb, 255:red, 74; green, 144; blue, 226 }  ,draw opacity=1 ][fill={rgb, 255:red, 74; green, 144; blue, 226 }  ,fill opacity=1 ] (498.33,157.67) .. controls (498.33,154.35) and (501.02,151.67) .. (504.33,151.67) .. controls (507.65,151.67) and (510.33,154.35) .. (510.33,157.67) .. controls (510.33,160.98) and (507.65,163.67) .. (504.33,163.67) .. controls (501.02,163.67) and (498.33,160.98) .. (498.33,157.67) -- cycle ;
%Shape: Circle [id:dp06506143308802792] 
\draw  [color={rgb, 255:red, 74; green, 144; blue, 226 }  ,draw opacity=1 ][fill={rgb, 255:red, 74; green, 144; blue, 226 }  ,fill opacity=1 ] (467.83,286.67) .. controls (467.83,283.35) and (470.52,280.67) .. (473.83,280.67) .. controls (477.15,280.67) and (479.83,283.35) .. (479.83,286.67) .. controls (479.83,289.98) and (477.15,292.67) .. (473.83,292.67) .. controls (470.52,292.67) and (467.83,289.98) .. (467.83,286.67) -- cycle ;
%Shape: Circle [id:dp9939141058309052] 
\draw  [color={rgb, 255:red, 74; green, 144; blue, 226 }  ,draw opacity=1 ][fill={rgb, 255:red, 74; green, 144; blue, 226 }  ,fill opacity=1 ] (531.83,279.67) .. controls (531.83,276.35) and (534.52,273.67) .. (537.83,273.67) .. controls (541.15,273.67) and (543.83,276.35) .. (543.83,279.67) .. controls (543.83,282.98) and (541.15,285.67) .. (537.83,285.67) .. controls (534.52,285.67) and (531.83,282.98) .. (531.83,279.67) -- cycle ;
%Shape: Circle [id:dp006791002980818472] 
\draw  [color={rgb, 255:red, 74; green, 144; blue, 226 }  ,draw opacity=1 ][fill={rgb, 255:red, 255; green, 255; blue, 255 }  ,fill opacity=1 ][line width=2.25]  (498.33,157.67) .. controls (498.33,154.35) and (501.02,151.67) .. (504.33,151.67) .. controls (507.65,151.67) and (510.33,154.35) .. (510.33,157.67) .. controls (510.33,160.98) and (507.65,163.67) .. (504.33,163.67) .. controls (501.02,163.67) and (498.33,160.98) .. (498.33,157.67) -- cycle ;
%Shape: Circle [id:dp919820840892326] 
\draw  [color={rgb, 255:red, 74; green, 144; blue, 226 }  ,draw opacity=1 ][fill={rgb, 255:red, 255; green, 255; blue, 255 }  ,fill opacity=1 ][line width=2.25]  (467.83,286.67) .. controls (467.83,283.35) and (470.52,280.67) .. (473.83,280.67) .. controls (477.15,280.67) and (479.83,283.35) .. (479.83,286.67) .. controls (479.83,289.98) and (477.15,292.67) .. (473.83,292.67) .. controls (470.52,292.67) and (467.83,289.98) .. (467.83,286.67) -- cycle ;
%Shape: Circle [id:dp5154014024748038] 
\draw  [color={rgb, 255:red, 74; green, 144; blue, 226 }  ,draw opacity=1 ][fill={rgb, 255:red, 255; green, 255; blue, 255 }  ,fill opacity=1 ][line width=2.25]  (531.83,279.67) .. controls (531.83,276.35) and (534.52,273.67) .. (537.83,273.67) .. controls (541.15,273.67) and (543.83,276.35) .. (543.83,279.67) .. controls (543.83,282.98) and (541.15,285.67) .. (537.83,285.67) .. controls (534.52,285.67) and (531.83,282.98) .. (531.83,279.67) -- cycle ;
%Shape: Free Drawing [id:dp9094306621747547] 
\draw  [color={rgb, 255:red, 74; green, 144; blue, 226 }  ,draw opacity=1 ][line width=3] [line join = round][line cap = round] (498.33,34.33) .. controls (498.33,38.78) and (496.58,43.58) .. (498.33,47.67) .. controls (498.76,48.66) and (509.78,47.67) .. (510.33,47.67) ;
%Shape: Free Drawing [id:dp09473093158323864] 
\draw  [color={rgb, 255:red, 74; green, 144; blue, 226 }  ,draw opacity=1 ][line width=3] [line join = round][line cap = round] (490.34,68.66) .. controls (489.66,64.27) and (490.65,59.25) .. (488.29,55.48) .. controls (487.71,54.57) and (476.98,57.25) .. (476.43,57.33) ;
%Shape: Free Drawing [id:dp9521758844656053] 
\draw  [color={rgb, 255:red, 74; green, 144; blue, 226 }  ,draw opacity=1 ][line width=3] [line join = round][line cap = round] (511.27,53.24) .. controls (506.85,53.75) and (501.88,52.56) .. (498.02,54.77) .. controls (497.08,55.3) and (499.33,66.13) .. (499.4,66.69) ;
%Shape: Free Drawing [id:dp10884727512660441] 
\draw  [color={rgb, 255:red, 74; green, 144; blue, 226 }  ,draw opacity=1 ][line width=3] [line join = round][line cap = round] (468.47,46.51) .. controls (472.9,46.18) and (477.82,47.56) .. (481.77,45.51) .. controls (482.72,45.01) and (480.91,34.09) .. (480.86,33.54) ;
%Shape: Free Drawing [id:dp39512480087630053] 
\draw  [color={rgb, 255:red, 74; green, 144; blue, 226 }  ,draw opacity=1 ][line width=3] [line join = round][line cap = round] (456.94,276.69) .. controls (461.34,277.28) and (465.87,279.65) .. (470.15,278.45) .. controls (471.19,278.16) and (471.66,267.11) .. (471.74,266.55) ;
%Shape: Free Drawing [id:dp3019270320594637] 
\draw  [color={rgb, 255:red, 74; green, 144; blue, 226 }  ,draw opacity=1 ][line width=3] [line join = round][line cap = round] (532.71,260.5) .. controls (535.31,264.1) and (536.7,269.02) .. (540.51,271.31) .. controls (541.43,271.87) and (549.79,264.62) .. (550.24,264.29) ;
%Shape: Free Drawing [id:dp14991459054603584] 
\draw  [color={rgb, 255:red, 74; green, 144; blue, 226 }  ,draw opacity=1 ][line width=3] [line join = round][line cap = round] (516.28,283.48) .. controls (520.04,281.1) and (525.04,280.02) .. (527.56,276.36) .. controls (528.17,275.47) and (521.45,266.68) .. (521.16,266.21) ;
%Shape: Free Drawing [id:dp6903073615923814] 
\draw  [color={rgb, 255:red, 74; green, 144; blue, 226 }  ,draw opacity=1 ][line width=3] [line join = round][line cap = round] (538.15,300.38) .. controls (537.1,296.07) and (537.67,290.98) .. (535,287.43) .. controls (534.35,286.56) and (523.88,290.13) .. (523.34,290.26) ;

% Text Node
\draw (67.5,154.73) node [anchor=north west][inner sep=0.75pt]    {$A$};
% Text Node
\draw (230,158.73) node [anchor=north west][inner sep=0.75pt]    {$B$};
% Text Node
\draw (149.5,221.23) node [anchor=north west][inner sep=0.75pt]    {$C$};
% Text Node
\draw (12.5,282.73) node [anchor=north west][inner sep=0.75pt]    {$D$};
% Text Node
\draw (143.5,98.73) node [anchor=north west][inner sep=0.75pt]    {$E$};
% Text Node
\draw (152,284.23) node [anchor=north west][inner sep=0.75pt]    {$F$};
% Text Node
\draw (115.5,53.23) node [anchor=north west][inner sep=0.75pt]  [font=\small]  {$0$};
% Text Node
\draw (130.83,146.23) node [anchor=north west][inner sep=0.75pt]  [font=\small]  {$1$};
% Text Node
\draw (103,270.23) node [anchor=north west][inner sep=0.75pt]  [font=\small]  {$2$};
% Text Node
\draw (189.5,295.73) node [anchor=north west][inner sep=0.75pt]  [font=\small]  {$3$};
% Text Node
\draw (415.83,155.07) node [anchor=north west][inner sep=0.75pt]    {$A$};
% Text Node
\draw (578.33,159.07) node [anchor=north west][inner sep=0.75pt]    {$B$};
% Text Node
\draw (497.83,221.57) node [anchor=north west][inner sep=0.75pt]    {$C$};
% Text Node
\draw (360.83,283.07) node [anchor=north west][inner sep=0.75pt]    {$D$};
% Text Node
\draw (491.83,99.07) node [anchor=north west][inner sep=0.75pt]    {$E$};
% Text Node
\draw (500.33,284.57) node [anchor=north west][inner sep=0.75pt]    {$F$};
% Text Node
\draw (603.17,52.9) node [anchor=north west][inner sep=0.75pt]  [font=\small]  {$0$};
% Text Node
\draw (523.83,90.57) node [anchor=north west][inner sep=0.75pt]  [font=\small]  {$1$};
% Text Node
\draw (501.33,257.9) node [anchor=north west][inner sep=0.75pt]  [font=\small]  {$2$};
% Text Node
\draw (547.83,225.4) node [anchor=north west][inner sep=0.75pt]  [font=\small]  {$3$};
% Text Node
\draw (449.17,229.23) node [anchor=north west][inner sep=0.75pt]  [font=\small]  {$1$};
% Text Node
\draw (459.83,89.57) node [anchor=north west][inner sep=0.75pt]  [font=\small]  {$0$};
% Text Node
\draw (386,62.57) node [anchor=north west][inner sep=0.75pt]  [font=\small]  {$2$};
% Text Node
\draw (502.5,317.4) node [anchor=north west][inner sep=0.75pt]  [font=\small]  {$3$};


\end{tikzpicture}



    \caption{Paso intermedio y paso 3 para el nudo de 8}
\end{figure}
En este caso vemos la esfera desde afuera por arriba como es la visualización natural de un poliedro desde la parte exterior.\\

De esta forma obtenemos la descompocision en dos poliedros ideales bajo nuestra definición del mismo. Es fácil notar que la asignación de caras y aristas de pegado que nos muestran ambos diagramas recuperan el complemento del nudo inicial, simplemente coloque la cara $A$ de uno sobre la otra y pegue ambas asegurándose que las aristas coincidan y lo mismo para las demás, por la construcción misma el pegado da el diagrama obtenido en el primer paso, salvo que en vez de tener dibujadas las 4 aristas isotopicas solo hay una.

\begin{note}
    En general en un poliedro usual con una definición de formas clásicas evitaríamos caras como la $E$ y $F$ donde solo hay dos lados, ya que lo usual es que nuestras caras básicas sean triángulos, para obtener algo como esto observe que por ejemplo en el diagrama del paso 1 si uno en la cara $F$ desliza la arista 2 a través de la región puede llegar de manera isotopica a la arista 3.
    \begin{figure}[H]
    \centering
    
%!TEX root = ./main.tex



\tikzset{every picture/.style={line width=0.75pt}} %set default line width to 0.75pt        

\begin{tikzpicture}[x=0.75pt,y=0.75pt,yscale=-1,xscale=1]
%uncomment if require: \path (0,877); %set diagram left start at 0, and has height of 877

%Curve Lines [id:da19609189587014564] 
\draw [line width=1.5]    (195.29,140) .. controls (237.61,27.33) and (303.88,0.67) .. (394.12,60) ;
%Straight Lines [id:da2771116463618114] 
\draw [line width=1.5]    (180,40) -- (210.59,80) ;
%Curve Lines [id:da5075059974941536] 
\draw [line width=1.5]    (225.88,100) .. controls (315.1,191.33) and (377.29,100.67) .. (440,20) ;
%Straight Lines [id:da01404980269438516] 
\draw [line width=1.5]    (409.41,80) -- (440,120) ;
%Straight Lines [id:da6027225475778994] 
\draw [color={rgb, 255:red, 74; green, 144; blue, 226 }  ,draw opacity=1 ][line width=1.5]    (250,120) -- (240.49,62.96) ;
\draw [shift={(240,60)}, rotate = 80.54] [color={rgb, 255:red, 74; green, 144; blue, 226 }  ,draw opacity=1 ][line width=1.5]    (14.21,-4.28) .. controls (9.04,-1.82) and (4.3,-0.39) .. (0,0) .. controls (4.3,0.39) and (9.04,1.82) .. (14.21,4.28)   ;
%Straight Lines [id:da602091940582576] 
\draw [color={rgb, 255:red, 74; green, 144; blue, 226 }  ,draw opacity=1 ][line width=1.5]    (370,100) -- (370,53) ;
\draw [shift={(370,50)}, rotate = 90] [color={rgb, 255:red, 74; green, 144; blue, 226 }  ,draw opacity=1 ][line width=1.5]    (14.21,-4.28) .. controls (9.04,-1.82) and (4.3,-0.39) .. (0,0) .. controls (4.3,0.39) and (9.04,1.82) .. (14.21,4.28)   ;
%Straight Lines [id:da7589682295521507] 
\draw [color={rgb, 255:red, 74; green, 144; blue, 226 }  ,draw opacity=1 ][line width=1.5]    (310,140) -- (310,33) ;
\draw [shift={(310,30)}, rotate = 90] [color={rgb, 255:red, 74; green, 144; blue, 226 }  ,draw opacity=1 ][line width=1.5]    (14.21,-4.28) .. controls (9.04,-1.82) and (4.3,-0.39) .. (0,0) .. controls (4.3,0.39) and (9.04,1.82) .. (14.21,4.28)   ;

% Text Node
\draw (281,72.4) node [anchor=north west][inner sep=0.75pt]    {$F$};
% Text Node
\draw (247,82.4) node [anchor=north west][inner sep=0.75pt]    {$2$};
% Text Node
\draw (357,72.4) node [anchor=north west][inner sep=0.75pt]    {$3$};


\end{tikzpicture}
        \caption{Isotopia de la Region F}
    \end{figure}
    En general en regiones así siempre las aristas entre cruces serán isotopicas, luego basta con tener una para tener toda la información, por lo que estas regiones pueden colapsarse en los diagramas finales a solo una arista y eliminar la cara completamente obteniendo la presentación dada por Thurston original, dos tetraedros, en este caso aplanados y uno visto desde adentro, donde la arista 0 se identifica con la 1 y la arista 2 con la 3, note ademas que cambiamos la numeración por lineas en las aristas, esto sera útil luego para no sobrecargar el dibujo.
    \begin{figure}[H]
    \centering
    \input{tetra8}
       \caption{Tetraedros del nudo de 8 luego de contraer las regiones E y F.} 
    \end{figure}
\end{note}
\subsection{Otros ejemplos de descompocision y una vista a la generalización}
El proceso realizado antes funciona muy bien tanto para enlaces como para nudos, para mostrar esto haremos uno de cada uno
\begin{eg}
    Considere la proyección del enlace de Whitehead estándar con regiones y aristas ya identificadas
    \begin{figure}[H]
    \centering
    
%!TEX root = ./main.tex

\tikzset{every picture/.style={line width=0.75pt}} %set default line width to 0.75pt        

\begin{tikzpicture}[x=0.75pt,y=0.75pt,yscale=-1,xscale=1]
%uncomment if require: \path (0,877); %set diagram left start at 0, and has height of 877

%Curve Lines [id:da1815470268489796] 
\draw [line width=1.5]    (190,150) .. controls (173.67,-40.33) and (360.33,-31) .. (360,90) ;
%Curve Lines [id:da1886044146181638] 
\draw [line width=1.5]    (190,180) .. controls (197,281) and (361,301.67) .. (360,120) ;
%Curve Lines [id:da5499420166410015] 
\draw [line width=1.5]    (270,150) .. controls (163.67,198.33) and (4.33,110.33) .. (180,100) ;
%Curve Lines [id:da628102200884786] 
\draw [line width=1.5]    (290,130) .. controls (468.33,29) and (552.33,154.33) .. (380,160) ;
%Curve Lines [id:da6746806126730861] 
\draw [line width=1.5]    (200,100) .. controls (269.67,101.67) and (281,177.67) .. (340,160) ;
%Straight Lines [id:da6086058360131362] 
\draw [color={rgb, 255:red, 74; green, 144; blue, 226 }  ,draw opacity=1 ][line width=1.5]    (200,60) -- (209.27,97.09) ;
\draw [shift={(210,100)}, rotate = 255.96] [color={rgb, 255:red, 74; green, 144; blue, 226 }  ,draw opacity=1 ][line width=1.5]    (14.21,-4.28) .. controls (9.04,-1.82) and (4.3,-0.39) .. (0,0) .. controls (4.3,0.39) and (9.04,1.82) .. (14.21,4.28)   ;
%Straight Lines [id:da014067308052953753] 
\draw [color={rgb, 255:red, 74; green, 144; blue, 226 }  ,draw opacity=1 ][line width=1.5]    (200,60) -- (171.8,97.6) ;
\draw [shift={(170,100)}, rotate = 306.87] [color={rgb, 255:red, 74; green, 144; blue, 226 }  ,draw opacity=1 ][line width=1.5]    (14.21,-4.28) .. controls (9.04,-1.82) and (4.3,-0.39) .. (0,0) .. controls (4.3,0.39) and (9.04,1.82) .. (14.21,4.28)   ;
%Straight Lines [id:da6519101916534229] 
\draw [color={rgb, 255:red, 74; green, 144; blue, 226 }  ,draw opacity=1 ][line width=1.5]    (190,120) -- (172.12,102.12) ;
\draw [shift={(170,100)}, rotate = 45] [color={rgb, 255:red, 74; green, 144; blue, 226 }  ,draw opacity=1 ][line width=1.5]    (14.21,-4.28) .. controls (9.04,-1.82) and (4.3,-0.39) .. (0,0) .. controls (4.3,0.39) and (9.04,1.82) .. (14.21,4.28)   ;
%Straight Lines [id:da8188548073525511] 
\draw [color={rgb, 255:red, 74; green, 144; blue, 226 }  ,draw opacity=1 ][line width=1.5]    (190,120) -- (207.88,102.12) ;
\draw [shift={(210,100)}, rotate = 135] [color={rgb, 255:red, 74; green, 144; blue, 226 }  ,draw opacity=1 ][line width=1.5]    (14.21,-4.28) .. controls (9.04,-1.82) and (4.3,-0.39) .. (0,0) .. controls (4.3,0.39) and (9.04,1.82) .. (14.21,4.28)   ;
%Straight Lines [id:da9049936130099601] 
\draw [color={rgb, 255:red, 74; green, 144; blue, 226 }  ,draw opacity=1 ][line width=1.5]    (220,160) -- (192.5,141.66) ;
\draw [shift={(190,140)}, rotate = 33.69] [color={rgb, 255:red, 74; green, 144; blue, 226 }  ,draw opacity=1 ][line width=1.5]    (14.21,-4.28) .. controls (9.04,-1.82) and (4.3,-0.39) .. (0,0) .. controls (4.3,0.39) and (9.04,1.82) .. (14.21,4.28)   ;
%Straight Lines [id:da0032696330086658953] 
\draw [color={rgb, 255:red, 74; green, 144; blue, 226 }  ,draw opacity=1 ][line width=1.5]    (220,160) -- (192.12,187.88) ;
\draw [shift={(190,190)}, rotate = 315] [color={rgb, 255:red, 74; green, 144; blue, 226 }  ,draw opacity=1 ][line width=1.5]    (14.21,-4.28) .. controls (9.04,-1.82) and (4.3,-0.39) .. (0,0) .. controls (4.3,0.39) and (9.04,1.82) .. (14.21,4.28)   ;
%Straight Lines [id:da6624458862731804] 
\draw [color={rgb, 255:red, 74; green, 144; blue, 226 }  ,draw opacity=1 ][line width=1.5]    (160,160) -- (187.88,187.88) ;
\draw [shift={(190,190)}, rotate = 225] [color={rgb, 255:red, 74; green, 144; blue, 226 }  ,draw opacity=1 ][line width=1.5]    (14.21,-4.28) .. controls (9.04,-1.82) and (4.3,-0.39) .. (0,0) .. controls (4.3,0.39) and (9.04,1.82) .. (14.21,4.28)   ;
%Straight Lines [id:da9943610350622754] 
\draw [color={rgb, 255:red, 74; green, 144; blue, 226 }  ,draw opacity=1 ][line width=1.5]    (160,160) -- (187.5,141.66) ;
\draw [shift={(190,140)}, rotate = 146.31] [color={rgb, 255:red, 74; green, 144; blue, 226 }  ,draw opacity=1 ][line width=1.5]    (14.21,-4.28) .. controls (9.04,-1.82) and (4.3,-0.39) .. (0,0) .. controls (4.3,0.39) and (9.04,1.82) .. (14.21,4.28)   ;
%Straight Lines [id:da7835353652276644] 
\draw [color={rgb, 255:red, 74; green, 144; blue, 226 }  ,draw opacity=1 ][line width=1.5]    (260,120) -- (250.73,157.09) ;
\draw [shift={(250,160)}, rotate = 284.04] [color={rgb, 255:red, 74; green, 144; blue, 226 }  ,draw opacity=1 ][line width=1.5]    (14.21,-4.28) .. controls (9.04,-1.82) and (4.3,-0.39) .. (0,0) .. controls (4.3,0.39) and (9.04,1.82) .. (14.21,4.28)   ;
%Straight Lines [id:da45301259693483675] 
\draw [color={rgb, 255:red, 74; green, 144; blue, 226 }  ,draw opacity=1 ][line width=1.5]    (290,150) -- (252.91,159.27) ;
\draw [shift={(250,160)}, rotate = 345.96] [color={rgb, 255:red, 74; green, 144; blue, 226 }  ,draw opacity=1 ][line width=1.5]    (14.21,-4.28) .. controls (9.04,-1.82) and (4.3,-0.39) .. (0,0) .. controls (4.3,0.39) and (9.04,1.82) .. (14.21,4.28)   ;
%Straight Lines [id:da4594824971051823] 
\draw [color={rgb, 255:red, 74; green, 144; blue, 226 }  ,draw opacity=1 ][line width=1.5]    (260,120) -- (307,120) ;
\draw [shift={(310,120)}, rotate = 180] [color={rgb, 255:red, 74; green, 144; blue, 226 }  ,draw opacity=1 ][line width=1.5]    (14.21,-4.28) .. controls (9.04,-1.82) and (4.3,-0.39) .. (0,0) .. controls (4.3,0.39) and (9.04,1.82) .. (14.21,4.28)   ;
%Straight Lines [id:da1631157459977307] 
\draw [color={rgb, 255:red, 74; green, 144; blue, 226 }  ,draw opacity=1 ][line width=1.5]    (290,150) -- (308.34,122.5) ;
\draw [shift={(310,120)}, rotate = 123.69] [color={rgb, 255:red, 74; green, 144; blue, 226 }  ,draw opacity=1 ][line width=1.5]    (14.21,-4.28) .. controls (9.04,-1.82) and (4.3,-0.39) .. (0,0) .. controls (4.3,0.39) and (9.04,1.82) .. (14.21,4.28)   ;
%Straight Lines [id:da8901598603000586] 
\draw [color={rgb, 255:red, 74; green, 144; blue, 226 }  ,draw opacity=1 ][line width=1.5]    (340,110) -- (358.34,137.5) ;
\draw [shift={(360,140)}, rotate = 236.31] [color={rgb, 255:red, 74; green, 144; blue, 226 }  ,draw opacity=1 ][line width=1.5]    (14.21,-4.28) .. controls (9.04,-1.82) and (4.3,-0.39) .. (0,0) .. controls (4.3,0.39) and (9.04,1.82) .. (14.21,4.28)   ;
%Straight Lines [id:da28864384570380464] 
\draw [color={rgb, 255:red, 74; green, 144; blue, 226 }  ,draw opacity=1 ][line width=1.5]    (390,90) -- (361.54,137.43) ;
\draw [shift={(360,140)}, rotate = 300.96] [color={rgb, 255:red, 74; green, 144; blue, 226 }  ,draw opacity=1 ][line width=1.5]    (14.21,-4.28) .. controls (9.04,-1.82) and (4.3,-0.39) .. (0,0) .. controls (4.3,0.39) and (9.04,1.82) .. (14.21,4.28)   ;
%Straight Lines [id:da7479932996118415] 
\draw [color={rgb, 255:red, 74; green, 144; blue, 226 }  ,draw opacity=1 ][line width=1.5]    (340,110) -- (358.89,62.79) ;
\draw [shift={(360,60)}, rotate = 111.8] [color={rgb, 255:red, 74; green, 144; blue, 226 }  ,draw opacity=1 ][line width=1.5]    (14.21,-4.28) .. controls (9.04,-1.82) and (4.3,-0.39) .. (0,0) .. controls (4.3,0.39) and (9.04,1.82) .. (14.21,4.28)   ;
%Straight Lines [id:da21649314531412311] 
\draw [color={rgb, 255:red, 74; green, 144; blue, 226 }  ,draw opacity=1 ][line width=1.5]    (390,90) -- (362.12,62.12) ;
\draw [shift={(360,60)}, rotate = 45] [color={rgb, 255:red, 74; green, 144; blue, 226 }  ,draw opacity=1 ][line width=1.5]    (14.21,-4.28) .. controls (9.04,-1.82) and (4.3,-0.39) .. (0,0) .. controls (4.3,0.39) and (9.04,1.82) .. (14.21,4.28)   ;
%Straight Lines [id:da07141714735873272] 
\draw [color={rgb, 255:red, 74; green, 144; blue, 226 }  ,draw opacity=1 ][line width=1.5]    (350,190) -- (322.12,162.12) ;
\draw [shift={(320,160)}, rotate = 45] [color={rgb, 255:red, 74; green, 144; blue, 226 }  ,draw opacity=1 ][line width=1.5]    (14.21,-4.28) .. controls (9.04,-1.82) and (4.3,-0.39) .. (0,0) .. controls (4.3,0.39) and (9.04,1.82) .. (14.21,4.28)   ;
%Straight Lines [id:da761624160799076] 
\draw [color={rgb, 255:red, 74; green, 144; blue, 226 }  ,draw opacity=1 ][line width=1.5]    (360,150) -- (322.91,159.27) ;
\draw [shift={(320,160)}, rotate = 345.96] [color={rgb, 255:red, 74; green, 144; blue, 226 }  ,draw opacity=1 ][line width=1.5]    (14.21,-4.28) .. controls (9.04,-1.82) and (4.3,-0.39) .. (0,0) .. controls (4.3,0.39) and (9.04,1.82) .. (14.21,4.28)   ;
%Straight Lines [id:da8783890100064231] 
\draw [color={rgb, 255:red, 74; green, 144; blue, 226 }  ,draw opacity=1 ][line width=1.5]    (350,190) -- (397.43,161.54) ;
\draw [shift={(400,160)}, rotate = 149.04] [color={rgb, 255:red, 74; green, 144; blue, 226 }  ,draw opacity=1 ][line width=1.5]    (14.21,-4.28) .. controls (9.04,-1.82) and (4.3,-0.39) .. (0,0) .. controls (4.3,0.39) and (9.04,1.82) .. (14.21,4.28)   ;
%Straight Lines [id:da8165383548278232] 
\draw [color={rgb, 255:red, 74; green, 144; blue, 226 }  ,draw opacity=1 ][line width=1.5]    (360,150) -- (397.09,159.27) ;
\draw [shift={(400,160)}, rotate = 194.04] [color={rgb, 255:red, 74; green, 144; blue, 226 }  ,draw opacity=1 ][line width=1.5]    (14.21,-4.28) .. controls (9.04,-1.82) and (4.3,-0.39) .. (0,0) .. controls (4.3,0.39) and (9.04,1.82) .. (14.21,4.28)   ;

% Text Node
\draw (264,52.4) node [anchor=north west][inner sep=0.75pt]    {$A$};
% Text Node
\draw (218,129.4) node [anchor=north west][inner sep=0.75pt]    {$B$};
% Text Node
\draw (321,132.4) node [anchor=north west][inner sep=0.75pt]    {$C$};
% Text Node
\draw (264,192.4) node [anchor=north west][inner sep=0.75pt]    {$D$};
% Text Node
\draw (147,122.4) node [anchor=north west][inner sep=0.75pt]    {$E$};
% Text Node
\draw (411,112.4) node [anchor=north west][inner sep=0.75pt]    {$F$};
% Text Node
\draw (171,62.4) node [anchor=north west][inner sep=0.75pt]    {$0$};
% Text Node
\draw (157,172.4) node [anchor=north west][inner sep=0.75pt]    {$1$};
% Text Node
\draw (271,102.4) node [anchor=north west][inner sep=0.75pt]    {$2$};
% Text Node
\draw (337,72.4) node [anchor=north west][inner sep=0.75pt]    {$3$};
% Text Node
\draw (377,178.4) node [anchor=north west][inner sep=0.75pt]    {$4$};


\end{tikzpicture}
        \caption{enlace de Whitehead}
    \end{figure}
    Siguiendo los pasos indicados antes obtenemos los siguientes diagramas que nos dan ambos poliedros, el de arriba visto desde adentro y el de abajo visto desde afuera.
    \begin{figure}[H]
    \centering
    %!TEX root = ./main.tex



\tikzset{every picture/.style={line width=0.75pt}} %set default line width to 0.75pt        

\begin{tikzpicture}[x=0.75pt,y=0.75pt,yscale=-1,xscale=1]
%uncomment if require: \path (0,877); %set diagram left start at 0, and has height of 877

%Straight Lines [id:da9943610350622754] 
\draw [color={rgb, 255:red, 74; green, 144; blue, 226 }  ,draw opacity=1 ][line width=1.5]    (97.18,155.15) -- (97.59,121.34) ;
\draw [shift={(97.64,117.34)}, rotate = 90.69] [fill={rgb, 255:red, 74; green, 144; blue, 226 }  ,fill opacity=1 ][line width=0.08]  [draw opacity=0] (11.61,-5.58) -- (0,0) -- (11.61,5.58) -- cycle    ;
%Straight Lines [id:da28864384570380464] 
\draw [color={rgb, 255:red, 74; green, 144; blue, 226 }  ,draw opacity=1 ][line width=1.5]    (226.57,119.58) -- (225.39,149.46) ;
\draw [shift={(225.23,153.46)}, rotate = 272.26] [fill={rgb, 255:red, 74; green, 144; blue, 226 }  ,fill opacity=1 ][line width=0.08]  [draw opacity=0] (11.61,-5.58) -- (0,0) -- (11.61,5.58) -- cycle    ;
%Shape: Ellipse [id:dp029827328935257524] 
\draw  [color={rgb, 255:red, 74; green, 144; blue, 226 }  ,draw opacity=1 ][fill={rgb, 255:red, 255; green, 255; blue, 255 }  ,fill opacity=1 ][line width=1.5]  (90.95,110.99) .. controls (90.95,107.48) and (93.95,104.64) .. (97.64,104.64) .. controls (101.34,104.64) and (104.33,107.48) .. (104.33,110.99) .. controls (104.33,114.5) and (101.34,117.34) .. (97.64,117.34) .. controls (93.95,117.34) and (90.95,114.5) .. (90.95,110.99) -- cycle ;
%Shape: Ellipse [id:dp4202696786928657] 
\draw  [color={rgb, 255:red, 74; green, 144; blue, 226 }  ,draw opacity=1 ][fill={rgb, 255:red, 255; green, 255; blue, 255 }  ,fill opacity=1 ][line width=1.5]  (90.49,161.51) .. controls (90.49,158) and (93.49,155.15) .. (97.18,155.15) .. controls (100.88,155.15) and (103.88,158) .. (103.88,161.51) .. controls (103.88,165.01) and (100.88,167.86) .. (97.18,167.86) .. controls (93.49,167.86) and (90.49,165.01) .. (90.49,161.51) -- cycle ;
%Shape: Ellipse [id:dp7878916678981432] 
\draw  [color={rgb, 255:red, 74; green, 144; blue, 226 }  ,draw opacity=1 ][fill={rgb, 255:red, 255; green, 255; blue, 255 }  ,fill opacity=1 ][line width=1.5]  (157.42,142.03) .. controls (157.42,138.52) and (160.41,135.67) .. (164.11,135.67) .. controls (167.8,135.67) and (170.8,138.52) .. (170.8,142.03) .. controls (170.8,145.53) and (167.8,148.38) .. (164.11,148.38) .. controls (160.41,148.38) and (157.42,145.53) .. (157.42,142.03) -- cycle ;
%Shape: Ellipse [id:dp9437208822572988] 
\draw  [color={rgb, 255:red, 74; green, 144; blue, 226 }  ,draw opacity=1 ][fill={rgb, 255:red, 255; green, 255; blue, 255 }  ,fill opacity=1 ][line width=1.5]  (219.88,113.23) .. controls (219.88,109.72) and (222.87,106.88) .. (226.57,106.88) .. controls (230.27,106.88) and (233.26,109.72) .. (233.26,113.23) .. controls (233.26,116.74) and (230.27,119.58) .. (226.57,119.58) .. controls (222.87,119.58) and (219.88,116.74) .. (219.88,113.23) -- cycle ;
%Shape: Ellipse [id:dp06226933559817249] 
\draw  [color={rgb, 255:red, 74; green, 144; blue, 226 }  ,draw opacity=1 ][fill={rgb, 255:red, 255; green, 255; blue, 255 }  ,fill opacity=1 ][line width=1.5]  (218.09,158.96) .. controls (218.09,155.46) and (221.09,152.61) .. (224.78,152.61) .. controls (228.48,152.61) and (231.48,155.46) .. (231.48,158.96) .. controls (231.48,162.47) and (228.48,165.32) .. (224.78,165.32) .. controls (221.09,165.32) and (218.09,162.47) .. (218.09,158.96) -- cycle ;
%Curve Lines [id:da2979171995828367] 
\draw [color={rgb, 255:red, 74; green, 144; blue, 226 }  ,draw opacity=1 ][line width=1.5]    (98.19,100.11) .. controls (111.44,2.47) and (212.65,4.55) .. (226.57,106.88) ;
\draw [shift={(97.64,104.64)}, rotate = 275.97] [fill={rgb, 255:red, 74; green, 144; blue, 226 }  ,fill opacity=1 ][line width=0.08]  [draw opacity=0] (11.61,-5.58) -- (0,0) -- (11.61,5.58) -- cycle    ;
%Curve Lines [id:da8764254520788726] 
\draw [color={rgb, 255:red, 74; green, 144; blue, 226 }  ,draw opacity=1 ][line width=1.5]    (104.33,110.99) .. controls (124.31,111.9) and (152.85,115.08) .. (161.39,132.92) ;
\draw [shift={(162.77,136.52)}, rotate = 253.46] [fill={rgb, 255:red, 74; green, 144; blue, 226 }  ,fill opacity=1 ][line width=0.08]  [draw opacity=0] (11.61,-5.58) -- (0,0) -- (11.61,5.58) -- cycle    ;
%Curve Lines [id:da809160987469365] 
\draw [color={rgb, 255:red, 74; green, 144; blue, 226 }  ,draw opacity=1 ][line width=1.5]    (86.4,161.52) .. controls (-1.46,161.13) and (26.63,107.5) .. (90.95,110.99) ;
\draw [shift={(90.49,161.51)}, rotate = 179.22] [fill={rgb, 255:red, 74; green, 144; blue, 226 }  ,fill opacity=1 ][line width=0.08]  [draw opacity=0] (11.61,-5.58) -- (0,0) -- (11.61,5.58) -- cycle    ;
%Curve Lines [id:da9659495047647262] 
\draw [color={rgb, 255:red, 74; green, 144; blue, 226 }  ,draw opacity=1 ][line width=1.5]    (97.18,167.86) .. controls (118.72,285.2) and (222.93,239.85) .. (224.77,168.59) ;
\draw [shift={(224.78,165.32)}, rotate = 88.95] [fill={rgb, 255:red, 74; green, 144; blue, 226 }  ,fill opacity=1 ][line width=0.08]  [draw opacity=0] (11.61,-5.58) -- (0,0) -- (11.61,5.58) -- cycle    ;
%Curve Lines [id:da5397459745235432] 
\draw [color={rgb, 255:red, 74; green, 144; blue, 226 }  ,draw opacity=1 ][line width=1.5]    (164.11,148.38) .. controls (153.97,163.28) and (137.12,164.92) .. (107.63,161.91) ;
\draw [shift={(103.88,161.51)}, rotate = 6.33] [fill={rgb, 255:red, 74; green, 144; blue, 226 }  ,fill opacity=1 ][line width=0.08]  [draw opacity=0] (11.61,-5.58) -- (0,0) -- (11.61,5.58) -- cycle    ;
%Curve Lines [id:da4211882280413548] 
\draw [color={rgb, 255:red, 74; green, 144; blue, 226 }  ,draw opacity=1 ][line width=1.5]    (164.11,135.67) .. controls (176.75,114.11) and (190.34,113.14) .. (216.15,113.22) ;
\draw [shift={(219.88,113.23)}, rotate = 180.29] [fill={rgb, 255:red, 74; green, 144; blue, 226 }  ,fill opacity=1 ][line width=0.08]  [draw opacity=0] (11.61,-5.58) -- (0,0) -- (11.61,5.58) -- cycle    ;
%Curve Lines [id:da956839943660936] 
\draw [color={rgb, 255:red, 74; green, 144; blue, 226 }  ,draw opacity=1 ][line width=1.5]    (218.09,158.96) .. controls (198.7,163.63) and (181.69,167.42) .. (166.7,151.38) ;
\draw [shift={(164.11,148.38)}, rotate = 51.28] [fill={rgb, 255:red, 74; green, 144; blue, 226 }  ,fill opacity=1 ][line width=0.08]  [draw opacity=0] (11.61,-5.58) -- (0,0) -- (11.61,5.58) -- cycle    ;
%Curve Lines [id:da8737526251165968] 
\draw [color={rgb, 255:red, 74; green, 144; blue, 226 }  ,draw opacity=1 ][line width=1.5]    (231.48,158.96) .. controls (325.12,164.41) and (327.8,103.05) .. (236.94,110.86) ;
\draw [shift={(234.15,111.11)}, rotate = 354.32] [fill={rgb, 255:red, 74; green, 144; blue, 226 }  ,fill opacity=1 ][line width=0.08]  [draw opacity=0] (11.61,-5.58) -- (0,0) -- (11.61,5.58) -- cycle    ;
%Straight Lines [id:da7109511854568149] 
\draw [color={rgb, 255:red, 74; green, 144; blue, 226 }  ,draw opacity=1 ][line width=1.5]    (411.08,150.53) -- (411.49,116.72) ;
\draw [shift={(411.53,112.72)}, rotate = 90.69] [fill={rgb, 255:red, 74; green, 144; blue, 226 }  ,fill opacity=1 ][line width=0.08]  [draw opacity=0] (11.61,-5.58) -- (0,0) -- (11.61,5.58) -- cycle    ;
%Straight Lines [id:da49493506455500713] 
\draw [color={rgb, 255:red, 74; green, 144; blue, 226 }  ,draw opacity=1 ][line width=1.5]    (540.46,114.96) -- (539.28,144.84) ;
\draw [shift={(539.12,148.84)}, rotate = 272.26] [fill={rgb, 255:red, 74; green, 144; blue, 226 }  ,fill opacity=1 ][line width=0.08]  [draw opacity=0] (11.61,-5.58) -- (0,0) -- (11.61,5.58) -- cycle    ;
%Shape: Ellipse [id:dp9403539709464994] 
\draw  [color={rgb, 255:red, 74; green, 144; blue, 226 }  ,draw opacity=1 ][fill={rgb, 255:red, 255; green, 255; blue, 255 }  ,fill opacity=1 ][line width=1.5]  (404.84,106.37) .. controls (404.84,102.86) and (407.84,100.02) .. (411.53,100.02) .. controls (415.23,100.02) and (418.23,102.86) .. (418.23,106.37) .. controls (418.23,109.88) and (415.23,112.72) .. (411.53,112.72) .. controls (407.84,112.72) and (404.84,109.88) .. (404.84,106.37) -- cycle ;
%Shape: Ellipse [id:dp9543762925369562] 
\draw  [color={rgb, 255:red, 74; green, 144; blue, 226 }  ,draw opacity=1 ][fill={rgb, 255:red, 255; green, 255; blue, 255 }  ,fill opacity=1 ][line width=1.5]  (404.38,156.89) .. controls (404.38,153.38) and (407.38,150.53) .. (411.08,150.53) .. controls (414.77,150.53) and (417.77,153.38) .. (417.77,156.89) .. controls (417.77,160.39) and (414.77,163.24) .. (411.08,163.24) .. controls (407.38,163.24) and (404.38,160.39) .. (404.38,156.89) -- cycle ;
%Shape: Ellipse [id:dp6094250043745837] 
\draw  [color={rgb, 255:red, 74; green, 144; blue, 226 }  ,draw opacity=1 ][fill={rgb, 255:red, 255; green, 255; blue, 255 }  ,fill opacity=1 ][line width=1.5]  (471.31,137.41) .. controls (471.31,133.9) and (474.3,131.06) .. (478,131.06) .. controls (481.7,131.06) and (484.69,133.9) .. (484.69,137.41) .. controls (484.69,140.92) and (481.7,143.76) .. (478,143.76) .. controls (474.3,143.76) and (471.31,140.92) .. (471.31,137.41) -- cycle ;
%Shape: Ellipse [id:dp7048265837255207] 
\draw  [color={rgb, 255:red, 74; green, 144; blue, 226 }  ,draw opacity=1 ][fill={rgb, 255:red, 255; green, 255; blue, 255 }  ,fill opacity=1 ][line width=1.5]  (533.77,108.61) .. controls (533.77,105.1) and (536.77,102.26) .. (540.46,102.26) .. controls (544.16,102.26) and (547.15,105.1) .. (547.15,108.61) .. controls (547.15,112.12) and (544.16,114.96) .. (540.46,114.96) .. controls (536.77,114.96) and (533.77,112.12) .. (533.77,108.61) -- cycle ;
%Shape: Ellipse [id:dp8750927184704229] 
\draw  [color={rgb, 255:red, 74; green, 144; blue, 226 }  ,draw opacity=1 ][fill={rgb, 255:red, 255; green, 255; blue, 255 }  ,fill opacity=1 ][line width=1.5]  (531.98,154.35) .. controls (531.98,150.84) and (534.98,147.99) .. (538.68,147.99) .. controls (542.37,147.99) and (545.37,150.84) .. (545.37,154.35) .. controls (545.37,157.85) and (542.37,160.7) .. (538.68,160.7) .. controls (534.98,160.7) and (531.98,157.85) .. (531.98,154.35) -- cycle ;
%Curve Lines [id:da5236484284848602] 
\draw [color={rgb, 255:red, 74; green, 144; blue, 226 }  ,draw opacity=1 ][line width=1.5]    (412.08,95.49) .. controls (425.33,-2.15) and (526.55,-0.07) .. (540.46,102.26) ;
\draw [shift={(411.53,100.02)}, rotate = 275.97] [fill={rgb, 255:red, 74; green, 144; blue, 226 }  ,fill opacity=1 ][line width=0.08]  [draw opacity=0] (11.61,-5.58) -- (0,0) -- (11.61,5.58) -- cycle    ;
%Curve Lines [id:da9662673859391505] 
\draw [color={rgb, 255:red, 74; green, 144; blue, 226 }  ,draw opacity=1 ][line width=1.5]    (418.23,106.37) .. controls (438.21,107.28) and (466.74,110.46) .. (475.28,128.3) ;
\draw [shift={(476.66,131.9)}, rotate = 253.46] [fill={rgb, 255:red, 74; green, 144; blue, 226 }  ,fill opacity=1 ][line width=0.08]  [draw opacity=0] (11.61,-5.58) -- (0,0) -- (11.61,5.58) -- cycle    ;
%Curve Lines [id:da2629698729180272] 
\draw [color={rgb, 255:red, 74; green, 144; blue, 226 }  ,draw opacity=1 ][line width=1.5]    (400.29,156.91) .. controls (312.43,156.51) and (340.52,102.88) .. (404.84,106.37) ;
\draw [shift={(404.38,156.89)}, rotate = 179.22] [fill={rgb, 255:red, 74; green, 144; blue, 226 }  ,fill opacity=1 ][line width=0.08]  [draw opacity=0] (11.61,-5.58) -- (0,0) -- (11.61,5.58) -- cycle    ;
%Curve Lines [id:da48834267482785687] 
\draw [color={rgb, 255:red, 74; green, 144; blue, 226 }  ,draw opacity=1 ][line width=1.5]    (411.08,163.24) .. controls (432.61,280.58) and (536.82,235.23) .. (538.66,163.97) ;
\draw [shift={(538.68,160.7)}, rotate = 88.95] [fill={rgb, 255:red, 74; green, 144; blue, 226 }  ,fill opacity=1 ][line width=0.08]  [draw opacity=0] (11.61,-5.58) -- (0,0) -- (11.61,5.58) -- cycle    ;
%Curve Lines [id:da8150566103908512] 
\draw [color={rgb, 255:red, 74; green, 144; blue, 226 }  ,draw opacity=1 ][line width=1.5]    (478,143.76) .. controls (467.86,158.66) and (451.01,160.3) .. (421.52,157.29) ;
\draw [shift={(417.77,156.89)}, rotate = 6.33] [fill={rgb, 255:red, 74; green, 144; blue, 226 }  ,fill opacity=1 ][line width=0.08]  [draw opacity=0] (11.61,-5.58) -- (0,0) -- (11.61,5.58) -- cycle    ;
%Curve Lines [id:da4757930219164036] 
\draw [color={rgb, 255:red, 74; green, 144; blue, 226 }  ,draw opacity=1 ][line width=1.5]    (478,131.06) .. controls (490.64,109.49) and (504.23,108.52) .. (530.04,108.6) ;
\draw [shift={(533.77,108.61)}, rotate = 180.29] [fill={rgb, 255:red, 74; green, 144; blue, 226 }  ,fill opacity=1 ][line width=0.08]  [draw opacity=0] (11.61,-5.58) -- (0,0) -- (11.61,5.58) -- cycle    ;
%Curve Lines [id:da9886701870479793] 
\draw [color={rgb, 255:red, 74; green, 144; blue, 226 }  ,draw opacity=1 ][line width=1.5]    (531.98,154.35) .. controls (512.59,159.01) and (495.59,162.8) .. (480.6,146.76) ;
\draw [shift={(478,143.76)}, rotate = 51.28] [fill={rgb, 255:red, 74; green, 144; blue, 226 }  ,fill opacity=1 ][line width=0.08]  [draw opacity=0] (11.61,-5.58) -- (0,0) -- (11.61,5.58) -- cycle    ;
%Curve Lines [id:da21521285354906916] 
\draw [color={rgb, 255:red, 74; green, 144; blue, 226 }  ,draw opacity=1 ][line width=1.5]    (545.37,154.35) .. controls (639.01,159.8) and (641.69,98.43) .. (550.83,106.24) ;
\draw [shift={(548.05,106.5)}, rotate = 354.32] [fill={rgb, 255:red, 74; green, 144; blue, 226 }  ,fill opacity=1 ][line width=0.08]  [draw opacity=0] (11.61,-5.58) -- (0,0) -- (11.61,5.58) -- cycle    ;

% Text Node
\draw (151.61,72.18) node [anchor=north west][inner sep=0.75pt]    {$A$};
% Text Node
\draw (116.72,132.26) node [anchor=north west][inner sep=0.75pt]    {$B$};
% Text Node
\draw (193.25,128.92) node [anchor=north west][inner sep=0.75pt]    {$C$};
% Text Node
\draw (151.61,181.41) node [anchor=north west][inner sep=0.75pt]    {$D$};
% Text Node
\draw (62.33,126.8) node [anchor=north west][inner sep=0.75pt]    {$E$};
% Text Node
\draw (264.56,118.99) node [anchor=north west][inner sep=0.75pt]    {$F$};
% Text Node
\draw (20.25,128.92) node [anchor=north west][inner sep=0.75pt]    {$0$};
% Text Node
\draw (150.52,239.02) node [anchor=north west][inner sep=0.75pt]    {$1$};
% Text Node
\draw (175.51,100.97) node [anchor=north west][inner sep=0.75pt]    {$2$};
% Text Node
\draw (154.09,12.4) node [anchor=north west][inner sep=0.75pt]    {$3$};
% Text Node
\draw (189.78,161.1) node [anchor=north west][inner sep=0.75pt]    {$4$};
% Text Node
\draw (126.43,97.59) node [anchor=north west][inner sep=0.75pt]    {$0$};
% Text Node
\draw (88.06,128.92) node [anchor=north west][inner sep=0.75pt]    {$1$};
% Text Node
\draw (132.68,162.8) node [anchor=north west][inner sep=0.75pt]    {$2$};
% Text Node
\draw (307,123.93) node [anchor=north west][inner sep=0.75pt]    {$4$};
% Text Node
\draw (225.48,122.5) node [anchor=north west][inner sep=0.75pt]    {$3$};
% Text Node
\draw (465.5,67.56) node [anchor=north west][inner sep=0.75pt]    {$A$};
% Text Node
\draw (430.61,127.64) node [anchor=north west][inner sep=0.75pt]    {$B$};
% Text Node
\draw (507.14,124.3) node [anchor=north west][inner sep=0.75pt]    {$C$};
% Text Node
\draw (465.5,176.79) node [anchor=north west][inner sep=0.75pt]    {$D$};
% Text Node
\draw (376.22,122.18) node [anchor=north west][inner sep=0.75pt]    {$E$};
% Text Node
\draw (578.45,114.37) node [anchor=north west][inner sep=0.75pt]    {$F$};
% Text Node
\draw (398.38,124.3) node [anchor=north west][inner sep=0.75pt]    {$0$};
% Text Node
\draw (443,158.18) node [anchor=north west][inner sep=0.75pt]    {$1$};
% Text Node
\draw (439.43,92.12) node [anchor=north west][inner sep=0.75pt]    {$2$};
% Text Node
\draw (617,116.68) node [anchor=north west][inner sep=0.75pt]    {$3$};
% Text Node
\draw (541.15,124.3) node [anchor=north west][inner sep=0.75pt]    {$4$};
% Text Node
\draw (466.2,8.28) node [anchor=north west][inner sep=0.75pt]    {$0$};
% Text Node
\draw (332.35,124.3) node [anchor=north west][inner sep=0.75pt]    {$1$};
% Text Node
\draw (501.89,158.18) node [anchor=north west][inner sep=0.75pt]    {$2$};
% Text Node
\draw (469.77,234.4) node [anchor=north west][inner sep=0.75pt]    {$4$};
% Text Node
\draw (492.97,97.2) node [anchor=north west][inner sep=0.75pt]    {$3$};
% Text Node
\draw (85,222.4) node [anchor=north west][inner sep=0.75pt]    {$G$};
% Text Node
\draw (391,220.4) node [anchor=north west][inner sep=0.75pt]    {$G$};


\end{tikzpicture}
        \caption{Descompocision del complemento del enlace}
    \end{figure}
    Observe que en ambos poliedros las regiones $E$ y $F$ solo poseen dos lados, en el caso de $E$ son $0$ y $1$ con orientaciones opuestas, mientras que en $F$ son $3$ y $4$ con orientaciones opuestas también, basta con recordar eso al momento de eliminar esas caras observe que la cara exterior posee 4 lados y es la única que es de este estilo, por lo que si identificamos esta obtenemos un octaedro ideal que fue justamente la descripción que se da en las notas de Thurston, aquí hacemos el dibujo plano que es mas general
    \begin{figure}[H]
     \centering
    %!TEX root = ./main.tex



\tikzset{every picture/.style={line width=0.75pt}} %set default line width to 0.75pt        

\begin{tikzpicture}[x=0.75pt,y=0.75pt,yscale=-1,xscale=1]
%uncomment if require: \path (0,877); %set diagram left start at 0, and has height of 877

%Straight Lines [id:da9943610350622754] 
\draw [color={rgb, 255:red, 74; green, 144; blue, 226 }  ,draw opacity=1 ][line width=1.5]    (143.18,155.15) -- (143.59,121.34) ;
\draw [shift={(143.64,117.34)}, rotate = 90.69] [fill={rgb, 255:red, 74; green, 144; blue, 226 }  ,fill opacity=1 ][line width=0.08]  [draw opacity=0] (11.61,-5.58) -- (0,0) -- (11.61,5.58) -- cycle    ;
%Straight Lines [id:da28864384570380464] 
\draw [color={rgb, 255:red, 74; green, 144; blue, 226 }  ,draw opacity=1 ][line width=1.5]    (272.57,119.58) -- (271.39,149.46) ;
\draw [shift={(271.23,153.46)}, rotate = 272.26] [fill={rgb, 255:red, 74; green, 144; blue, 226 }  ,fill opacity=1 ][line width=0.08]  [draw opacity=0] (11.61,-5.58) -- (0,0) -- (11.61,5.58) -- cycle    ;
%Shape: Ellipse [id:dp029827328935257524] 
\draw  [color={rgb, 255:red, 74; green, 144; blue, 226 }  ,draw opacity=1 ][fill={rgb, 255:red, 255; green, 255; blue, 255 }  ,fill opacity=1 ][line width=1.5]  (136.95,110.99) .. controls (136.95,107.48) and (139.95,104.64) .. (143.64,104.64) .. controls (147.34,104.64) and (150.33,107.48) .. (150.33,110.99) .. controls (150.33,114.5) and (147.34,117.34) .. (143.64,117.34) .. controls (139.95,117.34) and (136.95,114.5) .. (136.95,110.99) -- cycle ;
%Shape: Ellipse [id:dp4202696786928657] 
\draw  [color={rgb, 255:red, 74; green, 144; blue, 226 }  ,draw opacity=1 ][fill={rgb, 255:red, 255; green, 255; blue, 255 }  ,fill opacity=1 ][line width=1.5]  (136.49,161.51) .. controls (136.49,158) and (139.49,155.15) .. (143.18,155.15) .. controls (146.88,155.15) and (149.88,158) .. (149.88,161.51) .. controls (149.88,165.01) and (146.88,167.86) .. (143.18,167.86) .. controls (139.49,167.86) and (136.49,165.01) .. (136.49,161.51) -- cycle ;
%Shape: Ellipse [id:dp7878916678981432] 
\draw  [color={rgb, 255:red, 74; green, 144; blue, 226 }  ,draw opacity=1 ][fill={rgb, 255:red, 255; green, 255; blue, 255 }  ,fill opacity=1 ][line width=1.5]  (203.42,142.03) .. controls (203.42,138.52) and (206.41,135.67) .. (210.11,135.67) .. controls (213.8,135.67) and (216.8,138.52) .. (216.8,142.03) .. controls (216.8,145.53) and (213.8,148.38) .. (210.11,148.38) .. controls (206.41,148.38) and (203.42,145.53) .. (203.42,142.03) -- cycle ;
%Shape: Ellipse [id:dp9437208822572988] 
\draw  [color={rgb, 255:red, 74; green, 144; blue, 226 }  ,draw opacity=1 ][fill={rgb, 255:red, 255; green, 255; blue, 255 }  ,fill opacity=1 ][line width=1.5]  (265.88,113.23) .. controls (265.88,109.72) and (268.87,106.88) .. (272.57,106.88) .. controls (276.27,106.88) and (279.26,109.72) .. (279.26,113.23) .. controls (279.26,116.74) and (276.27,119.58) .. (272.57,119.58) .. controls (268.87,119.58) and (265.88,116.74) .. (265.88,113.23) -- cycle ;
%Shape: Ellipse [id:dp06226933559817249] 
\draw  [color={rgb, 255:red, 74; green, 144; blue, 226 }  ,draw opacity=1 ][fill={rgb, 255:red, 255; green, 255; blue, 255 }  ,fill opacity=1 ][line width=1.5]  (264.09,158.96) .. controls (264.09,155.46) and (267.09,152.61) .. (270.78,152.61) .. controls (274.48,152.61) and (277.48,155.46) .. (277.48,158.96) .. controls (277.48,162.47) and (274.48,165.32) .. (270.78,165.32) .. controls (267.09,165.32) and (264.09,162.47) .. (264.09,158.96) -- cycle ;
%Curve Lines [id:da2979171995828367] 
\draw [color={rgb, 255:red, 74; green, 144; blue, 226 }  ,draw opacity=1 ][line width=1.5]    (144.19,100.11) .. controls (157.44,2.47) and (258.65,4.55) .. (272.57,106.88) ;
\draw [shift={(143.64,104.64)}, rotate = 275.97] [fill={rgb, 255:red, 74; green, 144; blue, 226 }  ,fill opacity=1 ][line width=0.08]  [draw opacity=0] (11.61,-5.58) -- (0,0) -- (11.61,5.58) -- cycle    ;
%Curve Lines [id:da8764254520788726] 
\draw [color={rgb, 255:red, 74; green, 144; blue, 226 }  ,draw opacity=1 ][line width=1.5]    (154.59,111.21) .. controls (175.51,112.49) and (202.93,116.86) .. (208.77,136.52) ;
\draw [shift={(150.33,110.99)}, rotate = 2.62] [fill={rgb, 255:red, 74; green, 144; blue, 226 }  ,fill opacity=1 ][line width=0.08]  [draw opacity=0] (11.61,-5.58) -- (0,0) -- (11.61,5.58) -- cycle    ;
%Curve Lines [id:da9659495047647262] 
\draw [color={rgb, 255:red, 74; green, 144; blue, 226 }  ,draw opacity=1 ][line width=1.5]    (143.18,167.86) .. controls (164.72,285.2) and (268.93,239.85) .. (270.77,168.59) ;
\draw [shift={(270.78,165.32)}, rotate = 88.95] [fill={rgb, 255:red, 74; green, 144; blue, 226 }  ,fill opacity=1 ][line width=0.08]  [draw opacity=0] (11.61,-5.58) -- (0,0) -- (11.61,5.58) -- cycle    ;
%Curve Lines [id:da5397459745235432] 
\draw [color={rgb, 255:red, 74; green, 144; blue, 226 }  ,draw opacity=1 ][line width=1.5]    (210.11,148.38) .. controls (199.97,163.28) and (183.12,164.92) .. (153.63,161.91) ;
\draw [shift={(149.88,161.51)}, rotate = 6.33] [fill={rgb, 255:red, 74; green, 144; blue, 226 }  ,fill opacity=1 ][line width=0.08]  [draw opacity=0] (11.61,-5.58) -- (0,0) -- (11.61,5.58) -- cycle    ;
%Curve Lines [id:da4211882280413548] 
\draw [color={rgb, 255:red, 74; green, 144; blue, 226 }  ,draw opacity=1 ][line width=1.5]    (210.11,135.67) .. controls (222.75,114.11) and (236.34,113.14) .. (262.15,113.22) ;
\draw [shift={(265.88,113.23)}, rotate = 180.29] [fill={rgb, 255:red, 74; green, 144; blue, 226 }  ,fill opacity=1 ][line width=0.08]  [draw opacity=0] (11.61,-5.58) -- (0,0) -- (11.61,5.58) -- cycle    ;
%Curve Lines [id:da956839943660936] 
\draw [color={rgb, 255:red, 74; green, 144; blue, 226 }  ,draw opacity=1 ][line width=1.5]    (260.12,159.91) .. controls (241.26,164.32) and (224.71,166.59) .. (210.11,148.38) ;
\draw [shift={(264.09,158.96)}, rotate = 166.47] [fill={rgb, 255:red, 74; green, 144; blue, 226 }  ,fill opacity=1 ][line width=0.08]  [draw opacity=0] (11.61,-5.58) -- (0,0) -- (11.61,5.58) -- cycle    ;
%Straight Lines [id:da7109511854568149] 
\draw [color={rgb, 255:red, 74; green, 144; blue, 226 }  ,draw opacity=1 ][line width=1.5]    (366.97,144.53) -- (367.38,110.72) ;
\draw [shift={(366.92,148.53)}, rotate = 270.69] [fill={rgb, 255:red, 74; green, 144; blue, 226 }  ,fill opacity=1 ][line width=0.08]  [draw opacity=0] (11.61,-5.58) -- (0,0) -- (11.61,5.58) -- cycle    ;
%Straight Lines [id:da49493506455500713] 
\draw [color={rgb, 255:red, 74; green, 144; blue, 226 }  ,draw opacity=1 ][line width=1.5]    (496.15,116.96) -- (494.97,146.84) ;
\draw [shift={(496.31,112.96)}, rotate = 92.26] [fill={rgb, 255:red, 74; green, 144; blue, 226 }  ,fill opacity=1 ][line width=0.08]  [draw opacity=0] (11.61,-5.58) -- (0,0) -- (11.61,5.58) -- cycle    ;
%Shape: Ellipse [id:dp9403539709464994] 
\draw  [color={rgb, 255:red, 74; green, 144; blue, 226 }  ,draw opacity=1 ][fill={rgb, 255:red, 255; green, 255; blue, 255 }  ,fill opacity=1 ][line width=1.5]  (360.69,104.37) .. controls (360.69,100.86) and (363.68,98.02) .. (367.38,98.02) .. controls (371.08,98.02) and (374.07,100.86) .. (374.07,104.37) .. controls (374.07,107.88) and (371.08,110.72) .. (367.38,110.72) .. controls (363.68,110.72) and (360.69,107.88) .. (360.69,104.37) -- cycle ;
%Shape: Ellipse [id:dp9543762925369562] 
\draw  [color={rgb, 255:red, 74; green, 144; blue, 226 }  ,draw opacity=1 ][fill={rgb, 255:red, 255; green, 255; blue, 255 }  ,fill opacity=1 ][line width=1.5]  (360.23,154.89) .. controls (360.23,151.38) and (363.23,148.53) .. (366.92,148.53) .. controls (370.62,148.53) and (373.62,151.38) .. (373.62,154.89) .. controls (373.62,158.39) and (370.62,161.24) .. (366.92,161.24) .. controls (363.23,161.24) and (360.23,158.39) .. (360.23,154.89) -- cycle ;
%Shape: Ellipse [id:dp6094250043745837] 
\draw  [color={rgb, 255:red, 74; green, 144; blue, 226 }  ,draw opacity=1 ][fill={rgb, 255:red, 255; green, 255; blue, 255 }  ,fill opacity=1 ][line width=1.5]  (427.15,135.41) .. controls (427.15,131.9) and (430.15,129.06) .. (433.85,129.06) .. controls (437.54,129.06) and (440.54,131.9) .. (440.54,135.41) .. controls (440.54,138.92) and (437.54,141.76) .. (433.85,141.76) .. controls (430.15,141.76) and (427.15,138.92) .. (427.15,135.41) -- cycle ;
%Shape: Ellipse [id:dp7048265837255207] 
\draw  [color={rgb, 255:red, 74; green, 144; blue, 226 }  ,draw opacity=1 ][fill={rgb, 255:red, 255; green, 255; blue, 255 }  ,fill opacity=1 ][line width=1.5]  (489.62,106.61) .. controls (489.62,103.1) and (492.61,100.26) .. (496.31,100.26) .. controls (500,100.26) and (503,103.1) .. (503,106.61) .. controls (503,110.12) and (500,112.96) .. (496.31,112.96) .. controls (492.61,112.96) and (489.62,110.12) .. (489.62,106.61) -- cycle ;
%Shape: Ellipse [id:dp8750927184704229] 
\draw  [color={rgb, 255:red, 74; green, 144; blue, 226 }  ,draw opacity=1 ][fill={rgb, 255:red, 255; green, 255; blue, 255 }  ,fill opacity=1 ][line width=1.5]  (487.83,152.35) .. controls (487.83,148.84) and (490.83,145.99) .. (494.52,145.99) .. controls (498.22,145.99) and (501.22,148.84) .. (501.22,152.35) .. controls (501.22,155.85) and (498.22,158.7) .. (494.52,158.7) .. controls (490.83,158.7) and (487.83,155.85) .. (487.83,152.35) -- cycle ;
%Curve Lines [id:da5236484284848602] 
\draw [color={rgb, 255:red, 74; green, 144; blue, 226 }  ,draw opacity=1 ][line width=1.5]    (367.38,98.02) .. controls (377.97,-3.17) and (480.11,-3.63) .. (495.86,97.17) ;
\draw [shift={(496.31,100.26)}, rotate = 262.26] [fill={rgb, 255:red, 74; green, 144; blue, 226 }  ,fill opacity=1 ][line width=0.08]  [draw opacity=0] (11.61,-5.58) -- (0,0) -- (11.61,5.58) -- cycle    ;
%Curve Lines [id:da9662673859391505] 
\draw [color={rgb, 255:red, 74; green, 144; blue, 226 }  ,draw opacity=1 ][line width=1.5]    (374.07,104.37) .. controls (394.05,105.28) and (422.58,108.46) .. (431.12,126.3) ;
\draw [shift={(432.51,129.9)}, rotate = 253.46] [fill={rgb, 255:red, 74; green, 144; blue, 226 }  ,fill opacity=1 ][line width=0.08]  [draw opacity=0] (11.61,-5.58) -- (0,0) -- (11.61,5.58) -- cycle    ;
%Curve Lines [id:da48834267482785687] 
\draw [color={rgb, 255:red, 74; green, 144; blue, 226 }  ,draw opacity=1 ][line width=1.5]    (367.96,166.49) .. controls (391.97,278.91) and (495.84,230.72) .. (494.52,158.7) ;
\draw [shift={(366.92,161.24)}, rotate = 79.6] [fill={rgb, 255:red, 74; green, 144; blue, 226 }  ,fill opacity=1 ][line width=0.08]  [draw opacity=0] (11.61,-5.58) -- (0,0) -- (11.61,5.58) -- cycle    ;
%Curve Lines [id:da8150566103908512] 
\draw [color={rgb, 255:red, 74; green, 144; blue, 226 }  ,draw opacity=1 ][line width=1.5]    (433.85,141.76) .. controls (423.71,156.66) and (406.86,158.3) .. (377.37,155.29) ;
\draw [shift={(373.62,154.89)}, rotate = 6.33] [fill={rgb, 255:red, 74; green, 144; blue, 226 }  ,fill opacity=1 ][line width=0.08]  [draw opacity=0] (11.61,-5.58) -- (0,0) -- (11.61,5.58) -- cycle    ;
%Curve Lines [id:da4757930219164036] 
\draw [color={rgb, 255:red, 74; green, 144; blue, 226 }  ,draw opacity=1 ][line width=1.5]    (433.85,129.06) .. controls (446.49,107.49) and (460.08,106.52) .. (485.88,106.6) ;
\draw [shift={(489.62,106.61)}, rotate = 180.29] [fill={rgb, 255:red, 74; green, 144; blue, 226 }  ,fill opacity=1 ][line width=0.08]  [draw opacity=0] (11.61,-5.58) -- (0,0) -- (11.61,5.58) -- cycle    ;
%Curve Lines [id:da9886701870479793] 
\draw [color={rgb, 255:red, 74; green, 144; blue, 226 }  ,draw opacity=1 ][line width=1.5]    (487.83,152.35) .. controls (468.44,157.01) and (451.43,160.8) .. (436.44,144.76) ;
\draw [shift={(433.85,141.76)}, rotate = 51.28] [fill={rgb, 255:red, 74; green, 144; blue, 226 }  ,fill opacity=1 ][line width=0.08]  [draw opacity=0] (11.61,-5.58) -- (0,0) -- (11.61,5.58) -- cycle    ;

% Text Node
\draw (197.61,72.18) node [anchor=north west][inner sep=0.75pt]    {$A$};
% Text Node
\draw (162.72,132.26) node [anchor=north west][inner sep=0.75pt]    {$B$};
% Text Node
\draw (239.25,128.92) node [anchor=north west][inner sep=0.75pt]    {$C$};
% Text Node
\draw (197.61,181.41) node [anchor=north west][inner sep=0.75pt]    {$D$};
% Text Node
\draw (196.52,239.02) node [anchor=north west][inner sep=0.75pt]    {$1$};
% Text Node
\draw (221.51,100.97) node [anchor=north west][inner sep=0.75pt]    {$2$};
% Text Node
\draw (200.09,12.4) node [anchor=north west][inner sep=0.75pt]    {$3$};
% Text Node
\draw (235.78,161.1) node [anchor=north west][inner sep=0.75pt]    {$3$};
% Text Node
\draw (177,97.4) node [anchor=north west][inner sep=0.75pt]    {$1$};
% Text Node
\draw (134.06,128.92) node [anchor=north west][inner sep=0.75pt]    {$1$};
% Text Node
\draw (178.68,162.8) node [anchor=north west][inner sep=0.75pt]    {$2$};
% Text Node
\draw (271.48,122.5) node [anchor=north west][inner sep=0.75pt]    {$3$};
% Text Node
\draw (421.34,65.56) node [anchor=north west][inner sep=0.75pt]    {$A$};
% Text Node
\draw (386.46,125.64) node [anchor=north west][inner sep=0.75pt]    {$B$};
% Text Node
\draw (462.98,122.3) node [anchor=north west][inner sep=0.75pt]    {$C$};
% Text Node
\draw (421.34,174.79) node [anchor=north west][inner sep=0.75pt]    {$D$};
% Text Node
\draw (354.23,122.3) node [anchor=north west][inner sep=0.75pt]    {$1$};
% Text Node
\draw (398.85,156.18) node [anchor=north west][inner sep=0.75pt]    {$1$};
% Text Node
\draw (395.28,90.12) node [anchor=north west][inner sep=0.75pt]    {$2$};
% Text Node
\draw (497,122.3) node [anchor=north west][inner sep=0.75pt]    {$3$};
% Text Node
\draw (422.05,6.28) node [anchor=north west][inner sep=0.75pt]    {$1$};
% Text Node
\draw (457.74,156.18) node [anchor=north west][inner sep=0.75pt]    {$2$};
% Text Node
\draw (425.62,232.4) node [anchor=north west][inner sep=0.75pt]    {$3$};
% Text Node
\draw (448.82,95.2) node [anchor=north west][inner sep=0.75pt]    {$3$};
% Text Node
\draw (131,222.4) node [anchor=north west][inner sep=0.75pt]    {$G$};
% Text Node
\draw (346.85,218.4) node [anchor=north west][inner sep=0.75pt]    {$G$};


\end{tikzpicture}
        \caption{dos pirámides ideales del enlace de Whitehead}
    \end{figure}
\end{eg}
\begin{eg}
    Considere la siguiente proyeccion del nudo $6_1$ que no es precisamente la estandar mas tiene un parecido con la proyeccion del nudo de 8 y por eso la tomamos
    \begin{figure}[H]
    \centering
    %!TEX root = ./main.tex



\tikzset{every picture/.style={line width=0.75pt}} %set default line width to 0.75pt        

\begin{tikzpicture}[x=0.75pt,y=0.75pt,yscale=-1,xscale=1]
%uncomment if require: \path (0,877); %set diagram left start at 0, and has height of 877
\path[use as bounding box] (20,10) rectangle (660,300);
%Curve Lines [id:da7624936084054103] 
\draw [line width=1.5]    (332.42,57.3) .. controls (269.42,104.3) and (350.91,114.97) .. (386.58,137.3) .. controls (422.25,159.64) and (448.25,198.3) .. (416.58,237.3) ;
%Curve Lines [id:da19391436252011207] 
\draw [line width=1.5]    (316.58,237.3) .. controls (302.91,294.97) and (208.25,240.97) .. (326.58,117.3) ;
%Curve Lines [id:da04931308389248612] 
\draw [line width=1.5]    (276.58,237.3) .. controls (326.25,158.3) and (323.58,304.3) .. (366.58,247.3) ;
%Curve Lines [id:da7925895816985254] 
\draw [line width=1.5]    (346.58,107.3) .. controls (390.25,75.64) and (266.91,-60.36) .. (206.58,127.3) .. controls (146.25,314.97) and (270.25,258.97) .. (266.58,247.3) ;
%Curve Lines [id:da7432199945489849] 
\draw [line width=1.5]    (376.58,237.3) .. controls (406.91,159.64) and (402.25,313.64) .. (456.58,257.3) .. controls (510.91,200.97) and (461.58,12.3) .. (346.58,47.3) ;
%Curve Lines [id:da5685625074522969] 
\draw [line width=1.5]    (326.58,217.3) .. controls (364.91,172.3) and (381.91,308.97) .. (403.58,250.3) ;
%Straight Lines [id:da1141281197725319] 
\draw [color={rgb, 255:red, 74; green, 144; blue, 226 }  ,draw opacity=1 ][line width=1.5]    (318.68,70) -- (347.68,70) ;
\draw [shift={(351.68,70)}, rotate = 180] [fill={rgb, 255:red, 74; green, 144; blue, 226 }  ,fill opacity=1 ][line width=0.08]  [draw opacity=0] (11.61,-5.58) -- (0,0) -- (11.61,5.58) -- cycle    ;
%Straight Lines [id:da5128754818859278] 
\draw [color={rgb, 255:red, 74; green, 144; blue, 226 }  ,draw opacity=1 ][line width=1.5]    (318.68,70) -- (318.68,41) ;
\draw [shift={(318.68,37)}, rotate = 90] [fill={rgb, 255:red, 74; green, 144; blue, 226 }  ,fill opacity=1 ][line width=0.08]  [draw opacity=0] (11.61,-5.58) -- (0,0) -- (11.61,5.58) -- cycle    ;
%Straight Lines [id:da6937466521525161] 
\draw [color={rgb, 255:red, 74; green, 144; blue, 226 }  ,draw opacity=1 ][line width=1.5]    (352.68,42) -- (351.82,66) ;
\draw [shift={(351.68,70)}, rotate = 272.05] [fill={rgb, 255:red, 74; green, 144; blue, 226 }  ,fill opacity=1 ][line width=0.08]  [draw opacity=0] (11.61,-5.58) -- (0,0) -- (11.61,5.58) -- cycle    ;
%Straight Lines [id:da1076986821814464] 
\draw [color={rgb, 255:red, 74; green, 144; blue, 226 }  ,draw opacity=1 ][line width=1.5]    (352.68,42) -- (322.67,40.23) ;
\draw [shift={(318.68,40)}, rotate = 3.37] [fill={rgb, 255:red, 74; green, 144; blue, 226 }  ,fill opacity=1 ][line width=0.08]  [draw opacity=0] (11.61,-5.58) -- (0,0) -- (11.61,5.58) -- cycle    ;
%Straight Lines [id:da5957859563126462] 
\draw [color={rgb, 255:red, 74; green, 144; blue, 226 }  ,draw opacity=1 ][line width=1.5]    (355.68,100) -- (318.68,100) ;
\draw [shift={(314.68,100)}, rotate = 360] [fill={rgb, 255:red, 74; green, 144; blue, 226 }  ,fill opacity=1 ][line width=0.08]  [draw opacity=0] (11.61,-5.58) -- (0,0) -- (11.61,5.58) -- cycle    ;
%Straight Lines [id:da6302195089744549] 
\draw [color={rgb, 255:red, 74; green, 144; blue, 226 }  ,draw opacity=1 ][line width=1.5]    (316.68,130) -- (366.68,130) ;
\draw [shift={(370.68,130)}, rotate = 180] [fill={rgb, 255:red, 74; green, 144; blue, 226 }  ,fill opacity=1 ][line width=0.08]  [draw opacity=0] (11.61,-5.58) -- (0,0) -- (11.61,5.58) -- cycle    ;
%Straight Lines [id:da5412838641193578] 
\draw [color={rgb, 255:red, 74; green, 144; blue, 226 }  ,draw opacity=1 ][line width=1.5]    (316.68,130) -- (314.94,103.99) ;
\draw [shift={(314.68,100)}, rotate = 86.19] [fill={rgb, 255:red, 74; green, 144; blue, 226 }  ,fill opacity=1 ][line width=0.08]  [draw opacity=0] (11.61,-5.58) -- (0,0) -- (11.61,5.58) -- cycle    ;
%Straight Lines [id:da3701320323663966] 
\draw [color={rgb, 255:red, 74; green, 144; blue, 226 }  ,draw opacity=1 ][line width=1.5]    (370.68,130) -- (357.47,103.58) ;
\draw [shift={(355.68,100)}, rotate = 63.43] [fill={rgb, 255:red, 74; green, 144; blue, 226 }  ,fill opacity=1 ][line width=0.08]  [draw opacity=0] (11.61,-5.58) -- (0,0) -- (11.61,5.58) -- cycle    ;
%Straight Lines [id:da4037972105202199] 
\draw [color={rgb, 255:red, 74; green, 144; blue, 226 }  ,draw opacity=1 ][line width=1.5]    (250.68,260) -- (267.04,223.65) ;
\draw [shift={(268.68,220)}, rotate = 114.23] [fill={rgb, 255:red, 74; green, 144; blue, 226 }  ,fill opacity=1 ][line width=0.08]  [draw opacity=0] (11.61,-5.58) -- (0,0) -- (11.61,5.58) -- cycle    ;
%Straight Lines [id:da8718338183932847] 
\draw [color={rgb, 255:red, 74; green, 144; blue, 226 }  ,draw opacity=1 ][line width=1.5]    (288.68,220) -- (288.68,256) ;
\draw [shift={(288.68,260)}, rotate = 270] [fill={rgb, 255:red, 74; green, 144; blue, 226 }  ,fill opacity=1 ][line width=0.08]  [draw opacity=0] (11.61,-5.58) -- (0,0) -- (11.61,5.58) -- cycle    ;
%Straight Lines [id:da7553369000057689] 
\draw [color={rgb, 255:red, 74; green, 144; blue, 226 }  ,draw opacity=1 ][line width=1.5]    (306.68,256) -- (308.48,219) ;
\draw [shift={(308.68,215)}, rotate = 92.79] [fill={rgb, 255:red, 74; green, 144; blue, 226 }  ,fill opacity=1 ][line width=0.08]  [draw opacity=0] (11.61,-5.58) -- (0,0) -- (11.61,5.58) -- cycle    ;
%Straight Lines [id:da7979498405597526] 
\draw [color={rgb, 255:red, 74; green, 144; blue, 226 }  ,draw opacity=1 ][line width=1.5]    (338.68,210) -- (338.68,253) ;
\draw [shift={(338.68,257)}, rotate = 270] [fill={rgb, 255:red, 74; green, 144; blue, 226 }  ,fill opacity=1 ][line width=0.08]  [draw opacity=0] (11.61,-5.58) -- (0,0) -- (11.61,5.58) -- cycle    ;
%Straight Lines [id:da9990695315517517] 
\draw [color={rgb, 255:red, 74; green, 144; blue, 226 }  ,draw opacity=1 ][line width=1.5]    (358.68,256) -- (358.68,224) ;
\draw [shift={(358.68,220)}, rotate = 90] [fill={rgb, 255:red, 74; green, 144; blue, 226 }  ,fill opacity=1 ][line width=0.08]  [draw opacity=0] (11.61,-5.58) -- (0,0) -- (11.61,5.58) -- cycle    ;
%Straight Lines [id:da5789283453234753] 
\draw [color={rgb, 255:red, 74; green, 144; blue, 226 }  ,draw opacity=1 ][line width=1.5]    (382.68,226) -- (385.34,257.01) ;
\draw [shift={(385.68,261)}, rotate = 265.1] [fill={rgb, 255:red, 74; green, 144; blue, 226 }  ,fill opacity=1 ][line width=0.08]  [draw opacity=0] (11.61,-5.58) -- (0,0) -- (11.61,5.58) -- cycle    ;
%Straight Lines [id:da8699663998653359] 
\draw [color={rgb, 255:red, 74; green, 144; blue, 226 }  ,draw opacity=1 ][line width=1.5]    (398.68,260) -- (398.68,224) ;
\draw [shift={(398.68,220)}, rotate = 90] [fill={rgb, 255:red, 74; green, 144; blue, 226 }  ,fill opacity=1 ][line width=0.08]  [draw opacity=0] (11.61,-5.58) -- (0,0) -- (11.61,5.58) -- cycle    ;
%Straight Lines [id:da798656100582777] 
\draw [color={rgb, 255:red, 74; green, 144; blue, 226 }  ,draw opacity=1 ][line width=1.5]    (426.68,222) -- (421.26,259.04) ;
\draw [shift={(420.68,263)}, rotate = 278.33] [fill={rgb, 255:red, 74; green, 144; blue, 226 }  ,fill opacity=1 ][line width=0.08]  [draw opacity=0] (11.61,-5.58) -- (0,0) -- (11.61,5.58) -- cycle    ;
%Straight Lines [id:da33279694069000054] 
\draw [color={rgb, 255:red, 74; green, 144; blue, 226 }  ,draw opacity=1 ][line width=1.5]    (288.68,220) -- (272.68,220) ;
\draw [shift={(268.68,220)}, rotate = 360] [fill={rgb, 255:red, 74; green, 144; blue, 226 }  ,fill opacity=1 ][line width=0.08]  [draw opacity=0] (11.61,-5.58) -- (0,0) -- (11.61,5.58) -- cycle    ;
%Straight Lines [id:da43289558059406263] 
\draw [color={rgb, 255:red, 74; green, 144; blue, 226 }  ,draw opacity=1 ][line width=1.5]    (338.68,210) -- (312.62,214.34) ;
\draw [shift={(308.68,215)}, rotate = 350.54] [fill={rgb, 255:red, 74; green, 144; blue, 226 }  ,fill opacity=1 ][line width=0.08]  [draw opacity=0] (11.61,-5.58) -- (0,0) -- (11.61,5.58) -- cycle    ;
%Straight Lines [id:da6303991823506472] 
\draw [color={rgb, 255:red, 74; green, 144; blue, 226 }  ,draw opacity=1 ][line width=1.5]    (382.68,226) -- (362.56,220.97) ;
\draw [shift={(358.68,220)}, rotate = 14.04] [fill={rgb, 255:red, 74; green, 144; blue, 226 }  ,fill opacity=1 ][line width=0.08]  [draw opacity=0] (11.61,-5.58) -- (0,0) -- (11.61,5.58) -- cycle    ;
%Straight Lines [id:da8930356844444237] 
\draw [color={rgb, 255:red, 74; green, 144; blue, 226 }  ,draw opacity=1 ][line width=1.5]    (426.68,222) -- (402.67,220.28) ;
\draw [shift={(398.68,220)}, rotate = 4.09] [fill={rgb, 255:red, 74; green, 144; blue, 226 }  ,fill opacity=1 ][line width=0.08]  [draw opacity=0] (11.61,-5.58) -- (0,0) -- (11.61,5.58) -- cycle    ;
%Straight Lines [id:da9121806098865429] 
\draw [color={rgb, 255:red, 74; green, 144; blue, 226 }  ,draw opacity=1 ][line width=1.5]    (250.68,260) -- (284.68,260) ;
\draw [shift={(288.68,260)}, rotate = 180] [fill={rgb, 255:red, 74; green, 144; blue, 226 }  ,fill opacity=1 ][line width=0.08]  [draw opacity=0] (11.61,-5.58) -- (0,0) -- (11.61,5.58) -- cycle    ;
%Straight Lines [id:da632074891050223] 
\draw [color={rgb, 255:red, 74; green, 144; blue, 226 }  ,draw opacity=1 ][line width=1.5]    (306.68,256) -- (334.68,256.88) ;
\draw [shift={(338.68,257)}, rotate = 181.79] [fill={rgb, 255:red, 74; green, 144; blue, 226 }  ,fill opacity=1 ][line width=0.08]  [draw opacity=0] (11.61,-5.58) -- (0,0) -- (11.61,5.58) -- cycle    ;
%Straight Lines [id:da7854913192561737] 
\draw [color={rgb, 255:red, 74; green, 144; blue, 226 }  ,draw opacity=1 ][line width=1.5]    (358.68,256) -- (381.74,260.27) ;
\draw [shift={(385.68,261)}, rotate = 190.49] [fill={rgb, 255:red, 74; green, 144; blue, 226 }  ,fill opacity=1 ][line width=0.08]  [draw opacity=0] (11.61,-5.58) -- (0,0) -- (11.61,5.58) -- cycle    ;
%Straight Lines [id:da3077701960838797] 
\draw [color={rgb, 255:red, 74; green, 144; blue, 226 }  ,draw opacity=1 ][line width=1.5]    (398.68,260) -- (416.71,262.46) ;
\draw [shift={(420.68,263)}, rotate = 187.77] [fill={rgb, 255:red, 74; green, 144; blue, 226 }  ,fill opacity=1 ][line width=0.08]  [draw opacity=0] (11.61,-5.58) -- (0,0) -- (11.61,5.58) -- cycle    ;

% Text Node
\draw (242.68,132.4) node [anchor=north west][inner sep=0.75pt]    {$A$};
% Text Node
\draw (435.68,142.4) node [anchor=north west][inner sep=0.75pt]    {$B$};
% Text Node
\draw (339.68,162.4) node [anchor=north west][inner sep=0.75pt]    {$C$};
% Text Node
\draw (192.68,282.4) node [anchor=north west][inner sep=0.75pt]    {$D$};
% Text Node
\draw (289.68,232.4) node [anchor=north west][inner sep=0.75pt]    {$E$};
% Text Node
\draw (339.68,222.4) node [anchor=north west][inner sep=0.75pt]    {$F$};
% Text Node
\draw (383.68,229.4) node [anchor=north west][inner sep=0.75pt]    {$G$};
% Text Node
\draw (299.68,52.4) node [anchor=north west][inner sep=0.75pt]    {$1$};
% Text Node
\draw (299.68,112.4) node [anchor=north west][inner sep=0.75pt]    {$2$};
% Text Node
\draw (255.68,262.4) node [anchor=north west][inner sep=0.75pt]    {$3$};
% Text Node
\draw (310.68,258.4) node [anchor=north west][inner sep=0.75pt]    {$4$};
% Text Node
\draw (360.68,260.4) node [anchor=north west][inner sep=0.75pt]    {$5$};
% Text Node
\draw (399.68,262.4) node [anchor=north west][inner sep=0.75pt]    {$6$};


\end{tikzpicture}
        \caption{Nudo $6_1$}
    \end{figure}
   Nuevamente por los pasos de antes obtenemos los siguientes diagramas que nos dan ambos poliedros, el de arriba visto desde adentro y el de abajo visto desde afuera.
    \begin{figure}[H]
    %!TEX root = ./main.tex



\tikzset{every picture/.style={line width=0.75pt}} %set default line width to 0.75pt        

\begin{tikzpicture}[x=0.75pt,y=0.75pt,yscale=-1,xscale=1]
%uncomment if require: \path (0,877); %set diagram left start at 0, and has height of 877

%Shape: Ellipse [id:dp2062482206131232] 
\draw  [color={rgb, 255:red, 74; green, 144; blue, 226 }  ,draw opacity=1 ][line width=1.5]  (181.18,43.69) .. controls (181.18,40.42) and (183.82,37.77) .. (187.08,37.77) .. controls (190.33,37.77) and (192.97,40.42) .. (192.97,43.69) .. controls (192.97,46.95) and (190.33,49.6) .. (187.08,49.6) .. controls (183.82,49.6) and (181.18,46.95) .. (181.18,43.69) -- cycle ;
%Shape: Ellipse [id:dp42181631305663214] 
\draw  [color={rgb, 255:red, 74; green, 144; blue, 226 }  ,draw opacity=1 ][line width=1.5]  (181.18,125.31) .. controls (181.18,122.05) and (183.82,119.39) .. (187.08,119.39) .. controls (190.33,119.39) and (192.97,122.05) .. (192.97,125.31) .. controls (192.97,128.58) and (190.33,131.23) .. (187.08,131.23) .. controls (183.82,131.23) and (181.18,128.58) .. (181.18,125.31) -- cycle ;
%Shape: Ellipse [id:dp1893981875046039] 
\draw  [color={rgb, 255:red, 74; green, 144; blue, 226 }  ,draw opacity=1 ][line width=1.5]  (81.33,202.58) .. controls (81.33,199.31) and (83.97,196.66) .. (87.22,196.66) .. controls (90.48,196.66) and (93.11,199.31) .. (93.11,202.58) .. controls (93.11,205.85) and (90.48,208.5) .. (87.22,208.5) .. controls (83.97,208.5) and (81.33,205.85) .. (81.33,202.58) -- cycle ;
%Shape: Ellipse [id:dp3630097335542829] 
\draw  [color={rgb, 255:red, 74; green, 144; blue, 226 }  ,draw opacity=1 ][line width=1.5]  (276.69,200.09) .. controls (276.69,196.82) and (279.33,194.17) .. (282.59,194.17) .. controls (285.84,194.17) and (288.48,196.82) .. (288.48,200.09) .. controls (288.48,203.36) and (285.84,206.01) .. (282.59,206.01) .. controls (279.33,206.01) and (276.69,203.36) .. (276.69,200.09) -- cycle ;
%Shape: Ellipse [id:dp12213112120599945] 
\draw  [color={rgb, 255:red, 74; green, 144; blue, 226 }  ,draw opacity=1 ][line width=1.5]  (146.45,202.58) .. controls (146.45,199.31) and (149.09,196.66) .. (152.34,196.66) .. controls (155.6,196.66) and (158.24,199.31) .. (158.24,202.58) .. controls (158.24,205.85) and (155.6,208.5) .. (152.34,208.5) .. controls (149.09,208.5) and (146.45,205.85) .. (146.45,202.58) -- cycle ;
%Shape: Ellipse [id:dp14956651149377553] 
\draw  [color={rgb, 255:red, 74; green, 144; blue, 226 }  ,draw opacity=1 ][line width=1.5]  (218.4,202.58) .. controls (218.4,199.31) and (221.03,196.66) .. (224.29,196.66) .. controls (227.54,196.66) and (230.18,199.31) .. (230.18,202.58) .. controls (230.18,205.85) and (227.54,208.5) .. (224.29,208.5) .. controls (221.03,208.5) and (218.4,205.85) .. (218.4,202.58) -- cycle ;
%Curve Lines [id:da13558369698297235] 
\draw [color={rgb, 255:red, 74; green, 144; blue, 226 }  ,draw opacity=1 ][line width=1.5]    (190.32,116.91) .. controls (224.15,90.46) and (221.37,75.9) .. (187.08,49.6) ;
\draw [shift={(187.08,119.39)}, rotate = 323] [fill={rgb, 255:red, 74; green, 144; blue, 226 }  ,fill opacity=1 ][line width=0.08]  [draw opacity=0] (11.61,-5.58) -- (0,0) -- (11.61,5.58) -- cycle    ;
%Curve Lines [id:da08096805629350934] 
\draw [color={rgb, 255:red, 74; green, 144; blue, 226 }  ,draw opacity=1 ][line width=1.5]    (187.08,119.39) .. controls (153.21,100.86) and (142.93,76.71) .. (183.83,51.54) ;
\draw [shift={(187.08,49.6)}, rotate = 149.92] [fill={rgb, 255:red, 74; green, 144; blue, 226 }  ,fill opacity=1 ][line width=0.08]  [draw opacity=0] (11.61,-5.58) -- (0,0) -- (11.61,5.58) -- cycle    ;
%Curve Lines [id:da46118430818328493] 
\draw [color={rgb, 255:red, 74; green, 144; blue, 226 }  ,draw opacity=1 ][line width=1.5]    (81.33,202.58) .. controls (46.2,174.35) and (108.26,145.51) .. (183.62,131.85) ;
\draw [shift={(187.08,131.23)}, rotate = 170.12] [fill={rgb, 255:red, 74; green, 144; blue, 226 }  ,fill opacity=1 ][line width=0.08]  [draw opacity=0] (11.61,-5.58) -- (0,0) -- (11.61,5.58) -- cycle    ;
%Curve Lines [id:da12684350601217753] 
\draw [color={rgb, 255:red, 74; green, 144; blue, 226 }  ,draw opacity=1 ][line width=1.5]    (292.05,197.16) .. controls (321.77,170.32) and (267.27,146.57) .. (187.08,131.23) ;
\draw [shift={(288.48,200.09)}, rotate = 323] [fill={rgb, 255:red, 74; green, 144; blue, 226 }  ,fill opacity=1 ][line width=0.08]  [draw opacity=0] (11.61,-5.58) -- (0,0) -- (11.61,5.58) -- cycle    ;
%Curve Lines [id:da12533726813089074] 
\draw [color={rgb, 255:red, 74; green, 144; blue, 226 }  ,draw opacity=1 ][line width=1.5]    (288.48,200.09) .. controls (375.8,242.25) and (338.78,-38.85) .. (189.34,36.6) ;
\draw [shift={(187.08,37.77)}, rotate = 332.3] [fill={rgb, 255:red, 74; green, 144; blue, 226 }  ,fill opacity=1 ][line width=0.08]  [draw opacity=0] (11.61,-5.58) -- (0,0) -- (11.61,5.58) -- cycle    ;
%Curve Lines [id:da052098391594170956] 
\draw [color={rgb, 255:red, 74; green, 144; blue, 226 }  ,draw opacity=1 ][line width=1.5]    (76.95,203.7) .. controls (7.03,216.62) and (19.32,-31.54) .. (187.08,37.77) ;
\draw [shift={(81.33,202.58)}, rotate = 161.99] [fill={rgb, 255:red, 74; green, 144; blue, 226 }  ,fill opacity=1 ][line width=0.08]  [draw opacity=0] (11.61,-5.58) -- (0,0) -- (11.61,5.58) -- cycle    ;
%Curve Lines [id:da545296203025464] 
\draw [color={rgb, 255:red, 74; green, 144; blue, 226 }  ,draw opacity=1 ][line width=1.5]    (87.22,208.5) .. controls (95.26,249.18) and (133.74,236.79) .. (150.38,211.7) ;
\draw [shift={(152.34,208.5)}, rotate = 119.46] [fill={rgb, 255:red, 74; green, 144; blue, 226 }  ,fill opacity=1 ][line width=0.08]  [draw opacity=0] (11.61,-5.58) -- (0,0) -- (11.61,5.58) -- cycle    ;
%Curve Lines [id:da5810485905154714] 
\draw [color={rgb, 255:red, 74; green, 144; blue, 226 }  ,draw opacity=1 ][line width=1.5]    (152.34,208.5) .. controls (160.38,249.18) and (203.97,234.33) .. (222.15,211.41) ;
\draw [shift={(224.29,208.5)}, rotate = 124.16] [fill={rgb, 255:red, 74; green, 144; blue, 226 }  ,fill opacity=1 ][line width=0.08]  [draw opacity=0] (11.61,-5.58) -- (0,0) -- (11.61,5.58) -- cycle    ;
%Curve Lines [id:da14849415068653748] 
\draw [color={rgb, 255:red, 74; green, 144; blue, 226 }  ,draw opacity=1 ][line width=1.5]    (224.29,208.5) .. controls (232.28,248.97) and (276.2,234.31) .. (282.02,209.57) ;
\draw [shift={(282.59,206.01)}, rotate = 94.63] [fill={rgb, 255:red, 74; green, 144; blue, 226 }  ,fill opacity=1 ][line width=0.08]  [draw opacity=0] (11.61,-5.58) -- (0,0) -- (11.61,5.58) -- cycle    ;
%Curve Lines [id:da46264718836376073] 
\draw [color={rgb, 255:red, 74; green, 144; blue, 226 }  ,draw opacity=1 ][line width=1.5]    (89.34,192.72) .. controls (107.05,162.8) and (143.41,171.84) .. (152.34,196.66) ;
\draw [shift={(87.22,196.66)}, rotate = 295.84] [fill={rgb, 255:red, 74; green, 144; blue, 226 }  ,fill opacity=1 ][line width=0.08]  [draw opacity=0] (11.61,-5.58) -- (0,0) -- (11.61,5.58) -- cycle    ;
%Curve Lines [id:da44941926103098273] 
\draw [color={rgb, 255:red, 74; green, 144; blue, 226 }  ,draw opacity=1 ][line width=1.5]    (154.46,192.67) .. controls (172.13,161.99) and (208.21,161.67) .. (224.29,196.66) ;
\draw [shift={(152.34,196.66)}, rotate = 295.84] [fill={rgb, 255:red, 74; green, 144; blue, 226 }  ,fill opacity=1 ][line width=0.08]  [draw opacity=0] (11.61,-5.58) -- (0,0) -- (11.61,5.58) -- cycle    ;
%Curve Lines [id:da22225476038038006] 
\draw [color={rgb, 255:red, 74; green, 144; blue, 226 }  ,draw opacity=1 ][line width=1.5]    (225.66,192.76) .. controls (237.53,162.85) and (269.49,165.16) .. (282.59,194.17) ;
\draw [shift={(224.29,196.66)}, rotate = 287.25] [fill={rgb, 255:red, 74; green, 144; blue, 226 }  ,fill opacity=1 ][line width=0.08]  [draw opacity=0] (11.61,-5.58) -- (0,0) -- (11.61,5.58) -- cycle    ;
%Shape: Ellipse [id:dp42372727795787335] 
\draw  [color={rgb, 255:red, 74; green, 144; blue, 226 }  ,draw opacity=1 ][line width=1.5]  (495.85,43.02) .. controls (495.85,39.75) and (498.49,37.1) .. (501.74,37.1) .. controls (505,37.1) and (507.63,39.75) .. (507.63,43.02) .. controls (507.63,46.29) and (505,48.94) .. (501.74,48.94) .. controls (498.49,48.94) and (495.85,46.29) .. (495.85,43.02) -- cycle ;
%Shape: Ellipse [id:dp1870179348026778] 
\draw  [color={rgb, 255:red, 74; green, 144; blue, 226 }  ,draw opacity=1 ][line width=1.5]  (495.85,124.65) .. controls (495.85,121.38) and (498.49,118.73) .. (501.74,118.73) .. controls (505,118.73) and (507.63,121.38) .. (507.63,124.65) .. controls (507.63,127.92) and (505,130.57) .. (501.74,130.57) .. controls (498.49,130.57) and (495.85,127.92) .. (495.85,124.65) -- cycle ;
%Shape: Ellipse [id:dp6111256363608131] 
\draw  [color={rgb, 255:red, 74; green, 144; blue, 226 }  ,draw opacity=1 ][line width=1.5]  (396,201.92) .. controls (396,198.65) and (398.64,196) .. (401.89,196) .. controls (405.14,196) and (407.78,198.65) .. (407.78,201.92) .. controls (407.78,205.18) and (405.14,207.83) .. (401.89,207.83) .. controls (398.64,207.83) and (396,205.18) .. (396,201.92) -- cycle ;
%Shape: Ellipse [id:dp9315304169845421] 
\draw  [color={rgb, 255:red, 74; green, 144; blue, 226 }  ,draw opacity=1 ][line width=1.5]  (591.36,199.42) .. controls (591.36,196.15) and (594,193.5) .. (597.25,193.5) .. controls (600.51,193.5) and (603.15,196.15) .. (603.15,199.42) .. controls (603.15,202.69) and (600.51,205.34) .. (597.25,205.34) .. controls (594,205.34) and (591.36,202.69) .. (591.36,199.42) -- cycle ;
%Shape: Ellipse [id:dp322536109170936] 
\draw  [color={rgb, 255:red, 74; green, 144; blue, 226 }  ,draw opacity=1 ][line width=1.5]  (461.12,201.92) .. controls (461.12,198.65) and (463.76,196) .. (467.01,196) .. controls (470.27,196) and (472.9,198.65) .. (472.9,201.92) .. controls (472.9,205.18) and (470.27,207.83) .. (467.01,207.83) .. controls (463.76,207.83) and (461.12,205.18) .. (461.12,201.92) -- cycle ;
%Shape: Ellipse [id:dp9620679121888993] 
\draw  [color={rgb, 255:red, 74; green, 144; blue, 226 }  ,draw opacity=1 ][line width=1.5]  (533.06,201.92) .. controls (533.06,198.65) and (535.7,196) .. (538.95,196) .. controls (542.21,196) and (544.85,198.65) .. (544.85,201.92) .. controls (544.85,205.18) and (542.21,207.83) .. (538.95,207.83) .. controls (535.7,207.83) and (533.06,205.18) .. (533.06,201.92) -- cycle ;
%Curve Lines [id:da7735261830071775] 
\draw [color={rgb, 255:red, 74; green, 144; blue, 226 }  ,draw opacity=1 ][line width=1.5]    (504.99,116.24) .. controls (538.81,89.8) and (536.03,75.23) .. (501.74,48.94) ;
\draw [shift={(501.74,118.73)}, rotate = 323] [fill={rgb, 255:red, 74; green, 144; blue, 226 }  ,fill opacity=1 ][line width=0.08]  [draw opacity=0] (11.61,-5.58) -- (0,0) -- (11.61,5.58) -- cycle    ;
%Curve Lines [id:da5332538798838518] 
\draw [color={rgb, 255:red, 74; green, 144; blue, 226 }  ,draw opacity=1 ][line width=1.5]    (501.74,118.73) .. controls (467.88,100.2) and (457.6,76.04) .. (498.5,50.88) ;
\draw [shift={(501.74,48.94)}, rotate = 149.92] [fill={rgb, 255:red, 74; green, 144; blue, 226 }  ,fill opacity=1 ][line width=0.08]  [draw opacity=0] (11.61,-5.58) -- (0,0) -- (11.61,5.58) -- cycle    ;
%Curve Lines [id:da7801197789531258] 
\draw [color={rgb, 255:red, 74; green, 144; blue, 226 }  ,draw opacity=1 ][line width=1.5]    (396,201.92) .. controls (360.87,173.68) and (422.92,144.84) .. (498.29,131.18) ;
\draw [shift={(501.74,130.57)}, rotate = 170.12] [fill={rgb, 255:red, 74; green, 144; blue, 226 }  ,fill opacity=1 ][line width=0.08]  [draw opacity=0] (11.61,-5.58) -- (0,0) -- (11.61,5.58) -- cycle    ;
%Curve Lines [id:da022731159174359417] 
\draw [color={rgb, 255:red, 74; green, 144; blue, 226 }  ,draw opacity=1 ][line width=1.5]    (606.71,196.49) .. controls (636.44,169.65) and (581.94,145.9) .. (501.74,130.57) ;
\draw [shift={(603.15,199.42)}, rotate = 323] [fill={rgb, 255:red, 74; green, 144; blue, 226 }  ,fill opacity=1 ][line width=0.08]  [draw opacity=0] (11.61,-5.58) -- (0,0) -- (11.61,5.58) -- cycle    ;
%Curve Lines [id:da8965305497852382] 
\draw [color={rgb, 255:red, 74; green, 144; blue, 226 }  ,draw opacity=1 ][line width=1.5]    (603.15,199.42) .. controls (690.47,241.58) and (653.45,-39.52) .. (504,35.93) ;
\draw [shift={(501.74,37.1)}, rotate = 332.3] [fill={rgb, 255:red, 74; green, 144; blue, 226 }  ,fill opacity=1 ][line width=0.08]  [draw opacity=0] (11.61,-5.58) -- (0,0) -- (11.61,5.58) -- cycle    ;
%Curve Lines [id:da10010639301698165] 
\draw [color={rgb, 255:red, 74; green, 144; blue, 226 }  ,draw opacity=1 ][line width=1.5]    (391.62,203.04) .. controls (321.7,215.96) and (333.99,-32.21) .. (501.74,37.1) ;
\draw [shift={(396,201.92)}, rotate = 161.99] [fill={rgb, 255:red, 74; green, 144; blue, 226 }  ,fill opacity=1 ][line width=0.08]  [draw opacity=0] (11.61,-5.58) -- (0,0) -- (11.61,5.58) -- cycle    ;
%Curve Lines [id:da20611488989151905] 
\draw [color={rgb, 255:red, 74; green, 144; blue, 226 }  ,draw opacity=1 ][line width=1.5]    (401.89,207.83) .. controls (409.97,248.72) and (449.43,239.72) .. (465.37,211.03) ;
\draw [shift={(467.01,207.83)}, rotate = 115.22] [fill={rgb, 255:red, 74; green, 144; blue, 226 }  ,fill opacity=1 ][line width=0.08]  [draw opacity=0] (11.61,-5.58) -- (0,0) -- (11.61,5.58) -- cycle    ;
%Curve Lines [id:da6184979101341241] 
\draw [color={rgb, 255:red, 74; green, 144; blue, 226 }  ,draw opacity=1 ][line width=1.5]    (467.01,207.83) .. controls (475.09,248.72) and (524.06,244.69) .. (537.62,211.55) ;
\draw [shift={(538.95,207.83)}, rotate = 107.19] [fill={rgb, 255:red, 74; green, 144; blue, 226 }  ,fill opacity=1 ][line width=0.08]  [draw opacity=0] (11.61,-5.58) -- (0,0) -- (11.61,5.58) -- cycle    ;
%Curve Lines [id:da10958800394354862] 
\draw [color={rgb, 255:red, 74; green, 144; blue, 226 }  ,draw opacity=1 ][line width=1.5]    (538.95,207.83) .. controls (546.95,248.3) and (590.86,233.64) .. (596.69,208.9) ;
\draw [shift={(597.25,205.34)}, rotate = 94.63] [fill={rgb, 255:red, 74; green, 144; blue, 226 }  ,fill opacity=1 ][line width=0.08]  [draw opacity=0] (11.61,-5.58) -- (0,0) -- (11.61,5.58) -- cycle    ;
%Curve Lines [id:da7991255722818383] 
\draw [color={rgb, 255:red, 74; green, 144; blue, 226 }  ,draw opacity=1 ][line width=1.5]    (404,192.05) .. controls (421.71,162.14) and (458.08,171.17) .. (467.01,196) ;
\draw [shift={(401.89,196)}, rotate = 295.84] [fill={rgb, 255:red, 74; green, 144; blue, 226 }  ,fill opacity=1 ][line width=0.08]  [draw opacity=0] (11.61,-5.58) -- (0,0) -- (11.61,5.58) -- cycle    ;
%Curve Lines [id:da5824171246073673] 
\draw [color={rgb, 255:red, 74; green, 144; blue, 226 }  ,draw opacity=1 ][line width=1.5]    (469.12,192) .. controls (486.8,161.32) and (522.88,161) .. (538.95,196) ;
\draw [shift={(467.01,196)}, rotate = 295.84] [fill={rgb, 255:red, 74; green, 144; blue, 226 }  ,fill opacity=1 ][line width=0.08]  [draw opacity=0] (11.61,-5.58) -- (0,0) -- (11.61,5.58) -- cycle    ;
%Curve Lines [id:da15923659585236594] 
\draw [color={rgb, 255:red, 74; green, 144; blue, 226 }  ,draw opacity=1 ][line width=1.5]    (540.33,192.09) .. controls (552.19,162.18) and (584.15,164.49) .. (597.25,193.5) ;
\draw [shift={(538.95,196)}, rotate = 287.25] [fill={rgb, 255:red, 74; green, 144; blue, 226 }  ,fill opacity=1 ][line width=0.08]  [draw opacity=0] (11.61,-5.58) -- (0,0) -- (11.61,5.58) -- cycle    ;

% Text Node
\draw (103.33,96.4) node [anchor=north west][inner sep=0.75pt]    {$A$};
% Text Node
\draw (255.33,92.4) node [anchor=north west][inner sep=0.75pt]    {$B$};
% Text Node
\draw (180,145.07) node [anchor=north west][inner sep=0.75pt]    {$C$};
% Text Node
\draw (50.67,238.4) node [anchor=north west][inner sep=0.75pt]    {$D$};
% Text Node
\draw (110.67,197.73) node [anchor=north west][inner sep=0.75pt]    {$E$};
% Text Node
\draw (180,196.4) node [anchor=north west][inner sep=0.75pt]    {$F$};
% Text Node
\draw (245.33,196.07) node [anchor=north west][inner sep=0.75pt]    {$G$};
% Text Node
\draw (176.67,81.4) node [anchor=north west][inner sep=0.75pt]    {$H$};
% Text Node
\draw (411.33,96.07) node [anchor=north west][inner sep=0.75pt]    {$A$};
% Text Node
\draw (576,96.4) node [anchor=north west][inner sep=0.75pt]    {$B$};
% Text Node
\draw (492.67,78.07) node [anchor=north west][inner sep=0.75pt]    {$H$};
% Text Node
\draw (495.33,145.73) node [anchor=north west][inner sep=0.75pt]    {$C$};
% Text Node
\draw (427.33,197.07) node [anchor=north west][inner sep=0.75pt]    {$E$};
% Text Node
\draw (496,196.4) node [anchor=north west][inner sep=0.75pt]    {$F$};
% Text Node
\draw (561.33,194.07) node [anchor=north west][inner sep=0.75pt]    {$G$};
% Text Node
\draw (140.67,76.07) node [anchor=north west][inner sep=0.75pt]    {$1$};
% Text Node
\draw (303.33,37.4) node [anchor=north west][inner sep=0.75pt]    {$1$};
% Text Node
\draw (218.67,74.4) node [anchor=north west][inner sep=0.75pt]    {$2$};
% Text Node
\draw (113.33,127.73) node [anchor=north west][inner sep=0.75pt]    {$2$};
% Text Node
\draw (257.33,128.07) node [anchor=north west][inner sep=0.75pt]    {$6$};
% Text Node
\draw (246.67,236.07) node [anchor=north west][inner sep=0.75pt]    {$6$};
% Text Node
\draw (181.33,236.4) node [anchor=north west][inner sep=0.75pt]    {$5$};
% Text Node
\draw (242,153.07) node [anchor=north west][inner sep=0.75pt]    {$5$};
% Text Node
\draw (165.33,158.73) node [anchor=north west][inner sep=0.75pt]    {$4$};
% Text Node
\draw (110,233.4) node [anchor=north west][inner sep=0.75pt]    {$4$};
% Text Node
\draw (117.33,156.73) node [anchor=north west][inner sep=0.75pt]    {$3$};
% Text Node
\draw (42.67,46.07) node [anchor=north west][inner sep=0.75pt]    {$3$};
% Text Node
\draw (573.33,156.73) node [anchor=north west][inner sep=0.75pt]    {$6$};
% Text Node
\draw (516,155.07) node [anchor=north west][inner sep=0.75pt]    {$5$};
% Text Node
\draw (436.67,154.07) node [anchor=north west][inner sep=0.75pt]    {$4$};
% Text Node
\draw (424,236.07) node [anchor=north west][inner sep=0.75pt]    {$3$};
% Text Node
\draw (496,233.73) node [anchor=north west][inner sep=0.75pt]    {$4$};
% Text Node
\draw (562.67,234.4) node [anchor=north west][inner sep=0.75pt]    {$5$};
% Text Node
\draw (426,126.73) node [anchor=north west][inner sep=0.75pt]    {$3$};
% Text Node
\draw (564,126.4) node [anchor=north west][inner sep=0.75pt]    {$2$};
% Text Node
\draw (458.67,76.07) node [anchor=north west][inner sep=0.75pt]    {$2$};
% Text Node
\draw (531.33,74.73) node [anchor=north west][inner sep=0.75pt]    {$1$};
% Text Node
\draw (370.67,34.4) node [anchor=north west][inner sep=0.75pt]    {$1$};
% Text Node
\draw (624,44.07) node [anchor=north west][inner sep=0.75pt]    {$6$};


\end{tikzpicture}
        \caption{Descompocision del complemento del nudo $6_1$}
    \end{figure}
    Observe que podemos eliminar las caras $E,F,G$ y $H$ ya que solo tienen dos lados, así si identificamos adecuadamente obtenemos los siguientes dos poliedros que son muy similares a los tetraedros salvo por las caras $C$ y $D$.
    \begin{figure}[H]
    \centering
    %!TEX root = ./main.tex



\tikzset{every picture/.style={line width=0.75pt}} %set default line width to 0.75pt        

\begin{tikzpicture}[x=0.75pt,y=0.75pt,yscale=-1,xscale=1]
%uncomment if require: \path (0,877); %set diagram left start at 0, and has height of 877

%Shape: Ellipse [id:dp2062482206131232] 
\draw  [color={rgb, 255:red, 74; green, 144; blue, 226 }  ,draw opacity=1 ][line width=1.5]  (181.18,43.69) .. controls (181.18,40.42) and (183.82,37.77) .. (187.08,37.77) .. controls (190.33,37.77) and (192.97,40.42) .. (192.97,43.69) .. controls (192.97,46.95) and (190.33,49.6) .. (187.08,49.6) .. controls (183.82,49.6) and (181.18,46.95) .. (181.18,43.69) -- cycle ;
%Shape: Ellipse [id:dp42181631305663214] 
\draw  [color={rgb, 255:red, 74; green, 144; blue, 226 }  ,draw opacity=1 ][line width=1.5]  (181.18,125.31) .. controls (181.18,122.05) and (183.82,119.39) .. (187.08,119.39) .. controls (190.33,119.39) and (192.97,122.05) .. (192.97,125.31) .. controls (192.97,128.58) and (190.33,131.23) .. (187.08,131.23) .. controls (183.82,131.23) and (181.18,128.58) .. (181.18,125.31) -- cycle ;
%Shape: Ellipse [id:dp1893981875046039] 
\draw  [color={rgb, 255:red, 74; green, 144; blue, 226 }  ,draw opacity=1 ][line width=1.5]  (81.33,202.58) .. controls (81.33,199.31) and (83.97,196.66) .. (87.22,196.66) .. controls (90.48,196.66) and (93.11,199.31) .. (93.11,202.58) .. controls (93.11,205.85) and (90.48,208.5) .. (87.22,208.5) .. controls (83.97,208.5) and (81.33,205.85) .. (81.33,202.58) -- cycle ;
%Shape: Ellipse [id:dp3630097335542829] 
\draw  [color={rgb, 255:red, 74; green, 144; blue, 226 }  ,draw opacity=1 ][line width=1.5]  (276.69,200.09) .. controls (276.69,196.82) and (279.33,194.17) .. (282.59,194.17) .. controls (285.84,194.17) and (288.48,196.82) .. (288.48,200.09) .. controls (288.48,203.36) and (285.84,206.01) .. (282.59,206.01) .. controls (279.33,206.01) and (276.69,203.36) .. (276.69,200.09) -- cycle ;
%Shape: Ellipse [id:dp12213112120599945] 
\draw  [color={rgb, 255:red, 74; green, 144; blue, 226 }  ,draw opacity=1 ][line width=1.5]  (146.45,202.58) .. controls (146.45,199.31) and (149.09,196.66) .. (152.34,196.66) .. controls (155.6,196.66) and (158.24,199.31) .. (158.24,202.58) .. controls (158.24,205.85) and (155.6,208.5) .. (152.34,208.5) .. controls (149.09,208.5) and (146.45,205.85) .. (146.45,202.58) -- cycle ;
%Shape: Ellipse [id:dp14956651149377553] 
\draw  [color={rgb, 255:red, 74; green, 144; blue, 226 }  ,draw opacity=1 ][line width=1.5]  (218.4,202.58) .. controls (218.4,199.31) and (221.03,196.66) .. (224.29,196.66) .. controls (227.54,196.66) and (230.18,199.31) .. (230.18,202.58) .. controls (230.18,205.85) and (227.54,208.5) .. (224.29,208.5) .. controls (221.03,208.5) and (218.4,205.85) .. (218.4,202.58) -- cycle ;
%Curve Lines [id:da46118430818328493] 
\draw [color={rgb, 255:red, 74; green, 144; blue, 226 }  ,draw opacity=1 ][line width=1.5]    (80.43,198.63) .. controls (73.79,162.96) and (112.86,144.16) .. (187.08,131.23) ;
\draw [shift={(81.33,202.58)}, rotate = 254.88] [fill={rgb, 255:red, 74; green, 144; blue, 226 }  ,fill opacity=1 ][line width=0.08]  [draw opacity=0] (11.61,-5.58) -- (0,0) -- (11.61,5.58) -- cycle    ;
%Curve Lines [id:da12684350601217753] 
\draw [color={rgb, 255:red, 74; green, 144; blue, 226 }  ,draw opacity=1 ][line width=1.5]    (292.05,197.16) .. controls (321.77,170.32) and (267.27,146.57) .. (187.08,131.23) ;
\draw [shift={(288.48,200.09)}, rotate = 323] [fill={rgb, 255:red, 74; green, 144; blue, 226 }  ,fill opacity=1 ][line width=0.08]  [draw opacity=0] (11.61,-5.58) -- (0,0) -- (11.61,5.58) -- cycle    ;
%Curve Lines [id:da12533726813089074] 
\draw [color={rgb, 255:red, 74; green, 144; blue, 226 }  ,draw opacity=1 ][line width=1.5]    (288.48,200.09) .. controls (375.8,242.25) and (338.78,-38.85) .. (189.34,36.6) ;
\draw [shift={(187.08,37.77)}, rotate = 332.3] [fill={rgb, 255:red, 74; green, 144; blue, 226 }  ,fill opacity=1 ][line width=0.08]  [draw opacity=0] (11.61,-5.58) -- (0,0) -- (11.61,5.58) -- cycle    ;
%Curve Lines [id:da052098391594170956] 
\draw [color={rgb, 255:red, 74; green, 144; blue, 226 }  ,draw opacity=1 ][line width=1.5]    (81.33,202.58) .. controls (6.97,226.76) and (15.81,-30.37) .. (184.52,36.73) ;
\draw [shift={(187.08,37.77)}, rotate = 202.45] [fill={rgb, 255:red, 74; green, 144; blue, 226 }  ,fill opacity=1 ][line width=0.08]  [draw opacity=0] (11.61,-5.58) -- (0,0) -- (11.61,5.58) -- cycle    ;
%Straight Lines [id:da025851980169662614] 
\draw [color={rgb, 255:red, 74; green, 144; blue, 226 }  ,draw opacity=1 ][line width=1.5]    (93.11,202.58) -- (142.45,202.58) ;
\draw [shift={(146.45,202.58)}, rotate = 180] [fill={rgb, 255:red, 74; green, 144; blue, 226 }  ,fill opacity=1 ][line width=0.08]  [draw opacity=0] (11.61,-5.58) -- (0,0) -- (11.61,5.58) -- cycle    ;
%Straight Lines [id:da7108101232087555] 
\draw [color={rgb, 255:red, 74; green, 144; blue, 226 }  ,draw opacity=1 ][line width=1.5]    (162.24,202.58) -- (218.4,202.58) ;
\draw [shift={(158.24,202.58)}, rotate = 0] [fill={rgb, 255:red, 74; green, 144; blue, 226 }  ,fill opacity=1 ][line width=0.08]  [draw opacity=0] (11.61,-5.58) -- (0,0) -- (11.61,5.58) -- cycle    ;
%Straight Lines [id:da4593573245399082] 
\draw [color={rgb, 255:red, 74; green, 144; blue, 226 }  ,draw opacity=1 ][line width=1.5]    (230.18,202.58) -- (272.7,200.3) ;
\draw [shift={(276.69,200.09)}, rotate = 176.93] [fill={rgb, 255:red, 74; green, 144; blue, 226 }  ,fill opacity=1 ][line width=0.08]  [draw opacity=0] (11.61,-5.58) -- (0,0) -- (11.61,5.58) -- cycle    ;
%Straight Lines [id:da3678344551576813] 
\draw [color={rgb, 255:red, 74; green, 144; blue, 226 }  ,draw opacity=1 ][line width=1.5]    (187.08,119.39) -- (187.08,53.6) ;
\draw [shift={(187.08,49.6)}, rotate = 90] [fill={rgb, 255:red, 74; green, 144; blue, 226 }  ,fill opacity=1 ][line width=0.08]  [draw opacity=0] (11.61,-5.58) -- (0,0) -- (11.61,5.58) -- cycle    ;
%Shape: Ellipse [id:dp5834798930697471] 
\draw  [color={rgb, 255:red, 74; green, 144; blue, 226 }  ,draw opacity=1 ][line width=1.5]  (491.85,42.35) .. controls (491.85,39.08) and (494.49,36.43) .. (497.74,36.43) .. controls (501,36.43) and (503.63,39.08) .. (503.63,42.35) .. controls (503.63,45.62) and (501,48.27) .. (497.74,48.27) .. controls (494.49,48.27) and (491.85,45.62) .. (491.85,42.35) -- cycle ;
%Shape: Ellipse [id:dp9667981945481139] 
\draw  [color={rgb, 255:red, 74; green, 144; blue, 226 }  ,draw opacity=1 ][line width=1.5]  (491.85,123.98) .. controls (491.85,120.71) and (494.49,118.06) .. (497.74,118.06) .. controls (501,118.06) and (503.63,120.71) .. (503.63,123.98) .. controls (503.63,127.25) and (501,129.9) .. (497.74,129.9) .. controls (494.49,129.9) and (491.85,127.25) .. (491.85,123.98) -- cycle ;
%Shape: Ellipse [id:dp09059574964185924] 
\draw  [color={rgb, 255:red, 74; green, 144; blue, 226 }  ,draw opacity=1 ][line width=1.5]  (392,201.25) .. controls (392,197.98) and (394.64,195.33) .. (397.89,195.33) .. controls (401.14,195.33) and (403.78,197.98) .. (403.78,201.25) .. controls (403.78,204.52) and (401.14,207.17) .. (397.89,207.17) .. controls (394.64,207.17) and (392,204.52) .. (392,201.25) -- cycle ;
%Shape: Ellipse [id:dp3605131756242862] 
\draw  [color={rgb, 255:red, 74; green, 144; blue, 226 }  ,draw opacity=1 ][line width=1.5]  (587.36,198.76) .. controls (587.36,195.49) and (590,192.84) .. (593.25,192.84) .. controls (596.51,192.84) and (599.15,195.49) .. (599.15,198.76) .. controls (599.15,202.03) and (596.51,204.68) .. (593.25,204.68) .. controls (590,204.68) and (587.36,202.03) .. (587.36,198.76) -- cycle ;
%Shape: Ellipse [id:dp6418673252768715] 
\draw  [color={rgb, 255:red, 74; green, 144; blue, 226 }  ,draw opacity=1 ][line width=1.5]  (457.12,201.25) .. controls (457.12,197.98) and (459.76,195.33) .. (463.01,195.33) .. controls (466.27,195.33) and (468.9,197.98) .. (468.9,201.25) .. controls (468.9,204.52) and (466.27,207.17) .. (463.01,207.17) .. controls (459.76,207.17) and (457.12,204.52) .. (457.12,201.25) -- cycle ;
%Shape: Ellipse [id:dp39805151621059387] 
\draw  [color={rgb, 255:red, 74; green, 144; blue, 226 }  ,draw opacity=1 ][line width=1.5]  (529.06,201.25) .. controls (529.06,197.98) and (531.7,195.33) .. (534.95,195.33) .. controls (538.21,195.33) and (540.85,197.98) .. (540.85,201.25) .. controls (540.85,204.52) and (538.21,207.17) .. (534.95,207.17) .. controls (531.7,207.17) and (529.06,204.52) .. (529.06,201.25) -- cycle ;
%Curve Lines [id:da1957524342127539] 
\draw [color={rgb, 255:red, 74; green, 144; blue, 226 }  ,draw opacity=1 ][line width=1.5]    (391.09,197.3) .. controls (384.45,161.63) and (423.53,142.83) .. (497.74,129.9) ;
\draw [shift={(392,201.25)}, rotate = 254.88] [fill={rgb, 255:red, 74; green, 144; blue, 226 }  ,fill opacity=1 ][line width=0.08]  [draw opacity=0] (11.61,-5.58) -- (0,0) -- (11.61,5.58) -- cycle    ;
%Curve Lines [id:da3321764957558633] 
\draw [color={rgb, 255:red, 74; green, 144; blue, 226 }  ,draw opacity=1 ][line width=1.5]    (599.15,198.76) .. controls (635.8,171.14) and (582.49,146.54) .. (501.46,130.62) ;
\draw [shift={(497.74,129.9)}, rotate = 10.82] [fill={rgb, 255:red, 74; green, 144; blue, 226 }  ,fill opacity=1 ][line width=0.08]  [draw opacity=0] (11.61,-5.58) -- (0,0) -- (11.61,5.58) -- cycle    ;
%Curve Lines [id:da2799756577005189] 
\draw [color={rgb, 255:red, 74; green, 144; blue, 226 }  ,draw opacity=1 ][line width=1.5]    (599.15,198.76) .. controls (686.47,240.92) and (649.45,-40.18) .. (500,35.27) ;
\draw [shift={(497.74,36.43)}, rotate = 332.3] [fill={rgb, 255:red, 74; green, 144; blue, 226 }  ,fill opacity=1 ][line width=0.08]  [draw opacity=0] (11.61,-5.58) -- (0,0) -- (11.61,5.58) -- cycle    ;
%Curve Lines [id:da9026596524821855] 
\draw [color={rgb, 255:red, 74; green, 144; blue, 226 }  ,draw opacity=1 ][line width=1.5]    (387.62,202.37) .. controls (317.7,215.29) and (329.99,-32.88) .. (497.74,36.43) ;
\draw [shift={(392,201.25)}, rotate = 161.99] [fill={rgb, 255:red, 74; green, 144; blue, 226 }  ,fill opacity=1 ][line width=0.08]  [draw opacity=0] (11.61,-5.58) -- (0,0) -- (11.61,5.58) -- cycle    ;
%Straight Lines [id:da29172929203159703] 
\draw [color={rgb, 255:red, 74; green, 144; blue, 226 }  ,draw opacity=1 ][line width=1.5]    (407.78,201.25) -- (457.12,201.25) ;
\draw [shift={(403.78,201.25)}, rotate = 0] [fill={rgb, 255:red, 74; green, 144; blue, 226 }  ,fill opacity=1 ][line width=0.08]  [draw opacity=0] (11.61,-5.58) -- (0,0) -- (11.61,5.58) -- cycle    ;
%Straight Lines [id:da23336632170754767] 
\draw [color={rgb, 255:red, 74; green, 144; blue, 226 }  ,draw opacity=1 ][line width=1.5]    (468.9,201.25) -- (525.06,201.25) ;
\draw [shift={(529.06,201.25)}, rotate = 180] [fill={rgb, 255:red, 74; green, 144; blue, 226 }  ,fill opacity=1 ][line width=0.08]  [draw opacity=0] (11.61,-5.58) -- (0,0) -- (11.61,5.58) -- cycle    ;
%Straight Lines [id:da7582939485718254] 
\draw [color={rgb, 255:red, 74; green, 144; blue, 226 }  ,draw opacity=1 ][line width=1.5]    (544.84,201.03) -- (587.36,198.76) ;
\draw [shift={(540.85,201.25)}, rotate = 356.93] [fill={rgb, 255:red, 74; green, 144; blue, 226 }  ,fill opacity=1 ][line width=0.08]  [draw opacity=0] (11.61,-5.58) -- (0,0) -- (11.61,5.58) -- cycle    ;
%Straight Lines [id:da2633005133934979] 
\draw [color={rgb, 255:red, 74; green, 144; blue, 226 }  ,draw opacity=1 ][line width=1.5]    (497.74,114.06) -- (497.74,48.27) ;
\draw [shift={(497.74,118.06)}, rotate = 270] [fill={rgb, 255:red, 74; green, 144; blue, 226 }  ,fill opacity=1 ][line width=0.08]  [draw opacity=0] (11.61,-5.58) -- (0,0) -- (11.61,5.58) -- cycle    ;

% Text Node
\draw (103.33,96.4) node [anchor=north west][inner sep=0.75pt]    {$A$};
% Text Node
\draw (255.33,92.4) node [anchor=north west][inner sep=0.75pt]    {$B$};
% Text Node
\draw (180,145.07) node [anchor=north west][inner sep=0.75pt]    {$C$};
% Text Node
\draw (50.67,238.4) node [anchor=north west][inner sep=0.75pt]    {$D$};
% Text Node
\draw (303.33,37.4) node [anchor=north west][inner sep=0.75pt]    {$1$};
% Text Node
\draw (113.33,127.73) node [anchor=north west][inner sep=0.75pt]    {$1$};
% Text Node
\draw (257.33,128.07) node [anchor=north west][inner sep=0.75pt]    {$4$};
% Text Node
\draw (42.67,46.07) node [anchor=north west][inner sep=0.75pt]    {$4$};
% Text Node
\draw (114,206.07) node [anchor=north west][inner sep=0.75pt]    {$4$};
% Text Node
\draw (184.67,204.07) node [anchor=north west][inner sep=0.75pt]    {$4$};
% Text Node
\draw (242,205.4) node [anchor=north west][inner sep=0.75pt]    {$4$};
% Text Node
\draw (170,79.73) node [anchor=north west][inner sep=0.75pt]    {$1$};
% Text Node
\draw (414,95.07) node [anchor=north west][inner sep=0.75pt]    {$A$};
% Text Node
\draw (566,91.07) node [anchor=north west][inner sep=0.75pt]    {$B$};
% Text Node
\draw (490.67,143.73) node [anchor=north west][inner sep=0.75pt]    {$C$};
% Text Node
\draw (361.33,237.07) node [anchor=north west][inner sep=0.75pt]    {$D$};
% Text Node
\draw (614,36.07) node [anchor=north west][inner sep=0.75pt]    {$4$};
% Text Node
\draw (424,126.4) node [anchor=north west][inner sep=0.75pt]    {$4$};
% Text Node
\draw (568,126.73) node [anchor=north west][inner sep=0.75pt]    {$1$};
% Text Node
\draw (353.33,44.73) node [anchor=north west][inner sep=0.75pt]    {$1$};
% Text Node
\draw (424.67,204.73) node [anchor=north west][inner sep=0.75pt]    {$4$};
% Text Node
\draw (495.33,202.73) node [anchor=north west][inner sep=0.75pt]    {$4$};
% Text Node
\draw (559.33,204.07) node [anchor=north west][inner sep=0.75pt]    {$4$};
% Text Node
\draw (480.67,78.4) node [anchor=north west][inner sep=0.75pt]    {$1$};


\end{tikzpicture}
        \caption{Descompocision reduciendo las caras de dos aristas}
    \end{figure}
    
\end{eg}
Observe que particularmente este proceso es primero bastante mas fácil de realizar de lo que en un principio se podría pensar, pero ademas de esto, nos da una descompocision relativamente sencilla que nos permite visualizar la tres variedad, y no solo eso, nos da indicios de un análogo de la descomposición de dos variedades por medio de triangulaciones. Un comentario extra es que esta construcción también funciona en nudos y enlaces alternantes salvo que hay que hacer simplificaciones similares a la reducción de caras donde solo hay dos lados.
    \section{Estructura Hiperbolica del complemento de un nudo}
    Finalmente hemos llegado la parte donde con ayuda de las descomposiciones obtenidas podemos dotar a nuestras estructura, en un principio, netamente topológicas, de una estructura geométrica, en nuestro caso hiperbólica, naturalmente surge la duda de si esto se puede hacer para cualquier complemento de nudo y de que forma lo haremos, ademas nos interesa saber si la métrica inducida es completa, a continuación presentaremos una forma de hacerlo, mas se puede dar condiciones mas generales incluso especificar en el caso de la estructura hiperbólica dos dimensional. Similar a la sección anterior iremos haciendo nuestro ejemplo característico y algún otro posterior a introducir las condiciones
    \subsection{Ecuaciones de Pegado}
    %!TEX root = ./main.tex
\begin{definition}
   Sea $M$ una 3-variedad, decimos que $M$ tiene una estructura hiperbólica si para cada punto $x\in M$ existe una vecindad $U$ isometrica a una bola abierta en $\bb{H}^3$. Ademas decimos que esta estructura es completa, si la métrica inducida antes en $M$ es completa.
\end{definition}

A lo largo del curso hemos visto como las herramientas de cortado y pegado son nuestras grandes aliadas en ámbitos topológicos, esto es lo que nos permite visualizar de una manera diferente diferentes superficies por medio de identificaciones de polígonos, en nuestro caso mas estándar cuadrados. Ya hemos visto que nuestro complemento de nudos se puede descomponer en poliedros ideales, pero no necesariamente estos poliedros son los que queremos como descompocision sino que queremos el tipo de poliedro mas básico, los tetraedros
\begin{definition}
    Sea $M$ una 3-variedad, una triangulación topológica ideal de $M$ es una manera combinatoria de pegar tetraedros truncados (tetraedros ideales) tal que el resultado es homeomorfo a $M$, las partes truncadas corresponderán a la frontera de $M$. 
\end{definition}

Recordemos que previamente ya teníamos una descompocision así para el caso del nudo de 8, pero en nuestros otros dos ejemplos no tenemos algo así, mas posteriormente veremos que dada la descompocision en poliedros, no es difícil llegar a la de tetraedros. 

\begin{definition}
    Una triangulación ideal geométrica de $M$ es una triangulación topológica ideal donde cada tetraedro tiene una estructura hiperbólica positivamente orientada, y que el resultado de el pegado de justamente una variedad suave con una métrica completa.
\end{definition}

Así consideremos un tetraedro ideal cualquiera de una descompocision $T$, este podemos verlo embebido en $\bb{H}^3$, este tetraedro tiene 6 aristas sin preocuparnos por la identificación, si fijamos una $e$ podemos tomas una isometria que lleve los vértices ideales de la arista a $0$ y $\infty$, y un tercer vértice ideal a 1, recordemos que esta elección de tres puntos determina una única isometria, por lo que el cuarto vértice ideal es enviado a algún punto $z^\prime\in\bb{C}$, Observe que podemos asumir que este tiene parte imaginaria positiva, si no aplique la isometria que actuá sobre $\partial\bb{H}^3$ por medio de $z\mapsto \dfrac{z}{z^\prime}$, note que esta fija a $0$ y $\infty$, enviá a $z^\prime$ a $1$, y por ultimo a 1 lo enviá a un complejo con parte imaginaria positiva.

\begin{definition}
    Dada la convención anterior, defina el invariante de arista $e$, $z(e)\in\bb{C}$, como el numero complejo con parte imaginaria positiva obtenida de aplicar la isometria única de $\bb{H}^3$ construida antes.
\end{definition}

\begin{note}
    Es posible que este ultimo vértice se repita con alguno de los primeros, es decir $0,1,\infty$, en dado caso decimos que el tetraedro es degenerado, si se asigna a la recta real pero no a los puntos mencionados anteriormente decimos que es plano. en ninguno de los dos casos no podremos asignar la estructura hiperbólica deseada.
\end{note}
Observe que a partir de un tetraedro ideal ya ubicado uno podría cambiar los vértices ideales entre si para así obtener invariantes de aristas para las otras aristas del tetraedro, el siguiente lema muestra como se relacionan
\begin{lemma}
    Sea $T$ un tetraedro ideal con arista $e_1$, tal que su imagen a través de la isometria enviá los vértices de $T$ a $\infty,0,1$ y $z(e_1)$, tal que la imagen de $e_1$, se encuentra entre los vértices 0 y $\infty$, luego $T$ tiene los siguiente invariantes de aristas
    \begin{itemize}
        \item Si $e^\prime_1$ es la arista con 1 y $z(e_1)$, tenemos que $z(e^\prime_1)=z(e_1)$.
        \item La arista $e_2$ con vertices en $\infty$ y 1 su invariante es $z(e_2)=\dfrac{1}{1-z(e_1)}$. 
        \item La arista $e_3$ con vertices $\infty$ y $z(e_1)$ su invariante es $z(e_3)=\dfrac{z(e_1)-1}{z(e_1)}.$
    \end{itemize}
\end{lemma}
\begin{proof}
    Basta en cada caso tomar una isometria que mande los puntos finales de cada arista a la geodésica de $0$ a $\infty$ y mirar la imagen del punto restante. Por facilidad en la notación $z=z(e_1).$ 
    \begin{itemize}
        \item Si queremos mandar $1\mapsto 0$,  $z\mapsto\infty$ y $0\mapsto 1$ esto nos determina la transformación en la frontera $w\mapsto \dfrac{zw-z}{w-z}$, el punto $\infty$ lo enviá nuevamente a $z$ por lo que $z(e_1^\prime)=z.$
        \item Como ya $\infty$ esta fijo en nuestra arista queremos entonces enviar $1\mapsto 0$ y $z\mapsto 1$, esto nos determina la transformación $w\mapsto\dfrac{w-1}{z-1}$, así  $0\mapsto\dfrac{1}{1-z}$, por lo que $z(e_2)=\dfrac{1}{1-z}.$
        \item De manera similar fijando el infinito, y llevando $z\mapsto 0$ y $0\mapsto 1$ tenemos $w\mapsto\dfrac{w-z}{-z}$, así $1\mapsto \dfrac{z-1}{z}$, es decir $z(e_3)=\dfrac{z-1}{z}$
    \end{itemize}
\end{proof}
Observe que el hecho de que los demás invariantes dependan de solo uno a simple vista parece útil ya que quita la dependencia de la elección de la arista inicial. Veamos ahora como pegar los tetraedros. Concentrémonos en una arista $e$ identificada y fijemos esta como nuestro punto de partida. Sea $T_1$ el primer tetraedro con arista $e_1$ pegada a $e$, donde colocamos a la arista $e_1$ de $0$ a $\infty$, con un tercer vértice en uno y el cuarto en $z(e_1)$, luego sea $T_2$ el siguiente tetraedro con arista $e_2$, al pegarla a $e_1$, se pegara una cara $F_1$ de $T_1$ con vértices $0,z(e_1)$ y $\infty$ a una cara $F_2$, este nuevo tetraedro lo embebemos en primera instancia de la manera estándar con la cara $F_2$ en los vértices $0,\infty$ y $1$ y posterior con la isometria $w\mapsto z(e_1)w$ fijamos al $0$ y $\infty$, enviamos 1 a $z(e_1)$ y por tanto el vértice restante de $T_2$, $z(e_2)$ es enviado a $z(e_1)z(e_2)$, siguiendo este proceso iterativa mente en sentido antihorario de pegado tendremos algo como en la siguiente figura
\begin{figure}[H]
\centering

%!TEX root = ./main.tex



\tikzset{every picture/.style={line width=0.75pt}} %set default line width to 0.75pt        

\begin{tikzpicture}[x=0.75pt,y=0.75pt,yscale=-1,xscale=1]
%uncomment if require: \path (0,300); %set diagram left start at 0, and has height of 300

%Shape: Triangle [id:dp3094193515878545] 
\draw  [line width=1.5]  (306.18,69.64) -- (402.79,88.57) -- (326.67,178.33) -- cycle ;
%Shape: Triangle [id:dp975605725949933] 
\draw  [line width=1.5]  (403,88.67) -- (455.67,178.33) -- (326.67,178.33) -- cycle ;
%Shape: Triangle [id:dp6752076580728734] 
\draw  [line width=1.5]  (326.45,178.47) -- (202.62,134.91) -- (306.18,69.64) -- cycle ;

% Text Node
\draw (352,154.07) node [anchor=north west][inner sep=0.75pt]    {$e_{1}$};
% Text Node
\draw (327.33,133.4) node [anchor=north west][inner sep=0.75pt]    {$e_{2}$};
% Text Node
\draw (294.67,142.73) node [anchor=north west][inner sep=0.75pt]    {$e_{3}$};
% Text Node
\draw (324,184.4) node [anchor=north west][inner sep=0.75pt]    {$0$};
% Text Node
\draw (448.67,178.73) node [anchor=north west][inner sep=0.75pt]    {$1$};
% Text Node
\draw (408.67,70.07) node [anchor=north west][inner sep=0.75pt]    {$z( e_{1})$};
% Text Node
\draw (303.33,42.73) node [anchor=north west][inner sep=0.75pt]    {$z( e_{1}) z( e_{2})$};
% Text Node
\draw (93.33,118.07) node [anchor=north west][inner sep=0.75pt]    {$z( e_{1}) z( e_{2}) z( e_{3})$};


\end{tikzpicture}
    \caption{Vertices de triangulos pegados}
\end{figure}
Observe que cada vez multiplicamos el invariante y dependiendo del angulo entre las caras del tetraedro se ira sumando un angulo diferente que en particular queremos que coincida con dar una vuelta completa, esta intuición justifica el siguiente teorema
\begin{theorem}
    Sea $M$ una 3-variedad que admite una triangulación topológica ideal tal que cada tetraedro tiene una estructura hiperbólica. La estructura hiperbólica en los tetraedros induce una estructura hiperbólica en el pegado si y solo si para cada arista $e$ tenemos que 
    $$\prod z(e_i)=1\quad\text{y}\quad\sum\arg z(e_i)=2\pi.$$
\end{theorem}
En primera instancia podríamos pensar que tenemos muchas variables, pero como todos los invariantes de aristas dependen de uno solo por el lema 4.1, cada tetraedro aporta a lo mas una variable. Con esta ecuación dotemos a nuestro ejemplo principal.
\begin{eg}
    Recordemos la descompocision en dos tetraedros ideales que ya teníamos del nudo de $8$
    \begin{figure}[H]
    \centering
    
%!TEX root = ./main.tex



\tikzset{every picture/.style={line width=0.75pt}} %set default line width to 0.75pt        

\begin{tikzpicture}[x=0.75pt,y=0.75pt,yscale=-1,xscale=1]
%uncomment if require: \path (0,877); %set diagram left start at 0, and has height of 877

%Shape: Circle [id:dp84958386512667] 
\draw  [line width=1.5]  (189.33,55.33) .. controls (189.33,52.39) and (191.72,50) .. (194.67,50) .. controls (197.61,50) and (200,52.39) .. (200,55.33) .. controls (200,58.28) and (197.61,60.67) .. (194.67,60.67) .. controls (191.72,60.67) and (189.33,58.28) .. (189.33,55.33) -- cycle ;
%Shape: Circle [id:dp7070798701116545] 
\draw  [line width=1.5]  (190,145.33) .. controls (190,142.39) and (192.39,140) .. (195.33,140) .. controls (198.28,140) and (200.67,142.39) .. (200.67,145.33) .. controls (200.67,148.28) and (198.28,150.67) .. (195.33,150.67) .. controls (192.39,150.67) and (190,148.28) .. (190,145.33) -- cycle ;
%Shape: Circle [id:dp5401287447911013] 
\draw  [line width=1.5]  (114.67,210) .. controls (114.67,207.05) and (117.05,204.67) .. (120,204.67) .. controls (122.95,204.67) and (125.33,207.05) .. (125.33,210) .. controls (125.33,212.95) and (122.95,215.33) .. (120,215.33) .. controls (117.05,215.33) and (114.67,212.95) .. (114.67,210) -- cycle ;
%Shape: Circle [id:dp024366460256337152] 
\draw  [line width=1.5]  (264.67,210) .. controls (264.67,207.05) and (267.05,204.67) .. (270,204.67) .. controls (272.95,204.67) and (275.33,207.05) .. (275.33,210) .. controls (275.33,212.95) and (272.95,215.33) .. (270,215.33) .. controls (267.05,215.33) and (264.67,212.95) .. (264.67,210) -- cycle ;
%Straight Lines [id:da5819091680445531] 
\draw [line width=1.5]    (194.67,60.67) -- (195.3,136) ;
\draw [shift={(195.33,140)}, rotate = 269.52] [fill={rgb, 255:red, 0; green, 0; blue, 0 }  ][line width=0.08]  [draw opacity=0] (11.61,-5.58) -- (0,0) -- (11.61,5.58) -- cycle    ;
\draw [shift={(195,100.33)}, rotate = 269.52] [color={rgb, 255:red, 0; green, 0; blue, 0 }  ][line width=1.5]    (0,6.71) -- (0,-6.71)   ;
%Straight Lines [id:da6153154443999455] 
\draw [line width=1.5]    (125.33,210) -- (260.67,210) ;
\draw [shift={(264.67,210)}, rotate = 180] [fill={rgb, 255:red, 0; green, 0; blue, 0 }  ][line width=0.08]  [draw opacity=0] (11.61,-5.58) -- (0,0) -- (11.61,5.58) -- cycle    ;
\draw [shift={(191.5,210)}, rotate = 180] [color={rgb, 255:red, 0; green, 0; blue, 0 }  ][line width=1.5]    (0,6.71) -- (0,-6.71)(-6.03,6.71) -- (-6.03,-6.71)   ;
%Straight Lines [id:da5399323102122296] 
\draw [line width=1.5]    (120,204.67) -- (186.95,147.92) ;
\draw [shift={(190,145.33)}, rotate = 139.71] [fill={rgb, 255:red, 0; green, 0; blue, 0 }  ][line width=0.08]  [draw opacity=0] (11.61,-5.58) -- (0,0) -- (11.61,5.58) -- cycle    ;
\draw [shift={(155,175)}, rotate = 139.71] [color={rgb, 255:red, 0; green, 0; blue, 0 }  ][line width=1.5]    (0,6.71) -- (0,-6.71)   ;
%Straight Lines [id:da7687719054145034] 
\draw [line width=1.5]    (200.67,145.33) -- (266.96,202.07) ;
\draw [shift={(270,204.67)}, rotate = 220.56] [fill={rgb, 255:red, 0; green, 0; blue, 0 }  ][line width=0.08]  [draw opacity=0] (11.61,-5.58) -- (0,0) -- (11.61,5.58) -- cycle    ;
\draw [shift={(232.67,172.72)}, rotate = 220.56] [color={rgb, 255:red, 0; green, 0; blue, 0 }  ][line width=1.5]    (0,6.71) -- (0,-6.71)(-6.03,6.71) -- (-6.03,-6.71)   ;
%Curve Lines [id:da44100647811810356] 
\draw [line width=1.5]    (114.67,210) .. controls (115.64,114.62) and (150.17,70.98) .. (186.01,56.59) ;
\draw [shift={(189.33,55.33)}, rotate = 160.64] [fill={rgb, 255:red, 0; green, 0; blue, 0 }  ][line width=0.08]  [draw opacity=0] (11.61,-5.58) -- (0,0) -- (11.61,5.58) -- cycle    ;
\draw [shift={(127.57,125.14)}, rotate = 110.35] [color={rgb, 255:red, 0; green, 0; blue, 0 }  ][line width=1.5]    (0,6.71) -- (0,-6.71)(-6.03,6.71) -- (-6.03,-6.71)   ;
%Curve Lines [id:da6379686625487593] 
\draw [line width=1.5]    (200,55.33) .. controls (239.98,55.01) and (285.97,157.03) .. (276.2,206.32) ;
\draw [shift={(275.33,210)}, rotate = 285.26] [fill={rgb, 255:red, 0; green, 0; blue, 0 }  ][line width=0.08]  [draw opacity=0] (11.61,-5.58) -- (0,0) -- (11.61,5.58) -- cycle    ;
\draw [shift={(259.87,118.95)}, rotate = 246.87] [color={rgb, 255:red, 0; green, 0; blue, 0 }  ][line width=1.5]    (0,6.71) -- (0,-6.71)   ;
%Shape: Circle [id:dp034440490903057674] 
\draw  [line width=1.5]  (431.81,55.33) .. controls (431.81,52.39) and (434.2,50) .. (437.14,50) .. controls (440.09,50) and (442.47,52.39) .. (442.47,55.33) .. controls (442.47,58.28) and (440.09,60.67) .. (437.14,60.67) .. controls (434.2,60.67) and (431.81,58.28) .. (431.81,55.33) -- cycle ;
%Shape: Circle [id:dp4739299080174333] 
\draw  [line width=1.5]  (432.47,145.33) .. controls (432.47,142.39) and (434.86,140) .. (437.81,140) .. controls (440.75,140) and (443.14,142.39) .. (443.14,145.33) .. controls (443.14,148.28) and (440.75,150.67) .. (437.81,150.67) .. controls (434.86,150.67) and (432.47,148.28) .. (432.47,145.33) -- cycle ;
%Shape: Circle [id:dp2986589831675226] 
\draw  [line width=1.5]  (357.14,210) .. controls (357.14,207.05) and (359.53,204.67) .. (362.47,204.67) .. controls (365.42,204.67) and (367.81,207.05) .. (367.81,210) .. controls (367.81,212.95) and (365.42,215.33) .. (362.47,215.33) .. controls (359.53,215.33) and (357.14,212.95) .. (357.14,210) -- cycle ;
%Shape: Circle [id:dp6438231834616311] 
\draw  [line width=1.5]  (507.14,210) .. controls (507.14,207.05) and (509.53,204.67) .. (512.47,204.67) .. controls (515.42,204.67) and (517.81,207.05) .. (517.81,210) .. controls (517.81,212.95) and (515.42,215.33) .. (512.47,215.33) .. controls (509.53,215.33) and (507.14,212.95) .. (507.14,210) -- cycle ;
%Straight Lines [id:da019687389237552866] 
\draw [line width=1.5]    (437.17,64.67) -- (437.81,140) ;
\draw [shift={(437.47,100.33)}, rotate = 269.52] [color={rgb, 255:red, 0; green, 0; blue, 0 }  ][line width=1.5]    (0,6.71) -- (0,-6.71)   ;
\draw [shift={(437.14,60.67)}, rotate = 89.52] [fill={rgb, 255:red, 0; green, 0; blue, 0 }  ][line width=0.08]  [draw opacity=0] (11.61,-5.58) -- (0,0) -- (11.61,5.58) -- cycle    ;
%Straight Lines [id:da5743118307270256] 
\draw [line width=1.5]    (367.81,210) -- (503.14,210) ;
\draw [shift={(507.14,210)}, rotate = 180] [fill={rgb, 255:red, 0; green, 0; blue, 0 }  ][line width=0.08]  [draw opacity=0] (11.61,-5.58) -- (0,0) -- (11.61,5.58) -- cycle    ;
\draw [shift={(433.97,210)}, rotate = 180] [color={rgb, 255:red, 0; green, 0; blue, 0 }  ][line width=1.5]    (0,6.71) -- (0,-6.71)(-6.03,6.71) -- (-6.03,-6.71)   ;
%Straight Lines [id:da15352339522122915] 
\draw [line width=1.5]    (365.53,202.08) -- (432.47,145.33) ;
\draw [shift={(397.47,175)}, rotate = 139.71] [color={rgb, 255:red, 0; green, 0; blue, 0 }  ][line width=1.5]    (0,6.71) -- (0,-6.71)   ;
\draw [shift={(362.47,204.67)}, rotate = 319.71] [fill={rgb, 255:red, 0; green, 0; blue, 0 }  ][line width=0.08]  [draw opacity=0] (11.61,-5.58) -- (0,0) -- (11.61,5.58) -- cycle    ;
%Straight Lines [id:da15958460866160962] 
\draw [line width=1.5]    (443.14,145.33) -- (509.43,202.07) ;
\draw [shift={(512.47,204.67)}, rotate = 220.56] [fill={rgb, 255:red, 0; green, 0; blue, 0 }  ][line width=0.08]  [draw opacity=0] (11.61,-5.58) -- (0,0) -- (11.61,5.58) -- cycle    ;
\draw [shift={(475.15,172.72)}, rotate = 220.56] [color={rgb, 255:red, 0; green, 0; blue, 0 }  ][line width=1.5]    (0,6.71) -- (0,-6.71)(-6.03,6.71) -- (-6.03,-6.71)   ;
%Curve Lines [id:da2355671835877725] 
\draw [line width=1.5]    (357.14,210) .. controls (358.11,114.62) and (392.64,70.98) .. (428.48,56.59) ;
\draw [shift={(431.81,55.33)}, rotate = 160.64] [fill={rgb, 255:red, 0; green, 0; blue, 0 }  ][line width=0.08]  [draw opacity=0] (11.61,-5.58) -- (0,0) -- (11.61,5.58) -- cycle    ;
\draw [shift={(370.05,125.14)}, rotate = 110.35] [color={rgb, 255:red, 0; green, 0; blue, 0 }  ][line width=1.5]    (0,6.71) -- (0,-6.71)(-6.03,6.71) -- (-6.03,-6.71)   ;
%Curve Lines [id:da10269880002458354] 
\draw [line width=1.5]    (446.8,55.69) .. controls (486.71,62.44) and (530.35,164) .. (517.81,210) ;
\draw [shift={(504.08,122.96)}, rotate = 247.81] [color={rgb, 255:red, 0; green, 0; blue, 0 }  ][line width=1.5]    (0,6.71) -- (0,-6.71)   ;
\draw [shift={(442.47,55.33)}, rotate = 359.53] [fill={rgb, 255:red, 0; green, 0; blue, 0 }  ][line width=0.08]  [draw opacity=0] (11.61,-5.58) -- (0,0) -- (11.61,5.58) -- cycle    ;

% Text Node
\draw (394,112.4) node [anchor=north west][inner sep=0.75pt]    {$A$};
% Text Node
\draw (467,112.4) node [anchor=north west][inner sep=0.75pt]    {$B$};
% Text Node
\draw (431,172.4) node [anchor=north west][inner sep=0.75pt]    {$C$};
% Text Node
\draw (384,222.4) node [anchor=north west][inner sep=0.75pt]    {$D$};
% Text Node
\draw (224,122.4) node [anchor=north west][inner sep=0.75pt]    {$A$};
% Text Node
\draw (157,122.4) node [anchor=north west][inner sep=0.75pt]    {$B$};
% Text Node
\draw (191,172.4) node [anchor=north west][inner sep=0.75pt]    {$C$};
% Text Node
\draw (154,222.4) node [anchor=north west][inner sep=0.75pt]    {$D$};
% Text Node
\draw (132,160.4) node [anchor=north west][inner sep=0.75pt]    {$z_{1}$};
% Text Node
\draw (271,112.4) node [anchor=north west][inner sep=0.75pt]    {$z_{1}$};
% Text Node
\draw (244,162.4) node [anchor=north west][inner sep=0.75pt]    {$z_{2}$};
% Text Node
\draw (102,122.4) node [anchor=north west][inner sep=0.75pt]    {$z_{2}$};
% Text Node
\draw (201,212.4) node [anchor=north west][inner sep=0.75pt]    {$z_{3}$};
% Text Node
\draw (202,82.4) node [anchor=north west][inner sep=0.75pt]    {$z_{3}$};
% Text Node
\draw (385,148.4) node [anchor=north west][inner sep=0.75pt]    {$w_{1}$};
% Text Node
\draw (511,122.4) node [anchor=north west][inner sep=0.75pt]    {$w_{1}$};
% Text Node
\draw (461,142.4) node [anchor=north west][inner sep=0.75pt]    {$w_{2}$};
% Text Node
\draw (341,122.4) node [anchor=north west][inner sep=0.75pt]    {$w_{2}$};
% Text Node
\draw (441,82.4) node [anchor=north west][inner sep=0.75pt]    {$w_{3}$};
% Text Node
\draw (447,210.4) node [anchor=north west][inner sep=0.75pt]    {$w_{3}$};


\end{tikzpicture}
        \caption{tetraedros ideales del nudo de 8 con invariantes de aristas}
    \end{figure}
    Observe que por ser mas consistentes cambiamos la vista del poliedro de arriba a la parte exterior, lo que ocasiona una reflexión en el diagrama obtenido antes, ademas observe que a cada arista la hemos dotado de un numero complejo $z_i$ y $w_i$, justamente este es el invariante de arista asociado a enviar esa arista a la arista que va de $0$ a $\infty$, note que por el primer ítem del lema 4.1 solo se necesitan tres indices en primera instancia. Así para cada arista tenemos respectivamente
    $$z_1^2z_3w_1^2w_3=1\quad\text{y}\quad z_2^2z_3w_2^2w_3=1.$$
    Observe que podemos simplificar aun mas usando los otros dos ítems, si $z_1=z$ y $w_1=w$, tenemos que las ecuaciones se convierten en
    $$z^2\frac{z-1}{z}w^2\frac{w-1}{w}=1\quad\text{y}\quad \frac{1}{(1-z)^2}\frac{z-1}{z}\frac{1}{(1-w)^2}\frac{w-1}{w}=1,$$
    Observe que en ambos casos esto se reduce a la única ecuación
    $$z(z-1)w(w-1)=1,$$
    por lo que si solucionamos la ecuación cuadrática en términos de $z$ las soluciones son
    $$z=\frac{1\pm\sqrt{1+\dfrac{4}{w(w-1)}}}{2},$$
    Esto nos determina las soluciones condicionadas netamente a $w$, primero recordemos que $w$ tiene parte imaginaria positiva por hipótesis, queremos de manera similar que $z$ tenga parte imaginaria positiva, esto se consigue siempre que $1+\dfrac{4}{w(w-1)}$ no sea un numero real positivo. Primero observe que es un numero real positivo si $\Im(w(w-1))= 0$, de aquí uno concluye que $\Re(w)=\dfrac{1}{2}$, y segundo necesitaríamos que $1+\dfrac{4}{w(w-1)}>0$, luego de algunas simplificaciones usando la condición anterior se llega a la cuadrática $4y^2-15>0$, donde $y=\Im(w)$, de ahí se concluye que $\Im(w)>\dfrac{\sqrt{15}}{2}$, por lo que nuestro conjunto de soluciones para $w$ que nos dan que $z$ tiene parte imaginaria positiva y dotan del pegado de una estructura hiperbólica es
    $$\bb{H}^2\setminus\left\{w\in\bb{C}:\Re(w)=\frac{1}{2}\quad\text{y}\quad \Im(w)>\frac{\sqrt{15}}{2}\right\}$$ 
\end{eg}

\subsection{Ecuaciones de Completez}
Antes solo observábamos el complemento del nudo, mas ahora la idea sera observar el exterior del nudo, ya que esto genera una frontera toroidal a la tres variedad y ademas la volverá compacta,, y es importante considerar estas cúspides para entender si hay una estructura hiperbólica completa.
\begin{definition}
    Sea $M$ una 3-variedad con frontera toroidal, definimos una cúspide de $M$, como una vecindad de la frontera que es homeomorfa a $T^2\times [0,\infty]$
\end{definition}

En un principio no resulta tan claro el como encontrar estas cúspides, mas en el caso de nuestra descompocision se pueden llegar a entender de una mejor manera

\begin{definition}
    Sea $M$ una 3-variedad que admite una triangulación topológica ideal, si truncamos los vértices de cada tetraedro ideal se obtiene una colección de triángulos, donde cada uno  se encuentra en la frontera de una cúspide, las aristas de cada triangulo adquieren un pegado natural dado por las caras del tetraedro. Esto nos da una triangulación de cada frontera toroidal. Que llamamos una triangulación de cúspide
\end{definition}

\begin{eg}
    Si truncamos los vértices de nuestra triangulación del complemento del nudo de $8$ obtenemos los siguientes triángulos ya pegados que justamente nos dan la identificación del toro
    \begin{figure}[H]
    \centering
    
%!TEX root = ./main.tex



\tikzset{every picture/.style={line width=0.75pt}} %set default line width to 0.75pt        

\begin{tikzpicture}[x=0.75pt,y=0.75pt,yscale=-1,xscale=1]
%uncomment if require: \path (0,877); %set diagram left start at 0, and has height of 877

%Shape: Circle [id:dp84958386512667] 
\draw  [line width=1.5]  (189.33,55.33) .. controls (189.33,52.39) and (191.72,50) .. (194.67,50) .. controls (197.61,50) and (200,52.39) .. (200,55.33) .. controls (200,58.28) and (197.61,60.67) .. (194.67,60.67) .. controls (191.72,60.67) and (189.33,58.28) .. (189.33,55.33) -- cycle ;
%Shape: Circle [id:dp7070798701116545] 
\draw  [line width=1.5]  (190,145.33) .. controls (190,142.39) and (192.39,140) .. (195.33,140) .. controls (198.28,140) and (200.67,142.39) .. (200.67,145.33) .. controls (200.67,148.28) and (198.28,150.67) .. (195.33,150.67) .. controls (192.39,150.67) and (190,148.28) .. (190,145.33) -- cycle ;
%Shape: Circle [id:dp5401287447911013] 
\draw  [line width=1.5]  (114.67,210) .. controls (114.67,207.05) and (117.05,204.67) .. (120,204.67) .. controls (122.95,204.67) and (125.33,207.05) .. (125.33,210) .. controls (125.33,212.95) and (122.95,215.33) .. (120,215.33) .. controls (117.05,215.33) and (114.67,212.95) .. (114.67,210) -- cycle ;
%Shape: Circle [id:dp024366460256337152] 
\draw  [line width=1.5]  (264.67,210) .. controls (264.67,207.05) and (267.05,204.67) .. (270,204.67) .. controls (272.95,204.67) and (275.33,207.05) .. (275.33,210) .. controls (275.33,212.95) and (272.95,215.33) .. (270,215.33) .. controls (267.05,215.33) and (264.67,212.95) .. (264.67,210) -- cycle ;
%Straight Lines [id:da5819091680445531] 
\draw [line width=1.5]    (194.67,60.67) -- (195.3,136) ;
\draw [shift={(195.33,140)}, rotate = 269.52] [fill={rgb, 255:red, 0; green, 0; blue, 0 }  ][line width=0.08]  [draw opacity=0] (11.61,-5.58) -- (0,0) -- (11.61,5.58) -- cycle    ;
\draw [shift={(195,100.33)}, rotate = 269.52] [color={rgb, 255:red, 0; green, 0; blue, 0 }  ][line width=1.5]    (0,6.71) -- (0,-6.71)   ;
%Straight Lines [id:da6153154443999455] 
\draw [line width=1.5]    (125.33,210) -- (260.67,210) ;
\draw [shift={(264.67,210)}, rotate = 180] [fill={rgb, 255:red, 0; green, 0; blue, 0 }  ][line width=0.08]  [draw opacity=0] (11.61,-5.58) -- (0,0) -- (11.61,5.58) -- cycle    ;
\draw [shift={(191.5,210)}, rotate = 180] [color={rgb, 255:red, 0; green, 0; blue, 0 }  ][line width=1.5]    (0,6.71) -- (0,-6.71)(-6.03,6.71) -- (-6.03,-6.71)   ;
%Straight Lines [id:da5399323102122296] 
\draw [line width=1.5]    (120,204.67) -- (186.95,147.92) ;
\draw [shift={(190,145.33)}, rotate = 139.71] [fill={rgb, 255:red, 0; green, 0; blue, 0 }  ][line width=0.08]  [draw opacity=0] (11.61,-5.58) -- (0,0) -- (11.61,5.58) -- cycle    ;
\draw [shift={(155,175)}, rotate = 139.71] [color={rgb, 255:red, 0; green, 0; blue, 0 }  ][line width=1.5]    (0,6.71) -- (0,-6.71)   ;
%Straight Lines [id:da7687719054145034] 
\draw [line width=1.5]    (200.67,145.33) -- (266.96,202.07) ;
\draw [shift={(270,204.67)}, rotate = 220.56] [fill={rgb, 255:red, 0; green, 0; blue, 0 }  ][line width=0.08]  [draw opacity=0] (11.61,-5.58) -- (0,0) -- (11.61,5.58) -- cycle    ;
\draw [shift={(232.67,172.72)}, rotate = 220.56] [color={rgb, 255:red, 0; green, 0; blue, 0 }  ][line width=1.5]    (0,6.71) -- (0,-6.71)(-6.03,6.71) -- (-6.03,-6.71)   ;
%Curve Lines [id:da44100647811810356] 
\draw [line width=1.5]    (114.67,210) .. controls (115.64,114.62) and (150.17,70.98) .. (186.01,56.59) ;
\draw [shift={(189.33,55.33)}, rotate = 160.64] [fill={rgb, 255:red, 0; green, 0; blue, 0 }  ][line width=0.08]  [draw opacity=0] (11.61,-5.58) -- (0,0) -- (11.61,5.58) -- cycle    ;
\draw [shift={(127.57,125.14)}, rotate = 110.35] [color={rgb, 255:red, 0; green, 0; blue, 0 }  ][line width=1.5]    (0,6.71) -- (0,-6.71)(-6.03,6.71) -- (-6.03,-6.71)   ;
%Curve Lines [id:da6379686625487593] 
\draw [line width=1.5]    (200,55.33) .. controls (239.98,55.01) and (285.97,157.03) .. (276.2,206.32) ;
\draw [shift={(275.33,210)}, rotate = 285.26] [fill={rgb, 255:red, 0; green, 0; blue, 0 }  ][line width=0.08]  [draw opacity=0] (11.61,-5.58) -- (0,0) -- (11.61,5.58) -- cycle    ;
\draw [shift={(259.87,118.95)}, rotate = 246.87] [color={rgb, 255:red, 0; green, 0; blue, 0 }  ][line width=1.5]    (0,6.71) -- (0,-6.71)   ;
%Shape: Circle [id:dp034440490903057674] 
\draw  [line width=1.5]  (431.81,55.33) .. controls (431.81,52.39) and (434.2,50) .. (437.14,50) .. controls (440.09,50) and (442.47,52.39) .. (442.47,55.33) .. controls (442.47,58.28) and (440.09,60.67) .. (437.14,60.67) .. controls (434.2,60.67) and (431.81,58.28) .. (431.81,55.33) -- cycle ;
%Shape: Circle [id:dp4739299080174333] 
\draw  [line width=1.5]  (432.47,145.33) .. controls (432.47,142.39) and (434.86,140) .. (437.81,140) .. controls (440.75,140) and (443.14,142.39) .. (443.14,145.33) .. controls (443.14,148.28) and (440.75,150.67) .. (437.81,150.67) .. controls (434.86,150.67) and (432.47,148.28) .. (432.47,145.33) -- cycle ;
%Shape: Circle [id:dp2986589831675226] 
\draw  [line width=1.5]  (357.14,210) .. controls (357.14,207.05) and (359.53,204.67) .. (362.47,204.67) .. controls (365.42,204.67) and (367.81,207.05) .. (367.81,210) .. controls (367.81,212.95) and (365.42,215.33) .. (362.47,215.33) .. controls (359.53,215.33) and (357.14,212.95) .. (357.14,210) -- cycle ;
%Shape: Circle [id:dp6438231834616311] 
\draw  [line width=1.5]  (507.14,210) .. controls (507.14,207.05) and (509.53,204.67) .. (512.47,204.67) .. controls (515.42,204.67) and (517.81,207.05) .. (517.81,210) .. controls (517.81,212.95) and (515.42,215.33) .. (512.47,215.33) .. controls (509.53,215.33) and (507.14,212.95) .. (507.14,210) -- cycle ;
%Straight Lines [id:da019687389237552866] 
\draw [line width=1.5]    (437.17,64.67) -- (437.81,140) ;
\draw [shift={(437.47,100.33)}, rotate = 269.52] [color={rgb, 255:red, 0; green, 0; blue, 0 }  ][line width=1.5]    (0,6.71) -- (0,-6.71)   ;
\draw [shift={(437.14,60.67)}, rotate = 89.52] [fill={rgb, 255:red, 0; green, 0; blue, 0 }  ][line width=0.08]  [draw opacity=0] (11.61,-5.58) -- (0,0) -- (11.61,5.58) -- cycle    ;
%Straight Lines [id:da5743118307270256] 
\draw [line width=1.5]    (367.81,210) -- (503.14,210) ;
\draw [shift={(507.14,210)}, rotate = 180] [fill={rgb, 255:red, 0; green, 0; blue, 0 }  ][line width=0.08]  [draw opacity=0] (11.61,-5.58) -- (0,0) -- (11.61,5.58) -- cycle    ;
\draw [shift={(433.97,210)}, rotate = 180] [color={rgb, 255:red, 0; green, 0; blue, 0 }  ][line width=1.5]    (0,6.71) -- (0,-6.71)(-6.03,6.71) -- (-6.03,-6.71)   ;
%Straight Lines [id:da15352339522122915] 
\draw [line width=1.5]    (365.53,202.08) -- (432.47,145.33) ;
\draw [shift={(397.47,175)}, rotate = 139.71] [color={rgb, 255:red, 0; green, 0; blue, 0 }  ][line width=1.5]    (0,6.71) -- (0,-6.71)   ;
\draw [shift={(362.47,204.67)}, rotate = 319.71] [fill={rgb, 255:red, 0; green, 0; blue, 0 }  ][line width=0.08]  [draw opacity=0] (11.61,-5.58) -- (0,0) -- (11.61,5.58) -- cycle    ;
%Straight Lines [id:da15958460866160962] 
\draw [line width=1.5]    (443.14,145.33) -- (509.43,202.07) ;
\draw [shift={(512.47,204.67)}, rotate = 220.56] [fill={rgb, 255:red, 0; green, 0; blue, 0 }  ][line width=0.08]  [draw opacity=0] (11.61,-5.58) -- (0,0) -- (11.61,5.58) -- cycle    ;
\draw [shift={(475.15,172.72)}, rotate = 220.56] [color={rgb, 255:red, 0; green, 0; blue, 0 }  ][line width=1.5]    (0,6.71) -- (0,-6.71)(-6.03,6.71) -- (-6.03,-6.71)   ;
%Curve Lines [id:da2355671835877725] 
\draw [line width=1.5]    (357.14,210) .. controls (358.11,114.62) and (392.64,70.98) .. (428.48,56.59) ;
\draw [shift={(431.81,55.33)}, rotate = 160.64] [fill={rgb, 255:red, 0; green, 0; blue, 0 }  ][line width=0.08]  [draw opacity=0] (11.61,-5.58) -- (0,0) -- (11.61,5.58) -- cycle    ;
\draw [shift={(370.05,125.14)}, rotate = 110.35] [color={rgb, 255:red, 0; green, 0; blue, 0 }  ][line width=1.5]    (0,6.71) -- (0,-6.71)(-6.03,6.71) -- (-6.03,-6.71)   ;
%Curve Lines [id:da10269880002458354] 
\draw [line width=1.5]    (446.8,55.69) .. controls (486.71,62.44) and (530.35,164) .. (517.81,210) ;
\draw [shift={(504.08,122.96)}, rotate = 247.81] [color={rgb, 255:red, 0; green, 0; blue, 0 }  ][line width=1.5]    (0,6.71) -- (0,-6.71)   ;
\draw [shift={(442.47,55.33)}, rotate = 359.53] [fill={rgb, 255:red, 0; green, 0; blue, 0 }  ][line width=0.08]  [draw opacity=0] (11.61,-5.58) -- (0,0) -- (11.61,5.58) -- cycle    ;
%Shape: Triangle [id:dp9575795612267674] 
\draw  [color={rgb, 255:red, 128; green, 128; blue, 128 }  ,draw opacity=0.61 ][dash pattern={on 1.69pt off 2.76pt}][line width=1.5]  (195.5,100) -- (230,170) -- (160,170) -- cycle ;
%Shape: Triangle [id:dp8304666288631158] 
\draw  [color={rgb, 255:red, 128; green, 128; blue, 128 }  ,draw opacity=0.61 ][dash pattern={on 1.69pt off 2.76pt}][line width=1.5]  (193.82,95.66) -- (159.91,35.34) -- (229.91,36.02) -- cycle ;
%Shape: Triangle [id:dp4999507729437477] 
\draw  [color={rgb, 255:red, 128; green, 128; blue, 128 }  ,draw opacity=0.61 ][dash pattern={on 1.69pt off 2.76pt}][line width=1.5]  (436.49,96.99) -- (402.58,36.67) -- (472.58,37.35) -- cycle ;
%Shape: Triangle [id:dp6286531931567367] 
\draw  [color={rgb, 255:red, 128; green, 128; blue, 128 }  ,draw opacity=0.61 ][dash pattern={on 1.69pt off 2.76pt}][line width=1.5]  (437.47,100.33) -- (471.97,170.33) -- (401.97,170.33) -- cycle ;
%Shape: Triangle [id:dp9402152936660696] 
\draw  [color={rgb, 255:red, 128; green, 128; blue, 128 }  ,draw opacity=0.61 ][dash pattern={on 1.69pt off 2.76pt}][line width=1.5]  (397.47,175) -- (373.53,248.23) -- (320.46,193.79) -- cycle ;
%Shape: Triangle [id:dp7181925212339698] 
\draw  [color={rgb, 255:red, 128; green, 128; blue, 128 }  ,draw opacity=0.61 ][dash pattern={on 1.69pt off 2.76pt}][line width=1.5]  (155,175) -- (131.06,248.23) -- (77.99,193.79) -- cycle ;
%Shape: Triangle [id:dp23139624969110628] 
\draw  [color={rgb, 255:red, 128; green, 128; blue, 128 }  ,draw opacity=0.61 ][dash pattern={on 1.69pt off 2.76pt}][line width=1.5]  (235.33,175) -- (306.59,204.29) -- (248.38,253.19) -- cycle ;
%Shape: Triangle [id:dp7463486020980554] 
\draw  [color={rgb, 255:red, 128; green, 128; blue, 128 }  ,draw opacity=0.61 ][dash pattern={on 1.69pt off 2.76pt}][line width=1.5]  (477.81,175) -- (549.07,204.29) -- (490.85,253.19) -- cycle ;
%Shape: Triangle [id:dp670762646815206] 
\draw   (117.67,263.33) -- (175.33,392.67) -- (60,392.67) -- cycle ;
%Shape: Triangle [id:dp6096644924278205] 
\draw   (233,263.33) -- (290.67,392.67) -- (175.33,392.67) -- cycle ;
%Shape: Triangle [id:dp6132969216533539] 
\draw   (348.33,263.33) -- (406,392.67) -- (290.67,392.67) -- cycle ;
%Shape: Triangle [id:dp1957398112131299] 
\draw   (463.67,263.33) -- (521.33,392.67) -- (406,392.67) -- cycle ;
%Shape: Triangle [id:dp6822990842952337] 
\draw   (175.33,392.67) -- (117.67,263.33) -- (233,263.33) -- cycle ;
%Shape: Triangle [id:dp6962930689163636] 
\draw   (290.67,392.67) -- (233,263.33) -- (348.33,263.33) -- cycle ;
%Shape: Triangle [id:dp1424783735549131] 
\draw   (406,392.67) -- (348.33,263.33) -- (463.67,263.33) -- cycle ;
%Shape: Triangle [id:dp3424699931659767] 
\draw   (521.33,392.67) -- (463.67,263.33) -- (579,263.33) -- cycle ;

% Text Node
\draw (394,112.4) node [anchor=north west][inner sep=0.75pt]    {$A$};
% Text Node
\draw (467,112.4) node [anchor=north west][inner sep=0.75pt]    {$B$};
% Text Node
\draw (431,172.4) node [anchor=north west][inner sep=0.75pt]    {$C$};
% Text Node
\draw (384,222.4) node [anchor=north west][inner sep=0.75pt]    {$D$};
% Text Node
\draw (224,122.4) node [anchor=north west][inner sep=0.75pt]    {$A$};
% Text Node
\draw (157,122.4) node [anchor=north west][inner sep=0.75pt]    {$B$};
% Text Node
\draw (191,172.4) node [anchor=north west][inner sep=0.75pt]    {$C$};
% Text Node
\draw (154,222.4) node [anchor=north west][inner sep=0.75pt]    {$D$};
% Text Node
\draw (132,160.4) node [anchor=north west][inner sep=0.75pt]    {$z_{1}$};
% Text Node
\draw (271,112.4) node [anchor=north west][inner sep=0.75pt]    {$z_{1}$};
% Text Node
\draw (244,162.4) node [anchor=north west][inner sep=0.75pt]    {$z_{2}$};
% Text Node
\draw (102,122.4) node [anchor=north west][inner sep=0.75pt]    {$z_{2}$};
% Text Node
\draw (201,212.4) node [anchor=north west][inner sep=0.75pt]    {$z_{3}$};
% Text Node
\draw (202,82.4) node [anchor=north west][inner sep=0.75pt]    {$z_{3}$};
% Text Node
\draw (385,148.4) node [anchor=north west][inner sep=0.75pt]    {$w_{1}$};
% Text Node
\draw (511,122.4) node [anchor=north west][inner sep=0.75pt]    {$w_{1}$};
% Text Node
\draw (461,142.4) node [anchor=north west][inner sep=0.75pt]    {$w_{2}$};
% Text Node
\draw (341,122.4) node [anchor=north west][inner sep=0.75pt]    {$w_{2}$};
% Text Node
\draw (441,82.4) node [anchor=north west][inner sep=0.75pt]    {$w_{3}$};
% Text Node
\draw (447,210.4) node [anchor=north west][inner sep=0.75pt]    {$w_{3}$};
% Text Node
\draw (190,149.4) node [anchor=north west][inner sep=0.75pt]    {$a$};
% Text Node
\draw (119.67,217.4) node [anchor=north west][inner sep=0.75pt]    {$b$};
% Text Node
\draw (256.67,215.4) node [anchor=north west][inner sep=0.75pt]    {$d$};
% Text Node
\draw (188,33.73) node [anchor=north west][inner sep=0.75pt]    {$c$};
% Text Node
\draw (431,33.07) node [anchor=north west][inner sep=0.75pt]    {$e$};
% Text Node
\draw (501.67,219.07) node [anchor=north west][inner sep=0.75pt]    {$f$};
% Text Node
\draw (433,146.73) node [anchor=north west][inner sep=0.75pt]    {$g$};
% Text Node
\draw (360,216.73) node [anchor=north west][inner sep=0.75pt]    {$h$};
% Text Node
\draw (108.67,340.4) node [anchor=north west][inner sep=0.75pt]    {$a$};
% Text Node
\draw (168.67,300.4) node [anchor=north west][inner sep=0.75pt]    {$h$};
% Text Node
\draw (224.67,340.4) node [anchor=north west][inner sep=0.75pt]    {$b$};
% Text Node
\draw (284.67,300.4) node [anchor=north west][inner sep=0.75pt]    {$g$};
% Text Node
\draw (338.67,340.4) node [anchor=north west][inner sep=0.75pt]    {$c$};
% Text Node
\draw (458.67,340.4) node [anchor=north west][inner sep=0.75pt]    {$d$};
% Text Node
\draw (398.67,300.4) node [anchor=north west][inner sep=0.75pt]    {$f$};
% Text Node
\draw (514.67,300.4) node [anchor=north west][inner sep=0.75pt]    {$e$};
% Text Node
\draw (68.67,372.4) node [anchor=north west][inner sep=0.75pt]    {$z_{1}$};
% Text Node
\draw (150.67,371.4) node [anchor=north west][inner sep=0.75pt]    {$z_{2}$};
% Text Node
\draw (183.67,372.4) node [anchor=north west][inner sep=0.75pt]    {$z_{3}$};
% Text Node
\draw (265.67,371.4) node [anchor=north west][inner sep=0.75pt]    {$z_{1}$};
% Text Node
\draw (301.67,371.4) node [anchor=north west][inner sep=0.75pt]    {$z_{1}$};
% Text Node
\draw (383.67,370.4) node [anchor=north west][inner sep=0.75pt]    {$z_{2}$};
% Text Node
\draw (416.67,371.4) node [anchor=north west][inner sep=0.75pt]    {$z_{3}$};
% Text Node
\draw (498.67,370.4) node [anchor=north west][inner sep=0.75pt]    {$z_{1}$};
% Text Node
\draw (127.67,265.4) node [anchor=north west][inner sep=0.75pt]    {$w_{1}$};
% Text Node
\draw (203.67,265.4) node [anchor=north west][inner sep=0.75pt]    {$w_{3}$};
% Text Node
\draw (241.67,264.4) node [anchor=north west][inner sep=0.75pt]    {$w_{2}$};
% Text Node
\draw (318.67,263.4) node [anchor=north west][inner sep=0.75pt]    {$w_{1}$};
% Text Node
\draw (357.67,262.4) node [anchor=north west][inner sep=0.75pt]    {$w_{1}$};
% Text Node
\draw (436.67,262.4) node [anchor=north west][inner sep=0.75pt]    {$w_{3}$};
% Text Node
\draw (474.67,263.4) node [anchor=north west][inner sep=0.75pt]    {$w_{2}$};
% Text Node
\draw (550.67,262.4) node [anchor=north west][inner sep=0.75pt]    {$w_{1}$};
% Text Node
\draw (165.67,350.4) node [anchor=north west][inner sep=0.75pt]    {$w_{2}$};
% Text Node
\draw (279.67,350.4) node [anchor=north west][inner sep=0.75pt]    {$w_{3}$};
% Text Node
\draw (395.67,352.4) node [anchor=north west][inner sep=0.75pt]    {$w_{2}$};
% Text Node
\draw (511.67,349.4) node [anchor=north west][inner sep=0.75pt]    {$w_{3}$};
% Text Node
\draw (109.67,277.4) node [anchor=north west][inner sep=0.75pt]    {$z_{3}$};
% Text Node
\draw (341.67,280.4) node [anchor=north west][inner sep=0.75pt]    {$z_{3}$};
% Text Node
\draw (225.67,278.4) node [anchor=north west][inner sep=0.75pt]    {$z_{2}$};
% Text Node
\draw (455.67,280.4) node [anchor=north west][inner sep=0.75pt]    {$z_{2}$};


\end{tikzpicture}
        \caption{Triangulación de la cuspide}
    
    \end{figure}
\end{eg}

Necesitamos de estas cúspides para el siguiente teorema que nos habla acerca de la completez de nuestra estructura inducida

\begin{theorem}
    Sea $M$ una 3-variedad con frontera toroidal y una estructura hiperbólica, esta estructura es completa en $M$ si y solo si la estructura inducida en cada frontera de una cúspide es una estructura euclidea sobre el toro.
\end{theorem}

Esto quiere decir que en el momento del pegado solo basta por preocuparnos por que la estructura hiperbólica del complemento de nuestro nudo, induzca una estructura euclidea en la triangulación de la cúspide, la pregunta es como podemos asegurar esto.

\begin{definition}
    Suponga que la 3-variedad $M$ posee una triangulación topológica ideal, y sea $T$ la frontera toroidal de una cuspide de $M$, sea $[\alpha]\in\pi_1(T)$, asociamos el numero $H(\alpha)$ de la siguiente manera. Dotando a $\alpha$ de una orientación, este por homotopia puede ser colocado de tal manera que en cada triangulo solo separe una esquina una vez, donde cada vértice esta dotado de un invariante de arista $z_i$, y como tiene orientación se puede determinar si esta a la izquierda o la derecha, esto lo determinaremos con una potencia $\epsilon_i=\pm 1$, positivo si esta a la izquierda y negativo en otro caso. Así 
    $$H(\alpha)=\prod_{i=1}^n z_i^{\epsilon_i},$$
\end{definition}
Tan extraño y misterioso como puede parecer este numero nos permite determinar cuando nuestra estructura hiperbólica es completa
\begin{theorem}
    Sea $T$ la frontera toroidal de una cúspide de $M$, donde $M$ admite una triangulación topológica ideal y los tetraedros ideales admiten una estructura hiperbólica que satisface las ecuaciones de pegado. Sean $\alpha$ y $\beta$ curvas que generan $\pi_1(T)$, si $H(\alpha)=H(\beta)=1$ entonces la triangulación es una triangulación ideal geométrica.
\end{theorem}

En estos últimos dos conceptos no se ahondo en sus justificaciones debido a que estas requieren herramientas de holonomia y funciones de de desarrollo, las cuales se salen un poco del propósito de este proyecto. Con esta herramienta que llamaremos ecuaciones de completez descubramos para que invariantes de arista el complemento del nudo de 8 posee una estructura hiperbólica completa

\begin{eg}
    Observe primero que dada la triangulación anterior basta con tomar las siguientes dos curvas que claramente generan el grupo fundamental
    \begin{figure}[H]
    \centering
    
%!TEX root = ./main.tex



\tikzset{every picture/.style={line width=0.75pt}} %set default line width to 0.75pt        

\begin{tikzpicture}[x=0.75pt,y=0.75pt,yscale=-1,xscale=1]
%uncomment if require: \path (0,877); %set diagram left start at 0, and has height of 877

%Shape: Triangle [id:dp670762646815206] 
\draw   (117.67,43.33) -- (175.33,172.67) -- (60,172.67) -- cycle ;
%Shape: Triangle [id:dp6096644924278205] 
\draw   (233,43.33) -- (290.67,172.67) -- (175.33,172.67) -- cycle ;
%Shape: Triangle [id:dp6132969216533539] 
\draw   (348.33,43.33) -- (406,172.67) -- (290.67,172.67) -- cycle ;
%Shape: Triangle [id:dp1957398112131299] 
\draw   (463.67,43.33) -- (521.33,172.67) -- (406,172.67) -- cycle ;
%Shape: Triangle [id:dp6822990842952337] 
\draw   (175.33,172.67) -- (117.67,43.33) -- (233,43.33) -- cycle ;
%Shape: Triangle [id:dp6962930689163636] 
\draw   (290.67,172.67) -- (233,43.33) -- (348.33,43.33) -- cycle ;
%Shape: Triangle [id:dp1424783735549131] 
\draw   (406,172.67) -- (348.33,43.33) -- (463.67,43.33) -- cycle ;
%Shape: Triangle [id:dp3424699931659767] 
\draw   (521.33,172.67) -- (463.67,43.33) -- (579,43.33) -- cycle ;
%Curve Lines [id:da21318320003967128] 
\draw [color={rgb, 255:red, 208; green, 2; blue, 27 }  ,draw opacity=1 ][line width=1.5]    (88,110.67) .. controls (128,80.67) and (512.33,136.33) .. (552.33,106.33) ;
\draw [shift={(323.35,108.69)}, rotate = 183.46] [fill={rgb, 255:red, 208; green, 2; blue, 27 }  ,fill opacity=1 ][line width=0.08]  [draw opacity=0] (11.61,-5.58) -- (0,0) -- (11.61,5.58) -- cycle    ;
%Curve Lines [id:da5420489243511813] 
\draw [color={rgb, 255:red, 74; green, 144; blue, 226 }  ,draw opacity=1 ][line width=1.5]    (112.67,173.33) .. controls (152.67,143.33) and (137,73) .. (177,43) ;
\draw [shift={(146.49,101.4)}, rotate = 103.75] [fill={rgb, 255:red, 74; green, 144; blue, 226 }  ,fill opacity=1 ][line width=0.08]  [draw opacity=0] (11.61,-5.58) -- (0,0) -- (11.61,5.58) -- cycle    ;

% Text Node
\draw (68.67,152.4) node [anchor=north west][inner sep=0.75pt]    {$z_{1}$};
% Text Node
\draw (150.67,151.4) node [anchor=north west][inner sep=0.75pt]    {$z_{2}$};
% Text Node
\draw (183.67,152.4) node [anchor=north west][inner sep=0.75pt]    {$z_{3}$};
% Text Node
\draw (265.67,151.4) node [anchor=north west][inner sep=0.75pt]    {$z_{1}$};
% Text Node
\draw (301.67,151.4) node [anchor=north west][inner sep=0.75pt]    {$z_{1}$};
% Text Node
\draw (383.67,150.4) node [anchor=north west][inner sep=0.75pt]    {$z_{2}$};
% Text Node
\draw (416.67,151.4) node [anchor=north west][inner sep=0.75pt]    {$z_{3}$};
% Text Node
\draw (498.67,150.4) node [anchor=north west][inner sep=0.75pt]    {$z_{1}$};
% Text Node
\draw (127.67,45.4) node [anchor=north west][inner sep=0.75pt]    {$w_{1}$};
% Text Node
\draw (203.67,45.4) node [anchor=north west][inner sep=0.75pt]    {$w_{3}$};
% Text Node
\draw (241.67,44.4) node [anchor=north west][inner sep=0.75pt]    {$w_{2}$};
% Text Node
\draw (318.67,43.4) node [anchor=north west][inner sep=0.75pt]    {$w_{1}$};
% Text Node
\draw (357.67,42.4) node [anchor=north west][inner sep=0.75pt]    {$w_{1}$};
% Text Node
\draw (436.67,42.4) node [anchor=north west][inner sep=0.75pt]    {$w_{3}$};
% Text Node
\draw (474.67,43.4) node [anchor=north west][inner sep=0.75pt]    {$w_{2}$};
% Text Node
\draw (550.67,42.4) node [anchor=north west][inner sep=0.75pt]    {$w_{1}$};
% Text Node
\draw (165.67,130.4) node [anchor=north west][inner sep=0.75pt]    {$w_{2}$};
% Text Node
\draw (279.67,130.4) node [anchor=north west][inner sep=0.75pt]    {$w_{3}$};
% Text Node
\draw (395.67,132.4) node [anchor=north west][inner sep=0.75pt]    {$w_{2}$};
% Text Node
\draw (511.67,129.4) node [anchor=north west][inner sep=0.75pt]    {$w_{3}$};
% Text Node
\draw (109.67,57.4) node [anchor=north west][inner sep=0.75pt]    {$z_{3}$};
% Text Node
\draw (341.67,60.4) node [anchor=north west][inner sep=0.75pt]    {$z_{3}$};
% Text Node
\draw (225.67,58.4) node [anchor=north west][inner sep=0.75pt]    {$z_{2}$};
% Text Node
\draw (455.67,60.4) node [anchor=north west][inner sep=0.75pt]    {$z_{2}$};
% Text Node
\draw (286.67,86.07) node [anchor=north west][inner sep=0.75pt]    {$\textcolor[rgb]{0.82,0.01,0.11}{\alpha }$};
% Text Node
\draw (122,131.4) node [anchor=north west][inner sep=0.75pt]    {$\textcolor[rgb]{0.29,0.56,0.89}{\beta }$};


\end{tikzpicture}
        \caption{Curvas $\alpha$ y $\beta$}
    \end{figure}
    Luego por la descripción dada antes $$H(\alpha)=z_3w_2^{-1}z_2w_3^{-1}z_3w_2^{-1}z_2w_3^{-1}=\left(\dfrac{z_2z_3}{w_2w_3}\right)^2$$, y $$H(\beta)=z_2^{-1}w_1=\dfrac{w_1}{z_2}$$
    Siguiendo el lema 4.1 si el invariante de arista $z_1=z$ y $w_1=w$  obtenemos que 
    $$H(\alpha)=\left(\frac{1}{1-z}\cdot\frac{z-1}{z}\cdot(1-w)\cdot\frac{w}{w-1}\right)^2=\frac{w^2}{z^2},$$
    y ademas $H(\beta)=w(1-z)$, por las ecuaciones de completez, como $H(\alpha)=1$, tenemos que $w^2=z^2$, luego $(w-z)(w+z)=0$, pero como $w$ y $z$ tienen parte imaginaria positiva el segundo factor no puede ser 0, luego $w=z$, sustituyendo esto en $H(\beta)=1$, obtenemos que $z^2-z+1=0$, luego la única posibilidad es que
    $$w=z=\frac{1}{2}+i\frac{\sqrt{3}}{2}$$
    el cual es justamente un punto de los puntos que si dotan de estructura hiperbólica, esto nos dice que todos los demás pegados no tienen estructura hiperbólica completa.
\end{eg}

\subsection{Ecuaciones de pegado para el nudo $6_1$}

Recordemos que antes ya habíamos hecho la descompocision del nudo $6_1$ en dos poliedros ideales los cuales si bien no son tetraedros, la forma en que están proyectados se asemeja mucho a unos, por lo que es natural preguntarse de que forma podemos convertir estos en tetraedros ideales, la idea sera tomar un vértices de uno de los poliedros en las caras que no son triángulos y luego conectar este con aquellos vértices con los que no tenga aristas, para subdividir la cara en triángulos, entonces dadas las proyecciones obtenidas previamente tenemos las siguientes subdivisiones
\begin{figure}[H]
\centering
%!TEX root = ./main.tex



\tikzset{every picture/.style={line width=0.75pt}} %set default line width to 0.75pt        

\begin{tikzpicture}[x=0.75pt,y=0.75pt,yscale=-1,xscale=1]
%uncomment if require: \path (0,331); %set diagram left start at 0, and has height of 331
\path[use as bounding box] (20,10) rectangle (660,300);
%Shape: Ellipse [id:dp012401419565516214] 
\draw  [color={rgb, 255:red, 0; green, 0; blue, 0 }  ,draw opacity=1 ][line width=1.5]  (193.64,71.71) .. controls (193.64,69.26) and (195.4,67.27) .. (197.56,67.27) .. controls (199.72,67.27) and (201.47,69.26) .. (201.47,71.71) .. controls (201.47,74.17) and (199.72,76.16) .. (197.56,76.16) .. controls (195.4,76.16) and (193.64,74.17) .. (193.64,71.71) -- cycle ;
%Shape: Ellipse [id:dp44563585877978507] 
\draw  [color={rgb, 255:red, 0; green, 0; blue, 0 }  ,draw opacity=1 ][line width=1.5]  (193.64,133.02) .. controls (193.64,130.56) and (195.4,128.57) .. (197.56,128.57) .. controls (199.72,128.57) and (201.47,130.56) .. (201.47,133.02) .. controls (201.47,135.47) and (199.72,137.47) .. (197.56,137.47) .. controls (195.4,137.47) and (193.64,135.47) .. (193.64,133.02) -- cycle ;
%Shape: Ellipse [id:dp12283994981009283] 
\draw  [color={rgb, 255:red, 0; green, 0; blue, 0 }  ,draw opacity=1 ][line width=1.5]  (127.35,191.05) .. controls (127.35,188.6) and (129.1,186.61) .. (131.26,186.61) .. controls (133.42,186.61) and (135.17,188.6) .. (135.17,191.05) .. controls (135.17,193.51) and (133.42,195.5) .. (131.26,195.5) .. controls (129.1,195.5) and (127.35,193.51) .. (127.35,191.05) -- cycle ;
%Shape: Ellipse [id:dp43640670243388857] 
\draw  [color={rgb, 255:red, 0; green, 0; blue, 0 }  ,draw opacity=1 ][line width=1.5]  (257.06,189.18) .. controls (257.06,186.72) and (258.81,184.73) .. (260.97,184.73) .. controls (263.13,184.73) and (264.88,186.72) .. (264.88,189.18) .. controls (264.88,191.63) and (263.13,193.63) .. (260.97,193.63) .. controls (258.81,193.63) and (257.06,191.63) .. (257.06,189.18) -- cycle ;
%Shape: Ellipse [id:dp8232398640866994] 
\draw  [color={rgb, 255:red, 0; green, 0; blue, 0 }  ,draw opacity=1 ][line width=1.5]  (170.59,191.05) .. controls (170.59,188.6) and (172.34,186.61) .. (174.5,186.61) .. controls (176.66,186.61) and (178.41,188.6) .. (178.41,191.05) .. controls (178.41,193.51) and (176.66,195.5) .. (174.5,195.5) .. controls (172.34,195.5) and (170.59,193.51) .. (170.59,191.05) -- cycle ;
%Shape: Ellipse [id:dp43325876938835506] 
\draw  [color={rgb, 255:red, 0; green, 0; blue, 0 }  ,draw opacity=1 ][line width=1.5]  (218.35,191.05) .. controls (218.35,188.6) and (220.1,186.61) .. (222.26,186.61) .. controls (224.42,186.61) and (226.18,188.6) .. (226.18,191.05) .. controls (226.18,193.51) and (224.42,195.5) .. (222.26,195.5) .. controls (220.1,195.5) and (218.35,193.51) .. (218.35,191.05) -- cycle ;
%Curve Lines [id:da9923025827403492] 
\draw [color={rgb, 255:red, 0; green, 0; blue, 0 }  ,draw opacity=1 ][line width=1.5]    (126.57,186.85) .. controls (123.15,160.86) and (149.05,147.02) .. (197.56,137.47) ;
\draw [shift={(127.35,191.05)}, rotate = 256.57] [fill={rgb, 255:red, 0; green, 0; blue, 0 }  ,fill opacity=1 ][line width=0.08]  [draw opacity=0] (11.61,-5.58) -- (0,0) -- (11.61,5.58) -- cycle    ;
%Curve Lines [id:da25092915579525166] 
\draw [color={rgb, 255:red, 0; green, 0; blue, 0 }  ,draw opacity=1 ][line width=1.5]    (267.85,186.35) .. controls (286.23,166.45) and (250.25,148.86) .. (197.56,137.47) ;
\draw [shift={(264.88,189.18)}, rotate = 319.56] [fill={rgb, 255:red, 0; green, 0; blue, 0 }  ,fill opacity=1 ][line width=0.08]  [draw opacity=0] (11.61,-5.58) -- (0,0) -- (11.61,5.58) -- cycle    ;
%Curve Lines [id:da6555373806548633] 
\draw [color={rgb, 255:red, 0; green, 0; blue, 0 }  ,draw opacity=1 ][line width=1.5]    (264.88,189.18) .. controls (322.57,220.69) and (298.52,11.84) .. (200.55,65.56) ;
\draw [shift={(197.56,67.27)}, rotate = 329.29] [fill={rgb, 255:red, 0; green, 0; blue, 0 }  ,fill opacity=1 ][line width=0.08]  [draw opacity=0] (11.61,-5.58) -- (0,0) -- (11.61,5.58) -- cycle    ;
%Curve Lines [id:da7978745153273987] 
\draw [color={rgb, 255:red, 0; green, 0; blue, 0 }  ,draw opacity=1 ][line width=1.5]    (127.35,191.05) .. controls (78.23,209.12) and (83.79,18.03) .. (194.18,65.75) ;
\draw [shift={(197.56,67.27)}, rotate = 205.05] [fill={rgb, 255:red, 0; green, 0; blue, 0 }  ,fill opacity=1 ][line width=0.08]  [draw opacity=0] (11.61,-5.58) -- (0,0) -- (11.61,5.58) -- cycle    ;
%Straight Lines [id:da8527204625447116] 
\draw [color={rgb, 255:red, 0; green, 0; blue, 0 }  ,draw opacity=1 ][line width=1.5]    (135.17,191.05) -- (166.59,191.05) ;
\draw [shift={(170.59,191.05)}, rotate = 180] [fill={rgb, 255:red, 0; green, 0; blue, 0 }  ,fill opacity=1 ][line width=0.08]  [draw opacity=0] (11.61,-5.58) -- (0,0) -- (11.61,5.58) -- cycle    ;
%Straight Lines [id:da24688497873734994] 
\draw [color={rgb, 255:red, 0; green, 0; blue, 0 }  ,draw opacity=1 ][line width=1.5]    (182.41,191.05) -- (218.35,191.05) ;
\draw [shift={(178.41,191.05)}, rotate = 0] [fill={rgb, 255:red, 0; green, 0; blue, 0 }  ,fill opacity=1 ][line width=0.08]  [draw opacity=0] (11.61,-5.58) -- (0,0) -- (11.61,5.58) -- cycle    ;
%Straight Lines [id:da7578777207736729] 
\draw [color={rgb, 255:red, 0; green, 0; blue, 0 }  ,draw opacity=1 ][line width=1.5]    (226.18,191.05) -- (253.07,189.42) ;
\draw [shift={(257.06,189.18)}, rotate = 176.53] [fill={rgb, 255:red, 0; green, 0; blue, 0 }  ,fill opacity=1 ][line width=0.08]  [draw opacity=0] (11.61,-5.58) -- (0,0) -- (11.61,5.58) -- cycle    ;
%Straight Lines [id:da9029001126541812] 
\draw [color={rgb, 255:red, 0; green, 0; blue, 0 }  ,draw opacity=1 ][line width=1.5]    (197.56,128.57) -- (197.56,80.16) ;
\draw [shift={(197.56,76.16)}, rotate = 90] [fill={rgb, 255:red, 0; green, 0; blue, 0 }  ,fill opacity=1 ][line width=0.08]  [draw opacity=0] (11.61,-5.58) -- (0,0) -- (11.61,5.58) -- cycle    ;
%Shape: Ellipse [id:dp0659583824499258] 
\draw  [color={rgb, 255:red, 0; green, 0; blue, 0 }  ,draw opacity=1 ][line width=1.5]  (469.72,42.1) .. controls (469.72,38.74) and (471.79,36.01) .. (474.33,36.01) .. controls (476.88,36.01) and (478.94,38.74) .. (478.94,42.1) .. controls (478.94,45.47) and (476.88,48.2) .. (474.33,48.2) .. controls (471.79,48.2) and (469.72,45.47) .. (469.72,42.1) -- cycle ;
%Shape: Ellipse [id:dp3071305848970848] 
\draw  [color={rgb, 255:red, 0; green, 0; blue, 0 }  ,draw opacity=1 ][line width=1.5]  (469.72,126.18) .. controls (469.72,122.81) and (471.79,120.08) .. (474.33,120.08) .. controls (476.88,120.08) and (478.94,122.81) .. (478.94,126.18) .. controls (478.94,129.54) and (476.88,132.27) .. (474.33,132.27) .. controls (471.79,132.27) and (469.72,129.54) .. (469.72,126.18) -- cycle ;
%Shape: Ellipse [id:dp762429770691099] 
\draw  [color={rgb, 255:red, 0; green, 0; blue, 0 }  ,draw opacity=1 ][line width=1.5]  (391.65,205.76) .. controls (391.65,202.39) and (393.71,199.66) .. (396.25,199.66) .. controls (398.8,199.66) and (400.86,202.39) .. (400.86,205.76) .. controls (400.86,209.12) and (398.8,211.85) .. (396.25,211.85) .. controls (393.71,211.85) and (391.65,209.12) .. (391.65,205.76) -- cycle ;
%Shape: Ellipse [id:dp1838380971479946] 
\draw  [color={rgb, 255:red, 0; green, 0; blue, 0 }  ,draw opacity=1 ][line width=1.5]  (544.41,203.19) .. controls (544.41,199.82) and (546.47,197.09) .. (549.02,197.09) .. controls (551.56,197.09) and (553.62,199.82) .. (553.62,203.19) .. controls (553.62,206.56) and (551.56,209.29) .. (549.02,209.29) .. controls (546.47,209.29) and (544.41,206.56) .. (544.41,203.19) -- cycle ;
%Shape: Ellipse [id:dp6027195675802651] 
\draw  [color={rgb, 255:red, 0; green, 0; blue, 0 }  ,draw opacity=1 ][line width=1.5]  (442.57,205.76) .. controls (442.57,202.39) and (444.63,199.66) .. (447.17,199.66) .. controls (449.72,199.66) and (451.78,202.39) .. (451.78,205.76) .. controls (451.78,209.12) and (449.72,211.85) .. (447.17,211.85) .. controls (444.63,211.85) and (442.57,209.12) .. (442.57,205.76) -- cycle ;
%Shape: Ellipse [id:dp05709631909488433] 
\draw  [color={rgb, 255:red, 0; green, 0; blue, 0 }  ,draw opacity=1 ][line width=1.5]  (498.82,205.76) .. controls (498.82,202.39) and (500.89,199.66) .. (503.43,199.66) .. controls (505.97,199.66) and (508.04,202.39) .. (508.04,205.76) .. controls (508.04,209.12) and (505.97,211.85) .. (503.43,211.85) .. controls (500.89,211.85) and (498.82,209.12) .. (498.82,205.76) -- cycle ;
%Curve Lines [id:da41298006953224786] 
\draw [color={rgb, 255:red, 0; green, 0; blue, 0 }  ,draw opacity=1 ][line width=1.5]    (390.94,201.69) .. controls (385.75,164.95) and (416.3,145.59) .. (474.33,132.27) ;
\draw [shift={(391.65,205.76)}, rotate = 258.41] [fill={rgb, 255:red, 0; green, 0; blue, 0 }  ,fill opacity=1 ][line width=0.08]  [draw opacity=0] (11.61,-5.58) -- (0,0) -- (11.61,5.58) -- cycle    ;
%Curve Lines [id:da49093825994949347] 
\draw [color={rgb, 255:red, 0; green, 0; blue, 0 }  ,draw opacity=1 ][line width=1.5]    (553.62,203.19) .. controls (582.14,174.89) and (541.02,149.67) .. (478.2,133.27) ;
\draw [shift={(474.33,132.27)}, rotate = 14.13] [fill={rgb, 255:red, 0; green, 0; blue, 0 }  ,fill opacity=1 ][line width=0.08]  [draw opacity=0] (11.61,-5.58) -- (0,0) -- (11.61,5.58) -- cycle    ;
%Curve Lines [id:da8303588941464516] 
\draw [color={rgb, 255:red, 0; green, 0; blue, 0 }  ,draw opacity=1 ][line width=1.5]    (553.62,203.19) .. controls (621.9,246.61) and (592.96,-42.9) .. (476.1,34.81) ;
\draw [shift={(474.33,36.01)}, rotate = 325.34] [fill={rgb, 255:red, 0; green, 0; blue, 0 }  ,fill opacity=1 ][line width=0.08]  [draw opacity=0] (11.61,-5.58) -- (0,0) -- (11.61,5.58) -- cycle    ;
%Curve Lines [id:da39372336781948036] 
\draw [color={rgb, 255:red, 0; green, 0; blue, 0 }  ,draw opacity=1 ][line width=1.5]    (387.39,207.09) .. controls (333.65,217.62) and (343.83,-35.01) .. (474.33,36.01) ;
\draw [shift={(391.65,205.76)}, rotate = 156.81] [fill={rgb, 255:red, 0; green, 0; blue, 0 }  ,fill opacity=1 ][line width=0.08]  [draw opacity=0] (11.61,-5.58) -- (0,0) -- (11.61,5.58) -- cycle    ;
%Straight Lines [id:da5805044193484686] 
\draw [color={rgb, 255:red, 0; green, 0; blue, 0 }  ,draw opacity=1 ][line width=1.5]    (404.86,205.76) -- (442.57,205.76) ;
\draw [shift={(400.86,205.76)}, rotate = 0] [fill={rgb, 255:red, 0; green, 0; blue, 0 }  ,fill opacity=1 ][line width=0.08]  [draw opacity=0] (11.61,-5.58) -- (0,0) -- (11.61,5.58) -- cycle    ;
%Straight Lines [id:da8661077887659552] 
\draw [color={rgb, 255:red, 0; green, 0; blue, 0 }  ,draw opacity=1 ][line width=1.5]    (451.78,205.76) -- (494.82,205.76) ;
\draw [shift={(498.82,205.76)}, rotate = 180] [fill={rgb, 255:red, 0; green, 0; blue, 0 }  ,fill opacity=1 ][line width=0.08]  [draw opacity=0] (11.61,-5.58) -- (0,0) -- (11.61,5.58) -- cycle    ;
%Straight Lines [id:da050572247984595586] 
\draw [color={rgb, 255:red, 0; green, 0; blue, 0 }  ,draw opacity=1 ][line width=1.5]    (512.03,205.47) -- (544.41,203.19) ;
\draw [shift={(508.04,205.76)}, rotate = 355.96] [fill={rgb, 255:red, 0; green, 0; blue, 0 }  ,fill opacity=1 ][line width=0.08]  [draw opacity=0] (11.61,-5.58) -- (0,0) -- (11.61,5.58) -- cycle    ;
%Straight Lines [id:da2961895295115585] 
\draw [color={rgb, 255:red, 0; green, 0; blue, 0 }  ,draw opacity=1 ][line width=1.5]    (474.33,116.08) -- (474.33,48.2) ;
\draw [shift={(474.33,120.08)}, rotate = 270] [fill={rgb, 255:red, 0; green, 0; blue, 0 }  ,fill opacity=1 ][line width=0.08]  [draw opacity=0] (11.61,-5.58) -- (0,0) -- (11.61,5.58) -- cycle    ;
%Straight Lines [id:da22131359142153095] 
\draw [color={rgb, 255:red, 0; green, 0; blue, 0 }  ,draw opacity=1 ][line width=1.5]    (197.56,137.47) -- (176.2,182.98) ;
\draw [shift={(174.5,186.61)}, rotate = 295.14] [fill={rgb, 255:red, 0; green, 0; blue, 0 }  ,fill opacity=1 ][line width=0.08]  [draw opacity=0] (11.61,-5.58) -- (0,0) -- (11.61,5.58) -- cycle    ;
%Straight Lines [id:da6387841366636108] 
\draw [color={rgb, 255:red, 0; green, 0; blue, 0 }  ,draw opacity=1 ][line width=1.5]    (197.56,137.47) -- (220.47,183.03) ;
\draw [shift={(222.26,186.61)}, rotate = 243.31] [fill={rgb, 255:red, 0; green, 0; blue, 0 }  ,fill opacity=1 ][line width=0.08]  [draw opacity=0] (11.61,-5.58) -- (0,0) -- (11.61,5.58) -- cycle    ;
%Curve Lines [id:da15562696066606252] 
\draw [line width=1.5]    (197.56,67.27) .. controls (-20.14,-91.67) and (57.44,363.9) .. (172.76,198.05) ;
\draw [shift={(174.5,195.5)}, rotate = 123.91] [fill={rgb, 255:red, 0; green, 0; blue, 0 }  ][line width=0.08]  [draw opacity=0] (11.61,-5.58) -- (0,0) -- (11.61,5.58) -- cycle    ;
%Curve Lines [id:da37566750523718095] 
\draw [line width=1.5]    (197.56,67.27) .. controls (349.17,-122.29) and (403.31,351.18) .. (224.97,197.86) ;
\draw [shift={(222.26,195.5)}, rotate = 41.47] [fill={rgb, 255:red, 0; green, 0; blue, 0 }  ][line width=0.08]  [draw opacity=0] (11.61,-5.58) -- (0,0) -- (11.61,5.58) -- cycle    ;
%Curve Lines [id:da3329202190242504] 
\draw [line width=1.5]    (549.02,197.09) .. controls (564.29,157.47) and (415.67,137.92) .. (397.22,196) ;
\draw [shift={(396.25,199.66)}, rotate = 281.99] [fill={rgb, 255:red, 0; green, 0; blue, 0 }  ][line width=0.08]  [draw opacity=0] (11.61,-5.58) -- (0,0) -- (11.61,5.58) -- cycle    ;
%Curve Lines [id:da2054994932773433] 
\draw [line width=1.5]    (450.68,196.73) .. controls (483.13,170.76) and (520.84,171.65) .. (549.02,197.09) ;
\draw [shift={(447.17,199.66)}, rotate = 318.94] [fill={rgb, 255:red, 0; green, 0; blue, 0 }  ][line width=0.08]  [draw opacity=0] (11.61,-5.58) -- (0,0) -- (11.61,5.58) -- cycle    ;
%Curve Lines [id:da5787501033813364] 
\draw [line width=1.5]    (396.25,211.85) .. controls (415.7,279.44) and (487.8,264.88) .. (502.42,215.7) ;
\draw [shift={(503.43,211.85)}, rotate = 102.83] [fill={rgb, 255:red, 0; green, 0; blue, 0 }  ][line width=0.08]  [draw opacity=0] (11.61,-5.58) -- (0,0) -- (11.61,5.58) -- cycle    ;
%Curve Lines [id:da7560422734737169] 
\draw [line width=1.5]    (396.25,211.85) .. controls (405.83,335.69) and (555.06,283.65) .. (549.38,212.55) ;
\draw [shift={(549.02,209.29)}, rotate = 81.69] [fill={rgb, 255:red, 0; green, 0; blue, 0 }  ][line width=0.08]  [draw opacity=0] (11.61,-5.58) -- (0,0) -- (11.61,5.58) -- cycle    ;

% Text Node
\draw (139.44,109.41) node [anchor=north west][inner sep=0.75pt]    {$A$};
% Text Node
\draw (240.86,106.41) node [anchor=north west][inner sep=0.75pt]    {$B$};
% Text Node
\draw (65.28,104.24) node [anchor=north west][inner sep=0.75pt]    {$D_{1}$};
% Text Node
\draw (272.73,65.1) node [anchor=north west][inner sep=0.75pt]    {$1$};
% Text Node
\draw (146.58,132.94) node [anchor=north west][inner sep=0.75pt]    {$1$};
% Text Node
\draw (242.19,133.19) node [anchor=north west][inner sep=0.75pt]    {$4$};
% Text Node
\draw (99.66,71.61) node [anchor=north west][inner sep=0.75pt]    {$4$};
% Text Node
\draw (147.02,191.78) node [anchor=north west][inner sep=0.75pt]    {$4$};
% Text Node
\draw (193.94,190.27) node [anchor=north west][inner sep=0.75pt]    {$4$};
% Text Node
\draw (232.01,191.28) node [anchor=north west][inner sep=0.75pt]    {$4$};
% Text Node
\draw (184.2,96.89) node [anchor=north west][inner sep=0.75pt]    {$1$};
% Text Node
\draw (407.22,96.62) node [anchor=north west][inner sep=0.75pt]    {$A$};
% Text Node
\draw (526.4,92.5) node [anchor=north west][inner sep=0.75pt]    {$B$};
% Text Node
\draw (563.93,35.86) node [anchor=north west][inner sep=0.75pt]    {$4$};
% Text Node
\draw (415.36,128.89) node [anchor=north west][inner sep=0.75pt]    {$4$};
% Text Node
\draw (527.96,129.24) node [anchor=north west][inner sep=0.75pt]    {$1$};
% Text Node
\draw (360.1,44.78) node [anchor=north west][inner sep=0.75pt]    {$1$};
% Text Node
\draw (415.88,209.57) node [anchor=north west][inner sep=0.75pt]    {$4$};
% Text Node
\draw (471.14,207.51) node [anchor=north west][inner sep=0.75pt]    {$4$};
% Text Node
\draw (521.18,208.89) node [anchor=north west][inner sep=0.75pt]    {$4$};
% Text Node
\draw (459.67,79.46) node [anchor=north west][inner sep=0.75pt]    {$1$};
% Text Node
\draw (149.45,158.75) node [anchor=north west][inner sep=0.75pt]    {$C_{1}$};
% Text Node
\draw (186.09,163.84) node [anchor=north west][inner sep=0.75pt]    {$C2$};
% Text Node
\draw (228.91,158.47) node [anchor=north west][inner sep=0.75pt]    {$C_{3}$};
% Text Node
\draw (301.4,104.62) node [anchor=north west][inner sep=0.75pt]    {$D_{3}$};
% Text Node
\draw (186.76,223.14) node [anchor=north west][inner sep=0.75pt]    {$D_{2}$};
% Text Node
\draw (66.7,34.97) node [anchor=north west][inner sep=0.75pt]    {$9$};
% Text Node
\draw (288.03,9.61) node [anchor=north west][inner sep=0.75pt]    {$0$};
% Text Node
\draw (175.09,152.92) node [anchor=north west][inner sep=0.75pt]    {$7$};
% Text Node
\draw (211.7,152.35) node [anchor=north west][inner sep=0.75pt]    {$8$};
% Text Node
\draw (488.84,181.83) node [anchor=north west][inner sep=0.75pt]    {$C_{3}$};
% Text Node
\draw (422.44,173.57) node [anchor=north west][inner sep=0.75pt]    {$C_{2}$};
% Text Node
\draw (459.89,137.67) node [anchor=north west][inner sep=0.75pt]    {$C_{1}$};
% Text Node
\draw (491.35,145.22) node [anchor=north west][inner sep=0.75pt]    {$7$};
% Text Node
\draw (460.7,168.59) node [anchor=north west][inner sep=0.75pt]    {$8$};
% Text Node
\draw (436.02,223.04) node [anchor=north west][inner sep=0.75pt]    {$D_{1}$};
% Text Node
\draw (494.76,244.83) node [anchor=north west][inner sep=0.75pt]    {$D_{2}$};
% Text Node
\draw (556.91,242.27) node [anchor=north west][inner sep=0.75pt]    {$D_{3}$};
% Text Node
\draw (419.13,232.83) node [anchor=north west][inner sep=0.75pt]    {$9$};
% Text Node
\draw (413.6,274.85) node [anchor=north west][inner sep=0.75pt]    {$0$};


\end{tikzpicture}
    \caption{Subdivisión de caras $C$ y $D$}
\end{figure}
Recuerde que el pegado de las caras se hace de tal forma que las aristas coincidan por lo que una subdivisión de las caras $C$ y $D$ en un poliedro ideal ya dotan de una división al otro poliedro ideal. Ahora observe que hay secciones que poseen tres triangulo en cada poliedro, si cortamos estas regiones y asignamos una cara a la parte recortada por identificación, obtenemos precisamente tetraedros, esta subdivisión se hace tomando las regiones $D_3,B,C_3$ para uno y $C_1,A,D_1$, en el caso del poliedro de arriba y nombrando aquellas nuevas caras como $E_1$ y $E_2$ respectivamente
\begin{figure}[H]
\centering
%!TEX root = ./main.tex


\tikzset{every picture/.style={line width=0.75pt}} %set default line width to 0.75pt        

\begin{tikzpicture}[x=0.75pt,y=0.75pt,yscale=-0.9,xscale=0.9]
%uncomment if require: \path (0,498); %set diagram left start at 0, and has height of 498
\path[use as bounding box] (20,10) rectangle (660,250);
%Shape: Ellipse [id:dp012401419565516214] 
\draw  [color={rgb, 255:red, 0; green, 0; blue, 0 }  ,draw opacity=1 ][line width=1.5]  (310.98,80.38) .. controls (310.98,77.92) and (312.73,75.93) .. (314.89,75.93) .. controls (317.05,75.93) and (318.8,77.92) .. (318.8,80.38) .. controls (318.8,82.83) and (317.05,84.82) .. (314.89,84.82) .. controls (312.73,84.82) and (310.98,82.83) .. (310.98,80.38) -- cycle ;
%Shape: Ellipse [id:dp44563585877978507] 
\draw  [color={rgb, 255:red, 0; green, 0; blue, 0 }  ,draw opacity=1 ][line width=1.5]  (310.98,141.69) .. controls (310.98,139.23) and (312.73,137.24) .. (314.89,137.24) .. controls (317.05,137.24) and (318.8,139.23) .. (318.8,141.69) .. controls (318.8,144.14) and (317.05,146.13) .. (314.89,146.13) .. controls (312.73,146.13) and (310.98,144.14) .. (310.98,141.69) -- cycle ;
%Shape: Ellipse [id:dp8232398640866994] 
\draw  [color={rgb, 255:red, 0; green, 0; blue, 0 }  ,draw opacity=1 ][line width=1.5]  (287.92,199.72) .. controls (287.92,197.26) and (289.67,195.27) .. (291.83,195.27) .. controls (293.99,195.27) and (295.74,197.26) .. (295.74,199.72) .. controls (295.74,202.17) and (293.99,204.16) .. (291.83,204.16) .. controls (289.67,204.16) and (287.92,202.17) .. (287.92,199.72) -- cycle ;
%Shape: Ellipse [id:dp43325876938835506] 
\draw  [color={rgb, 255:red, 0; green, 0; blue, 0 }  ,draw opacity=1 ][line width=1.5]  (335.68,199.72) .. controls (335.68,197.26) and (337.44,195.27) .. (339.6,195.27) .. controls (341.76,195.27) and (343.51,197.26) .. (343.51,199.72) .. controls (343.51,202.17) and (341.76,204.16) .. (339.6,204.16) .. controls (337.44,204.16) and (335.68,202.17) .. (335.68,199.72) -- cycle ;
%Straight Lines [id:da24688497873734994] 
\draw [color={rgb, 255:red, 0; green, 0; blue, 0 }  ,draw opacity=1 ][line width=1.5]    (299.74,199.72) -- (335.68,199.72) ;
\draw [shift={(295.74,199.72)}, rotate = 0] [fill={rgb, 255:red, 0; green, 0; blue, 0 }  ,fill opacity=1 ][line width=0.08]  [draw opacity=0] (11.61,-5.58) -- (0,0) -- (11.61,5.58) -- cycle    ;
%Straight Lines [id:da9029001126541812] 
\draw [color={rgb, 255:red, 0; green, 0; blue, 0 }  ,draw opacity=1 ][line width=1.5]    (314.89,137.24) -- (314.89,88.82) ;
\draw [shift={(314.89,84.82)}, rotate = 90] [fill={rgb, 255:red, 0; green, 0; blue, 0 }  ,fill opacity=1 ][line width=0.08]  [draw opacity=0] (11.61,-5.58) -- (0,0) -- (11.61,5.58) -- cycle    ;
%Straight Lines [id:da22131359142153095] 
\draw [color={rgb, 255:red, 0; green, 0; blue, 0 }  ,draw opacity=1 ][line width=1.5]    (314.89,146.13) -- (293.53,191.65) ;
\draw [shift={(291.83,195.27)}, rotate = 295.14] [fill={rgb, 255:red, 0; green, 0; blue, 0 }  ,fill opacity=1 ][line width=0.08]  [draw opacity=0] (11.61,-5.58) -- (0,0) -- (11.61,5.58) -- cycle    ;
%Straight Lines [id:da6387841366636108] 
\draw [color={rgb, 255:red, 0; green, 0; blue, 0 }  ,draw opacity=1 ][line width=1.5]    (314.89,146.13) -- (337.8,191.7) ;
\draw [shift={(339.6,195.27)}, rotate = 243.31] [fill={rgb, 255:red, 0; green, 0; blue, 0 }  ,fill opacity=1 ][line width=0.08]  [draw opacity=0] (11.61,-5.58) -- (0,0) -- (11.61,5.58) -- cycle    ;
%Curve Lines [id:da15562696066606252] 
\draw [line width=1.5]    (314.89,75.93) .. controls (97.2,-83.01) and (174.77,372.56) .. (290.09,206.71) ;
\draw [shift={(291.83,204.16)}, rotate = 123.91] [fill={rgb, 255:red, 0; green, 0; blue, 0 }  ][line width=0.08]  [draw opacity=0] (11.61,-5.58) -- (0,0) -- (11.61,5.58) -- cycle    ;
%Curve Lines [id:da37566750523718095] 
\draw [line width=1.5]    (314.89,75.93) .. controls (466.51,-113.62) and (520.64,359.85) .. (342.3,206.52) ;
\draw [shift={(339.6,204.16)}, rotate = 41.47] [fill={rgb, 255:red, 0; green, 0; blue, 0 }  ][line width=0.08]  [draw opacity=0] (11.61,-5.58) -- (0,0) -- (11.61,5.58) -- cycle    ;
%Shape: Ellipse [id:dp2127358898052878] 
\draw  [color={rgb, 255:red, 0; green, 0; blue, 0 }  ,draw opacity=1 ][line width=1.5]  (85.35,201.05) .. controls (85.35,198.6) and (87.1,196.61) .. (89.26,196.61) .. controls (91.42,196.61) and (93.17,198.6) .. (93.17,201.05) .. controls (93.17,203.51) and (91.42,205.5) .. (89.26,205.5) .. controls (87.1,205.5) and (85.35,203.51) .. (85.35,201.05) -- cycle ;
%Curve Lines [id:da6892910775973166] 
\draw [color={rgb, 255:red, 0; green, 0; blue, 0 }  ,draw opacity=1 ][line width=1.5]    (84.57,196.85) .. controls (81.15,170.86) and (107.05,157.02) .. (155.56,147.47) ;
\draw [shift={(85.35,201.05)}, rotate = 256.57] [fill={rgb, 255:red, 0; green, 0; blue, 0 }  ,fill opacity=1 ][line width=0.08]  [draw opacity=0] (11.61,-5.58) -- (0,0) -- (11.61,5.58) -- cycle    ;
%Curve Lines [id:da5288751216338294] 
\draw [color={rgb, 255:red, 0; green, 0; blue, 0 }  ,draw opacity=1 ][line width=1.5]    (85.35,201.05) .. controls (36.23,219.12) and (41.79,28.03) .. (152.18,75.75) ;
\draw [shift={(155.56,77.27)}, rotate = 205.05] [fill={rgb, 255:red, 0; green, 0; blue, 0 }  ,fill opacity=1 ][line width=0.08]  [draw opacity=0] (11.61,-5.58) -- (0,0) -- (11.61,5.58) -- cycle    ;
%Straight Lines [id:da24276479087471536] 
\draw [color={rgb, 255:red, 0; green, 0; blue, 0 }  ,draw opacity=1 ][line width=1.5]    (93.17,201.05) -- (124.59,201.05) ;
\draw [shift={(128.59,201.05)}, rotate = 180] [fill={rgb, 255:red, 0; green, 0; blue, 0 }  ,fill opacity=1 ][line width=0.08]  [draw opacity=0] (11.61,-5.58) -- (0,0) -- (11.61,5.58) -- cycle    ;
%Curve Lines [id:da32788257215994565] 
\draw [line width=1.5]    (155.56,77.27) .. controls (-62.14,-81.67) and (15.44,373.9) .. (130.76,208.05) ;
\draw [shift={(132.5,205.5)}, rotate = 123.91] [fill={rgb, 255:red, 0; green, 0; blue, 0 }  ][line width=0.08]  [draw opacity=0] (11.61,-5.58) -- (0,0) -- (11.61,5.58) -- cycle    ;
%Shape: Ellipse [id:dp4018621662081351] 
\draw  [color={rgb, 255:red, 0; green, 0; blue, 0 }  ,draw opacity=1 ][line width=1.5]  (127.25,201.72) .. controls (127.25,199.26) and (129,197.27) .. (131.16,197.27) .. controls (133.32,197.27) and (135.08,199.26) .. (135.08,201.72) .. controls (135.08,204.17) and (133.32,206.16) .. (131.16,206.16) .. controls (129,206.16) and (127.25,204.17) .. (127.25,201.72) -- cycle ;
%Straight Lines [id:da06341837878705825] 
\draw [color={rgb, 255:red, 0; green, 0; blue, 0 }  ,draw opacity=1 ][line width=1.5]    (154.22,148.13) -- (132.86,193.65) ;
\draw [shift={(131.16,197.27)}, rotate = 295.14] [fill={rgb, 255:red, 0; green, 0; blue, 0 }  ,fill opacity=1 ][line width=0.08]  [draw opacity=0] (11.61,-5.58) -- (0,0) -- (11.61,5.58) -- cycle    ;
%Shape: Ellipse [id:dp38906306008016145] 
\draw  [color={rgb, 255:red, 0; green, 0; blue, 0 }  ,draw opacity=1 ][line width=1.5]  (154.31,82.38) .. controls (154.31,79.92) and (156.06,77.93) .. (158.22,77.93) .. controls (160.38,77.93) and (162.14,79.92) .. (162.14,82.38) .. controls (162.14,84.83) and (160.38,86.82) .. (158.22,86.82) .. controls (156.06,86.82) and (154.31,84.83) .. (154.31,82.38) -- cycle ;
%Shape: Ellipse [id:dp08743735948745213] 
\draw  [color={rgb, 255:red, 0; green, 0; blue, 0 }  ,draw opacity=1 ][line width=1.5]  (154.31,143.69) .. controls (154.31,141.23) and (156.06,139.24) .. (158.22,139.24) .. controls (160.38,139.24) and (162.14,141.23) .. (162.14,143.69) .. controls (162.14,146.14) and (160.38,148.13) .. (158.22,148.13) .. controls (156.06,148.13) and (154.31,146.14) .. (154.31,143.69) -- cycle ;
%Straight Lines [id:da6699514353142748] 
\draw [color={rgb, 255:red, 0; green, 0; blue, 0 }  ,draw opacity=1 ][line width=1.5]    (158.22,139.24) -- (158.22,90.82) ;
\draw [shift={(158.22,86.82)}, rotate = 90] [fill={rgb, 255:red, 0; green, 0; blue, 0 }  ,fill opacity=1 ][line width=0.08]  [draw opacity=0] (11.61,-5.58) -- (0,0) -- (11.61,5.58) -- cycle    ;
%Shape: Ellipse [id:dp8050348596658142] 
\draw  [color={rgb, 255:red, 0; green, 0; blue, 0 }  ,draw opacity=1 ][line width=1.5]  (547.06,202.3) .. controls (547.06,199.84) and (548.81,197.85) .. (550.97,197.85) .. controls (553.13,197.85) and (554.88,199.84) .. (554.88,202.3) .. controls (554.88,204.75) and (553.13,206.75) .. (550.97,206.75) .. controls (548.81,206.75) and (547.06,204.75) .. (547.06,202.3) -- cycle ;
%Shape: Ellipse [id:dp4814843770420125] 
\draw  [color={rgb, 255:red, 0; green, 0; blue, 0 }  ,draw opacity=1 ][line width=1.5]  (508.35,204.17) .. controls (508.35,201.72) and (510.1,199.73) .. (512.26,199.73) .. controls (514.42,199.73) and (516.18,201.72) .. (516.18,204.17) .. controls (516.18,206.63) and (514.42,208.62) .. (512.26,208.62) .. controls (510.1,208.62) and (508.35,206.63) .. (508.35,204.17) -- cycle ;
%Curve Lines [id:da6756687162833178] 
\draw [color={rgb, 255:red, 0; green, 0; blue, 0 }  ,draw opacity=1 ][line width=1.5]    (557.85,199.47) .. controls (576.23,179.57) and (540.25,161.98) .. (487.56,150.59) ;
\draw [shift={(554.88,202.3)}, rotate = 319.56] [fill={rgb, 255:red, 0; green, 0; blue, 0 }  ,fill opacity=1 ][line width=0.08]  [draw opacity=0] (11.61,-5.58) -- (0,0) -- (11.61,5.58) -- cycle    ;
%Curve Lines [id:da3167596918642579] 
\draw [color={rgb, 255:red, 0; green, 0; blue, 0 }  ,draw opacity=1 ][line width=1.5]    (554.88,202.3) .. controls (612.57,233.81) and (588.52,24.96) .. (490.55,78.68) ;
\draw [shift={(487.56,80.38)}, rotate = 329.29] [fill={rgb, 255:red, 0; green, 0; blue, 0 }  ,fill opacity=1 ][line width=0.08]  [draw opacity=0] (11.61,-5.58) -- (0,0) -- (11.61,5.58) -- cycle    ;
%Straight Lines [id:da5914875377963786] 
\draw [color={rgb, 255:red, 0; green, 0; blue, 0 }  ,draw opacity=1 ][line width=1.5]    (516.18,204.17) -- (543.07,202.54) ;
\draw [shift={(547.06,202.3)}, rotate = 176.53] [fill={rgb, 255:red, 0; green, 0; blue, 0 }  ,fill opacity=1 ][line width=0.08]  [draw opacity=0] (11.61,-5.58) -- (0,0) -- (11.61,5.58) -- cycle    ;
%Straight Lines [id:da07161147378506905] 
\draw [color={rgb, 255:red, 0; green, 0; blue, 0 }  ,draw opacity=1 ][line width=1.5]    (487.56,150.59) -- (510.47,196.15) ;
\draw [shift={(512.26,199.73)}, rotate = 243.31] [fill={rgb, 255:red, 0; green, 0; blue, 0 }  ,fill opacity=1 ][line width=0.08]  [draw opacity=0] (11.61,-5.58) -- (0,0) -- (11.61,5.58) -- cycle    ;
%Curve Lines [id:da5080118287083738] 
\draw [line width=1.5]    (487.56,80.38) .. controls (639.17,-109.17) and (693.31,364.3) .. (514.97,210.98) ;
\draw [shift={(512.26,208.62)}, rotate = 41.47] [fill={rgb, 255:red, 0; green, 0; blue, 0 }  ][line width=0.08]  [draw opacity=0] (11.61,-5.58) -- (0,0) -- (11.61,5.58) -- cycle    ;
%Shape: Ellipse [id:dp8095971328672295] 
\draw  [color={rgb, 255:red, 0; green, 0; blue, 0 }  ,draw opacity=1 ][line width=1.5]  (482.31,85.71) .. controls (482.31,83.26) and (484.06,81.27) .. (486.22,81.27) .. controls (488.38,81.27) and (490.14,83.26) .. (490.14,85.71) .. controls (490.14,88.17) and (488.38,90.16) .. (486.22,90.16) .. controls (484.06,90.16) and (482.31,88.17) .. (482.31,85.71) -- cycle ;
%Shape: Ellipse [id:dp915558520007347] 
\draw  [color={rgb, 255:red, 0; green, 0; blue, 0 }  ,draw opacity=1 ][line width=1.5]  (482.31,147.02) .. controls (482.31,144.56) and (484.06,142.57) .. (486.22,142.57) .. controls (488.38,142.57) and (490.14,144.56) .. (490.14,147.02) .. controls (490.14,149.47) and (488.38,151.47) .. (486.22,151.47) .. controls (484.06,151.47) and (482.31,149.47) .. (482.31,147.02) -- cycle ;
%Straight Lines [id:da8569739762561568] 
\draw [color={rgb, 255:red, 0; green, 0; blue, 0 }  ,draw opacity=1 ][line width=1.5]    (486.22,142.57) -- (486.22,94.16) ;
\draw [shift={(486.22,90.16)}, rotate = 90] [fill={rgb, 255:red, 0; green, 0; blue, 0 }  ,fill opacity=1 ][line width=0.08]  [draw opacity=0] (11.61,-5.58) -- (0,0) -- (11.61,5.58) -- cycle    ;

% Text Node
\draw (311.27,198.94) node [anchor=north west][inner sep=0.75pt]    {$4$};
% Text Node
\draw (301.54,105.56) node [anchor=north west][inner sep=0.75pt]    {$1$};
% Text Node
\draw (303.43,172.5) node [anchor=north west][inner sep=0.75pt]    {$C2$};
% Text Node
\draw (304.1,231.81) node [anchor=north west][inner sep=0.75pt]    {$D_{2}$};
% Text Node
\draw (184.03,43.64) node [anchor=north west][inner sep=0.75pt]    {$9$};
% Text Node
\draw (405.37,18.28) node [anchor=north west][inner sep=0.75pt]    {$0$};
% Text Node
\draw (292.43,161.58) node [anchor=north west][inner sep=0.75pt]    {$7$};
% Text Node
\draw (97.44,119.41) node [anchor=north west][inner sep=0.75pt]    {$A$};
% Text Node
\draw (23.28,114.24) node [anchor=north west][inner sep=0.75pt]    {$D_{1}$};
% Text Node
\draw (104.58,142.94) node [anchor=north west][inner sep=0.75pt]    {$1$};
% Text Node
\draw (57.66,81.61) node [anchor=north west][inner sep=0.75pt]    {$4$};
% Text Node
\draw (105.02,201.78) node [anchor=north west][inner sep=0.75pt]    {$4$};
% Text Node
\draw (107.45,168.75) node [anchor=north west][inner sep=0.75pt]    {$C_{1}$};
% Text Node
\draw (24.7,44.97) node [anchor=north west][inner sep=0.75pt]    {$9$};
% Text Node
\draw (131.76,163.58) node [anchor=north west][inner sep=0.75pt]    {$7$};
% Text Node
\draw (144.87,107.56) node [anchor=north west][inner sep=0.75pt]    {$1$};
% Text Node
\draw (633.33,207.07) node [anchor=north west][inner sep=0.75pt]    {$E_{2}$};
% Text Node
\draw (380.67,126.4) node [anchor=north west][inner sep=0.75pt]    {$E_{2}$};
% Text Node
\draw (530.86,119.53) node [anchor=north west][inner sep=0.75pt]    {$B$};
% Text Node
\draw (562.73,78.22) node [anchor=north west][inner sep=0.75pt]    {$1$};
% Text Node
\draw (532.19,146.31) node [anchor=north west][inner sep=0.75pt]    {$4$};
% Text Node
\draw (522.01,204.4) node [anchor=north west][inner sep=0.75pt]    {$4$};
% Text Node
\draw (518.91,171.59) node [anchor=north west][inner sep=0.75pt]    {$C_{3}$};
% Text Node
\draw (591.4,117.74) node [anchor=north west][inner sep=0.75pt]    {$D_{3}$};
% Text Node
\draw (578.03,22.73) node [anchor=north west][inner sep=0.75pt]    {$0$};
% Text Node
\draw (501.7,165.47) node [anchor=north west][inner sep=0.75pt]    {$8$};
% Text Node
\draw (472.87,110.89) node [anchor=north west][inner sep=0.75pt]    {$1$};
% Text Node
\draw (232.67,125.07) node [anchor=north west][inner sep=0.75pt]    {$E_{1}$};
% Text Node
\draw (20,229.73) node [anchor=north west][inner sep=0.75pt]    {$E_{1}$};


\end{tikzpicture}
    \caption{Subdivisión en tres tetraedros del poliedro de arriba}
\end{figure}
De una manera similar podemos hacer una subdivisión del poliedro de abajo como se muestra a continuación
\begin{figure}[H]
\centering
%!TEX root = ./main.tex




\tikzset{every picture/.style={line width=0.75pt}} %set default line width to 0.75pt        

\begin{tikzpicture}[x=0.75pt,y=0.75pt,yscale=-1,xscale=1]
%uncomment if require: \path (0,498); %set diagram left start at 0, and has height of 498
\path[use as bounding box] (20,10) rectangle (660,380);
%Shape: Ellipse [id:dp35589263376995484] 
\draw  [color={rgb, 255:red, 0; green, 0; blue, 0 }  ,draw opacity=1 ][line width=1.5]  (507.06,44.65) .. controls (507.06,41.28) and (509.12,38.55) .. (511.67,38.55) .. controls (514.21,38.55) and (516.27,41.28) .. (516.27,44.65) .. controls (516.27,48.02) and (514.21,50.75) .. (511.67,50.75) .. controls (509.12,50.75) and (507.06,48.02) .. (507.06,44.65) -- cycle ;
%Shape: Ellipse [id:dp5273685537709042] 
\draw  [color={rgb, 255:red, 0; green, 0; blue, 0 }  ,draw opacity=1 ][line width=1.5]  (428.98,208.3) .. controls (428.98,204.94) and (431.04,202.21) .. (433.59,202.21) .. controls (436.13,202.21) and (438.19,204.94) .. (438.19,208.3) .. controls (438.19,211.67) and (436.13,214.4) .. (433.59,214.4) .. controls (431.04,214.4) and (428.98,211.67) .. (428.98,208.3) -- cycle ;
%Shape: Ellipse [id:dp9933087224039621] 
\draw  [color={rgb, 255:red, 0; green, 0; blue, 0 }  ,draw opacity=1 ][line width=1.5]  (581.74,205.74) .. controls (581.74,202.37) and (583.8,199.64) .. (586.35,199.64) .. controls (588.89,199.64) and (590.96,202.37) .. (590.96,205.74) .. controls (590.96,209.1) and (588.89,211.83) .. (586.35,211.83) .. controls (583.8,211.83) and (581.74,209.1) .. (581.74,205.74) -- cycle ;
%Shape: Ellipse [id:dp4223187867652599] 
\draw  [color={rgb, 255:red, 0; green, 0; blue, 0 }  ,draw opacity=1 ][line width=1.5]  (536.16,208.3) .. controls (536.16,204.94) and (538.22,202.21) .. (540.76,202.21) .. controls (543.31,202.21) and (545.37,204.94) .. (545.37,208.3) .. controls (545.37,211.67) and (543.31,214.4) .. (540.76,214.4) .. controls (538.22,214.4) and (536.16,211.67) .. (536.16,208.3) -- cycle ;
%Curve Lines [id:da32707407117974907] 
\draw [color={rgb, 255:red, 0; green, 0; blue, 0 }  ,draw opacity=1 ][line width=1.5]    (590.96,205.74) .. controls (659.24,249.16) and (630.29,-40.35) .. (513.43,37.36) ;
\draw [shift={(511.67,38.55)}, rotate = 325.34] [fill={rgb, 255:red, 0; green, 0; blue, 0 }  ,fill opacity=1 ][line width=0.08]  [draw opacity=0] (11.61,-5.58) -- (0,0) -- (11.61,5.58) -- cycle    ;
%Curve Lines [id:da7020868150054066] 
\draw [color={rgb, 255:red, 0; green, 0; blue, 0 }  ,draw opacity=1 ][line width=1.5]    (424.72,209.64) .. controls (370.98,220.16) and (381.16,-32.47) .. (511.67,38.55) ;
\draw [shift={(428.98,208.3)}, rotate = 156.81] [fill={rgb, 255:red, 0; green, 0; blue, 0 }  ,fill opacity=1 ][line width=0.08]  [draw opacity=0] (11.61,-5.58) -- (0,0) -- (11.61,5.58) -- cycle    ;
%Straight Lines [id:da6734901748040498] 
\draw [color={rgb, 255:red, 0; green, 0; blue, 0 }  ,draw opacity=1 ][line width=1.5]    (549.36,208.02) -- (581.74,205.74) ;
\draw [shift={(545.37,208.3)}, rotate = 355.96] [fill={rgb, 255:red, 0; green, 0; blue, 0 }  ,fill opacity=1 ][line width=0.08]  [draw opacity=0] (11.61,-5.58) -- (0,0) -- (11.61,5.58) -- cycle    ;
%Curve Lines [id:da03703637046060837] 
\draw [line width=1.5]    (586.35,199.64) .. controls (601.62,160.02) and (453,140.46) .. (434.55,198.55) ;
\draw [shift={(433.59,202.21)}, rotate = 281.99] [fill={rgb, 255:red, 0; green, 0; blue, 0 }  ][line width=0.08]  [draw opacity=0] (11.61,-5.58) -- (0,0) -- (11.61,5.58) -- cycle    ;
%Curve Lines [id:da798612936169141] 
\draw [line width=1.5]    (433.59,214.4) .. controls (453.03,281.98) and (525.13,267.42) .. (539.75,218.25) ;
\draw [shift={(540.76,214.4)}, rotate = 102.83] [fill={rgb, 255:red, 0; green, 0; blue, 0 }  ][line width=0.08]  [draw opacity=0] (11.61,-5.58) -- (0,0) -- (11.61,5.58) -- cycle    ;
%Curve Lines [id:da8888761441506123] 
\draw [line width=1.5]    (433.59,214.4) .. controls (443.17,338.24) and (592.39,286.2) .. (586.72,215.09) ;
\draw [shift={(586.35,211.83)}, rotate = 81.69] [fill={rgb, 255:red, 0; green, 0; blue, 0 }  ][line width=0.08]  [draw opacity=0] (11.61,-5.58) -- (0,0) -- (11.61,5.58) -- cycle    ;
%Shape: Ellipse [id:dp8150638560518382] 
\draw  [color={rgb, 255:red, 0; green, 0; blue, 0 }  ,draw opacity=1 ][line width=1.5]  (196.39,53.32) .. controls (196.39,49.95) and (198.45,47.22) .. (201,47.22) .. controls (203.54,47.22) and (205.61,49.95) .. (205.61,53.32) .. controls (205.61,56.69) and (203.54,59.41) .. (201,59.41) .. controls (198.45,59.41) and (196.39,56.69) .. (196.39,53.32) -- cycle ;
%Shape: Ellipse [id:dp07405350698131474] 
\draw  [color={rgb, 255:red, 0; green, 0; blue, 0 }  ,draw opacity=1 ][line width=1.5]  (196.39,137.39) .. controls (196.39,134.02) and (198.45,131.29) .. (201,131.29) .. controls (203.54,131.29) and (205.61,134.02) .. (205.61,137.39) .. controls (205.61,140.76) and (203.54,143.49) .. (201,143.49) .. controls (198.45,143.49) and (196.39,140.76) .. (196.39,137.39) -- cycle ;
%Curve Lines [id:da5729143394323514] 
\draw [color={rgb, 255:red, 0; green, 0; blue, 0 }  ,draw opacity=1 ][line width=1.5]    (117.61,212.9) .. controls (112.41,176.16) and (142.97,156.8) .. (201,143.49) ;
\draw [shift={(118.31,216.97)}, rotate = 258.41] [fill={rgb, 255:red, 0; green, 0; blue, 0 }  ,fill opacity=1 ][line width=0.08]  [draw opacity=0] (11.61,-5.58) -- (0,0) -- (11.61,5.58) -- cycle    ;
%Curve Lines [id:da985022859832252] 
\draw [color={rgb, 255:red, 0; green, 0; blue, 0 }  ,draw opacity=1 ][line width=1.5]    (280.29,214.4) .. controls (308.81,186.1) and (267.69,160.88) .. (204.87,144.48) ;
\draw [shift={(201,143.49)}, rotate = 14.13] [fill={rgb, 255:red, 0; green, 0; blue, 0 }  ,fill opacity=1 ][line width=0.08]  [draw opacity=0] (11.61,-5.58) -- (0,0) -- (11.61,5.58) -- cycle    ;
%Curve Lines [id:da42144751436003447] 
\draw [color={rgb, 255:red, 0; green, 0; blue, 0 }  ,draw opacity=1 ][line width=1.5]    (280.29,214.4) .. controls (348.57,257.82) and (319.62,-31.69) .. (202.77,46.02) ;
\draw [shift={(201,47.22)}, rotate = 325.34] [fill={rgb, 255:red, 0; green, 0; blue, 0 }  ,fill opacity=1 ][line width=0.08]  [draw opacity=0] (11.61,-5.58) -- (0,0) -- (11.61,5.58) -- cycle    ;
%Curve Lines [id:da7830144424190421] 
\draw [color={rgb, 255:red, 0; green, 0; blue, 0 }  ,draw opacity=1 ][line width=1.5]    (114.05,218.31) .. controls (60.32,228.83) and (70.5,-23.8) .. (201,47.22) ;
\draw [shift={(118.31,216.97)}, rotate = 156.81] [fill={rgb, 255:red, 0; green, 0; blue, 0 }  ,fill opacity=1 ][line width=0.08]  [draw opacity=0] (11.61,-5.58) -- (0,0) -- (11.61,5.58) -- cycle    ;
%Straight Lines [id:da11394903068946982] 
\draw [color={rgb, 255:red, 0; green, 0; blue, 0 }  ,draw opacity=1 ][line width=1.5]    (201,127.29) -- (201,59.41) ;
\draw [shift={(201,131.29)}, rotate = 270] [fill={rgb, 255:red, 0; green, 0; blue, 0 }  ,fill opacity=1 ][line width=0.08]  [draw opacity=0] (11.61,-5.58) -- (0,0) -- (11.61,5.58) -- cycle    ;
%Curve Lines [id:da9079411866287189] 
\draw [line width=1.5]    (275.68,208.31) .. controls (290.96,168.68) and (142.33,149.13) .. (123.88,207.22) ;
\draw [shift={(122.92,210.87)}, rotate = 281.99] [fill={rgb, 255:red, 0; green, 0; blue, 0 }  ][line width=0.08]  [draw opacity=0] (11.61,-5.58) -- (0,0) -- (11.61,5.58) -- cycle    ;
%Shape: Ellipse [id:dp9797878316847205] 
\draw  [color={rgb, 255:red, 0; green, 0; blue, 0 }  ,draw opacity=1 ][line width=1.5]  (271.08,214.4) .. controls (271.08,211.03) and (273.14,208.31) .. (275.68,208.31) .. controls (278.23,208.31) and (280.29,211.03) .. (280.29,214.4) .. controls (280.29,217.77) and (278.23,220.5) .. (275.68,220.5) .. controls (273.14,220.5) and (271.08,217.77) .. (271.08,214.4) -- cycle ;
%Shape: Ellipse [id:dp11924297827138353] 
\draw  [color={rgb, 255:red, 0; green, 0; blue, 0 }  ,draw opacity=1 ][line width=1.5]  (120.31,308.3) .. controls (120.31,304.94) and (122.38,302.21) .. (124.92,302.21) .. controls (127.46,302.21) and (129.53,304.94) .. (129.53,308.3) .. controls (129.53,311.67) and (127.46,314.4) .. (124.92,314.4) .. controls (122.38,314.4) and (120.31,311.67) .. (120.31,308.3) -- cycle ;
%Shape: Ellipse [id:dp18019428491254885] 
\draw  [color={rgb, 255:red, 0; green, 0; blue, 0 }  ,draw opacity=1 ][line width=1.5]  (118.31,216.97) .. controls (118.31,213.6) and (120.38,210.87) .. (122.92,210.87) .. controls (125.46,210.87) and (127.53,213.6) .. (127.53,216.97) .. controls (127.53,220.34) and (125.46,223.07) .. (122.92,223.07) .. controls (120.38,223.07) and (118.31,220.34) .. (118.31,216.97) -- cycle ;
%Shape: Ellipse [id:dp2447839509836378] 
\draw  [color={rgb, 255:red, 0; green, 0; blue, 0 }  ,draw opacity=1 ][line width=1.5]  (171.23,308.3) .. controls (171.23,304.94) and (173.3,302.21) .. (175.84,302.21) .. controls (178.39,302.21) and (180.45,304.94) .. (180.45,308.3) .. controls (180.45,311.67) and (178.39,314.4) .. (175.84,314.4) .. controls (173.3,314.4) and (171.23,311.67) .. (171.23,308.3) -- cycle ;
%Shape: Ellipse [id:dp33013962263442154] 
\draw  [color={rgb, 255:red, 0; green, 0; blue, 0 }  ,draw opacity=1 ][line width=1.5]  (227.49,308.3) .. controls (227.49,304.94) and (229.55,302.21) .. (232.1,302.21) .. controls (234.64,302.21) and (236.7,304.94) .. (236.7,308.3) .. controls (236.7,311.67) and (234.64,314.4) .. (232.1,314.4) .. controls (229.55,314.4) and (227.49,311.67) .. (227.49,308.3) -- cycle ;
%Straight Lines [id:da31396117550667135] 
\draw [color={rgb, 255:red, 0; green, 0; blue, 0 }  ,draw opacity=1 ][line width=1.5]    (133.53,308.3) -- (171.23,308.3) ;
\draw [shift={(129.53,308.3)}, rotate = 0] [fill={rgb, 255:red, 0; green, 0; blue, 0 }  ,fill opacity=1 ][line width=0.08]  [draw opacity=0] (11.61,-5.58) -- (0,0) -- (11.61,5.58) -- cycle    ;
%Straight Lines [id:da6005265572468635] 
\draw [color={rgb, 255:red, 0; green, 0; blue, 0 }  ,draw opacity=1 ][line width=1.5]    (180.45,308.3) -- (223.49,308.3) ;
\draw [shift={(227.49,308.3)}, rotate = 180] [fill={rgb, 255:red, 0; green, 0; blue, 0 }  ,fill opacity=1 ][line width=0.08]  [draw opacity=0] (11.61,-5.58) -- (0,0) -- (11.61,5.58) -- cycle    ;
%Straight Lines [id:da1808031690077344] 
\draw [color={rgb, 255:red, 0; green, 0; blue, 0 }  ,draw opacity=1 ][line width=1.5]    (240.69,308.02) -- (273.08,305.74) ;
\draw [shift={(236.7,308.3)}, rotate = 355.96] [fill={rgb, 255:red, 0; green, 0; blue, 0 }  ,fill opacity=1 ][line width=0.08]  [draw opacity=0] (11.61,-5.58) -- (0,0) -- (11.61,5.58) -- cycle    ;
%Curve Lines [id:da3357677157773651] 
\draw [line width=1.5]    (277.68,299.64) .. controls (292.96,260.02) and (144.33,240.46) .. (125.88,298.55) ;
\draw [shift={(124.92,302.21)}, rotate = 281.99] [fill={rgb, 255:red, 0; green, 0; blue, 0 }  ][line width=0.08]  [draw opacity=0] (11.61,-5.58) -- (0,0) -- (11.61,5.58) -- cycle    ;
%Curve Lines [id:da670250481056926] 
\draw [line width=1.5]    (179.35,299.27) .. controls (211.8,273.31) and (249.51,274.2) .. (277.68,299.64) ;
\draw [shift={(175.84,302.21)}, rotate = 318.94] [fill={rgb, 255:red, 0; green, 0; blue, 0 }  ][line width=0.08]  [draw opacity=0] (11.61,-5.58) -- (0,0) -- (11.61,5.58) -- cycle    ;
%Curve Lines [id:da1585249402391844] 
\draw [line width=1.5]    (124.92,314.4) .. controls (144.36,381.98) and (216.47,367.42) .. (231.09,318.25) ;
\draw [shift={(232.1,314.4)}, rotate = 102.83] [fill={rgb, 255:red, 0; green, 0; blue, 0 }  ][line width=0.08]  [draw opacity=0] (11.61,-5.58) -- (0,0) -- (11.61,5.58) -- cycle    ;
%Shape: Ellipse [id:dp8849254018327529] 
\draw  [color={rgb, 255:red, 0; green, 0; blue, 0 }  ,draw opacity=1 ][line width=1.5]  (273.08,305.74) .. controls (273.08,302.37) and (275.14,299.64) .. (277.68,299.64) .. controls (280.23,299.64) and (282.29,302.37) .. (282.29,305.74) .. controls (282.29,309.1) and (280.23,311.83) .. (277.68,311.83) .. controls (275.14,311.83) and (273.08,309.1) .. (273.08,305.74) -- cycle ;

% Text Node
\draw (601.26,38.41) node [anchor=north west][inner sep=0.75pt]    {$4$};
% Text Node
\draw (397.44,47.33) node [anchor=north west][inner sep=0.75pt]    {$1$};
% Text Node
\draw (558.52,211.43) node [anchor=north west][inner sep=0.75pt]    {$4$};
% Text Node
\draw (532.09,247.38) node [anchor=north west][inner sep=0.75pt]    {$D_{2}$};
% Text Node
\draw (594.24,244.82) node [anchor=north west][inner sep=0.75pt]    {$D_{3}$};
% Text Node
\draw (442.27,268.73) node [anchor=north west][inner sep=0.75pt]    {$0$};
% Text Node
\draw (133.88,107.84) node [anchor=north west][inner sep=0.75pt]    {$A$};
% Text Node
\draw (253.06,103.72) node [anchor=north west][inner sep=0.75pt]    {$B$};
% Text Node
\draw (290.6,47.07) node [anchor=north west][inner sep=0.75pt]    {$4$};
% Text Node
\draw (142.03,140.11) node [anchor=north west][inner sep=0.75pt]    {$4$};
% Text Node
\draw (254.63,140.45) node [anchor=north west][inner sep=0.75pt]    {$1$};
% Text Node
\draw (86.77,56) node [anchor=north west][inner sep=0.75pt]    {$1$};
% Text Node
\draw (186.34,90.67) node [anchor=north west][inner sep=0.75pt]    {$1$};
% Text Node
\draw (186.56,148.89) node [anchor=north west][inner sep=0.75pt]    {$C_{1}$};
% Text Node
\draw (218.02,156.44) node [anchor=north west][inner sep=0.75pt]    {$7$};
% Text Node
\draw (190,188.4) node [anchor=north west][inner sep=0.75pt]    {$E_{3}$};
% Text Node
\draw (144.55,312.12) node [anchor=north west][inner sep=0.75pt]    {$4$};
% Text Node
\draw (199.81,310.06) node [anchor=north west][inner sep=0.75pt]    {$4$};
% Text Node
\draw (249.85,311.43) node [anchor=north west][inner sep=0.75pt]    {$4$};
% Text Node
\draw (217.5,284.38) node [anchor=north west][inner sep=0.75pt]    {$C_{3}$};
% Text Node
\draw (151.1,276.12) node [anchor=north west][inner sep=0.75pt]    {$C_{2}$};
% Text Node
\draw (220.02,247.77) node [anchor=north west][inner sep=0.75pt]    {$7$};
% Text Node
\draw (189.37,271.13) node [anchor=north west][inner sep=0.75pt]    {$8$};
% Text Node
\draw (164.69,325.59) node [anchor=north west][inner sep=0.75pt]    {$D_{1}$};
% Text Node
\draw (223.43,347.38) node [anchor=north west][inner sep=0.75pt]    {$E_{4}$};
% Text Node
\draw (147.8,335.38) node [anchor=north west][inner sep=0.75pt]    {$9$};
% Text Node
\draw (512.68,143.77) node [anchor=north west][inner sep=0.75pt]    {$7$};
% Text Node
\draw (502.67,95.07) node [anchor=north west][inner sep=0.75pt]    {$E_{3}$};
% Text Node
\draw (475.13,233.38) node [anchor=north west][inner sep=0.75pt]    {$9$};
% Text Node
\draw (474.76,204.05) node [anchor=north west][inner sep=0.75pt]    {$E_{4}$};


\end{tikzpicture}
    \caption{Subdivisión del poliedro de abajo}
\end{figure}
Observe particularmente que el poliedro de la derecha del todo no es exactamente un tetraedro, ya que la unión de las caras $E_4$ y $D_2$ forman un digono y no un cuadrilátero, como pasaría en un tetraedro, Por lo que si identificamos 0 con 7, cuando hacemos esto las caras $E_4$ y $D_2$ se identifican entre si, lo mismo ocurre con $E_3$ y $D_3$, por lo que esto no es otra cosa que el grafo de un triangulo sobre la 3-bola, por lo que en el pegado no tienen incidencia, así nuestra triangulación topológica ideal del complemento del nudo $6_1$ la podemos ver como el pegado de los siguientes 5 tetraedros todos vistos ahora desde afuera.
\begin{figure}[H]
\centering
%!TEX root = ./main.tex



\tikzset{every picture/.style={line width=0.75pt}} %set default line width to 0.75pt        

\begin{tikzpicture}[x=0.75pt,y=0.75pt,yscale=-1,xscale=1]
%uncomment if require: \path (0,498); %set diagram left start at 0, and has height of 498

%Shape: Ellipse [id:dp722498185471097] 
\draw  [line width=1.5]  (140.33,28.88) .. controls (140.33,26.46) and (142.38,24.5) .. (144.91,24.5) .. controls (147.43,24.5) and (149.48,26.46) .. (149.48,28.88) .. controls (149.48,31.29) and (147.43,33.25) .. (144.91,33.25) .. controls (142.38,33.25) and (140.33,31.29) .. (140.33,28.88) -- cycle ;
%Shape: Ellipse [id:dp09419658359942129] 
\draw  [line width=1.5]  (140.9,102.76) .. controls (140.9,100.34) and (142.95,98.38) .. (145.48,98.38) .. controls (148.01,98.38) and (150.05,100.34) .. (150.05,102.76) .. controls (150.05,105.18) and (148.01,107.14) .. (145.48,107.14) .. controls (142.95,107.14) and (140.9,105.18) .. (140.9,102.76) -- cycle ;
%Shape: Ellipse [id:dp4574028971026852] 
\draw  [line width=1.5]  (76.27,155.84) .. controls (76.27,153.42) and (78.32,151.46) .. (80.85,151.46) .. controls (83.37,151.46) and (85.42,153.42) .. (85.42,155.84) .. controls (85.42,158.26) and (83.37,160.22) .. (80.85,160.22) .. controls (78.32,160.22) and (76.27,158.26) .. (76.27,155.84) -- cycle ;
%Shape: Ellipse [id:dp012434768366862081] 
\draw  [line width=1.5]  (204.96,155.84) .. controls (204.96,153.42) and (207.01,151.46) .. (209.54,151.46) .. controls (212.06,151.46) and (214.11,153.42) .. (214.11,155.84) .. controls (214.11,158.26) and (212.06,160.22) .. (209.54,160.22) .. controls (207.01,160.22) and (204.96,158.26) .. (204.96,155.84) -- cycle ;
%Straight Lines [id:da7104639893284079] 
\draw [line width=1.5]    (144.94,37.25) -- (145.48,98.38) ;
\draw [shift={(144.91,33.25)}, rotate = 89.5] [fill={rgb, 255:red, 0; green, 0; blue, 0 }  ][line width=0.08]  [draw opacity=0] (11.61,-5.58) -- (0,0) -- (11.61,5.58) -- cycle    ;
%Straight Lines [id:da06252865798857044] 
\draw [line width=1.5]    (89.42,155.84) -- (204.96,155.84) ;
\draw [shift={(85.42,155.84)}, rotate = 0] [fill={rgb, 255:red, 0; green, 0; blue, 0 }  ][line width=0.08]  [draw opacity=0] (11.61,-5.58) -- (0,0) -- (11.61,5.58) -- cycle    ;
%Straight Lines [id:da8407976043494823] 
\draw [line width=1.5]    (83.95,148.94) -- (140.9,102.76) ;
\draw [shift={(80.85,151.46)}, rotate = 320.96] [fill={rgb, 255:red, 0; green, 0; blue, 0 }  ][line width=0.08]  [draw opacity=0] (11.61,-5.58) -- (0,0) -- (11.61,5.58) -- cycle    ;
%Straight Lines [id:da93252160050729] 
\draw [line width=1.5]    (150.05,102.76) -- (206.44,148.93) ;
\draw [shift={(209.54,151.46)}, rotate = 219.31] [fill={rgb, 255:red, 0; green, 0; blue, 0 }  ][line width=0.08]  [draw opacity=0] (11.61,-5.58) -- (0,0) -- (11.61,5.58) -- cycle    ;
%Curve Lines [id:da914075530376739] 
\draw [line width=1.5]    (76.27,155.84) .. controls (77.1,77.95) and (106.42,42.09) .. (137,30.09) ;
\draw [shift={(140.33,28.88)}, rotate = 161.42] [fill={rgb, 255:red, 0; green, 0; blue, 0 }  ][line width=0.08]  [draw opacity=0] (11.61,-5.58) -- (0,0) -- (11.61,5.58) -- cycle    ;
%Curve Lines [id:da5685611468391036] 
\draw [line width=1.5]    (153.73,29.26) .. controls (187.82,35.57) and (224.82,118.28) .. (214.11,155.84) ;
\draw [shift={(149.48,28.88)}, rotate = 359.55] [fill={rgb, 255:red, 0; green, 0; blue, 0 }  ][line width=0.08]  [draw opacity=0] (11.61,-5.58) -- (0,0) -- (11.61,5.58) -- cycle    ;
%Shape: Ellipse [id:dp9871169656532187] 
\draw  [line width=1.5]  (320.49,28.33) .. controls (320.49,25.91) and (322.54,23.95) .. (325.07,23.95) .. controls (327.6,23.95) and (329.65,25.91) .. (329.65,28.33) .. controls (329.65,30.75) and (327.6,32.71) .. (325.07,32.71) .. controls (322.54,32.71) and (320.49,30.75) .. (320.49,28.33) -- cycle ;
%Shape: Ellipse [id:dp4998971196831187] 
\draw  [line width=1.5]  (321.07,102.21) .. controls (321.07,99.79) and (323.12,97.83) .. (325.64,97.83) .. controls (328.17,97.83) and (330.22,99.79) .. (330.22,102.21) .. controls (330.22,104.63) and (328.17,106.59) .. (325.64,106.59) .. controls (323.12,106.59) and (321.07,104.63) .. (321.07,102.21) -- cycle ;
%Shape: Ellipse [id:dp5334863650227393] 
\draw  [line width=1.5]  (256.44,155.3) .. controls (256.44,152.88) and (258.49,150.92) .. (261.01,150.92) .. controls (263.54,150.92) and (265.59,152.88) .. (265.59,155.3) .. controls (265.59,157.71) and (263.54,159.67) .. (261.01,159.67) .. controls (258.49,159.67) and (256.44,157.71) .. (256.44,155.3) -- cycle ;
%Shape: Ellipse [id:dp8295563632609373] 
\draw  [line width=1.5]  (385.13,155.3) .. controls (385.13,152.88) and (387.17,150.92) .. (389.7,150.92) .. controls (392.23,150.92) and (394.28,152.88) .. (394.28,155.3) .. controls (394.28,157.71) and (392.23,159.67) .. (389.7,159.67) .. controls (387.17,159.67) and (385.13,157.71) .. (385.13,155.3) -- cycle ;
%Straight Lines [id:da6853527229248998] 
\draw [line width=1.5]    (325.11,36.71) -- (325.64,97.83) ;
\draw [shift={(325.07,32.71)}, rotate = 89.5] [fill={rgb, 255:red, 0; green, 0; blue, 0 }  ][line width=0.08]  [draw opacity=0] (11.61,-5.58) -- (0,0) -- (11.61,5.58) -- cycle    ;
%Straight Lines [id:da37663298733859474] 
\draw [line width=1.5]    (265.59,155.3) -- (381.13,155.3) ;
\draw [shift={(385.13,155.3)}, rotate = 180] [fill={rgb, 255:red, 0; green, 0; blue, 0 }  ][line width=0.08]  [draw opacity=0] (11.61,-5.58) -- (0,0) -- (11.61,5.58) -- cycle    ;
%Straight Lines [id:da1331834266330163] 
\draw [line width=1.5]    (264.12,148.4) -- (321.07,102.21) ;
\draw [shift={(261.01,150.92)}, rotate = 320.96] [fill={rgb, 255:red, 0; green, 0; blue, 0 }  ][line width=0.08]  [draw opacity=0] (11.61,-5.58) -- (0,0) -- (11.61,5.58) -- cycle    ;
%Straight Lines [id:da93580380730357] 
\draw [line width=1.5]    (330.22,102.21) -- (386.61,148.38) ;
\draw [shift={(389.7,150.92)}, rotate = 219.31] [fill={rgb, 255:red, 0; green, 0; blue, 0 }  ][line width=0.08]  [draw opacity=0] (11.61,-5.58) -- (0,0) -- (11.61,5.58) -- cycle    ;
%Curve Lines [id:da9473646706061418] 
\draw [line width=1.5]    (256.44,155.3) .. controls (257.26,77.4) and (286.59,41.55) .. (317.16,29.54) ;
\draw [shift={(320.49,28.33)}, rotate = 161.42] [fill={rgb, 255:red, 0; green, 0; blue, 0 }  ][line width=0.08]  [draw opacity=0] (11.61,-5.58) -- (0,0) -- (11.61,5.58) -- cycle    ;
%Curve Lines [id:da4025171529168252] 
\draw [line width=1.5]    (329.65,28.33) .. controls (363.77,28.06) and (403,110.96) .. (395.14,151.65) ;
\draw [shift={(394.28,155.3)}, rotate = 285.91] [fill={rgb, 255:red, 0; green, 0; blue, 0 }  ][line width=0.08]  [draw opacity=0] (11.61,-5.58) -- (0,0) -- (11.61,5.58) -- cycle    ;
%Shape: Ellipse [id:dp4598430040551579] 
\draw  [line width=1.5]  (503.52,25.04) .. controls (503.52,22.63) and (505.57,20.67) .. (508.09,20.67) .. controls (510.62,20.67) and (512.67,22.63) .. (512.67,25.04) .. controls (512.67,27.46) and (510.62,29.42) .. (508.09,29.42) .. controls (505.57,29.42) and (503.52,27.46) .. (503.52,25.04) -- cycle ;
%Shape: Ellipse [id:dp5035867782009985] 
\draw  [line width=1.5]  (504.09,98.93) .. controls (504.09,96.51) and (506.14,94.55) .. (508.67,94.55) .. controls (511.19,94.55) and (513.24,96.51) .. (513.24,98.93) .. controls (513.24,101.34) and (511.19,103.3) .. (508.67,103.3) .. controls (506.14,103.3) and (504.09,101.34) .. (504.09,98.93) -- cycle ;
%Shape: Ellipse [id:dp6235578561815645] 
\draw  [line width=1.5]  (439.46,152.01) .. controls (439.46,149.59) and (441.51,147.63) .. (444.04,147.63) .. controls (446.56,147.63) and (448.61,149.59) .. (448.61,152.01) .. controls (448.61,154.43) and (446.56,156.39) .. (444.04,156.39) .. controls (441.51,156.39) and (439.46,154.43) .. (439.46,152.01) -- cycle ;
%Shape: Ellipse [id:dp47621421692962784] 
\draw  [line width=1.5]  (568.15,152.01) .. controls (568.15,149.59) and (570.2,147.63) .. (572.73,147.63) .. controls (575.25,147.63) and (577.3,149.59) .. (577.3,152.01) .. controls (577.3,154.43) and (575.25,156.39) .. (572.73,156.39) .. controls (570.2,156.39) and (568.15,154.43) .. (568.15,152.01) -- cycle ;
%Straight Lines [id:da24752576794389303] 
\draw [line width=1.5]    (508.13,33.42) -- (508.67,94.55) ;
\draw [shift={(508.09,29.42)}, rotate = 89.5] [fill={rgb, 255:red, 0; green, 0; blue, 0 }  ][line width=0.08]  [draw opacity=0] (11.61,-5.58) -- (0,0) -- (11.61,5.58) -- cycle    ;
%Straight Lines [id:da2745160206955022] 
\draw [line width=1.5]    (452.61,152.01) -- (568.15,152.01) ;
\draw [shift={(448.61,152.01)}, rotate = 0] [fill={rgb, 255:red, 0; green, 0; blue, 0 }  ][line width=0.08]  [draw opacity=0] (11.61,-5.58) -- (0,0) -- (11.61,5.58) -- cycle    ;
%Straight Lines [id:da1251063457941054] 
\draw [line width=1.5]    (447.14,145.11) -- (504.09,98.93) ;
\draw [shift={(444.04,147.63)}, rotate = 320.96] [fill={rgb, 255:red, 0; green, 0; blue, 0 }  ][line width=0.08]  [draw opacity=0] (11.61,-5.58) -- (0,0) -- (11.61,5.58) -- cycle    ;
%Straight Lines [id:da8560302252250104] 
\draw [line width=1.5]    (513.24,98.93) -- (569.63,145.1) ;
\draw [shift={(572.73,147.63)}, rotate = 219.31] [fill={rgb, 255:red, 0; green, 0; blue, 0 }  ][line width=0.08]  [draw opacity=0] (11.61,-5.58) -- (0,0) -- (11.61,5.58) -- cycle    ;
%Curve Lines [id:da5942103699772806] 
\draw [line width=1.5]    (439.55,147.22) .. controls (441.58,69.88) and (472.41,35.5) .. (503.52,25.04) ;
\draw [shift={(439.46,152.01)}, rotate = 270.61] [fill={rgb, 255:red, 0; green, 0; blue, 0 }  ][line width=0.08]  [draw opacity=0] (11.61,-5.58) -- (0,0) -- (11.61,5.58) -- cycle    ;
%Curve Lines [id:da7202124541602778] 
\draw [line width=1.5]    (516.91,25.43) .. controls (551.01,31.74) and (588.01,114.45) .. (577.3,152.01) ;
\draw [shift={(512.67,25.04)}, rotate = 359.55] [fill={rgb, 255:red, 0; green, 0; blue, 0 }  ][line width=0.08]  [draw opacity=0] (11.61,-5.58) -- (0,0) -- (11.61,5.58) -- cycle    ;
%Shape: Ellipse [id:dp7971440912052996] 
\draw  [line width=1.5]  (228.98,199.62) .. controls (228.98,197.21) and (231.03,195.25) .. (233.56,195.25) .. controls (236.09,195.25) and (238.13,197.21) .. (238.13,199.62) .. controls (238.13,202.04) and (236.09,204) .. (233.56,204) .. controls (231.03,204) and (228.98,202.04) .. (228.98,199.62) -- cycle ;
%Shape: Ellipse [id:dp4840470955224502] 
\draw  [line width=1.5]  (229.55,273.51) .. controls (229.55,271.09) and (231.6,269.13) .. (234.13,269.13) .. controls (236.66,269.13) and (238.71,271.09) .. (238.71,273.51) .. controls (238.71,275.92) and (236.66,277.88) .. (234.13,277.88) .. controls (231.6,277.88) and (229.55,275.92) .. (229.55,273.51) -- cycle ;
%Shape: Ellipse [id:dp7205406460641789] 
\draw  [line width=1.5]  (164.92,326.59) .. controls (164.92,324.17) and (166.97,322.21) .. (169.5,322.21) .. controls (172.03,322.21) and (174.08,324.17) .. (174.08,326.59) .. controls (174.08,329.01) and (172.03,330.97) .. (169.5,330.97) .. controls (166.97,330.97) and (164.92,329.01) .. (164.92,326.59) -- cycle ;
%Shape: Ellipse [id:dp16212402515113167] 
\draw  [line width=1.5]  (293.61,326.59) .. controls (293.61,324.17) and (295.66,322.21) .. (298.19,322.21) .. controls (300.72,322.21) and (302.76,324.17) .. (302.76,326.59) .. controls (302.76,329.01) and (300.72,330.97) .. (298.19,330.97) .. controls (295.66,330.97) and (293.61,329.01) .. (293.61,326.59) -- cycle ;
%Straight Lines [id:da38679418807864374] 
\draw [line width=1.5]    (233.56,204) -- (234.1,265.13) ;
\draw [shift={(234.13,269.13)}, rotate = 269.5] [fill={rgb, 255:red, 0; green, 0; blue, 0 }  ][line width=0.08]  [draw opacity=0] (11.61,-5.58) -- (0,0) -- (11.61,5.58) -- cycle    ;
%Straight Lines [id:da7340741733057103] 
\draw [line width=1.5]    (178.08,326.59) -- (293.61,326.59) ;
\draw [shift={(174.08,326.59)}, rotate = 0] [fill={rgb, 255:red, 0; green, 0; blue, 0 }  ][line width=0.08]  [draw opacity=0] (11.61,-5.58) -- (0,0) -- (11.61,5.58) -- cycle    ;
%Straight Lines [id:da6492127454486433] 
\draw [line width=1.5]    (172.61,319.69) -- (229.55,273.51) ;
\draw [shift={(169.5,322.21)}, rotate = 320.96] [fill={rgb, 255:red, 0; green, 0; blue, 0 }  ][line width=0.08]  [draw opacity=0] (11.61,-5.58) -- (0,0) -- (11.61,5.58) -- cycle    ;
%Straight Lines [id:da1534318649264561] 
\draw [line width=1.5]    (241.8,276.04) -- (298.19,322.21) ;
\draw [shift={(238.71,273.51)}, rotate = 39.31] [fill={rgb, 255:red, 0; green, 0; blue, 0 }  ][line width=0.08]  [draw opacity=0] (11.61,-5.58) -- (0,0) -- (11.61,5.58) -- cycle    ;
%Curve Lines [id:da15708072815584284] 
\draw [line width=1.5]    (165.01,321.8) .. controls (167.04,244.46) and (197.87,210.08) .. (228.98,199.62) ;
\draw [shift={(164.92,326.59)}, rotate = 270.61] [fill={rgb, 255:red, 0; green, 0; blue, 0 }  ][line width=0.08]  [draw opacity=0] (11.61,-5.58) -- (0,0) -- (11.61,5.58) -- cycle    ;
%Curve Lines [id:da4532354822589155] 
\draw [line width=1.5]    (242.38,200.01) .. controls (276.47,206.32) and (313.47,289.03) .. (302.76,326.59) ;
\draw [shift={(238.13,199.62)}, rotate = 359.55] [fill={rgb, 255:red, 0; green, 0; blue, 0 }  ][line width=0.08]  [draw opacity=0] (11.61,-5.58) -- (0,0) -- (11.61,5.58) -- cycle    ;
%Shape: Ellipse [id:dp707487152357257] 
\draw  [line width=1.5]  (420.01,199.08) .. controls (420.01,196.66) and (422.06,194.7) .. (424.59,194.7) .. controls (427.12,194.7) and (429.17,196.66) .. (429.17,199.08) .. controls (429.17,201.5) and (427.12,203.46) .. (424.59,203.46) .. controls (422.06,203.46) and (420.01,201.5) .. (420.01,199.08) -- cycle ;
%Shape: Ellipse [id:dp2675712667635096] 
\draw  [line width=1.5]  (420.59,272.96) .. controls (420.59,270.54) and (422.63,268.58) .. (425.16,268.58) .. controls (427.69,268.58) and (429.74,270.54) .. (429.74,272.96) .. controls (429.74,275.38) and (427.69,277.34) .. (425.16,277.34) .. controls (422.63,277.34) and (420.59,275.38) .. (420.59,272.96) -- cycle ;
%Shape: Ellipse [id:dp04604583726861322] 
\draw  [line width=1.5]  (355.96,326.04) .. controls (355.96,323.63) and (358,321.67) .. (360.53,321.67) .. controls (363.06,321.67) and (365.11,323.63) .. (365.11,326.04) .. controls (365.11,328.46) and (363.06,330.42) .. (360.53,330.42) .. controls (358,330.42) and (355.96,328.46) .. (355.96,326.04) -- cycle ;
%Shape: Ellipse [id:dp527192499801234] 
\draw  [line width=1.5]  (484.64,326.04) .. controls (484.64,323.63) and (486.69,321.67) .. (489.22,321.67) .. controls (491.75,321.67) and (493.8,323.63) .. (493.8,326.04) .. controls (493.8,328.46) and (491.75,330.42) .. (489.22,330.42) .. controls (486.69,330.42) and (484.64,328.46) .. (484.64,326.04) -- cycle ;
%Straight Lines [id:da2748072495693461] 
\draw [line width=1.5]    (424.63,207.46) -- (425.16,268.58) ;
\draw [shift={(424.59,203.46)}, rotate = 89.5] [fill={rgb, 255:red, 0; green, 0; blue, 0 }  ][line width=0.08]  [draw opacity=0] (11.61,-5.58) -- (0,0) -- (11.61,5.58) -- cycle    ;
%Straight Lines [id:da8571687174187388] 
\draw [line width=1.5]    (369.11,326.04) -- (484.64,326.04) ;
\draw [shift={(365.11,326.04)}, rotate = 0] [fill={rgb, 255:red, 0; green, 0; blue, 0 }  ][line width=0.08]  [draw opacity=0] (11.61,-5.58) -- (0,0) -- (11.61,5.58) -- cycle    ;
%Straight Lines [id:da009179926411796768] 
\draw [line width=1.5]    (363.64,319.15) -- (420.59,272.96) ;
\draw [shift={(360.53,321.67)}, rotate = 320.96] [fill={rgb, 255:red, 0; green, 0; blue, 0 }  ][line width=0.08]  [draw opacity=0] (11.61,-5.58) -- (0,0) -- (11.61,5.58) -- cycle    ;
%Straight Lines [id:da09268248986267924] 
\draw [line width=1.5]    (432.83,275.49) -- (489.22,321.67) ;
\draw [shift={(429.74,272.96)}, rotate = 39.31] [fill={rgb, 255:red, 0; green, 0; blue, 0 }  ][line width=0.08]  [draw opacity=0] (11.61,-5.58) -- (0,0) -- (11.61,5.58) -- cycle    ;
%Curve Lines [id:da7275208733646754] 
\draw [line width=1.5]    (356.04,321.25) .. controls (358.07,243.91) and (388.91,209.54) .. (420.01,199.08) ;
\draw [shift={(355.96,326.04)}, rotate = 270.61] [fill={rgb, 255:red, 0; green, 0; blue, 0 }  ][line width=0.08]  [draw opacity=0] (11.61,-5.58) -- (0,0) -- (11.61,5.58) -- cycle    ;
%Curve Lines [id:da8929863303964553] 
\draw [line width=1.5]    (433.41,199.46) .. controls (467.51,205.77) and (504.5,288.48) .. (493.8,326.04) ;
\draw [shift={(429.17,199.08)}, rotate = 359.55] [fill={rgb, 255:red, 0; green, 0; blue, 0 }  ][line width=0.08]  [draw opacity=0] (11.61,-5.58) -- (0,0) -- (11.61,5.58) -- cycle    ;

% Text Node
\draw (135.57,124.53) node [anchor=north west][inner sep=0.75pt]    {$C_{3}$};
% Text Node
\draw (108.19,78.19) node [anchor=north west][inner sep=0.75pt]    {$B$};
% Text Node
\draw (168.24,80.2) node [anchor=north west][inner sep=0.75pt]    {$E_{2}$};
% Text Node
\draw (109.12,171.05) node [anchor=north west][inner sep=0.75pt]    {$D_{3}$};
% Text Node
\draw (104.18,109.93) node [anchor=north west][inner sep=0.75pt]    {$4$};
% Text Node
\draw (180.82,114.59) node [anchor=north west][inner sep=0.75pt]    {$8$};
% Text Node
\draw (140.22,155.36) node [anchor=north west][inner sep=0.75pt]    {$4$};
% Text Node
\draw (134.5,59.31) node [anchor=north west][inner sep=0.75pt]    {$1$};
% Text Node
\draw (80.73,57.67) node [anchor=north west][inner sep=0.75pt]    {$1$};
% Text Node
\draw (197.41,54.93) node [anchor=north west][inner sep=0.75pt]    {$7$};
% Text Node
\draw (284.35,80.2) node [anchor=north west][inner sep=0.75pt]    {$E_{2}$};
% Text Node
\draw (346.12,79.65) node [anchor=north west][inner sep=0.75pt]    {$E_{1}$};
% Text Node
\draw (319.17,124.53) node [anchor=north west][inner sep=0.75pt]    {$C_{2}$};
% Text Node
\draw (272.13,163.66) node [anchor=north west][inner sep=0.75pt]    {$D_{2}$};
% Text Node
\draw (280.92,111.85) node [anchor=north west][inner sep=0.75pt]    {$8$};
% Text Node
\draw (358.13,109.11) node [anchor=north west][inner sep=0.75pt]    {$7$};
% Text Node
\draw (319.24,155.63) node [anchor=north west][inner sep=0.75pt]    {$4$};
% Text Node
\draw (327.82,60.95) node [anchor=north west][inner sep=0.75pt]    {$1$};
% Text Node
\draw (270.62,40.16) node [anchor=north west][inner sep=0.75pt]    {$7$};
% Text Node
\draw (369.57,40.7) node [anchor=north west][inner sep=0.75pt]    {$9$};
% Text Node
\draw (473.09,71.17) node [anchor=north west][inner sep=0.75pt]    {$E_{1}$};
% Text Node
\draw (530.07,71.9) node [anchor=north west][inner sep=0.75pt]    {$A$};
% Text Node
\draw (502.19,120.97) node [anchor=north west][inner sep=0.75pt]    {$C_{1}$};
% Text Node
\draw (451.72,161.47) node [anchor=north west][inner sep=0.75pt]    {$D_{1}$};
% Text Node
\draw (467.94,103.09) node [anchor=north west][inner sep=0.75pt]    {$7$};
% Text Node
\draw (508.55,57.67) node [anchor=north west][inner sep=0.75pt]    {$1$};
% Text Node
\draw (541.72,103.09) node [anchor=north west][inner sep=0.75pt]    {$1$};
% Text Node
\draw (502.26,153.99) node [anchor=north west][inner sep=0.75pt]    {$4$};
% Text Node
\draw (556.02,41.8) node [anchor=north west][inner sep=0.75pt]    {$4$};
% Text Node
\draw (446.21,42.35) node [anchor=north west][inner sep=0.75pt]    {$9$};
% Text Node
\draw (198.91,242.37) node [anchor=north west][inner sep=0.75pt]    {$A$};
% Text Node
\draw (255.18,241.28) node [anchor=north west][inner sep=0.75pt]    {$B$};
% Text Node
\draw (225.37,294.73) node [anchor=north west][inner sep=0.75pt]    {$C_{1}$};
% Text Node
\draw (182.33,336.87) node [anchor=north west][inner sep=0.75pt]    {$D_{3}$};
% Text Node
\draw (390.52,241.1) node [anchor=north west][inner sep=0.75pt]    {$D_{1}$};
% Text Node
\draw (445,242.19) node [anchor=north west][inner sep=0.75pt]    {$C_{2}$};
% Text Node
\draw (415.83,290.9) node [anchor=north west][inner sep=0.75pt]    {$C_{3}$};
% Text Node
\draw (363.07,340.7) node [anchor=north west][inner sep=0.75pt]    {$D_{2}$};
% Text Node
\draw (164.24,237.72) node [anchor=north west][inner sep=0.75pt]    {$1$};
% Text Node
\draw (234.02,230.06) node [anchor=north west][inner sep=0.75pt]    {$1$};
% Text Node
\draw (266.04,280.41) node [anchor=north west][inner sep=0.75pt]    {$1$};
% Text Node
\draw (193.98,277.67) node [anchor=north west][inner sep=0.75pt]    {$4$};
% Text Node
\draw (294.07,239.91) node [anchor=north west][inner sep=0.75pt]    {$4$};
% Text Node
\draw (231.16,331.85) node [anchor=north west][inner sep=0.75pt]    {$7$};
% Text Node
\draw (382.72,278.22) node [anchor=north west][inner sep=0.75pt]    {$4$};
% Text Node
\draw (427.33,233.89) node [anchor=north west][inner sep=0.75pt]    {$4$};
% Text Node
\draw (414.75,326.38) node [anchor=north west][inner sep=0.75pt]    {$4$};
% Text Node
\draw (356.41,230.06) node [anchor=north west][inner sep=0.75pt]    {$9$};
% Text Node
\draw (457.08,279.86) node [anchor=north west][inner sep=0.75pt]    {$8$};
% Text Node
\draw (480.53,229.51) node [anchor=north west][inner sep=0.75pt]    {$7$};
% Text Node
\draw (110.69,95.8) node [anchor=north west][inner sep=0.75pt]    {$z_{1}$};
% Text Node
\draw (206.21,72.81) node [anchor=north west][inner sep=0.75pt]    {$z_{1}$};
% Text Node
\draw (168.46,101.27) node [anchor=north west][inner sep=0.75pt]    {$z_{2}$};
% Text Node
\draw (66.65,79.38) node [anchor=north west][inner sep=0.75pt]    {$z_{2}$};
% Text Node
\draw (153.01,157.09) node [anchor=north west][inner sep=0.75pt]    {$z_{3}$};
% Text Node
\draw (146.72,60.77) node [anchor=north west][inner sep=0.75pt]    {$z_{3}$};
% Text Node
\draw (285.92,99.08) node [anchor=north west][inner sep=0.75pt]    {$w_{1}$};
% Text Node
\draw (383.16,73.91) node [anchor=north west][inner sep=0.75pt]    {$w_{1}$};
% Text Node
\draw (340.83,96.89) node [anchor=north west][inner sep=0.75pt]    {$w_{2}$};
% Text Node
\draw (247.03,68.98) node [anchor=north west][inner sep=0.75pt]    {$w_{2}$};
% Text Node
\draw (335.68,151.62) node [anchor=north west][inner sep=0.75pt]    {$w_{3}$};
% Text Node
\draw (304.23,58.04) node [anchor=north west][inner sep=0.75pt]    {$w_{3}$};
% Text Node
\draw (474.31,93.34) node [anchor=north west][inner sep=0.75pt]    {$u_{1}$};
% Text Node
\draw (569.83,74.73) node [anchor=north west][inner sep=0.75pt]    {$u_{1}$};
% Text Node
\draw (525.78,97.17) node [anchor=north west][inner sep=0.75pt]    {$u_{2}$};
% Text Node
\draw (434.27,59.41) node [anchor=north west][inner sep=0.75pt]    {$u_{2}$};
% Text Node
\draw (517.78,152.99) node [anchor=north west][inner sep=0.75pt]    {$u_{3}$};
% Text Node
\draw (489.75,53.93) node [anchor=north west][inner sep=0.75pt]    {$u_{3}$};
% Text Node
\draw (202.85,268.46) node [anchor=north west][inner sep=0.75pt]    {$s_{1}$};
% Text Node
\draw (284.06,222.49) node [anchor=north west][inner sep=0.75pt]    {$s_{1}$};
% Text Node
\draw (252.03,268.46) node [anchor=north west][inner sep=0.75pt]    {$s_{2}$};
% Text Node
\draw (169.1,222.49) node [anchor=north west][inner sep=0.75pt]    {$s_{2}$};
% Text Node
\draw (219.43,226.87) node [anchor=north west][inner sep=0.75pt]    {$s_{3}$};
% Text Node
\draw (245.17,325.38) node [anchor=north west][inner sep=0.75pt]    {$s_{3}$};
% Text Node
\draw (391.66,268.19) node [anchor=north west][inner sep=0.75pt]    {$t_{1}$};
% Text Node
\draw (487.18,243.56) node [anchor=north west][inner sep=0.75pt]    {$t_{1}$};
% Text Node
\draw (447.14,268.74) node [anchor=north west][inner sep=0.75pt]    {$t_{2}$};
% Text Node
\draw (346.48,255.06) node [anchor=north west][inner sep=0.75pt]    {$t_{2}$};
% Text Node
\draw (431.7,327.3) node [anchor=north west][inner sep=0.75pt]    {$t_{3}$};
% Text Node
\draw (411.68,237) node [anchor=north west][inner sep=0.75pt]    {$t_{3}$};


\end{tikzpicture}
    \caption{Triangulación topológica ideal del nudo $6_1$}
\end{figure}
Con esto ya podemos empezar a pensar primero en ecuaciones de pegado, Observe que dotamos a cada tetraedro de variables $z_i,w_i,u_i,s_i$ y $t_i$, con $i=1,2,3$, al igual que en el caso del nudo de ocho primero solo mantenemos el invariante de arista para aristas opuestas, la ecuaciones de pegado por arista son
\begin{itemize}
    \item \textbf{Arista 1:} $z_2z_3w_3u_2u_3s_2^2s_3=1.$
    \item \textbf{Arista 4:} $z_1z_3w_3u_1u_3s_1^2t_1t_3^2=1.$
    \item \textbf{Arista 7:} $z_1w_2^2u_1s_3t_1=1.$
    \item \textbf{Arista 8:} $z_2w_1t_2=1.$
    \item \textbf{Arista 9:} $w_1u_2t_2=1.$
\end{itemize}
Haciendo las sustituciones $z_1=z,w_1=w,u_1=u,s_1=s$ y $t_1=t$ y usando el lema 4.1 obtenemos el siguiente sistema de ecuaciones no lineal
$$\begin{cases}
    \dfrac{1-w}{zu(1-s)sw}=1\\
    \\
    \dfrac{(z-1)(w-1)(u-1)s^2(t-1)^2}{wt}=1\\
    \\
    \dfrac{zu(s-1)t}{(1-w)^2s}=1\\
    \\
    \dfrac{w}{(1-z)(1-t)}=1\\
    \\
    \dfrac{w}{(1-u)(1-t)}=1
\end{cases}$$
Algunas condiciones son fáciles de deducir, como por ejemplo que $z=u$, pero resolver este sistema e incluso otros sobrepasa la capacidad de calculo en general cuando se trata de determinar la parametrizacion de soluciones, por lo que  en general lo mejor es usar software para solucionar. Cuando se trata de las ecuaciones de completez pasa algo incluso mas complicado y es al momento de truncar los tetraedros en este caso se generan 20 triángulos que se necesitan pegar de manera adecuada siguiendo las identificaciones de las caras, mas esta tarea se salio de mis manos y entra en juego la ayuda de herramientas computacionales.

\subsection{Ecuaciones de pegado y completes para el enlace de Whitehead}

Al igual que antes ya teníamos dos poliedros ideales (pirámides) que reflejaban el complemento del enlace de Whitehead, en este caso la división en tetraedros es casi inmediata, pero seguiremos la idea del anterior ejemplo por consistencia. Primero dividimos en ambos poliedros la cara $G$ por medio de una arista extra
\begin{figure}[H]
\centering
%!TEX root = ./main.tex



\tikzset{every picture/.style={line width=0.75pt}} %set default line width to 0.75pt        

\begin{tikzpicture}[x=0.75pt,y=0.75pt,yscale=-1,xscale=1]
%uncomment if require: \path (0,498); %set diagram left start at 0, and has height of 498
\path[use as bounding box] (20,10) rectangle (660,300);
%Shape: Circle [id:dp8563981245220138] 
\draw  [line width=1.5]  (109.32,66.82) .. controls (109.5,63.88) and (112.03,61.64) .. (114.97,61.82) .. controls (117.91,62.01) and (120.15,64.54) .. (119.96,67.48) .. controls (119.78,70.42) and (117.25,72.65) .. (114.31,72.47) .. controls (111.37,72.29) and (109.14,69.76) .. (109.32,66.82) -- cycle ;
%Shape: Circle [id:dp3520905591144071] 
\draw  [line width=1.5]  (175.47,131.04) .. controls (175.65,128.1) and (178.18,125.86) .. (181.12,126.04) .. controls (184.06,126.23) and (186.3,128.76) .. (186.12,131.7) .. controls (185.93,134.64) and (183.4,136.87) .. (180.46,136.69) .. controls (177.52,136.51) and (175.29,133.98) .. (175.47,131.04) -- cycle ;
%Straight Lines [id:da29309615585626436] 
\draw [line width=1.5]    (121.25,73.55) -- (175.65,128.04) ;
\draw [shift={(118.43,70.72)}, rotate = 45.05] [fill={rgb, 255:red, 0; green, 0; blue, 0 }  ][line width=0.08]  [draw opacity=0] (11.61,-5.58) -- (0,0) -- (11.61,5.58) -- cycle    ;
%Shape: Circle [id:dp7112383425029831] 
\draw  [line width=1.5]  (247.1,65.43) .. controls (250.04,65.61) and (252.27,68.15) .. (252.08,71.09) .. controls (251.9,74.03) and (249.36,76.26) .. (246.42,76.07) .. controls (243.49,75.88) and (241.25,73.35) .. (241.44,70.41) .. controls (241.63,67.47) and (244.16,65.24) .. (247.1,65.43) -- cycle ;
%Straight Lines [id:da41995557597045774] 
\draw [line width=1.5]    (240.35,77.35) -- (185.78,131.68) ;
\draw [shift={(243.18,74.53)}, rotate = 135.13] [fill={rgb, 255:red, 0; green, 0; blue, 0 }  ][line width=0.08]  [draw opacity=0] (11.61,-5.58) -- (0,0) -- (11.61,5.58) -- cycle    ;
%Shape: Circle [id:dp5575025818197289] 
\draw  [line width=1.5]  (117.75,201.6) .. controls (114.81,201.47) and (112.53,198.97) .. (112.66,196.03) .. controls (112.79,193.09) and (115.28,190.81) .. (118.22,190.94) .. controls (121.17,191.07) and (123.45,193.57) .. (123.31,196.51) .. controls (123.18,199.45) and (120.69,201.73) .. (117.75,201.6) -- cycle ;
%Straight Lines [id:da4725169852169093] 
\draw [line width=1.5]    (124.28,189.55) -- (177.82,134.21) ;
\draw [shift={(121.49,192.42)}, rotate = 314.06] [fill={rgb, 255:red, 0; green, 0; blue, 0 }  ][line width=0.08]  [draw opacity=0] (11.61,-5.58) -- (0,0) -- (11.61,5.58) -- cycle    ;
%Shape: Circle [id:dp9549973004430545] 
\draw  [line width=1.5]  (251.2,197.58) .. controls (250.98,200.52) and (248.43,202.72) .. (245.49,202.51) .. controls (242.55,202.29) and (240.34,199.74) .. (240.56,196.8) .. controls (240.77,193.86) and (243.33,191.66) .. (246.27,191.87) .. controls (249.2,192.09) and (251.41,194.64) .. (251.2,197.58) -- cycle ;
%Straight Lines [id:da14548168825025642] 
\draw [line width=1.5]    (239.34,190.71) -- (185.54,135.63) ;
\draw [shift={(242.13,193.57)}, rotate = 225.68] [fill={rgb, 255:red, 0; green, 0; blue, 0 }  ][line width=0.08]  [draw opacity=0] (11.61,-5.58) -- (0,0) -- (11.61,5.58) -- cycle    ;
%Straight Lines [id:da04138952431338294] 
\draw [line width=1.5]    (123.31,196.51) -- (236.56,196.79) ;
\draw [shift={(240.56,196.8)}, rotate = 180.14] [fill={rgb, 255:red, 0; green, 0; blue, 0 }  ][line width=0.08]  [draw opacity=0] (11.61,-5.58) -- (0,0) -- (11.61,5.58) -- cycle    ;
%Straight Lines [id:da6360614459550309] 
\draw [line width=1.5]    (114.44,76.47) -- (118.22,190.94) ;
\draw [shift={(114.31,72.47)}, rotate = 88.11] [fill={rgb, 255:red, 0; green, 0; blue, 0 }  ][line width=0.08]  [draw opacity=0] (11.61,-5.58) -- (0,0) -- (11.61,5.58) -- cycle    ;
%Straight Lines [id:da22900843394258874] 
\draw [line width=1.5]    (123.96,67.57) -- (241.44,70.41) ;
\draw [shift={(119.96,67.48)}, rotate = 1.38] [fill={rgb, 255:red, 0; green, 0; blue, 0 }  ][line width=0.08]  [draw opacity=0] (11.61,-5.58) -- (0,0) -- (11.61,5.58) -- cycle    ;
%Straight Lines [id:da9297989788033489] 
\draw [line width=1.5]    (246.42,76.07) -- (246.27,187.87) ;
\draw [shift={(246.27,191.87)}, rotate = 270.08] [fill={rgb, 255:red, 0; green, 0; blue, 0 }  ][line width=0.08]  [draw opacity=0] (11.61,-5.58) -- (0,0) -- (11.61,5.58) -- cycle    ;
%Shape: Circle [id:dp6983363055718416] 
\draw  [line width=1.5]  (389.99,69.48) .. controls (390.17,66.54) and (392.7,64.31) .. (395.64,64.49) .. controls (398.58,64.67) and (400.81,67.2) .. (400.63,70.14) .. controls (400.45,73.08) and (397.92,75.32) .. (394.98,75.14) .. controls (392.04,74.96) and (389.8,72.42) .. (389.99,69.48) -- cycle ;
%Shape: Circle [id:dp34637775441133845] 
\draw  [line width=1.5]  (456.14,133.7) .. controls (456.32,130.76) and (458.85,128.53) .. (461.79,128.71) .. controls (464.73,128.89) and (466.96,131.42) .. (466.78,134.36) .. controls (466.6,137.3) and (464.07,139.54) .. (461.13,139.36) .. controls (458.19,139.17) and (455.95,136.64) .. (456.14,133.7) -- cycle ;
%Straight Lines [id:da32489783251721505] 
\draw [line width=1.5]    (399.09,73.39) -- (453.5,127.88) ;
\draw [shift={(456.32,130.71)}, rotate = 225.05] [fill={rgb, 255:red, 0; green, 0; blue, 0 }  ][line width=0.08]  [draw opacity=0] (11.61,-5.58) -- (0,0) -- (11.61,5.58) -- cycle    ;
%Shape: Circle [id:dp04140762256709696] 
\draw  [line width=1.5]  (527.77,68.09) .. controls (530.71,68.28) and (532.94,70.81) .. (532.75,73.75) .. controls (532.57,76.69) and (530.03,78.92) .. (527.09,78.74) .. controls (524.15,78.55) and (521.92,76.02) .. (522.11,73.08) .. controls (522.29,70.14) and (524.83,67.91) .. (527.77,68.09) -- cycle ;
%Straight Lines [id:da0022087921748601413] 
\draw [line width=1.5]    (521.01,80.02) -- (466.45,134.34) ;
\draw [shift={(523.85,77.2)}, rotate = 135.13] [fill={rgb, 255:red, 0; green, 0; blue, 0 }  ][line width=0.08]  [draw opacity=0] (11.61,-5.58) -- (0,0) -- (11.61,5.58) -- cycle    ;
%Shape: Circle [id:dp13672754996963388] 
\draw  [line width=1.5]  (398.41,204.26) .. controls (395.47,204.13) and (393.19,201.64) .. (393.32,198.7) .. controls (393.46,195.76) and (395.95,193.48) .. (398.89,193.61) .. controls (401.83,193.74) and (404.11,196.23) .. (403.98,199.17) .. controls (403.85,202.12) and (401.36,204.4) .. (398.41,204.26) -- cycle ;
%Straight Lines [id:da052100170009540814] 
\draw [line width=1.5]    (404.94,192.21) -- (458.48,136.88) ;
\draw [shift={(402.16,195.09)}, rotate = 314.06] [fill={rgb, 255:red, 0; green, 0; blue, 0 }  ][line width=0.08]  [draw opacity=0] (11.61,-5.58) -- (0,0) -- (11.61,5.58) -- cycle    ;
%Shape: Circle [id:dp8069072010758881] 
\draw  [line width=1.5]  (531.86,200.25) .. controls (531.65,203.18) and (529.09,205.39) .. (526.16,205.18) .. controls (523.22,204.96) and (521.01,202.41) .. (521.22,199.47) .. controls (521.44,196.53) and (523.99,194.32) .. (526.93,194.54) .. controls (529.87,194.75) and (532.08,197.31) .. (531.86,200.25) -- cycle ;
%Straight Lines [id:da801530605139214] 
\draw [line width=1.5]    (522.8,196.24) -- (469,141.16) ;
\draw [shift={(466.2,138.29)}, rotate = 45.68] [fill={rgb, 255:red, 0; green, 0; blue, 0 }  ][line width=0.08]  [draw opacity=0] (11.61,-5.58) -- (0,0) -- (11.61,5.58) -- cycle    ;
%Straight Lines [id:da7355128204598104] 
\draw [line width=1.5]    (407.98,199.18) -- (521.22,199.47) ;
\draw [shift={(403.98,199.17)}, rotate = 0.14] [fill={rgb, 255:red, 0; green, 0; blue, 0 }  ][line width=0.08]  [draw opacity=0] (11.61,-5.58) -- (0,0) -- (11.61,5.58) -- cycle    ;
%Straight Lines [id:da7937596572289117] 
\draw [line width=1.5]    (394.98,75.14) -- (398.76,189.61) ;
\draw [shift={(398.89,193.61)}, rotate = 268.11] [fill={rgb, 255:red, 0; green, 0; blue, 0 }  ][line width=0.08]  [draw opacity=0] (11.61,-5.58) -- (0,0) -- (11.61,5.58) -- cycle    ;
%Straight Lines [id:da2955657521942342] 
\draw [line width=1.5]    (400.63,70.14) -- (518.11,72.98) ;
\draw [shift={(522.11,73.08)}, rotate = 181.38] [fill={rgb, 255:red, 0; green, 0; blue, 0 }  ][line width=0.08]  [draw opacity=0] (11.61,-5.58) -- (0,0) -- (11.61,5.58) -- cycle    ;
%Straight Lines [id:da5139882653554566] 
\draw [line width=1.5]    (527.09,82.74) -- (526.93,194.54) ;
\draw [shift={(527.09,78.74)}, rotate = 90.08] [fill={rgb, 255:red, 0; green, 0; blue, 0 }  ][line width=0.08]  [draw opacity=0] (11.61,-5.58) -- (0,0) -- (11.61,5.58) -- cycle    ;
%Curve Lines [id:da6493881684844067] 
\draw [line width=1.5]    (114.97,61.82) .. controls (-87.98,189.03) and (169.07,357.97) .. (244.37,204.84) ;
\draw [shift={(245.49,202.51)}, rotate = 115.16] [fill={rgb, 255:red, 0; green, 0; blue, 0 }  ][line width=0.08]  [draw opacity=0] (11.61,-5.58) -- (0,0) -- (11.61,5.58) -- cycle    ;
%Curve Lines [id:da7408599747443351] 
\draw [line width=1.5]    (398.41,204.26) .. controls (582.74,396.7) and (670.78,118.47) .. (529.91,68.82) ;
\draw [shift={(527.77,68.09)}, rotate = 18.29] [fill={rgb, 255:red, 0; green, 0; blue, 0 }  ][line width=0.08]  [draw opacity=0] (11.61,-5.58) -- (0,0) -- (11.61,5.58) -- cycle    ;

% Text Node
\draw (174.67,86.07) node [anchor=north west][inner sep=0.75pt]    {$A$};
% Text Node
\draw (452,86.73) node [anchor=north west][inner sep=0.75pt]    {$A$};
% Text Node
\draw (135.33,126.07) node [anchor=north west][inner sep=0.75pt]    {$B$};
% Text Node
\draw (414,127.4) node [anchor=north west][inner sep=0.75pt]    {$B$};
% Text Node
\draw (216,122.07) node [anchor=north west][inner sep=0.75pt]    {$C$};
% Text Node
\draw (496.67,126.07) node [anchor=north west][inner sep=0.75pt]    {$C$};
% Text Node
\draw (173.33,165.4) node [anchor=north west][inner sep=0.75pt]    {$D$};
% Text Node
\draw (453.33,163.4) node [anchor=north west][inner sep=0.75pt]    {$D$};
% Text Node
\draw (100.67,129.4) node [anchor=north west][inner sep=0.75pt]    {$1$};
% Text Node
\draw (135.33,97.4) node [anchor=north west][inner sep=0.75pt]    {$1$};
% Text Node
\draw (172,198.73) node [anchor=north west][inner sep=0.75pt]    {$1$};
% Text Node
\draw (178.67,49.73) node [anchor=north west][inner sep=0.75pt]    {$3$};
% Text Node
\draw (251.33,122.4) node [anchor=north west][inner sep=0.75pt]    {$3$};
% Text Node
\draw (216.67,149.73) node [anchor=north west][inner sep=0.75pt]    {$3$};
% Text Node
\draw (135.33,149.73) node [anchor=north west][inner sep=0.75pt]    {$2$};
% Text Node
\draw (200,86.4) node [anchor=north west][inner sep=0.75pt]    {$2$};
% Text Node
\draw (378.67,125.4) node [anchor=north west][inner sep=0.75pt]    {$1$};
% Text Node
\draw (416.67,152.73) node [anchor=north west][inner sep=0.75pt]    {$1$};
% Text Node
\draw (441.33,52.73) node [anchor=north west][inner sep=0.75pt]    {$1$};
% Text Node
\draw (531.33,129.07) node [anchor=north west][inner sep=0.75pt]    {$3$};
% Text Node
\draw (481.33,92.4) node [anchor=north west][inner sep=0.75pt]    {$3$};
% Text Node
\draw (460.67,205.07) node [anchor=north west][inner sep=0.75pt]    {$3$};
% Text Node
\draw (415.33,103.07) node [anchor=north west][inner sep=0.75pt]    {$2$};
% Text Node
\draw (496,149.73) node [anchor=north west][inner sep=0.75pt]    {$2$};
% Text Node
\draw (71.33,206.73) node [anchor=north west][inner sep=0.75pt]    {$G_{1}$};
% Text Node
\draw (221.33,247.4) node [anchor=north west][inner sep=0.75pt]    {$G_{2}$};
% Text Node
\draw (103.33,260.07) node [anchor=north west][inner sep=0.75pt]    {$5$};
% Text Node
\draw (537.33,221.4) node [anchor=north west][inner sep=0.75pt]    {$G_{2}$};
% Text Node
\draw (400.67,249.4) node [anchor=north west][inner sep=0.75pt]    {$G_{1}$};
% Text Node
\draw (508.67,277.4) node [anchor=north west][inner sep=0.75pt]    {$5$};


\end{tikzpicture}
    \caption{División cara $G$ en los poliedros}
\end{figure}
Nuevamente es incluso mas facil que antes, basta con que en el poliedro de arriba se separen las caras $B,D,G_1$ del diagrama y se agregue una nueva cara $H_1$ y similarmente en el de abajo separar $C,D,G_2$ y agregar una nueva cara $H_2$ como se muestra a continuación
\begin{figure}[H]
\centering
%!TEX root = ./main.tex


\tikzset{every picture/.style={line width=0.75pt}} %set default line width to 0.75pt        

\begin{tikzpicture}[x=0.75pt,y=0.75pt,yscale=-1,xscale=1]
%uncomment if require: \path (0,498); %set diagram left start at 0, and has height of 498
\path[use as bounding box] (20,10) rectangle (660,500);
%Shape: Circle [id:dp8563981245220138] 
\draw  [line width=1.5]  (108.65,33.48) .. controls (108.83,30.54) and (111.37,28.31) .. (114.31,28.49) .. controls (117.24,28.67) and (119.48,31.2) .. (119.3,34.14) .. controls (119.12,37.08) and (116.58,39.32) .. (113.64,39.14) .. controls (110.7,38.96) and (108.47,36.42) .. (108.65,33.48) -- cycle ;
%Shape: Circle [id:dp3520905591144071] 
\draw  [line width=1.5]  (174.8,97.7) .. controls (174.98,94.76) and (177.52,92.53) .. (180.46,92.71) .. controls (183.4,92.89) and (185.63,95.42) .. (185.45,98.36) .. controls (185.27,101.3) and (182.74,103.54) .. (179.8,103.36) .. controls (176.86,103.17) and (174.62,100.64) .. (174.8,97.7) -- cycle ;
%Straight Lines [id:da29309615585626436] 
\draw [line width=1.5]    (120.59,40.22) -- (174.99,94.71) ;
\draw [shift={(117.76,37.39)}, rotate = 45.05] [fill={rgb, 255:red, 0; green, 0; blue, 0 }  ][line width=0.08]  [draw opacity=0] (11.61,-5.58) -- (0,0) -- (11.61,5.58) -- cycle    ;
%Shape: Circle [id:dp7112383425029831] 
\draw  [line width=1.5]  (246.43,32.09) .. controls (249.37,32.28) and (251.6,34.81) .. (251.42,37.75) .. controls (251.23,40.69) and (248.7,42.92) .. (245.76,42.74) .. controls (242.82,42.55) and (240.59,40.02) .. (240.77,37.08) .. controls (240.96,34.14) and (243.49,31.91) .. (246.43,32.09) -- cycle ;
%Straight Lines [id:da41995557597045774] 
\draw [line width=1.5]    (239.68,44.02) -- (185.12,98.34) ;
\draw [shift={(242.52,41.2)}, rotate = 135.13] [fill={rgb, 255:red, 0; green, 0; blue, 0 }  ][line width=0.08]  [draw opacity=0] (11.61,-5.58) -- (0,0) -- (11.61,5.58) -- cycle    ;
%Shape: Circle [id:dp9549973004430545] 
\draw  [line width=1.5]  (250.53,164.25) .. controls (250.32,167.18) and (247.76,169.39) .. (244.82,169.18) .. controls (241.88,168.96) and (239.68,166.41) .. (239.89,163.47) .. controls (240.11,160.53) and (242.66,158.32) .. (245.6,158.54) .. controls (248.54,158.75) and (250.74,161.31) .. (250.53,164.25) -- cycle ;
%Straight Lines [id:da14548168825025642] 
\draw [line width=1.5]    (238.67,157.38) -- (184.87,102.29) ;
\draw [shift={(241.46,160.24)}, rotate = 225.68] [fill={rgb, 255:red, 0; green, 0; blue, 0 }  ][line width=0.08]  [draw opacity=0] (11.61,-5.58) -- (0,0) -- (11.61,5.58) -- cycle    ;
%Straight Lines [id:da22900843394258874] 
\draw [line width=1.5]    (123.3,34.24) -- (240.77,37.08) ;
\draw [shift={(119.3,34.14)}, rotate = 1.38] [fill={rgb, 255:red, 0; green, 0; blue, 0 }  ][line width=0.08]  [draw opacity=0] (11.61,-5.58) -- (0,0) -- (11.61,5.58) -- cycle    ;
%Straight Lines [id:da9297989788033489] 
\draw [line width=1.5]    (245.76,42.74) -- (245.6,154.54) ;
\draw [shift={(245.6,158.54)}, rotate = 270.08] [fill={rgb, 255:red, 0; green, 0; blue, 0 }  ][line width=0.08]  [draw opacity=0] (11.61,-5.58) -- (0,0) -- (11.61,5.58) -- cycle    ;
%Shape: Circle [id:dp6983363055718416] 
\draw  [line width=1.5]  (390.65,32.82) .. controls (390.83,29.88) and (393.37,27.64) .. (396.31,27.82) .. controls (399.24,28.01) and (401.48,30.54) .. (401.3,33.48) .. controls (401.12,36.42) and (398.58,38.65) .. (395.64,38.47) .. controls (392.7,38.29) and (390.47,35.76) .. (390.65,32.82) -- cycle ;
%Shape: Circle [id:dp34637775441133845] 
\draw  [line width=1.5]  (456.8,97.04) .. controls (456.98,94.1) and (459.52,91.86) .. (462.46,92.04) .. controls (465.4,92.23) and (467.63,94.76) .. (467.45,97.7) .. controls (467.27,100.64) and (464.74,102.87) .. (461.8,102.69) .. controls (458.86,102.51) and (456.62,99.98) .. (456.8,97.04) -- cycle ;
%Straight Lines [id:da32489783251721505] 
\draw [line width=1.5]    (399.76,36.72) -- (454.16,91.21) ;
\draw [shift={(456.99,94.04)}, rotate = 225.05] [fill={rgb, 255:red, 0; green, 0; blue, 0 }  ][line width=0.08]  [draw opacity=0] (11.61,-5.58) -- (0,0) -- (11.61,5.58) -- cycle    ;
%Shape: Circle [id:dp04140762256709696] 
\draw  [line width=1.5]  (528.43,31.43) .. controls (531.37,31.61) and (533.6,34.15) .. (533.42,37.09) .. controls (533.23,40.03) and (530.7,42.26) .. (527.76,42.07) .. controls (524.82,41.88) and (522.59,39.35) .. (522.77,36.41) .. controls (522.96,33.47) and (525.49,31.24) .. (528.43,31.43) -- cycle ;
%Straight Lines [id:da0022087921748601413] 
\draw [line width=1.5]    (521.68,43.35) -- (467.12,97.68) ;
\draw [shift={(524.52,40.53)}, rotate = 135.13] [fill={rgb, 255:red, 0; green, 0; blue, 0 }  ][line width=0.08]  [draw opacity=0] (11.61,-5.58) -- (0,0) -- (11.61,5.58) -- cycle    ;
%Shape: Circle [id:dp13672754996963388] 
\draw  [line width=1.5]  (399.08,167.6) .. controls (396.14,167.47) and (393.86,164.97) .. (393.99,162.03) .. controls (394.12,159.09) and (396.62,156.81) .. (399.56,156.94) .. controls (402.5,157.07) and (404.78,159.57) .. (404.65,162.51) .. controls (404.52,165.45) and (402.02,167.73) .. (399.08,167.6) -- cycle ;
%Straight Lines [id:da052100170009540814] 
\draw [line width=1.5]    (405.61,155.55) -- (459.15,100.21) ;
\draw [shift={(402.83,158.42)}, rotate = 314.06] [fill={rgb, 255:red, 0; green, 0; blue, 0 }  ][line width=0.08]  [draw opacity=0] (11.61,-5.58) -- (0,0) -- (11.61,5.58) -- cycle    ;
%Straight Lines [id:da7937596572289117] 
\draw [line width=1.5]    (395.64,38.47) -- (399.43,152.94) ;
\draw [shift={(399.56,156.94)}, rotate = 268.11] [fill={rgb, 255:red, 0; green, 0; blue, 0 }  ][line width=0.08]  [draw opacity=0] (11.61,-5.58) -- (0,0) -- (11.61,5.58) -- cycle    ;
%Straight Lines [id:da2955657521942342] 
\draw [line width=1.5]    (401.3,33.48) -- (518.77,36.31) ;
\draw [shift={(522.77,36.41)}, rotate = 181.38] [fill={rgb, 255:red, 0; green, 0; blue, 0 }  ][line width=0.08]  [draw opacity=0] (11.61,-5.58) -- (0,0) -- (11.61,5.58) -- cycle    ;
%Curve Lines [id:da6493881684844067] 
\draw [line width=1.5]    (114.31,28.49) .. controls (-88.65,155.69) and (168.4,324.63) .. (243.7,171.51) ;
\draw [shift={(244.82,169.18)}, rotate = 115.16] [fill={rgb, 255:red, 0; green, 0; blue, 0 }  ][line width=0.08]  [draw opacity=0] (11.61,-5.58) -- (0,0) -- (11.61,5.58) -- cycle    ;
%Curve Lines [id:da7408599747443351] 
\draw [line width=1.5]    (399.08,167.6) .. controls (583.41,360.03) and (671.45,81.81) .. (530.57,32.16) ;
\draw [shift={(528.43,31.43)}, rotate = 18.29] [fill={rgb, 255:red, 0; green, 0; blue, 0 }  ][line width=0.08]  [draw opacity=0] (11.61,-5.58) -- (0,0) -- (11.61,5.58) -- cycle    ;
%Shape: Circle [id:dp7329325513471445] 
\draw  [line width=1.5]  (200.8,324.37) .. controls (200.98,321.43) and (203.52,319.19) .. (206.46,319.38) .. controls (209.4,319.56) and (211.63,322.09) .. (211.45,325.03) .. controls (211.27,327.97) and (208.74,330.21) .. (205.8,330.02) .. controls (202.86,329.84) and (200.62,327.31) .. (200.8,324.37) -- cycle ;
%Straight Lines [id:da24734427845057683] 
\draw [line width=1.5]    (146.59,266.89) -- (200.99,321.38) ;
\draw [shift={(143.76,264.06)}, rotate = 45.05] [fill={rgb, 255:red, 0; green, 0; blue, 0 }  ][line width=0.08]  [draw opacity=0] (11.61,-5.58) -- (0,0) -- (11.61,5.58) -- cycle    ;
%Shape: Circle [id:dp04065717725850437] 
\draw  [line width=1.5]  (143.08,394.93) .. controls (140.14,394.8) and (137.86,392.31) .. (137.99,389.36) .. controls (138.12,386.42) and (140.62,384.14) .. (143.56,384.27) .. controls (146.5,384.41) and (148.78,386.9) .. (148.65,389.84) .. controls (148.52,392.78) and (146.02,395.06) .. (143.08,394.93) -- cycle ;
%Straight Lines [id:da30678934938564406] 
\draw [line width=1.5]    (149.61,382.88) -- (203.15,327.55) ;
\draw [shift={(146.83,385.76)}, rotate = 314.06] [fill={rgb, 255:red, 0; green, 0; blue, 0 }  ][line width=0.08]  [draw opacity=0] (11.61,-5.58) -- (0,0) -- (11.61,5.58) -- cycle    ;
%Straight Lines [id:da749898484078513] 
\draw [line width=1.5]    (264.67,384.05) -- (210.87,328.96) ;
\draw [shift={(267.46,386.91)}, rotate = 225.68] [fill={rgb, 255:red, 0; green, 0; blue, 0 }  ][line width=0.08]  [draw opacity=0] (11.61,-5.58) -- (0,0) -- (11.61,5.58) -- cycle    ;
%Straight Lines [id:da9107200542859878] 
\draw [line width=1.5]    (148.65,389.84) -- (261.89,390.12) ;
\draw [shift={(265.89,390.13)}, rotate = 180.14] [fill={rgb, 255:red, 0; green, 0; blue, 0 }  ][line width=0.08]  [draw opacity=0] (11.61,-5.58) -- (0,0) -- (11.61,5.58) -- cycle    ;
%Straight Lines [id:da592816778318295] 
\draw [line width=1.5]    (139.78,269.8) -- (143.56,384.27) ;
\draw [shift={(139.64,265.8)}, rotate = 88.11] [fill={rgb, 255:red, 0; green, 0; blue, 0 }  ][line width=0.08]  [draw opacity=0] (11.61,-5.58) -- (0,0) -- (11.61,5.58) -- cycle    ;
%Curve Lines [id:da002061286055452305] 
\draw [line width=1.5]    (140.31,255.16) .. controls (-62.65,382.36) and (194.4,551.3) .. (269.7,398.18) ;
\draw [shift={(270.82,395.84)}, rotate = 115.16] [fill={rgb, 255:red, 0; green, 0; blue, 0 }  ][line width=0.08]  [draw opacity=0] (11.61,-5.58) -- (0,0) -- (11.61,5.58) -- cycle    ;
%Shape: Circle [id:dp8757397345630595] 
\draw  [line width=1.5]  (134.65,260.15) .. controls (134.83,257.21) and (137.37,254.98) .. (140.31,255.16) .. controls (143.24,255.34) and (145.48,257.87) .. (145.3,260.81) .. controls (145.12,263.75) and (142.58,265.99) .. (139.64,265.8) .. controls (136.7,265.62) and (134.47,263.09) .. (134.65,260.15) -- cycle ;
%Shape: Circle [id:dp0002885352406026831] 
\draw  [line width=1.5]  (261.81,391.9) .. controls (261.99,388.96) and (264.52,386.73) .. (267.46,386.91) .. controls (270.4,387.09) and (272.64,389.62) .. (272.46,392.56) .. controls (272.27,395.5) and (269.74,397.74) .. (266.8,397.55) .. controls (263.86,397.37) and (261.63,394.84) .. (261.81,391.9) -- cycle ;
%Straight Lines [id:da0900476780775239] 
\draw [line width=1.5]    (406.28,384.88) -- (459.82,329.55) ;
\draw [shift={(403.49,387.76)}, rotate = 314.06] [fill={rgb, 255:red, 0; green, 0; blue, 0 }  ][line width=0.08]  [draw opacity=0] (11.61,-5.58) -- (0,0) -- (11.61,5.58) -- cycle    ;
%Shape: Circle [id:dp6293143216952294] 
\draw  [line width=1.5]  (533.2,392.91) .. controls (532.98,395.85) and (530.43,398.06) .. (527.49,397.84) .. controls (524.55,397.63) and (522.34,395.07) .. (522.56,392.13) .. controls (522.77,389.2) and (525.33,386.99) .. (528.27,387.2) .. controls (531.2,387.42) and (533.41,389.97) .. (533.2,392.91) -- cycle ;
%Straight Lines [id:da1297418655137339] 
\draw [line width=1.5]    (524.13,388.91) -- (470.33,333.82) ;
\draw [shift={(467.54,330.96)}, rotate = 45.68] [fill={rgb, 255:red, 0; green, 0; blue, 0 }  ][line width=0.08]  [draw opacity=0] (11.61,-5.58) -- (0,0) -- (11.61,5.58) -- cycle    ;
%Straight Lines [id:da3079032519134891] 
\draw [line width=1.5]    (409.31,391.85) -- (522.56,392.13) ;
\draw [shift={(405.31,391.84)}, rotate = 0.14] [fill={rgb, 255:red, 0; green, 0; blue, 0 }  ][line width=0.08]  [draw opacity=0] (11.61,-5.58) -- (0,0) -- (11.61,5.58) -- cycle    ;
%Straight Lines [id:da9761283042486899] 
\draw [line width=1.5]    (528.42,275.4) -- (528.27,387.2) ;
\draw [shift={(528.42,271.4)}, rotate = 90.08] [fill={rgb, 255:red, 0; green, 0; blue, 0 }  ][line width=0.08]  [draw opacity=0] (11.61,-5.58) -- (0,0) -- (11.61,5.58) -- cycle    ;
%Curve Lines [id:da31523783789874726] 
\draw [line width=1.5]    (399.75,396.93) .. controls (584.07,589.37) and (672.12,311.14) .. (531.24,261.49) ;
\draw [shift={(529.1,260.76)}, rotate = 18.29] [fill={rgb, 255:red, 0; green, 0; blue, 0 }  ][line width=0.08]  [draw opacity=0] (11.61,-5.58) -- (0,0) -- (11.61,5.58) -- cycle    ;
%Shape: Circle [id:dp14865181765242919] 
\draw  [line width=1.5]  (403.02,398.41) .. controls (400.08,398.28) and (397.8,395.79) .. (397.93,392.85) .. controls (398.06,389.9) and (400.55,387.62) .. (403.49,387.76) .. controls (406.44,387.89) and (408.72,390.38) .. (408.58,393.32) .. controls (408.45,396.26) and (405.96,398.54) .. (403.02,398.41) -- cycle ;
%Shape: Circle [id:dp8864568871703659] 
\draw  [line width=1.5]  (464.91,335.11) .. controls (461.96,334.98) and (459.69,332.49) .. (459.82,329.55) .. controls (459.95,326.6) and (462.44,324.33) .. (465.38,324.46) .. controls (468.33,324.59) and (470.61,327.08) .. (470.47,330.02) .. controls (470.34,332.97) and (467.85,335.24) .. (464.91,335.11) -- cycle ;
%Shape: Circle [id:dp10976801910100586] 
\draw  [line width=1.5]  (528.42,271.4) .. controls (525.48,271.27) and (523.2,268.78) .. (523.33,265.84) .. controls (523.47,262.9) and (525.96,260.62) .. (528.9,260.75) .. controls (531.84,260.88) and (534.12,263.37) .. (533.99,266.31) .. controls (533.86,269.26) and (531.37,271.54) .. (528.42,271.4) -- cycle ;
%Straight Lines [id:da7063896897703602] 
\draw [line width=1.5]    (522.35,271.35) -- (467.78,325.68) ;
\draw [shift={(525.18,268.53)}, rotate = 135.13] [fill={rgb, 255:red, 0; green, 0; blue, 0 }  ][line width=0.08]  [draw opacity=0] (11.61,-5.58) -- (0,0) -- (11.61,5.58) -- cycle    ;

% Text Node
\draw (174,52.73) node [anchor=north west][inner sep=0.75pt]    {$A$};
% Text Node
\draw (452.67,50.07) node [anchor=north west][inner sep=0.75pt]    {$A$};
% Text Node
\draw (414.67,90.73) node [anchor=north west][inner sep=0.75pt]    {$B$};
% Text Node
\draw (215.33,88.73) node [anchor=north west][inner sep=0.75pt]    {$C$};
% Text Node
\draw (134.67,64.07) node [anchor=north west][inner sep=0.75pt]    {$1$};
% Text Node
\draw (178,16.4) node [anchor=north west][inner sep=0.75pt]    {$3$};
% Text Node
\draw (250.67,89.07) node [anchor=north west][inner sep=0.75pt]    {$3$};
% Text Node
\draw (196,129.73) node [anchor=north west][inner sep=0.75pt]    {$3$};
% Text Node
\draw (199.33,53.07) node [anchor=north west][inner sep=0.75pt]    {$2$};
% Text Node
\draw (379.33,88.73) node [anchor=north west][inner sep=0.75pt]    {$1$};
% Text Node
\draw (437.33,124.07) node [anchor=north west][inner sep=0.75pt]    {$1$};
% Text Node
\draw (442,16.07) node [anchor=north west][inner sep=0.75pt]    {$1$};
% Text Node
\draw (498,69.07) node [anchor=north west][inner sep=0.75pt]    {$3$};
% Text Node
\draw (416,66.4) node [anchor=north west][inner sep=0.75pt]    {$2$};
% Text Node
\draw (220.67,214.07) node [anchor=north west][inner sep=0.75pt]    {$G_{2}$};
% Text Node
\draw (102.67,226.73) node [anchor=north west][inner sep=0.75pt]    {$5$};
% Text Node
\draw (401.33,212.73) node [anchor=north west][inner sep=0.75pt]    {$G_{1}$};
% Text Node
\draw (509.33,240.73) node [anchor=north west][inner sep=0.75pt]    {$5$};
% Text Node
\draw (160.67,319.4) node [anchor=north west][inner sep=0.75pt]    {$B$};
% Text Node
\draw (198.67,358.73) node [anchor=north west][inner sep=0.75pt]    {$D$};
% Text Node
\draw (126,322.73) node [anchor=north west][inner sep=0.75pt]    {$1$};
% Text Node
\draw (178.67,279.4) node [anchor=north west][inner sep=0.75pt]    {$1$};
% Text Node
\draw (197.33,392.07) node [anchor=north west][inner sep=0.75pt]    {$1$};
% Text Node
\draw (160.67,343.07) node [anchor=north west][inner sep=0.75pt]    {$2$};
% Text Node
\draw (96.67,400.07) node [anchor=north west][inner sep=0.75pt]    {$G_{1}$};
% Text Node
\draw (128.67,453.4) node [anchor=north west][inner sep=0.75pt]    {$5$};
% Text Node
\draw (239.33,338.4) node [anchor=north west][inner sep=0.75pt]    {$3$};
% Text Node
\draw (218.67,304.07) node [anchor=north west][inner sep=0.75pt]    {$H_{1}$};
% Text Node
\draw (124,137.4) node [anchor=north west][inner sep=0.75pt]    {$H_{1}$};
% Text Node
\draw (454.67,356.07) node [anchor=north west][inner sep=0.75pt]    {$D$};
% Text Node
\draw (418,345.4) node [anchor=north west][inner sep=0.75pt]    {$1$};
% Text Node
\draw (462,397.73) node [anchor=north west][inner sep=0.75pt]    {$3$};
% Text Node
\draw (497.33,342.4) node [anchor=north west][inner sep=0.75pt]    {$2$};
% Text Node
\draw (538.67,414.07) node [anchor=north west][inner sep=0.75pt]    {$G_{2}$};
% Text Node
\draw (498,317.4) node [anchor=north west][inner sep=0.75pt]    {$C$};
% Text Node
\draw (482.67,283.73) node [anchor=north west][inner sep=0.75pt]    {$3$};
% Text Node
\draw (532.67,329.07) node [anchor=north west][inner sep=0.75pt]    {$3$};
% Text Node
\draw (425.33,296.73) node [anchor=north west][inner sep=0.75pt]    {$H_{2}$};
% Text Node
\draw (504,137.4) node [anchor=north west][inner sep=0.75pt]    {$H_{2}$};


\end{tikzpicture}
    \caption{Poliedros divididos, arriba a la izquierda y abajo a la derecha}
\end{figure}
Reorganizando las aristas para obtener el dibujo estándar y cambiando la vista del poliedro de arriba al exterior, obtenemos la triangulación topológica ideal
\begin{figure}[H]
\centering
\input{whiteheadtetra}
    \caption{Triangulación topológica ideal de el enlace de Whitehead}
\end{figure}
Observe que como ha sido costumbre ya dotamos de los respectivos invariantes de aristas a los diagramas, dotando a las opuestas por el lema 4.1 de la misma variable, así por las ecuaciones de pegado tenemos que
\begin{itemize}
    \item \textbf{Arista 1:} $z_1w_1w_2^2u_1u_2^2t_1=1.$
    \item \textbf{Arista 2:} $z_3w_3u_3t_3=1.$
    \item \textbf{Arista 3:} $z_1z_2^2w_1u_1t_1t_2^2=1$
    \item \textbf{Arista 5:} $z_3w_3u_3t_3=1.$
\end{itemize}
Observe que una ecuación se repite, luego por el lema 4.1 y tomando las sustituciones $z_1=z,w_1=w,u_1=u$ y $t_1=t$ obtenemos las ecuaciones de pegado
$$\begin{cases}
    \dfrac{zwut}{(1-w)^2(1-u)^2}=1\\
    \\
    \dfrac{(z-1)(w-1)(u-1)(t-1)}{zwut}=1\\
    \\
    \dfrac{zwut}{(1-z)^2(1-t)^2}=1
\end{cases}$$
Nuevamente hay deducciones que se pueden hacer solo mirando las ecuaciones y recordando la condición de la parte imaginaria positiva se puede llegar a que $(w-1)(u-1)=(z-1)(t-1)$, mas avanzar mas es complicado, por lo que haremos uso de las ecuaciones de completez. Para estas note que este es un enlace por lo que como son dos componentes deberíamos obtener dos cúspides y por tanto dos triangulaciones del toro, primero truncamos los vértices y damos nombres a las caras 
\begin{figure}[H]
\centering
%!TEX root = ./main.tex



\tikzset{every picture/.style={line width=0.75pt}} %set default line width to 0.75pt        

\begin{tikzpicture}[x=0.75pt,y=0.75pt,yscale=-1,xscale=1]
%uncomment if require: \path (0,498); %set diagram left start at 0, and has height of 498

%Shape: Circle [id:dp5729830365949617] 
\draw  [line width=1.5]  (161.81,25.33) .. controls (161.81,22.39) and (164.2,20) .. (167.14,20) .. controls (170.09,20) and (172.47,22.39) .. (172.47,25.33) .. controls (172.47,28.28) and (170.09,30.67) .. (167.14,30.67) .. controls (164.2,30.67) and (161.81,28.28) .. (161.81,25.33) -- cycle ;
%Shape: Circle [id:dp37612834718030397] 
\draw  [line width=1.5]  (162.47,115.33) .. controls (162.47,112.39) and (164.86,110) .. (167.81,110) .. controls (170.75,110) and (173.14,112.39) .. (173.14,115.33) .. controls (173.14,118.28) and (170.75,120.67) .. (167.81,120.67) .. controls (164.86,120.67) and (162.47,118.28) .. (162.47,115.33) -- cycle ;
%Straight Lines [id:da9158510271719965] 
\draw [line width=1.5]    (167.17,34.67) -- (167.81,110) ;
\draw [shift={(167.14,30.67)}, rotate = 89.52] [fill={rgb, 255:red, 0; green, 0; blue, 0 }  ][line width=0.08]  [draw opacity=0] (11.61,-5.58) -- (0,0) -- (11.61,5.58) -- cycle    ;
%Straight Lines [id:da691738049353374] 
\draw [line width=1.5]    (101.81,180) -- (237.14,180) ;
\draw [shift={(97.81,180)}, rotate = 0] [fill={rgb, 255:red, 0; green, 0; blue, 0 }  ][line width=0.08]  [draw opacity=0] (11.61,-5.58) -- (0,0) -- (11.61,5.58) -- cycle    ;
%Straight Lines [id:da21012216073255352] 
\draw [line width=1.5]    (95.53,172.08) -- (162.47,115.33) ;
\draw [shift={(92.47,174.67)}, rotate = 319.71] [fill={rgb, 255:red, 0; green, 0; blue, 0 }  ][line width=0.08]  [draw opacity=0] (11.61,-5.58) -- (0,0) -- (11.61,5.58) -- cycle    ;
%Straight Lines [id:da7410869854550676] 
\draw [line width=1.5]    (173.14,115.33) -- (239.43,172.07) ;
\draw [shift={(242.47,174.67)}, rotate = 220.56] [fill={rgb, 255:red, 0; green, 0; blue, 0 }  ][line width=0.08]  [draw opacity=0] (11.61,-5.58) -- (0,0) -- (11.61,5.58) -- cycle    ;
%Curve Lines [id:da21492669607570647] 
\draw [line width=1.5]    (87.21,175.61) .. controls (89.24,80.37) and (125.36,38.14) .. (161.81,25.33) ;
\draw [shift={(87.14,180)}, rotate = 270.58] [fill={rgb, 255:red, 0; green, 0; blue, 0 }  ][line width=0.08]  [draw opacity=0] (11.61,-5.58) -- (0,0) -- (11.61,5.58) -- cycle    ;
%Curve Lines [id:da9036325885810722] 
\draw [line width=1.5]    (172.47,25.33) .. controls (212.45,25.01) and (258.44,127.03) .. (248.67,176.32) ;
\draw [shift={(247.81,180)}, rotate = 285.26] [fill={rgb, 255:red, 0; green, 0; blue, 0 }  ][line width=0.08]  [draw opacity=0] (11.61,-5.58) -- (0,0) -- (11.61,5.58) -- cycle    ;
%Shape: Circle [id:dp6426609962699884] 
\draw  [line width=1.5]  (87.14,180) .. controls (87.14,177.05) and (89.53,174.67) .. (92.47,174.67) .. controls (95.42,174.67) and (97.81,177.05) .. (97.81,180) .. controls (97.81,182.95) and (95.42,185.33) .. (92.47,185.33) .. controls (89.53,185.33) and (87.14,182.95) .. (87.14,180) -- cycle ;
%Shape: Circle [id:dp45710635129265853] 
\draw  [line width=1.5]  (237.14,180) .. controls (237.14,177.05) and (239.53,174.67) .. (242.47,174.67) .. controls (245.42,174.67) and (247.81,177.05) .. (247.81,180) .. controls (247.81,182.95) and (245.42,185.33) .. (242.47,185.33) .. controls (239.53,185.33) and (237.14,182.95) .. (237.14,180) -- cycle ;
%Shape: Circle [id:dp9010615426685297] 
\draw  [line width=1.5]  (165.81,230) .. controls (165.81,227.05) and (168.2,224.67) .. (171.14,224.67) .. controls (174.09,224.67) and (176.47,227.05) .. (176.47,230) .. controls (176.47,232.95) and (174.09,235.33) .. (171.14,235.33) .. controls (168.2,235.33) and (165.81,232.95) .. (165.81,230) -- cycle ;
%Shape: Circle [id:dp8341928324956512] 
\draw  [line width=1.5]  (166.47,320) .. controls (166.47,317.05) and (168.86,314.67) .. (171.81,314.67) .. controls (174.75,314.67) and (177.14,317.05) .. (177.14,320) .. controls (177.14,322.95) and (174.75,325.33) .. (171.81,325.33) .. controls (168.86,325.33) and (166.47,322.95) .. (166.47,320) -- cycle ;
%Shape: Circle [id:dp7263331520607497] 
\draw  [line width=1.5]  (241.14,384.67) .. controls (241.14,381.72) and (243.53,379.33) .. (246.47,379.33) .. controls (249.42,379.33) and (251.81,381.72) .. (251.81,384.67) .. controls (251.81,387.61) and (249.42,390) .. (246.47,390) .. controls (243.53,390) and (241.14,387.61) .. (241.14,384.67) -- cycle ;
%Straight Lines [id:da9896464315387669] 
\draw [line width=1.5]    (171.17,239.33) -- (171.81,314.67) ;
\draw [shift={(171.14,235.33)}, rotate = 89.52] [fill={rgb, 255:red, 0; green, 0; blue, 0 }  ][line width=0.08]  [draw opacity=0] (11.61,-5.58) -- (0,0) -- (11.61,5.58) -- cycle    ;
%Straight Lines [id:da3068077560217968] 
\draw [line width=1.5]    (101.81,384.67) -- (237.14,384.67) ;
\draw [shift={(241.14,384.67)}, rotate = 180] [fill={rgb, 255:red, 0; green, 0; blue, 0 }  ][line width=0.08]  [draw opacity=0] (11.61,-5.58) -- (0,0) -- (11.61,5.58) -- cycle    ;
%Straight Lines [id:da1192544795528433] 
\draw [line width=1.5]    (99.53,376.75) -- (166.47,320) ;
\draw [shift={(96.47,379.33)}, rotate = 319.71] [fill={rgb, 255:red, 0; green, 0; blue, 0 }  ][line width=0.08]  [draw opacity=0] (11.61,-5.58) -- (0,0) -- (11.61,5.58) -- cycle    ;
%Straight Lines [id:da5014217572994905] 
\draw [line width=1.5]    (177.14,320) -- (243.43,376.73) ;
\draw [shift={(246.47,379.33)}, rotate = 220.56] [fill={rgb, 255:red, 0; green, 0; blue, 0 }  ][line width=0.08]  [draw opacity=0] (11.61,-5.58) -- (0,0) -- (11.61,5.58) -- cycle    ;
%Curve Lines [id:da14877238718382224] 
\draw [line width=1.5]    (91.21,380.28) .. controls (93.24,285.04) and (129.36,242.81) .. (165.81,230) ;
\draw [shift={(91.14,384.67)}, rotate = 270.58] [fill={rgb, 255:red, 0; green, 0; blue, 0 }  ][line width=0.08]  [draw opacity=0] (11.61,-5.58) -- (0,0) -- (11.61,5.58) -- cycle    ;
%Curve Lines [id:da6788551769771028] 
\draw [line width=1.5]    (176.47,230) .. controls (216.45,229.68) and (262.44,331.7) .. (252.67,380.98) ;
\draw [shift={(251.81,384.67)}, rotate = 285.26] [fill={rgb, 255:red, 0; green, 0; blue, 0 }  ][line width=0.08]  [draw opacity=0] (11.61,-5.58) -- (0,0) -- (11.61,5.58) -- cycle    ;
%Shape: Circle [id:dp8968629678267172] 
\draw  [line width=1.5]  (91.14,384.67) .. controls (91.14,381.72) and (93.53,379.33) .. (96.47,379.33) .. controls (99.42,379.33) and (101.81,381.72) .. (101.81,384.67) .. controls (101.81,387.61) and (99.42,390) .. (96.47,390) .. controls (93.53,390) and (91.14,387.61) .. (91.14,384.67) -- cycle ;
%Shape: Circle [id:dp7197685772792822] 
\draw  [line width=1.5]  (452.14,20) .. controls (452.14,17.05) and (454.53,14.67) .. (457.47,14.67) .. controls (460.42,14.67) and (462.81,17.05) .. (462.81,20) .. controls (462.81,22.95) and (460.42,25.33) .. (457.47,25.33) .. controls (454.53,25.33) and (452.14,22.95) .. (452.14,20) -- cycle ;
%Shape: Circle [id:dp49004728051988067] 
\draw  [line width=1.5]  (452.81,110) .. controls (452.81,107.05) and (455.2,104.67) .. (458.14,104.67) .. controls (461.09,104.67) and (463.47,107.05) .. (463.47,110) .. controls (463.47,112.95) and (461.09,115.33) .. (458.14,115.33) .. controls (455.2,115.33) and (452.81,112.95) .. (452.81,110) -- cycle ;
%Shape: Circle [id:dp9996086110664077] 
\draw  [line width=1.5]  (527.47,174.67) .. controls (527.47,171.72) and (529.86,169.33) .. (532.81,169.33) .. controls (535.75,169.33) and (538.14,171.72) .. (538.14,174.67) .. controls (538.14,177.61) and (535.75,180) .. (532.81,180) .. controls (529.86,180) and (527.47,177.61) .. (527.47,174.67) -- cycle ;
%Straight Lines [id:da5916120908249641] 
\draw [line width=1.5]    (457.47,25.33) -- (458.11,100.67) ;
\draw [shift={(458.14,104.67)}, rotate = 269.52] [fill={rgb, 255:red, 0; green, 0; blue, 0 }  ][line width=0.08]  [draw opacity=0] (11.61,-5.58) -- (0,0) -- (11.61,5.58) -- cycle    ;
%Straight Lines [id:da2675543814820319] 
\draw [line width=1.5]    (388.14,174.67) -- (523.47,174.67) ;
\draw [shift={(527.47,174.67)}, rotate = 180] [fill={rgb, 255:red, 0; green, 0; blue, 0 }  ][line width=0.08]  [draw opacity=0] (11.61,-5.58) -- (0,0) -- (11.61,5.58) -- cycle    ;
%Straight Lines [id:da8619340282694029] 
\draw [line width=1.5]    (385.86,166.75) -- (452.81,110) ;
\draw [shift={(382.81,169.33)}, rotate = 319.71] [fill={rgb, 255:red, 0; green, 0; blue, 0 }  ][line width=0.08]  [draw opacity=0] (11.61,-5.58) -- (0,0) -- (11.61,5.58) -- cycle    ;
%Straight Lines [id:da3311013263120153] 
\draw [line width=1.5]    (463.47,110) -- (529.77,166.73) ;
\draw [shift={(532.81,169.33)}, rotate = 220.56] [fill={rgb, 255:red, 0; green, 0; blue, 0 }  ][line width=0.08]  [draw opacity=0] (11.61,-5.58) -- (0,0) -- (11.61,5.58) -- cycle    ;
%Curve Lines [id:da6945139417973814] 
\draw [line width=1.5]    (377.54,170.28) .. controls (379.57,75.04) and (415.7,32.8) .. (452.14,20) ;
\draw [shift={(377.47,174.67)}, rotate = 270.58] [fill={rgb, 255:red, 0; green, 0; blue, 0 }  ][line width=0.08]  [draw opacity=0] (11.61,-5.58) -- (0,0) -- (11.61,5.58) -- cycle    ;
%Curve Lines [id:da7822794249467356] 
\draw [line width=1.5]    (462.81,20) .. controls (502.78,19.67) and (548.78,121.7) .. (539,170.98) ;
\draw [shift={(538.14,174.67)}, rotate = 285.26] [fill={rgb, 255:red, 0; green, 0; blue, 0 }  ][line width=0.08]  [draw opacity=0] (11.61,-5.58) -- (0,0) -- (11.61,5.58) -- cycle    ;
%Shape: Circle [id:dp6495661214839258] 
\draw  [line width=1.5]  (377.47,174.67) .. controls (377.47,171.72) and (379.86,169.33) .. (382.81,169.33) .. controls (385.75,169.33) and (388.14,171.72) .. (388.14,174.67) .. controls (388.14,177.61) and (385.75,180) .. (382.81,180) .. controls (379.86,180) and (377.47,177.61) .. (377.47,174.67) -- cycle ;
%Shape: Circle [id:dp10015321763196539] 
\draw  [line width=1.5]  (453.14,230) .. controls (453.14,227.05) and (455.53,224.67) .. (458.47,224.67) .. controls (461.42,224.67) and (463.81,227.05) .. (463.81,230) .. controls (463.81,232.95) and (461.42,235.33) .. (458.47,235.33) .. controls (455.53,235.33) and (453.14,232.95) .. (453.14,230) -- cycle ;
%Shape: Circle [id:dp03418543258112594] 
\draw  [line width=1.5]  (453.81,320) .. controls (453.81,317.05) and (456.2,314.67) .. (459.14,314.67) .. controls (462.09,314.67) and (464.47,317.05) .. (464.47,320) .. controls (464.47,322.95) and (462.09,325.33) .. (459.14,325.33) .. controls (456.2,325.33) and (453.81,322.95) .. (453.81,320) -- cycle ;
%Shape: Circle [id:dp7359218836228637] 
\draw  [line width=1.5]  (528.47,384.67) .. controls (528.47,381.72) and (530.86,379.33) .. (533.81,379.33) .. controls (536.75,379.33) and (539.14,381.72) .. (539.14,384.67) .. controls (539.14,387.61) and (536.75,390) .. (533.81,390) .. controls (530.86,390) and (528.47,387.61) .. (528.47,384.67) -- cycle ;
%Straight Lines [id:da3785724466533624] 
\draw [line width=1.5]    (458.47,235.33) -- (459.11,310.67) ;
\draw [shift={(459.14,314.67)}, rotate = 269.52] [fill={rgb, 255:red, 0; green, 0; blue, 0 }  ][line width=0.08]  [draw opacity=0] (11.61,-5.58) -- (0,0) -- (11.61,5.58) -- cycle    ;
%Straight Lines [id:da2697861935808339] 
\draw [line width=1.5]    (393.14,384.67) -- (528.47,384.67) ;
\draw [shift={(389.14,384.67)}, rotate = 0] [fill={rgb, 255:red, 0; green, 0; blue, 0 }  ][line width=0.08]  [draw opacity=0] (11.61,-5.58) -- (0,0) -- (11.61,5.58) -- cycle    ;
%Straight Lines [id:da17453965147392314] 
\draw [line width=1.5]    (386.86,376.75) -- (453.81,320) ;
\draw [shift={(383.81,379.33)}, rotate = 319.71] [fill={rgb, 255:red, 0; green, 0; blue, 0 }  ][line width=0.08]  [draw opacity=0] (11.61,-5.58) -- (0,0) -- (11.61,5.58) -- cycle    ;
%Straight Lines [id:da07864825916900786] 
\draw [line width=1.5]    (464.47,320) -- (530.77,376.73) ;
\draw [shift={(533.81,379.33)}, rotate = 220.56] [fill={rgb, 255:red, 0; green, 0; blue, 0 }  ][line width=0.08]  [draw opacity=0] (11.61,-5.58) -- (0,0) -- (11.61,5.58) -- cycle    ;
%Curve Lines [id:da9315806742362415] 
\draw [line width=1.5]    (378.54,380.28) .. controls (380.57,285.04) and (416.7,242.81) .. (453.14,230) ;
\draw [shift={(378.47,384.67)}, rotate = 270.58] [fill={rgb, 255:red, 0; green, 0; blue, 0 }  ][line width=0.08]  [draw opacity=0] (11.61,-5.58) -- (0,0) -- (11.61,5.58) -- cycle    ;
%Curve Lines [id:da17189628696395065] 
\draw [line width=1.5]    (463.81,230) .. controls (503.78,229.68) and (549.78,331.7) .. (540,380.98) ;
\draw [shift={(539.14,384.67)}, rotate = 285.26] [fill={rgb, 255:red, 0; green, 0; blue, 0 }  ][line width=0.08]  [draw opacity=0] (11.61,-5.58) -- (0,0) -- (11.61,5.58) -- cycle    ;
%Shape: Circle [id:dp47115208809575815] 
\draw  [line width=1.5]  (378.47,384.67) .. controls (378.47,381.72) and (380.86,379.33) .. (383.81,379.33) .. controls (386.75,379.33) and (389.14,381.72) .. (389.14,384.67) .. controls (389.14,387.61) and (386.75,390) .. (383.81,390) .. controls (380.86,390) and (378.47,387.61) .. (378.47,384.67) -- cycle ;
%Shape: Triangle [id:dp6389874267176925] 
\draw  [color={rgb, 255:red, 128; green, 128; blue, 128 }  ,draw opacity=0.61 ][dash pattern={on 1.69pt off 2.76pt}][line width=1.5]  (166.76,67.43) -- (132.85,7.11) -- (202.84,7.79) -- cycle ;
%Shape: Triangle [id:dp36108366327372987] 
\draw  [color={rgb, 255:red, 128; green, 128; blue, 128 }  ,draw opacity=0.61 ][dash pattern={on 1.69pt off 2.76pt}][line width=1.5]  (167.74,70.77) -- (202.24,140.77) -- (132.24,140.77) -- cycle ;
%Shape: Triangle [id:dp7364938247780082] 
\draw  [color={rgb, 255:red, 128; green, 128; blue, 128 }  ,draw opacity=0.61 ][dash pattern={on 1.69pt off 2.76pt}][line width=1.5]  (127.74,145.44) -- (103.8,218.67) -- (50.73,164.22) -- cycle ;
%Shape: Triangle [id:dp8688796137791671] 
\draw  [color={rgb, 255:red, 128; green, 128; blue, 128 }  ,draw opacity=0.61 ][dash pattern={on 1.69pt off 2.76pt}][line width=1.5]  (208.07,145.44) -- (279.33,174.73) -- (221.12,223.63) -- cycle ;
%Shape: Triangle [id:dp5510471902819721] 
\draw  [color={rgb, 255:red, 128; green, 128; blue, 128 }  ,draw opacity=0.61 ][dash pattern={on 1.69pt off 2.76pt}][line width=1.5]  (172.76,270.1) -- (138.85,209.77) -- (208.84,210.46) -- cycle ;
%Shape: Triangle [id:dp5832218749091068] 
\draw  [color={rgb, 255:red, 128; green, 128; blue, 128 }  ,draw opacity=0.61 ][dash pattern={on 1.69pt off 2.76pt}][line width=1.5]  (173.74,273.44) -- (208.24,343.44) -- (138.24,343.44) -- cycle ;
%Shape: Triangle [id:dp17185534363242116] 
\draw  [color={rgb, 255:red, 128; green, 128; blue, 128 }  ,draw opacity=0.61 ][dash pattern={on 1.69pt off 2.76pt}][line width=1.5]  (133.74,348.1) -- (109.8,421.33) -- (56.73,366.89) -- cycle ;
%Shape: Triangle [id:dp5927451881338996] 
\draw  [color={rgb, 255:red, 128; green, 128; blue, 128 }  ,draw opacity=0.61 ][dash pattern={on 1.69pt off 2.76pt}][line width=1.5]  (214.07,348.1) -- (285.33,377.39) -- (227.12,426.29) -- cycle ;
%Shape: Triangle [id:dp22961376462006855] 
\draw  [color={rgb, 255:red, 128; green, 128; blue, 128 }  ,draw opacity=0.61 ][dash pattern={on 1.69pt off 2.76pt}][line width=1.5]  (456.76,275.43) -- (422.85,215.11) -- (492.84,215.79) -- cycle ;
%Shape: Triangle [id:dp7412114645452874] 
\draw  [color={rgb, 255:red, 128; green, 128; blue, 128 }  ,draw opacity=0.61 ][dash pattern={on 1.69pt off 2.76pt}][line width=1.5]  (457.74,278.77) -- (492.24,348.77) -- (422.24,348.77) -- cycle ;
%Shape: Triangle [id:dp5335245159862751] 
\draw  [color={rgb, 255:red, 128; green, 128; blue, 128 }  ,draw opacity=0.61 ][dash pattern={on 1.69pt off 2.76pt}][line width=1.5]  (417.74,353.44) -- (393.8,426.67) -- (340.73,372.22) -- cycle ;
%Shape: Triangle [id:dp6675828844645042] 
\draw  [color={rgb, 255:red, 128; green, 128; blue, 128 }  ,draw opacity=0.61 ][dash pattern={on 1.69pt off 2.76pt}][line width=1.5]  (498.07,353.44) -- (569.33,382.73) -- (511.12,431.63) -- cycle ;
%Shape: Triangle [id:dp7953792568010355] 
\draw  [color={rgb, 255:red, 128; green, 128; blue, 128 }  ,draw opacity=0.61 ][dash pattern={on 1.69pt off 2.76pt}][line width=1.5]  (454.76,66.76) -- (420.85,6.44) -- (490.84,7.13) -- cycle ;
%Shape: Triangle [id:dp43961271511290967] 
\draw  [color={rgb, 255:red, 128; green, 128; blue, 128 }  ,draw opacity=0.61 ][dash pattern={on 1.69pt off 2.76pt}][line width=1.5]  (455.74,70.1) -- (490.24,140.1) -- (420.24,140.1) -- cycle ;
%Shape: Triangle [id:dp24827760567983648] 
\draw  [color={rgb, 255:red, 128; green, 128; blue, 128 }  ,draw opacity=0.61 ][dash pattern={on 1.69pt off 2.76pt}][line width=1.5]  (415.74,144.77) -- (391.8,218) -- (338.73,163.56) -- cycle ;
%Shape: Triangle [id:dp7191718328020039] 
\draw  [color={rgb, 255:red, 128; green, 128; blue, 128 }  ,draw opacity=0.61 ][dash pattern={on 1.69pt off 2.76pt}][line width=1.5]  (496.07,144.77) -- (567.33,174.06) -- (509.12,222.96) -- cycle ;

% Text Node
\draw (159.33,141.4) node [anchor=north west][inner sep=0.75pt]    {$H_{1}$};
% Text Node
\draw (194,86.73) node [anchor=north west][inner sep=0.75pt]    {$A$};
% Text Node
\draw (116,88.73) node [anchor=north west][inner sep=0.75pt]    {$C$};
% Text Node
\draw (115.33,188.4) node [anchor=north west][inner sep=0.75pt]    {$G_{2}$};
% Text Node
\draw (204,120.4) node [anchor=north west][inner sep=0.75pt]    {$1$};
% Text Node
\draw (123.33,120.4) node [anchor=north west][inner sep=0.75pt]    {$3$};
% Text Node
\draw (233.33,62.4) node [anchor=north west][inner sep=0.75pt]    {$3$};
% Text Node
\draw (96,51.07) node [anchor=north west][inner sep=0.75pt]    {$3$};
% Text Node
\draw (152,71.07) node [anchor=north west][inner sep=0.75pt]    {$2$};
% Text Node
\draw (164.67,185.07) node [anchor=north west][inner sep=0.75pt]    {$5$};
% Text Node
\draw (154.9,270.75) node [anchor=north west][inner sep=0.75pt]  [rotate=-2.55]  {$2$};
% Text Node
\draw (201.5,296.09) node [anchor=north west][inner sep=0.75pt]  [rotate=-2.55]  {$D$};
% Text Node
\draw (128.62,298.44) node [anchor=north west][inner sep=0.75pt]  [rotate=-2.55]  {$B$};
% Text Node
\draw (81.84,294.23) node [anchor=north west][inner sep=0.75pt]  [rotate=-2.55]  {$1$};
% Text Node
\draw (246.5,292.23) node [anchor=north west][inner sep=0.75pt]  [rotate=-2.55]  {$1$};
% Text Node
\draw (216.36,332.26) node [anchor=north west][inner sep=0.75pt]  [rotate=-2.55]  {$3$};
% Text Node
\draw (125.05,326.62) node [anchor=north west][inner sep=0.75pt]  [rotate=-2.55]  {$1$};
% Text Node
\draw (163.33,345.4) node [anchor=north west][inner sep=0.75pt]    {$H_{1}$};
% Text Node
\draw (168.67,392.73) node [anchor=north west][inner sep=0.75pt]    {$5$};
% Text Node
\draw (114.67,399.07) node [anchor=north west][inner sep=0.75pt]    {$G_{1}$};
% Text Node
\draw (417.33,82.4) node [anchor=north west][inner sep=0.75pt]    {$B$};
% Text Node
\draw (487.33,82.4) node [anchor=north west][inner sep=0.75pt]    {$A$};
% Text Node
\draw (451.33,140.4) node [anchor=north west][inner sep=0.75pt]    {$H_{2}$};
% Text Node
\draw (417.33,182.4) node [anchor=north west][inner sep=0.75pt]    {$G_{1}$};
% Text Node
\draw (381.33,72.4) node [anchor=north west][inner sep=0.75pt]    {$1$};
% Text Node
\draw (527.33,72.4) node [anchor=north west][inner sep=0.75pt]    {$1$};
% Text Node
\draw (407.33,122.4) node [anchor=north west][inner sep=0.75pt]    {$1$};
% Text Node
\draw (461.33,52.4) node [anchor=north west][inner sep=0.75pt]    {$2$};
% Text Node
\draw (497.33,122.4) node [anchor=north west][inner sep=0.75pt]    {$3$};
% Text Node
\draw (457.33,182.4) node [anchor=north west][inner sep=0.75pt]    {$5$};
% Text Node
\draw (416.33,292.4) node [anchor=north west][inner sep=0.75pt]    {$C$};
% Text Node
\draw (482.33,292.4) node [anchor=north west][inner sep=0.75pt]    {$D$};
% Text Node
\draw (447.33,342.4) node [anchor=north west][inner sep=0.75pt]    {$H_{2}$};
% Text Node
\draw (412.33,390.4) node [anchor=north west][inner sep=0.75pt]    {$G_{2}$};
% Text Node
\draw (382.33,282.4) node [anchor=north west][inner sep=0.75pt]    {$3$};
% Text Node
\draw (528.33,282.4) node [anchor=north west][inner sep=0.75pt]    {$3$};
% Text Node
\draw (408.33,332.4) node [anchor=north west][inner sep=0.75pt]    {$3$};
% Text Node
\draw (460.81,388.07) node [anchor=north west][inner sep=0.75pt]    {$5$};
% Text Node
\draw (462.33,272.4) node [anchor=north west][inner sep=0.75pt]    {$2$};
% Text Node
\draw (498.33,332.4) node [anchor=north west][inner sep=0.75pt]    {$1$};
% Text Node
\draw (212.67,130.4) node [anchor=north west][inner sep=0.75pt]    {$z_{1}$};
% Text Node
\draw (80.67,77.07) node [anchor=north west][inner sep=0.75pt]    {$z_{1}$};
% Text Node
\draw (106,131.07) node [anchor=north west][inner sep=0.75pt]    {$z_{2}$};
% Text Node
\draw (240,83.73) node [anchor=north west][inner sep=0.75pt]    {$z_{2}$};
% Text Node
\draw (170.67,67.4) node [anchor=north west][inner sep=0.75pt]    {$z_{3}$};
% Text Node
\draw (182,180.07) node [anchor=north west][inner sep=0.75pt]    {$z_{3}$};
% Text Node
\draw (506.67,129.4) node [anchor=north west][inner sep=0.75pt]    {$w_{1}$};
% Text Node
\draw (362.67,87.4) node [anchor=north west][inner sep=0.75pt]    {$w_{1}$};
% Text Node
\draw (413.33,109.4) node [anchor=north west][inner sep=0.75pt]    {$w_{2}$};
% Text Node
\draw (534,91.4) node [anchor=north west][inner sep=0.75pt]    {$w_{2}$};
% Text Node
\draw (436.47,59.73) node [anchor=north west][inner sep=0.75pt]    {$w_{3}$};
% Text Node
\draw (473.81,173.73) node [anchor=north west][inner sep=0.75pt]    {$w_{3}$};
% Text Node
\draw (222,342.4) node [anchor=north west][inner sep=0.75pt]    {$u_{1}$};
% Text Node
\draw (86,281.07) node [anchor=north west][inner sep=0.75pt]    {$u_{1}$};
% Text Node
\draw (107.33,339.07) node [anchor=north west][inner sep=0.75pt]    {$u_{2}$};
% Text Node
\draw (236,271.73) node [anchor=north west][inner sep=0.75pt]    {$u_{2}$};
% Text Node
\draw (173.33,271.73) node [anchor=north west][inner sep=0.75pt]    {$u_{3}$};
% Text Node
\draw (182,387.73) node [anchor=north west][inner sep=0.75pt]    {$u_{3}$};
% Text Node
\draw (511.33,339.07) node [anchor=north west][inner sep=0.75pt]    {$t_{1}$};
% Text Node
\draw (382.67,260.4) node [anchor=north west][inner sep=0.75pt]    {$t_{1}$};
% Text Node
\draw (394,342.73) node [anchor=north west][inner sep=0.75pt]    {$t_{2}$};
% Text Node
\draw (522,259.4) node [anchor=north west][inner sep=0.75pt]    {$t_{2}$};
% Text Node
\draw (443.33,264.4) node [anchor=north west][inner sep=0.75pt]    {$t_{3}$};
% Text Node
\draw (477.33,385.73) node [anchor=north west][inner sep=0.75pt]    {$t_{3}$};
% Text Node
\draw (161.33,117.4) node [anchor=north west][inner sep=0.75pt]    {$a$};
% Text Node
\draw (177.33,6.07) node [anchor=north west][inner sep=0.75pt]    {$b$};
% Text Node
\draw (454,118.73) node [anchor=north west][inner sep=0.75pt]    {$c$};
% Text Node
\draw (434,8.07) node [anchor=north west][inner sep=0.75pt]    {$d$};
% Text Node
\draw (176.47,210.73) node [anchor=north west][inner sep=0.75pt]    {$f$};
% Text Node
\draw (455.81,323.4) node [anchor=north west][inner sep=0.75pt]    {$g$};
% Text Node
\draw (467.14,212.07) node [anchor=north west][inner sep=0.75pt]    {$h$};
% Text Node
\draw (168.47,323.4) node [anchor=north west][inner sep=0.75pt]    {$e$};
% Text Node
\draw (64.67,161.07) node [anchor=north west][inner sep=0.75pt]    {$a^{'}$};
% Text Node
\draw (228,392.4) node [anchor=north west][inner sep=0.75pt]    {$b^{'}$};
% Text Node
\draw (517.33,393.73) node [anchor=north west][inner sep=0.75pt]    {$c^{'}$};
% Text Node
\draw (378.67,181.73) node [anchor=north west][inner sep=0.75pt]    {$d^{'}$};
% Text Node
\draw (98.47,388.07) node [anchor=north west][inner sep=0.75pt]    {$e'$};
% Text Node
\draw (224.67,187.73) node [anchor=north west][inner sep=0.75pt]    {$f'$};
% Text Node
\draw (512.67,186.73) node [anchor=north west][inner sep=0.75pt]    {$g'$};
% Text Node
\draw (380.67,392.4) node [anchor=north west][inner sep=0.75pt]    {$h'$};


\end{tikzpicture}
    \caption{Triangulos en los tetraedros truncados que representan las cuspides}
\end{figure}
Pegando de manera adecuada obtenemos dos paralelogramos que nos representan cada frontera de cúspide como se muestra a continuación
\begin{figure}[H]
\centering
%!TEX root = ./main.tex



\tikzset{every picture/.style={line width=0.75pt}} %set default line width to 0.75pt        

\begin{tikzpicture}[x=0.75pt,y=0.75pt,yscale=-1,xscale=1]
%uncomment if require: \path (0,542); %set diagram left start at 0, and has height of 542

%Straight Lines [id:da13006069871416692] 
\draw [line width=1.5]    (142,195) -- (630,195) ;
%Straight Lines [id:da13561558542782937] 
\draw [line width=1.5]    (142,195) -- (20,295) ;
%Straight Lines [id:da6048067045021074] 
\draw [line width=1.5]    (630,195) -- (508,295) ;
%Straight Lines [id:da7479845165592521] 
\draw [line width=1.5]    (20,295) -- (508,295) ;
%Straight Lines [id:da5146924711453463] 
\draw [line width=1.5]    (142,195) -- (142,295) ;
%Straight Lines [id:da7683491262763936] 
\draw [line width=1.5]    (264,195) -- (264,295) ;
%Straight Lines [id:da15527473375540246] 
\draw [line width=1.5]    (386,195) -- (386,295) ;
%Straight Lines [id:da45625934758008424] 
\draw [line width=1.5]    (508,195) -- (508,295) ;
%Straight Lines [id:da3616387902013627] 
\draw [line width=1.5]    (386,195) -- (264,295) ;
%Straight Lines [id:da044739648304707336] 
\draw [line width=1.5]    (264,295) -- (142,195) ;
%Straight Lines [id:da2836418136045904] 
\draw [line width=1.5]    (508,295) -- (386,195) ;
%Straight Lines [id:da867664282695949] 
\draw [line width=1.5]    (152,55) -- (640,55) ;
%Straight Lines [id:da9918221704170823] 
\draw [line width=1.5]    (152,55) -- (30,155) ;
%Straight Lines [id:da5944725906481361] 
\draw [line width=1.5]    (640,55) -- (518,155) ;
%Straight Lines [id:da8733020176164918] 
\draw [line width=1.5]    (30,155) -- (518,155) ;
%Straight Lines [id:da8720752195824351] 
\draw [line width=1.5]    (152,55) -- (152,155) ;
%Straight Lines [id:da9117542313308579] 
\draw [line width=1.5]    (274,55) -- (274,155) ;
%Straight Lines [id:da28083917441028194] 
\draw [line width=1.5]    (396,55) -- (396,155) ;
%Straight Lines [id:da06592295243917656] 
\draw [line width=1.5]    (518,55) -- (518,155) ;
%Straight Lines [id:da9134925142291349] 
\draw [line width=1.5]    (396,55) -- (274,155) ;
%Straight Lines [id:da022240721894983384] 
\draw [line width=1.5]    (274,155) -- (152,55) ;
%Straight Lines [id:da6293934000840508] 
\draw [line width=1.5]    (518,155) -- (396,55) ;
%Curve Lines [id:da5316729016494588] 
\draw [color={rgb, 255:red, 208; green, 2; blue, 27 }  ,draw opacity=1 ][line width=1.5]    (85,110) .. controls (125,80) and (539,135) .. (579,105) ;
\draw [shift={(338.81,107.87)}, rotate = 183.15] [fill={rgb, 255:red, 208; green, 2; blue, 27 }  ,fill opacity=1 ][line width=0.08]  [draw opacity=0] (11.61,-5.58) -- (0,0) -- (11.61,5.58) -- cycle    ;
%Curve Lines [id:da09367095573809125] 
\draw [color={rgb, 255:red, 208; green, 2; blue, 27 }  ,draw opacity=1 ][line width=1.5]    (75,250) .. controls (115,220) and (529,275) .. (569,245) ;
\draw [shift={(328.81,247.87)}, rotate = 183.15] [fill={rgb, 255:red, 208; green, 2; blue, 27 }  ,fill opacity=1 ][line width=0.08]  [draw opacity=0] (11.61,-5.58) -- (0,0) -- (11.61,5.58) -- cycle    ;
%Curve Lines [id:da08706634222859888] 
\draw [color={rgb, 255:red, 74; green, 144; blue, 226 }  ,draw opacity=1 ][line width=1.5]    (110,154) .. controls (150,124) and (183,84) .. (223,54) ;
\draw [shift={(171.43,99.28)}, rotate = 136.23] [fill={rgb, 255:red, 74; green, 144; blue, 226 }  ,fill opacity=1 ][line width=0.08]  [draw opacity=0] (11.61,-5.58) -- (0,0) -- (11.61,5.58) -- cycle    ;
%Curve Lines [id:da3000463539903655] 
\draw [color={rgb, 255:red, 74; green, 144; blue, 226 }  ,draw opacity=1 ][line width=1.5]    (97,295) .. controls (137,265) and (170,225) .. (210,195) ;
\draw [shift={(158.43,240.28)}, rotate = 136.23] [fill={rgb, 255:red, 74; green, 144; blue, 226 }  ,fill opacity=1 ][line width=0.08]  [draw opacity=0] (11.61,-5.58) -- (0,0) -- (11.61,5.58) -- cycle    ;

% Text Node
\draw (99,250.4) node [anchor=north west][inner sep=0.75pt]    {$a^{'}$};
% Text Node
\draw (122,274.4) node [anchor=north west][inner sep=0.75pt]    {$z_{3}$};
% Text Node
\draw (124,208.4) node [anchor=north west][inner sep=0.75pt]    {$z_{2}$};
% Text Node
\draw (43,274.4) node [anchor=north west][inner sep=0.75pt]    {$z_{1}$};
% Text Node
\draw (171,253.4) node [anchor=north west][inner sep=0.75pt]    {$b'$};
% Text Node
\draw (147,275.4) node [anchor=north west][inner sep=0.75pt]    {$u_{3}$};
% Text Node
\draw (143,209.4) node [anchor=north west][inner sep=0.75pt]    {$u_{1}$};
% Text Node
\draw (227,274.4) node [anchor=north west][inner sep=0.75pt]    {$u_{2}$};
% Text Node
\draw (211,217.4) node [anchor=north west][inner sep=0.75pt]    {$c'$};
% Text Node
\draw (246,261.4) node [anchor=north west][inner sep=0.75pt]    {$t_{1}$};
% Text Node
\draw (247,197.4) node [anchor=north west][inner sep=0.75pt]    {$t_{3}$};
% Text Node
\draw (165,195.4) node [anchor=north west][inner sep=0.75pt]    {$t_{2}$};
% Text Node
\draw (294,217.4) node [anchor=north west][inner sep=0.75pt]    {$d'$};
% Text Node
\draw (266,198.4) node [anchor=north west][inner sep=0.75pt]    {$w_{3}$};
% Text Node
\draw (266,256.4) node [anchor=north west][inner sep=0.75pt]    {$w_{2}$};
% Text Node
\draw (345,192.4) node [anchor=north west][inner sep=0.75pt]    {$w_{1}$};
% Text Node
\draw (336,254.4) node [anchor=north west][inner sep=0.75pt]    {$e'$};
% Text Node
\draw (365,207.4) node [anchor=north west][inner sep=0.75pt]    {$u_{2}$};
% Text Node
\draw (367,273.4) node [anchor=north west][inner sep=0.75pt]    {$u_{3}$};
% Text Node
\draw (286,275.4) node [anchor=north west][inner sep=0.75pt]    {$u_{1}$};
% Text Node
\draw (414,254.4) node [anchor=north west][inner sep=0.75pt]    {$f'$};
% Text Node
\draw (391,207.4) node [anchor=north west][inner sep=0.75pt]    {$z_{1}$};
% Text Node
\draw (390,273.4) node [anchor=north west][inner sep=0.75pt]    {$z_{3}$};
% Text Node
\draw (473,275.4) node [anchor=north west][inner sep=0.75pt]    {$z_{2}$};
% Text Node
\draw (458,217.4) node [anchor=north west][inner sep=0.75pt]    {$g'$};
% Text Node
\draw (487,260.4) node [anchor=north west][inner sep=0.75pt]    {$w_{1}$};
% Text Node
\draw (407,194.4) node [anchor=north west][inner sep=0.75pt]    {$w_{2}$};
% Text Node
\draw (487,197.4) node [anchor=north west][inner sep=0.75pt]    {$w_{3}$};
% Text Node
\draw (535,223.4) node [anchor=north west][inner sep=0.75pt]    {$h'$};
% Text Node
\draw (511,261.4) node [anchor=north west][inner sep=0.75pt]    {$t_{2}$};
% Text Node
\draw (511,197.4) node [anchor=north west][inner sep=0.75pt]    {$t_{3}$};
% Text Node
\draw (589,196.4) node [anchor=north west][inner sep=0.75pt]    {$t_{1}$};
% Text Node
\draw (109,110.4) node [anchor=north west][inner sep=0.75pt]    {$a$};
% Text Node
\draw (132,134.4) node [anchor=north west][inner sep=0.75pt]    {$z_{3}$};
% Text Node
\draw (134,68.4) node [anchor=north west][inner sep=0.75pt]    {$z_{2}$};
% Text Node
\draw (53,134.4) node [anchor=north west][inner sep=0.75pt]    {$z_{1}$};
% Text Node
\draw (181,113.4) node [anchor=north west][inner sep=0.75pt]    {$h$};
% Text Node
\draw (400,130.4) node [anchor=north west][inner sep=0.75pt]    {$u_{3}$};
% Text Node
\draw (400,70.4) node [anchor=north west][inner sep=0.75pt]    {$u_{1}$};
% Text Node
\draw (480,132.4) node [anchor=north west][inner sep=0.75pt]    {$u_{2}$};
% Text Node
\draw (221,77.4) node [anchor=north west][inner sep=0.75pt]    {$b$};
% Text Node
\draw (154,62.4) node [anchor=north west][inner sep=0.75pt]    {$t_{1}$};
% Text Node
\draw (155,133.4) node [anchor=north west][inner sep=0.75pt]    {$t_{3}$};
% Text Node
\draw (237,133.4) node [anchor=north west][inner sep=0.75pt]    {$t_{2}$};
% Text Node
\draw (304,77.4) node [anchor=north west][inner sep=0.75pt]    {$g$};
% Text Node
\draw (371,130.4) node [anchor=north west][inner sep=0.75pt]    {$w_{3}$};
% Text Node
\draw (372,69.4) node [anchor=north west][inner sep=0.75pt]    {$w_{2}$};
% Text Node
\draw (297,132.4) node [anchor=north west][inner sep=0.75pt]    {$w_{1}$};
% Text Node
\draw (346,114.4) node [anchor=north west][inner sep=0.75pt]    {$c$};
% Text Node
\draw (520,120.4) node [anchor=north west][inner sep=0.75pt]    {$u_{2}$};
% Text Node
\draw (523,57.4) node [anchor=north west][inner sep=0.75pt]    {$u_{3}$};
% Text Node
\draw (600,52.4) node [anchor=north west][inner sep=0.75pt]    {$u_{1}$};
% Text Node
\draw (424,114.4) node [anchor=north west][inner sep=0.75pt]    {$f$};
% Text Node
\draw (258,120.4) node [anchor=north west][inner sep=0.75pt]    {$z_{1}$};
% Text Node
\draw (252,52.4) node [anchor=north west][inner sep=0.75pt]    {$z_{3}$};
% Text Node
\draw (172,52.4) node [anchor=north west][inner sep=0.75pt]    {$z_{2}$};
% Text Node
\draw (468,77.4) node [anchor=north west][inner sep=0.75pt]    {$d$};
% Text Node
\draw (497,120.4) node [anchor=north west][inner sep=0.75pt]    {$w_{1}$};
% Text Node
\draw (417,54.4) node [anchor=north west][inner sep=0.75pt]    {$w_{2}$};
% Text Node
\draw (497,57.4) node [anchor=north west][inner sep=0.75pt]    {$w_{3}$};
% Text Node
\draw (545,83.4) node [anchor=north west][inner sep=0.75pt]    {$e$};
% Text Node
\draw (276,118.4) node [anchor=north west][inner sep=0.75pt]    {$t_{2}$};
% Text Node
\draw (275,54.4) node [anchor=north west][inner sep=0.75pt]    {$t_{3}$};
% Text Node
\draw (360,54.4) node [anchor=north west][inner sep=0.75pt]    {$t_{1}$};


\end{tikzpicture}
    \caption{Triangulación de las cúspides y curvas generadoras}
\end{figure}
Observe que de una vez en cada cúspide trazamos las curvas que generan el grupo fundamental, en rojo $\alpha,\alpha^\prime$ y en azul $\beta,\beta^\prime$. Así el numero asociado a cada curva es
\begin{itemize}
    \item $H(\alpha)=z_2t_1z_1^{-1}t_2^{-1}w_2u_1w_1^{-1}u_2^{-1}=\dfrac{z_2t_1w_2u_1}{z_1t_2w_1u_2}.$
    \item $H(\beta)=z_3^{-1}z_2t_1=\dfrac{z_2t_1}{z_3}.$
    \item $H(\alpha^\prime)=z_2u_1t_1^{-1}w_2^{-1}u_2z_1w_1^{-1}t_2^{-1}=\dfrac{z_2u_1u_2z_1}{t_1w_2w_1t_2}.$
    \item $H(\beta^\prime)=z_3^{-1}u_1t_2=\dfrac{u_1t_2}{z_3}$
\end{itemize}
Usando las relaciones de de invariantes de aristas obtenemos las ecuaciones de completez que complementan a las de pegado halladas antes
$$\begin{cases}
    \dfrac{ut(1-u)(1-t)}{zw(1-z)(1-w)}=1\\
    \\
    -\dfrac{zt}{(z-1)^2}=1\\
    \\
    \dfrac{uz(1-w)(1-t)}{tw(1-z)(1-u)}=1\\
    \\
    \dfrac{uz}{(z-1)(1-t)}=1
\end{cases}$$
Nuevamente hallar soluciones a mano seria excesivamente engorroso y para esto es mejor usar herramientas computacionales, pero note que esto nos describe en su totalidad los valores que puede tomar cada invariante de arista, por lo que si uno quisiera unos números en particular basta con remplazar y verificar que todo se cumple.
    \section{La conjetura del volumen}
    Hemos visto una forma de dotar a complementos de un nudo de una estructura hiperbólica completa, al igual que en los preliminares hallamos distancias y áreas dada la métrica hiperbólica, si dotamos a una 3-variedad de esta, de manera local en coordenadas la métrica es la hiperbólica por lo que podríamos calcular el volumen de esas regiones y por tanto el volumen de complementos de nudos
    \begin{definition}
        Dado un nudo $K\subset S^3$, definimos el volumen hiperbolico como
        $$Vol(K):=Vol(S^3\setminus K)$$
    \end{definition}
    Resulta que en la métrica hiperbólica el volumen de estas tres variedades siempre es finita, por lo que en un principio resultaría superflua la condición de completez de la métrica, pero el siguiente teorema de Mostow nos dice el por que es necesaria
    \begin{theorem}
        Suponga que $M$ es una 3 variedad con volumen finito, si $M$ posee una estructura hiperbólica completa, entonces esta estructura es única en $M$ salvo isometrias.
    \end{theorem}
    Este teorema nos asegura que la estructura hiperbólica completa en el complemento de un nudo es única salvo isometria y por tanto es un invariante de nudos, y por tanto todo aquello que calculemos de esa 3-variedad, como por el ejemplo el volumen, también sera un invariante. Hasta la fecha la mayoría de volúmenes hiperbólicos son hallados de manera computacional y por medio de cotas superiores e inferiores, mas no hay formula cerrada salvo casos contados como el nudo de $8$. Aproximaciones de algunos de estos los podemos visualizar en la siguiente tabla.
    \begin{table}[H]
\centering
\begin{tabular}{|c|c|c|}
\hline
\textbf{Nudo} & \textbf{Volumen hiperbólico} \\
\hline
$4_1$  & $2.02988$ \\
\hline
$5_2$  & $2.82812$ \\
\hline
$6_1$  & $3.16396$ \\
\hline
$6_2$   & $3.16396$ \\
\hline
$6_3$   & $5.69302$ \\
\hline
$7_2$   & $3.33174$ \\
\hline
$7_3$   & $4.59296$ \\
\hline
$7_4$   & $5.13794$ \\
\hline
$7_5$  & $5.33349$ \\
\hline
$7_6$   & $6.44353$ \\
\hline
$7_7$   & $6.89032$ \\
\hline
\end{tabular}
\caption{Volúmenes hiperbólicos de algunos nudos hiperbólicos.}
\end{table}
    Seria genial poder encontrar una formula cerrada para el volumen de cualquier nudo, mas esto es un problema abierto aun que se conoce como la conjetura del volumen, la cual nos dice que el volumen esta dado por la siguiente formula
    $$Vol(K)=\lim_{N\to\infty}\frac{2\pi\log|J_N(K;e^{\frac{\pi i}{N}})|}{N},$$
    Donde $J_N(K;e^{\frac{\pi i}{N}})$ es el $N-$esimo polinomio de Jones coloreado evaluado en $e^{\frac{\pi i}{N}}$. el cual es la generalización del polinomio de Jones clásico dada por los invariantes de Reshetikhin–Turaev.
\begin{thebibliography}{9}

%!TEX root = ./main.tex

\bibitem{purcell_hyperbolic_knot_theory}
Jessica S. Purcell,
\textit{Hyperbolic Knot Theory}.
School of Mathematics, Monash University,
Clayton VIC 3800, Australia.

\bibitem{thurston_3d_geometry_topology}
William P. Thurston,
\textit{Three-Dimensional Geometry and Topology, Vol. 1},
edited by Silvio Levy.
Princeton University Press,
Princeton, New Jersey, 1997.

\bibitem{KnotAtlas}
D. Bar-Natan et al.,
\emph{The Knot Atlas}.
Available at \url{http://katlas.org/},
accessed June 2026.

\bibitem{MurakamiVolumeConjecture}
Hitoshi Murakami,
\emph{An Introduction to the Volume Conjecture},
arXiv:1002.0120 [math.GT], 2010.

\bibitem{SnapPy}
Marc Culler, Nathan M. Dunfield, Matthias Goerner, and Jeffrey R. Weeks,
\emph{SnapPy, a Computer Program for Studying the Geometry and Topology of 3-Manifolds},
available at \url{https://snappy.computop.org}.

\end{thebibliography}

\end{document}


